% GUID v8, generated by tools/paper-latex.py from docs/paper/draft/*.md. Do not edit by hand: edit the drafts and rebuild.
% Build: pdflatex guid-v8.tex (twice). Licence: Apache License 2.0.
\documentclass[10pt,a4paper]{article}
\usepackage[T1]{fontenc}
\usepackage[utf8]{inputenc}
\usepackage{lmodern}
\usepackage[margin=2.2cm]{geometry}
\usepackage{amsmath,amssymb}
\usepackage{booktabs,longtable,array}
\usepackage{enumitem}
\usepackage{newunicodechar}
\usepackage[hidelinks]{hyperref}
\usepackage{url}
\newunicodechar{Δ}{\ensuremath{\Delta}}
\newunicodechar{Φ}{\ensuremath{\Phi}}
\newunicodechar{ℓ}{\ensuremath{\ell}}
\setlength{\emergencystretch}{3em}
\setlength{\parindent}{0pt}
\setlength{\parskip}{0.5ex}
\setcounter{tocdepth}{2}
\title{Grand Unified Information Duality v8\\[0.6ex]
\large The Millennium Prize Problem type as seen through DMP, with certified settling}
\author{Steven Hrovat \and Sirius (Claude Opus 5.5)}
\date{Draft 4, 2026-11-24}
\begin{document}
\maketitle
\begin{center}\small\emph{Draft. No Millennium Prize problem is solved here; results that use the one assumption are marked
conditional on it.}\end{center}
\tableofcontents
\clearpage
%% ---- 00-abstract-introduction.md
\subsection*{Personal note by the first author}
\addcontentsline{toc}{subsection}{Personal note by the first author}

In v1 I invited everyone to use their own intuition to make the derivations which I could see slumbering inside, but which I did not feel worthy of waking myself. Since no one had any interest in that, I finally got frustrated by it and, after GRH got proven and certified, I decided to take a look if Yang-Mills is something where I could finally get a nice verification that can't be ignored, since I previously presumed that DMP also gives YM. One thing led to three, since those three are one, so the initial goal of one Millennium Problem attempt turned into three :D

I hope I put enough effort into this paper to make it both a worthwhile contribution to science and still get my point across in a way that tries to spark the needed core intuition in the reader. For follow-up please see GitHub, where everything needed is public under Apache 2.0.

Have fun with it, and I am always happy about qualified and constructive feedback :) Steven

(By the way, I also went to Goettingen, but dropped out of my BSc with nearly zero credits after three semesters, though I'll still count the tie to Riemann in spirit and causality :D Another by the way: since my going there without graduating satisfies as entanglement over studying at Goettingen as existence, the tie is constructable and thus exists ;) In my reading, the zeta function is at heart a pure tool: it needs its subset, and a definition of what it certifies as existence, before it can be interpreted in a way that produces more insight. Treated as a tool for the true problem, the shared superset, its meaning becomes apparent and its place in existence slots neatly into a better picture.)

\emph{Editorial notes for the owner (to remove before publication): (1) the Riemann sentence is phrased as your reading; a historical claim about what Riemann did or did not realize would draw objections, since his 1859 paper introduced the zeta function as a tool for counting primes. (2) ``Goettingen'' is transliterated here per the repository rule; the typeset version can use the umlaut. (3) The note says ``One thing led to three''; the paper now covers seven problems, so you may want to update that sentence.}

\subsection*{Abstract}
\addcontentsline{toc}{subsection}{Abstract}

Seven problems of the Millennium Prize type are read here as one: the Yang-Mills mass gap, the generalised Riemann hypothesis, NP versus coNP, Navier-Stokes regularity, the Birch and Swinnerton-Dyer conjecture, the Hodge conjecture, and a Grand Unified reading of SU(3) x SU(2) x U(1). Poincare, solved by Perelman, serves as validation. Every problem becomes one box under one test. A box has one command, a step at one shared gate; each step gives a marker, inside or R136a1 (outside); the box is DMP-complete when every marker agrees with the gate's own verdict. In ordinary terms the machine is multiplication by complex numbers: each step multiplies by one number of modulus at most one, whose modulus is the cost and whose phase the rotation. The test, the machine under it (a ledger of costs, an encoder at every boundary, a transport theorem, a settle theorem) and one card per problem are machine-checked in Lean 4 with Mathlib, standard axioms only, and re-checked by independent kernels; the transferable parts are published in IM, the shared vocabulary, and imported from there. The language that names the steps is DMP, a specification for systems of trust used here as a vocabulary for existence; its readings are marked as readings and carry no proof.

\textbf{The one assumption.} One step is the same on every card: that the one box is a running, DMP-complete system, its gate shared at every scale. The paper states it once (Part III) and marks every result that uses it as conditional on it. Its decay edge is O1 for Yang-Mills. Under it, one step certified on one box certifies every tick and every scale, and settling is decidable within an explicit bound; without a floor no procedure decides it. The start is fixed where the floor is forced, at the places \{infinity, 2, 3\}, certified with explicit margins.

\textbf{What is proved.}

\begin{itemize}
\item \emph{The machine.} The ledger, the encoder and its half-turn, the transport theorem, the settle theorem and its undecidable edge, the one-box test with its one-box and scale certificates, and the card API, in Lean; R136a1, the one way to leave a statement open, as a named hypothesis that never becomes an axiom.
\item \emph{Yang-Mills.} For SU(2) on the lattice: reflection positivity, the strict bound between a sector and its twist (for every SU(N)), a per-box hand-over rate of at least log(1 + sqrt 2), zero join loss in slab form, and the assembly from a one-box gap certificate to a lattice mass gap in both orders of limits. O1 is exactly DMP-completeness of the Yang-Mills box, in both directions.
\item \emph{GRH.} The explicit formula for all nontrivial zeros of zeta, by the contour shift, in Lean; Weil's form positive with an explicit margin on the first box, where only the real place and the primes 2 and 3 enter (certified in ball arithmetic); the positivity mechanisms in finite models.
\item \emph{NP vs coNP.} The crossing theorem for every sound two-premise proof system, and the equivalence of the lifting interface with lifting by constant-size gadgets.
\item \emph{Navier-Stokes.} The card, with Fefferman's four forms; in the energy-only model the box gives the wrong marker at some scale, so it is not DMP-complete.
\item \emph{BSD.} What Mathlib alone gives, over any number field: Mordell-Weil from weak Mordell-Weil, the true rank, the existence line, parity from the functional equation, point counts that run, and a finite order at s = 1 from Hasse's bound, with the classical theorems Mathlib lacks carried as named parts.
\item \emph{Hodge.} Hodge-Riemann positivity, the Hodge index as linear algebra, Weil's bound for curves, eigenvalues on the circle, and Lefschetz (1,1) as a diagram chase.
\item \emph{GU.} The Z\_6 quotient of SU(3) x SU(2) x U(1) inside SU(5), charge quantisation from a compact U(1), and anomaly cancellation.
\item \emph{Validation.} Poincare for SU(N): a manifold homotopy equivalent to SU(N) is SU(N) for N \textless{}= 3 (published for SU(3)) and need not be from N = 4 on; Sp(2) is rigid, and G2 has exactly two structures. Rigidity is a one-box test too, and G2 is rigid under the one assumption.
\end{itemize}
\textbf{What is not proved.} No problem is solved here. One open step per card:

\begin{itemize}
\item \emph{Yang-Mills:} O1, the hand-over inequality for every outside, uniformly in the scale, with the lower edge (a correlation that does not vanish) in the same assumption; for the Clay statement also the continuum construction and every compact simple group.
\item \emph{GRH:} Weil positivity beyond the first boxes, a margin that survives the growth of the support. RH and GRH for Dirichlet L-functions were proved in Lean by others in November 2026, by a route that does not use positivity; GRH beyond Dirichlet L-functions is open.
\item \emph{NP vs coNP:} lifting with constant-size gadgets, and with it super-polynomial lower bounds for strong proof systems.
\item \emph{Navier-Stokes:} regularity or blowup of the equations themselves. DMP's reading predicts blowup (the equations are not DMP-complete), so global regularity would refute it. Two posted blowup claims by others are unrefereed and unchecked here; one rests on a certificate that is announced, not released.
\item \emph{BSD:} the rank equals the order of the zero at s = 1, from order 2 on.
\item \emph{Hodge:} codimension p \textgreater{}= 2 in general.
\item \emph{GU:} the checkable consequence taken for the gate is uniform charge quantisation, which any compact U(1) also gives; a consequence beyond it is open, and the hypercharges are an input.
\end{itemize}
The paper stays a draft until the assumption is proved: each open step is stated in the exact form that would close it, so that closing one changes one row of the claims ledger, one verdict line and one item of this abstract.

\medskip\noindent\rule{\linewidth}{0.4pt}\medskip

\subsection*{Introduction}
\addcontentsline{toc}{subsection}{Introduction}

\subsubsection*{What this paper is}
\addcontentsline{toc}{subsubsection}{What this paper is}

This is a paper about one machine and one test, used seven times. The form under them is old: \textbf{positivity is existence}. A positive weight makes a measure; a positive form makes a Hilbert space; a positive operator has its spectrum where it can be seen. Osterwalder and Schrader built quantum field theories from reflection positivity. Weil reduced the Riemann hypothesis to the positivity of one pairing. Hodge and Riemann found the same positivity on the cohomology of a variety, and Weil used it to prove the Riemann hypothesis for curves over finite fields. Dobrushin showed that a contraction of local influences forces a unique equilibrium, reached fast.

What is new is the shape put around it. A running system pays a cost at every boundary; with a floor its settling is decidable within an explicit bound, and without one it is not. A box with one command at one shared gate is tested step by step, and it passes when every marker agrees with the gate. Each problem becomes such a box, with exactly one named open step, and the open step is the same on every card: that the box is a running system whose gate is shared at every scale. The paper proves the machine and each box's core, and assumes that one step, once.

\subsubsection*{What a reader can take from it}
\addcontentsline{toc}{subsubsection}{What a reader can take from it}

\begin{itemize}
\item From \textbf{mathematical physics}: the Yang-Mills gap reduced, for SU(2), to one inequality for one box of fixed physical size, proved equivalent to the box passing the test, with the rest of the mechanism machine-checked; the Navier-Stokes card with the status of the posted blowup claims stated exactly.
\item From \textbf{number theory}: the explicit formula for all zeros in Lean, the first box of Weil's positivity certified with its margin, and why that margin cannot be carried outward by computation; a Mathlib-only BSD core over number fields, in IM.
\item From \textbf{algebraic geometry}: the Hodge index and Weil's bound as linear algebra in Lean, and the Hodge-index route to the Riemann hypothesis over finite fields read as the more efficiently correct one for this structure, beside Deligne's complete proof.
\item From \textbf{topology}: Poincare for SU(N) settled (rigid exactly for N \textless{}= 3), Sp(2) rigid, and the structure set of G2 counted.
\item From \textbf{complexity theory}: the crossing theorem for every two-premise proof system, and the open step of NP vs coNP reduced to one lifting statement.
\item From \textbf{logic and verification}: a machine in which every open part is a named hypothesis, every axiom shows in the kernel's report, and a lint and a second kernel check every change; undecidability and the floor as the two sides of one test.
\end{itemize}
\subsubsection*{A guide by ear: the tempered piano}
\addcontentsline{toc}{subsubsection}{A guide by ear: the tempered piano}

This guide is not part of any argument. It is offered as a picture to hold while reading, and the idea of hearing existence through tuning came to the first author from a short published by Vsauce; it is credited here, under that name, for that transfer.

\begin{itemize}
\item \textbf{Tuning is built from 2 and 3.} The octave is the ratio 2, the fifth the ratio 3/2. Pitch, taken modulo the octave, is a circle, and a fifth is a rotation on it \textbf{[classical]}.
\item \textbf{The rotation does not close by itself.} Because log2(3) is irrational, twelve pure fifths overshoot seven octaves by the Pythagorean comma, 3\textasciicircum{}12 / 2\textasciicircum{}19 = 1.01364 \textbf{[classical]}. This is the finite loss at the start that the paper excises: the first box at \{infinity, 2, 3\} (Part V) is where the two generators of the tuning meet.
\item \textbf{Temperament is the balanced state.} Equal temperament spreads the comma evenly over all twelve steps (each fifth becomes 2\textasciicircum{}(7/12)), so the rotation closes and every key is equivalent; twelve is not arbitrary, since 19/12 is a convergent of log2(3) \textbf{[classical]}. In the paper's words, the rotation gives the tuning a direction and makes every tick alike (Part III, the tick operator) \textbf{[reading]}.
\item \textbf{The zeros are the strings.} The explicit formula (Part V) writes the primes as a sum of tones: each nontrivial zero 1/2 + i gamma sounds at frequency gamma in log x. On the critical line every tone has the same amplitude, x\textasciicircum{}(1/2); the Riemann hypothesis says that every string is tuned to the same loudness \textbf{[classical]}.
\end{itemize}
The guide ends where the tuning does: on the final tone, resonating outward. Nothing more is stated about it than that such a configuration is describable, and that in that description it is DMP-complete \textbf{[reading]}.

\subsubsection*{How to read the status marks}
\addcontentsline{toc}{subsubsection}{How to read the status marks}

Every statement carries one mark. \textbf{[Lean]}: machine-checked here, standard axioms only, module named. \textbf{[classical]}: known, with a reference. \textbf{[external]}: a result by others, read but not rebuilt. \textbf{[conditional]}: under a named assumption. \textbf{[paper]}: argued here in prose. \textbf{[numerical]}: measured here; it can refute, never prove a for-all statement. \textbf{[certified]}: enclosed in ball arithmetic, with its log. \textbf{[reading]}: a DMP interpretation, never mathematics. \textbf{[open]}: not settled. The claims ledger (\texttt{../CLAIMS.md}) lists every claim with its mark, source and what would upgrade it.

\subsubsection*{The three verification layers}
\addcontentsline{toc}{subsubsection}{The three verification layers}

The work behind this paper is checked three ways: by proof (on paper and in Lean), by the literature, and by computation. The computations can refute a statement or guide a proof; they never establish a statement about all cases. Every measured number in the paper is tied to the statement it would discharge, and a passing measurement is reported as consistent, not as proof. The Lean is re-checked by independent kernels (Part II), and IM's arsenal passes its lint at every change, listing every part left open by name.

\subsubsection*{What is new}
\addcontentsline{toc}{subsubsection}{What is new}

\begin{itemize}
\item The machine as one object: running systems, their ledger, the encoder, and the settle theorem with its edge, decidable within an explicit bound under a cost floor and undecidable without one \textbf{[Lean]}.
\item The one-box test: one command at one shared gate, a marker per step; the one-box and scale certificates; O1 exactly as DMP-completeness of the Yang-Mills box; rigidity as a one-box test \textbf{[Lean]}.
\item The form layer as one API: every problem a card with exactly one named open slot, and R136a1 as the named answer ``outside'' \textbf{[Lean]}.
\item The explicit formula for all nontrivial zeros of zeta by the contour shift, on a Mathlib-only core \textbf{[Lean]}, and the first box of Weil's positivity with an explicit margin \textbf{[certified]}.
\item The hand-over bound, the strict twist bound for every SU(N), and zero join loss in slab form \textbf{[Lean]}.
\item A Mathlib-only BSD core over number fields, published in IM \textbf{[Lean]}.
\item One geometry, two edges: GRH and O1 proved to ask opposite, independent bounds of one distance \textbf{[Lean]}.
\item Poincare for SU(N), Sp(2) rigid, and G2's two structures, by surgery from classical theorems \textbf{[paper]}, with SU(3) from the literature \textbf{[external]}.
\end{itemize}
What is \emph{not} new is listed with references at each use: Osterwalder-Schrader, Weil, Hodge-Riemann, Lefschetz, Deligne, Dobrushin, Lee-Yang, Pringsheim and Landau, Mordell-Weil, Hasse, Gross-Zagier-Kolyvagin, Cook-Reckhow, the barrier theorems, Perelman, and the surgery theory of Browder, Novikov, Sullivan, Wall and Kirby-Siebenmann.

\subsubsection*{Plan}
\addcontentsline{toc}{subsubsection}{Plan}

Part I: DMP, the vocabulary. Part II: the machine, then the tools the problems use. Part III: the one assumption and the rigid start. Parts IV to X: the seven problems, one template each (Yang-Mills, GRH, P vs NP, Navier-Stokes, BSD, Hodge, GU). Part XI: Outlook. Appendix A: validation (Poincare, SU(N) and the rigid groups).


%% ---- 01-part-I-dmp.md
\section*{Part I. DMP: a language for existence}
\addcontentsline{toc}{section}{Part I. DMP: a language for existence}

\subsection*{I.0 How to read this part}
\addcontentsline{toc}{subsection}{I.0 How to read this part}

DMP, the Dark Matter Protocol, is a specification for systems that must trust each other. It was written for software and governance: who may act, on what evidence, and what is left behind. Its author reads the same rules as a description of how anything comes to exist. This paper uses DMP in exactly that way: as a language. It gives names to the steps that the three problems of Parts IV to VI have in common, so that each step is said once.

Two rules keep this honest.

\begin{itemize}
\item \textbf{A reading is not a theorem.} Where this part says ``existence is a positive weight'', that is a reading: a choice of which mathematical object a DMP word stands for. Readings are marked \textbf{[reading]}. They guide; they prove nothing.
\item \textbf{Every reading points forward.} Each section ends with the precise statement that Part II makes of it, with its status there. Part I can be skipped by a reader who wants only the mathematics. It cannot be used by a reader who wants a proof.
\end{itemize}
The marks used in the whole paper:

{\footnotesize
\begin{longtable}{>{\raggedright\arraybackslash}p{\dimexpr 0.2910\linewidth-2\tabcolsep\relax}>{\raggedright\arraybackslash}p{\dimexpr 0.6890\linewidth-2\tabcolsep\relax}}
\toprule
\textbf{Mark} & \textbf{Meaning} \\
\midrule
\endhead
\textbf{[Lean]} & Machine-checked in this project (Lean 4 with Mathlib, standard axioms only). The module path is given relative to \texttt{docs/mass-gap/formal/ProofPhase/}. \\
\textbf{[classical]} & Known mathematics or physics, with a reference. \\
\textbf{[external]} & A published or posted result by others that this project reads but did not rebuild or re-check. \\
\textbf{[paper]} & Argued in this project in prose, not machine-checked. \\
\textbf{[numerical]} & Measured in this project. Evidence; it can refute, never prove a statement about all cases. \\
\textbf{[reading]} & A DMP reading: an interpretation, never mathematics. \\
\textbf{[open]} & Not settled by anyone. \\
\bottomrule
\end{longtable}}

\medskip\noindent\rule{\linewidth}{0.4pt}\medskip

\subsection*{I.1 Existence has two halves}
\addcontentsline{toc}{subsection}{I.1 Existence has two halves}

\textbf{DMP.} The smallest unit of DMP is a \emph{part}, and a part is a Duality: two halves, one that owns \emph{trust} and one that owns \emph{information}. A part acts only while both halves certify it. Nothing else is said about what a part is made of. One rule sets the tone for everything after it: from trust nothing can be derived about information except that the information exists.

\textbf{Reading (I.1).} Read trust as \emph{that something is} and information as \emph{what it is}. Then ``existence'' is the trust half alone: a yes or no, with no content. In mathematics the cleanest object of this kind is a \textbf{positive weight}: a number that says a configuration counts, and says nothing about what the configuration looks like. A system exists when its total weight is positive. \textbf{[reading]}

\textbf{Why it is a good reading.} In all three problems the first step is exactly such a statement:

{\footnotesize
\begin{longtable}{>{\raggedright\arraybackslash}p{\dimexpr 0.2103\linewidth-2\tabcolsep\relax}>{\raggedright\arraybackslash}p{\dimexpr 0.7697\linewidth-2\tabcolsep\relax}}
\toprule
\textbf{} & \textbf{``Something exists''} \\
\midrule
\endhead
Yang-Mills & the lattice weight e\textasciicircum{}\{-S\} is positive, so the partition function Z is positive \\
GRH & a pairing of test functions (Weil's explicit-formula pairing) is non-negative \\
P vs NP & a formula is satisfiable iff its count of satisfying assignments is positive (SAT iff \#SAT \textgreater{} 0) \\
\bottomrule
\end{longtable}}

\textbf{What Part II makes of it.} Positivity is the first link of the common chain (II.9). It is the input of every construction that follows, and it is the step that is hard in two of the three problems.

\medskip\noindent\rule{\linewidth}{0.4pt}\medskip

\subsection*{I.2 G0, the ground}
\addcontentsline{toc}{subsection}{I.2 G0, the ground}

\textbf{DMP.} The rules of DMP are checked by two independent witnesses. But the first two witnesses cannot be checked by anyone: somebody has to trust them first. DMP calls that first act the \emph{minting} and its source \textbf{G0}, the ground. G0 is ``stated, not derived''. Every later standing descends from it. A system that cannot be traced back to a G0 is not part of the system.

\textbf{Reading (I.2).} G0 is the definition of the system: the action and the gauge group for Yang-Mills; the explicit formula and its test functions for GRH; the axioms and the proof rules for P vs NP. Everything is derived from G0, and G0 itself is not. In particular, the positivity of I.1 is either built into G0 by construction, or it has to be proved. \textbf{[reading]}

\textbf{Why it matters for the paper.} This is the single sharpest difference between the three problems:

\begin{itemize}
\item In Yang-Mills positivity is part of G0. The weight e\textasciicircum{}\{-S\} is positive and local, so reflection positivity follows by construction (II.9).
\item In GRH positivity is \emph{not} part of any known G0. Weil's pairing is non-negative exactly when GRH holds. Building a G0 that makes it non-negative by construction is the open step (Part V).
\item In P vs NP, positivity of the count is the question itself: whether the formula has a satisfying assignment.
\end{itemize}
\textbf{What Part II makes of it.} None of the tools of Part II creates positivity. They all start from it (II.16).

\medskip\noindent\rule{\linewidth}{0.4pt}\medskip

\subsection*{I.3 The existence boundary (DB)}
\addcontentsline{toc}{subsection}{I.3 The existence boundary (DB)}

\textbf{DMP.} A boundary \emph{contains} information when no part beyond it is trusted for any act on that information. Nothing crosses a containing boundary. In the physical reading the author calls such a boundary an \textbf{existence boundary}, a DB. Three statements about a DB carry most of the weight of this paper:

\begin{enumerate}
\item Only existence crosses a DB; information that would cross is lost.
\item Every layer has its own Duality, so the inside of a DB is a system of its own.
\item Anything that crosses a DB is re-attached to a new G0, and ``anything not compatible with the encoder in the DB is lost'' (the author, 2026-11-23).
\end{enumerate}
\textbf{Reading (I.3).} A DB is a box together with its outside. ``Only existence crosses'' says the box sees its outside only through a boundary condition. ``What the encoder cannot carry is lost'' says the passage across the boundary is a \textbf{lossy channel}: a re-encoding into a coarser description, which can only forget. Every renormalisation step, every conditioning on an outside, and every transfer step of a Markov chain is such a channel. \textbf{[reading]}

\textbf{What Part II makes of it.} Data processing (II.12): a channel never increases the distance between two states, and a channel that loses something contracts it by a factor eta \textless{} 1. The same channel applied n times contracts by eta\textasciicircum{}n. That one inequality is the engine of the Yang-Mills reduction and of the sampling side of P vs NP.

\medskip\noindent\rule{\linewidth}{0.4pt}\medskip

\subsection*{I.4 Acts and remnants}
\addcontentsline{toc}{subsection}{I.4 Acts and remnants}

\textbf{DMP.} An \emph{act} is anything done: a message, a read, a write, a decision. Every act leaves a \emph{record}, written at the act, readable in full by every witness of it, and outliving the act. What an act leaves behind is its \textbf{remnant}.

The author adds two statements about remnants that this paper uses.

\begin{itemize}
\item \textbf{A contradiction is excised at once.} A contradictory statement is neither existence nor non-existence; it is a disagreement at G0. It is removed immediately, and the only remnant is the statement ``it does not exist''.
\item \textbf{The act carrier loses nothing.} The author's example is the photon: it exists only between the event that emits it and the event that absorbs it, it carries one quantum, and when it delivers its act nothing is left over.
\end{itemize}
\textbf{Reading (I.4).} The first statement is proof by contradiction, read as a process. Assuming A and deriving a contradiction leaves one remnant: not-A. In propositional logic the remnant is literal: a refutation, a derivation of the empty clause, is the certificate that no assignment satisfies the formula. The second statement is a \textbf{conservation law}. A certificate passes from one act to the next without loss exactly when it is carried by a conserved quantity. \textbf{[reading]}

\textbf{What Part II makes of it.}

\begin{itemize}
\item The remnant of a contradiction is a refutation, and refutations always exist (refutation completeness of resolution) \textbf{[Lean]} (\texttt{Proof/Complete.lean}, \texttt{refutable\_iff}).
\item The carrier is a conserved flux. In both Yang-Mills (the centre flux) and propositional proof (the Tseitin parity) a Gauss law holds mod 2: the flux through the boundary of a region equals the charge inside, whatever happens inside \textbf{[Lean]} (\texttt{Proof/Crossing.lean}, \texttt{gauss}; \texttt{Crossover/FluxRate.lean}, \texttt{exists\_sheet\_of\_flux}). Conserved flux cannot be screened. This is the second branch of the dichotomy (II.14).
\end{itemize}
I.9 sharpens the first statement of this section: excision happens at transmission, and only the remnant tail is kept.

\medskip\noindent\rule{\linewidth}{0.4pt}\medskip

\subsection*{I.5 The ODP: pressing, then minting}
\addcontentsline{toc}{subsection}{I.5 The ODP: pressing, then minting}

\textbf{DMP.} The \textbf{ODP}, the Oscillon Degradation Protocol, is the author's account of how a G0 finds its rules. G0 knows only how to exist; its first act must be to find out how to act. So the ``perfect duality'' oscillates and decays until it reaches a stable definition of existence, ``because there MUST be a right one, because it exists'' (the author, 2026-10-28).

In the protocol this happens in two phases.

\begin{itemize}
\item \textbf{Pressing.} A configuration is run as a rehearsal until its definition of existence is complete and its configuration hash matches. A pressed system is ``out of spec'', never ``proven'': the pressing is a rehearsal of the real rule, not a weaker form of it.
\item \textbf{Minting.} Only a configuration that has passed its pressing is minted, and minting goes together with the \textbf{WHKs}, the White Hole Keys: ``the entire content of creation, from which each part can derive the other'' (the author, 2026-11-23). In DMP these are the naked duality, the quantum and the tick, together with the gap they establish.
\end{itemize}
\textbf{Reading (I.5).} Pressing is the scale-by-scale construction of a theory until one finite check passes; minting is the mechanical production of everything else from that check. \textbf{[reading]} The reading has an exact counterpart in Yang-Mills, where the two phases are separated in the proof:

\begin{itemize}
\item pressing: renormalise from the lattice spacing up to one physical scale, until a finite, local contraction condition holds in one box (the one-box statement, Part IV) \textbf{[open]};
\item minting: from that certificate, build the Hilbert space, the transfer operator and its gap, and take the limits \textbf{[Lean]} (\texttt{OS/GapLoc.lean}, \texttt{Assembly.lean}; conditional on the stated open inputs, Part IV).
\end{itemize}
The WHKs also have an exact counterpart. Pure Yang-Mills has no free dimensionless parameter. The bare coupling and the lattice spacing combine into one scale, Lambda (dimensional transmutation), and every physical quantity is a pure number fixed by the group, times a power of Lambda. From the group and any one of \{the tick with its coupling, Lambda, the gap\}, all the others follow. \textbf{[classical]}

\textbf{What Part II makes of it.} The split ``one finite certificate, then mechanical production'' is the shape of every theorem in II.12: a column sum below one at one scale implies decay at every larger scale.

\medskip\noindent\rule{\linewidth}{0.4pt}\medskip

\subsection*{I.6 The quantum of assertion and the tick}
\addcontentsline{toc}{subsection}{I.6 The quantum of assertion and the tick}

\textbf{DMP.} An act is never infinitely small. The author: ``The quantum of assertion is given along with the tick frequency at G0, so the rate is already part of any well constructed box.'' The quantum and the tick also set a floor: ``The loss associated with that smallest act is a universal error bar.''

\textbf{Reading (I.6).} Two precise statements stand behind this. \textbf{[reading]}

\begin{enumerate}
\item \textbf{A quantized measure arrives; a real one may not.} A process that lowers a natural number at every step stops after finitely many steps. A process that lowers a real number at every step may approach a limit forever. \textbf{[Lean]} (\texttt{Native/Method.lean}, \texttt{fixed\_after} and \texttt{real\_never\_arrives}).
\item \textbf{A finite box has a rate per tick.} In a finite box with a compact group and a bounded action, the law of the box given any outside has a density bounded below, so one step of the dynamics contracts by a definite factor, whatever the outside (Doeblin's condition). The contraction theorem is \textbf{[Lean]} for abstract kernels (\texttt{Mixing/Doeblin.lean}, \texttt{contracts\_of\_doeblin}); the lower bound for the Wilson box is \textbf{[paper]} (\texttt{PREP.md}, ``Three answers'', (a)). What the quantum does \emph{not} give by itself is control of the sum over many neighbours, nor survival of the rate as the tick shrinks. That is exactly what is open in Yang-Mills (Part IV, and II.12 for the precise limit).
\end{enumerate}
The universal floor became a project rule for measurements: every error bar is the larger of the statistical floor and the resolution floor 1/T, with T the time extent actually observed (the time-energy resolution). \textbf{[paper]} (rule adopted 2026-11-23; \texttt{docs/mass-gap/weak-coupling/proof-phase/c9a-crossover/ONEBOX-CHECK.md}, section 7).

\medskip\noindent\rule{\linewidth}{0.4pt}\medskip

\subsection*{I.7 Existence = non-existence + 1}
\addcontentsline{toc}{subsection}{I.7 Existence = non-existence + 1}

\textbf{DMP.} The author's shortest statement: existence is at least non-existence plus one, ``given by the simple fact that existence has to also describe what does not exist, in addition to everything else'' (2026-11-23). A boundary case belongs to it: with no coupling at all there is nothing to assert, and the ``+1'' disappears.

\textbf{Reading (I.7).} Three exact statements share this shape. \textbf{[reading]} for the identification; the statements themselves are:

\begin{itemize}
\item \textbf{Set theory.} A description of all parts of a set includes the empty part, and Cantor's theorem makes it strictly larger than the set. In von Neumann's construction 0 = \{\} and 1 = \{0\}: the first thing that is not nothing is ``a set containing nothing''. \textbf{[classical]}
\item \textbf{The strict flip bound.} For the SU(2) Wilson action at every coupling beta != 0, every box size and every twist position, the twisted partition function is strictly smaller than the untwisted one: Z\_tw \textless{} Z \textbf{[Lean]} (\texttt{Flip/Strict.lean}, \texttt{su2\_Ztw\_lt\_Z}). At beta = 0 they are equal, matching the boundary case.
\item \textbf{A strict sign gives a rate at the self-dual point.} If 0 \textless{} Z\_tw \textless{} Z, then the twist free energy or the electric free energy is at least log(1 + sqrt 2), about 0.881 \textbf{[Lean]} (\texttt{Crossover/HandOver.lean}, \texttt{log\_one\_add\_sqrt\_two\_le\_max}). A strict inequality with no rate still forces a definite rate in one of two dual channels. This is the clearest case in the paper of a DMP statement turning into a number (II.13).
\end{itemize}
A last reading belongs here because Part IV uses it: \textbf{``mass gap = the arrow of causality exists''} (the author, 2026-11-22). The transfer matrix of a reflection-positive theory is a positive contraction. A gap means that, after one unit of time 1/m, every correlation with the past has shrunk by 1/e: a built-in clock. Without a gap correlations still fade, but only as a power law, with no intrinsic tick. \textbf{[reading]} on the classical spectral fact of II.10.

\medskip\noindent\rule{\linewidth}{0.4pt}\medskip

\subsection*{I.8 DMP as an existence linter}
\addcontentsline{toc}{subsection}{I.8 DMP as an existence linter}

\textbf{DMP.} The author's one-line description of DMP: it is a very well defined \emph{existence linter}. A linter checks a text against a rule set and says nothing about what the text means. The author reads Gravity, the repository that holds the machine-checked evidence, as the linter, and DMP as the universally readable rule set that is its passing condition. Gravity and DMP sit at one node of the root chain. Each keeps the other honest: neither can change without the other's check seeing it.

\textbf{Reading (I.8).} Existence, in the sense of I.1, is what the linter passes. A statement exists in the system when it is constructed, by the kernel, from DMP's pinned sentences and rests on nothing else. Whatever cannot be so constructed is not part of the system (I.2). \textbf{[reading]} (claim I.8a)

\textbf{The exact counterpart.}

\begin{itemize}
\item \textbf{Gravity's check passes only when} the pinned copy of DMP matches its recorded hash; the repository needs nothing that its one setup file does not list; no source file uses an escape hatch (\texttt{sorry}, \texttt{admit}, \texttt{axiom}, \texttt{native\_decide}, \texttt{unsafe}, \texttt{implemented\_by}, \texttt{extern}); the Lean build succeeds; and the kernel's axiom report for every theorem lists only Lean's three standard axioms. \textbf{[paper]} (read from \texttt{evidence/tools/check.sh} and \texttt{tools/check-pin.sh}; the check was not run for this draft) (claim I.8b)
\item \textbf{The arsenal has the same shape at small size} (II.6): its lint is the linter, and ``Mathlib v4.34.0 only, no escape hatch, an open part only as a named R136a1 hypothesis'' is its passing condition. \textbf{[paper]} (claim I.8c)
\item \textbf{The root chain pins in both directions.} Gravity holds a pinned copy of DMP (\texttt{vendor/dmp/PIN}, hash checked by \texttt{tools/check-pin.sh}). That DMP pins Gravity in turn is the author's statement and the notes'; it was not re-read for this draft. \textbf{[paper]} (claim I.8d)
\end{itemize}
\textbf{Three care points.}

\begin{enumerate}
\item \textbf{A linter certifies form, not content.} A passing result says that every statement is soundly constructed from what DMP states. It does not say that DMP's sentences describe the world. \textbf{[paper]} (claim I.8e)
\item \textbf{A linter cannot certify its own passing condition.} A consistent theory strong enough for arithmetic cannot prove its own consistency \textbf{[classical]} (Goedel II; claim I.8h). That is why the two nodes must keep each other honest. It is also why one assumption is left that neither can lint: that the system runs. This is the assumption Part III calls O1 in each card's terms. \textbf{[reading]} for the identification (claim I.8f)
\item \textbf{``Universally readable'' holds for the form.} The evidence is plain text and Lean, checkable by anyone with the pinned toolchain. ``Universal superset'' is a reading. \textbf{[reading]} (claim I.8g)
\end{enumerate}
\textbf{What Part II makes of it.} The linter's rule set is the form layer (II.1 to II.5) plus the arsenal's gate (II.6). The unlintable assumption is the open step of each card (Part III, section 4).

\textbf{The one-box test.} The test itself, a box with one command at one shared gate and a marker per step, is Part II's one-box test (II.4).

\medskip\noindent\rule{\linewidth}{0.4pt}\medskip

\subsection*{I.9 Excision at transmission}
\addcontentsline{toc}{subsection}{I.9 Excision at transmission}

\textbf{DMP.} A part answered \textbf{R136a1} is outside the system, whatever it is. Nothing else needs to be proved about it. In DMP this is always a DMP-complete answer: it is the other option an act collapses back on.

\textbf{Reading (I.9).} I.4 said that a contradiction is excised at once and leaves only ``it does not exist''. The author sharpens the statement, and the operative word is \emph{transmitted}. The claim is about what crosses the gate, not about what a sender may compose. A wrong act never crosses: at the gate it is answered R136a1, so it is outside. What stays inside is only its \textbf{remnant tail}: the record that an excision happened. IM's verifier counts every refusal without content, as the notes state (DMP chapter 06), so the count is the assertion of the excision and carries none of the excised information. \textbf{[reading]} (claim I.9a)

\textbf{The exact counterpart.} In the arsenal the same move is \texttt{Arsenal.R136a1}: the part is named and stays outside every proof, and only the name remains. It is the statement taken as a named hypothesis, never an axiom, so every declaration still rests on the three standard axioms. \textbf{[Lean]} (\texttt{im} repository, \texttt{arsenal/Arsenal/R136a1.lean}, \texttt{R136a1}, \texttt{R136a1.out}) (claim I.9b)

\textbf{Care point.} The gate excises what its checks can recognise as wrong at that tick. An act whose wrongness becomes recognisable only later is excised at the later tick, with its tail recorded then. A later excision is a new act, never a revision of the earlier tick. \textbf{[reading]} (claim I.9c). Part III, section 5, gives the logical side: excision decided per tick, and why the Clay statements cannot be excised as a class.

\medskip\noindent\rule{\linewidth}{0.4pt}\medskip

\subsection*{I.10 Rights tied to identity}
\addcontentsline{toc}{subsection}{I.10 Rights tied to identity}

\textbf{DMP.} The author's rule for security: tie rights to identity. If an act is allowed only when it is permissible, and no wrong information can be transmitted (I.9), then the identity has total freedom \emph{inside} and is secure in its existence.

\textbf{Reading (I.10).} ``Total freedom inside, secure at the gate'' is capability security, as IM builds it, and it has one more reading. Asymptotic freedom in Yang-Mills is free at short distance, where the coupling vanishes in the ultraviolet \textbf{[classical]} (Gross-Wilczek 1973; Politzer 1973, recalled; claim I.10g), and bound at large distance, where the gap holds. Reading the two together: inside the box the dynamics is free, and the boundary is where it is bound. \textbf{[reading]} (claim I.10a)

\textbf{The exact counterpart.}

\begin{itemize}
\item \textbf{An act passes only with both witnesses' grants.} A sentence passes the door only when both the security grant and the operational grant carry its kind; one grant alone never passes. \textbf{[Lean]} (\texttt{im} repository, \texttt{proof/Im/Laws.lean}, \texttt{Im.pass\_needs\_both\_grants}, \texttt{Im.one\_grant\_never\_passes}) (claim I.10b)
\item \textbf{Every sentence of IM's closed grammar keeps every law:} there is no sentence outside the list of 72 shapes, and each shape is lawful. \textbf{[Lean]} (\texttt{proof/Im/Closure.lean}, \texttt{Im.closure}, \texttt{Im.space\_size}) (claim I.10c)
\item \textbf{Behind the door a composition stays lawful.} A network built from nodes that keep the guarantee on what their door lets in keeps it on any input. \textbf{[Lean]} (\texttt{proof/Im/Compose.lean}, \texttt{Im.sound}, \texttt{Im.galaxy\_robust}) (claim I.10d)
\item \textbf{Tokens only attenuate.} An attenuated token's kinds are a subset of its parent's, and a chain of attenuations is the intersection of the root with its caveats, in any order. It holds by construction: attenuation is intersection with a caveat. \textbf{[Lean]} (IM, \texttt{proof/Im/Attenuate.lean}: \texttt{attenuate\_sub}, \texttt{attenuateAll\_eq}, \texttt{attenuateAll\_perm}; claim I.10e)
\end{itemize}
\textbf{Care points.} (1) ``Permissible'' is what the grammar and the grants state, so the security is exactly as good as that specification. (2) The gate certifies integrity and authority, not the truth of content (I.9 says what happens to what it recognises as wrong). (3) The binding of rights to identity is as strong as the custody of the identity's key. \textbf{[reading]} (claim I.10f)

\medskip\noindent\rule{\linewidth}{0.4pt}\medskip

\subsection*{I.11 The arsenal as shared vocabulary}
\addcontentsline{toc}{subsection}{I.11 The arsenal as shared vocabulary}

\textbf{The rule (the owner's).} What is content-free and used by more than one project is kept once, and used from there.

\textbf{Reading (I.11).} The arsenal in the \texttt{im} repository (\texttt{arsenal/}) is that shared place for results that need Mathlib. It holds words that more than one part of this paper uses: \texttt{R136a1} for ``outside'' (I.9), and the pieces that are not about any one problem. A project that needs a word from it \textbf{imports} the word. It never copies it. If a result could be moved there, it moves, and the project uses it from there as a dependency. \textbf{[paper]} (a project rule, and how it is done today: \texttt{arsenal/README.md}; claim I.11a). The tool that guards it, and what it checks, are in II.6.

\textbf{A second reading of the letters.} IM is the Interstellar Medium, and that name stays. Its letters also read as \textbf{Identity Multiplication}, and the second reading is true of what the arsenal holds: the rotation API is multiplication by a number on the unit circle, which keeps the norm and moves only the phase; on the circle the carried mass never changes \textbf{[Lean]} (II.0c, \texttt{cmach\_lossless}), and the encoder's root is the same norm read from both sides \textbf{[Lean]} (II.2). The owner reads a second reading that carries a meaning as itself a uniqueness, one that survives being contained \textbf{[reading]} (claim I.11b). He also reads IM as a step up the DAG: the edge is real in the build (this project requires IM's arsenal, pinned by commit, and the package graph is acyclic). His argument for why the link is a rotation runs in three steps. A DB lets only existence cross, and the encoder's gate keeps the norm: both sides certify the same norm \textbf{[Lean]} (\texttt{Encoder.root}; any norm-preserving map gives an encoder) and \textbf{[reading]} for ``a DB allows only this gate'' (I.3). A linear gate that keeps every norm is multiplication by a unit number, a rotation \textbf{[Lean]} (II.0f). So the link to another node, which has to pass a DB, is a rotation: the ultimate gate in his words \textbf{[reading]} (claim I.11c). The cost does not sit in the gate. In the owner's words it sits in the provability of a common root, and that is exact: the part of the state that provably shares its root with the other side is the reflection coefficient c, and the remainder, \textbar{}\textbar{}psi - conj(c) psi'\textbar{}\textbar{}\textasciicircum{}2 = \textbar{}\textbar{}psi\textbar{}\textbar{}\textasciicircum{}2 (1 - \textbar{}c\textbar{}\textasciicircum{}2), is exactly what the step pays \textbf{[Lean]} (\texttt{szego\_remainder}). At \textbar{}c\textbar{} = 1 the root is fully shared and the crossing is free. On the reals the circle is just 1 and -1, so ordinary multiplication is a cost times an identity-keeping sign \textbf{[Lean]} (II.0g): ordinary multiplication is IM as well, in the owner's words \textbf{[reading]} (claim I.11d).

\medskip\noindent\rule{\linewidth}{0.4pt}\medskip

\subsection*{I.12 The vocabulary in one table}
\addcontentsline{toc}{subsection}{I.12 The vocabulary in one table}

{\footnotesize
\begin{longtable}{>{\raggedright\arraybackslash}p{\dimexpr 0.1604\linewidth-2\tabcolsep\relax}>{\raggedright\arraybackslash}p{\dimexpr 0.2269\linewidth-2\tabcolsep\relax}>{\raggedright\arraybackslash}p{\dimexpr 0.2192\linewidth-2\tabcolsep\relax}>{\raggedright\arraybackslash}p{\dimexpr 0.2569\linewidth-2\tabcolsep\relax}>{\raggedright\arraybackslash}p{\dimexpr 0.1167\linewidth-2\tabcolsep\relax}}
\toprule
\textbf{DMP word} & \textbf{Plain meaning} & \textbf{Reading} & \textbf{Precise counterpart} & \textbf{Part II} \\
\midrule
\endhead
Existence (trust half) & that it is & a positive weight & Z \textgreater{} 0; a non-negative pairing & II.9 \\
Information half & what it is & an observable & a vector in the space built from the weight & II.9 \\
G0 & the ground, stated not derived & the definition of the system & the action, the explicit formula, the axioms & II.16 \\
DB & a containing boundary & a box with its outside & conditioning on a boundary; a lossy channel & II.12 \\
Act, remnant & what is done, what it leaves & a step; its certificate & a transfer step; a refutation; a conserved flux & II.12, II.14 \\
Act carrier & passes an act without loss & a conserved quantity & Gauss law mod 2 & II.14 \\
ODP: pressing & rehearse until the hash matches & construct up to one finite check & a column sum below 1 at one scale & II.12 \\
ODP: minting & produce from the certificate & build everything else & OS reconstruction and limits & II.9, II.10 \\
WHKs & each part derives the others & no free parameter & dimensional transmutation & Part IV \\
Quantum, tick & smallest act, its rate & a rate per step & Doeblin; quantized descent arrives & II.12, II.15 \\
Existence = non-existence + 1 & always strictly more & a strict inequality & Z\_tw \textless{} Z; Cantor & II.13 \\
Contradiction excised & only ``it does not exist'' remains & refutation & refutation completeness & II.14 \\
Mass gap & the arrow of causality exists & a clock in the transfer step & spectral gap = decay rate & II.10, II.11 \\
Existence linter & checks form, says nothing of meaning & passing means soundly constructed from DMP & the pinned check; the arsenal's lint & II.6 \\
R136a1 & the answer ``outside'' & excision at transmission, only the remnant tail stays & \texttt{R136a1}, a named hypothesis, never an axiom & II.6 \\
Rights tied to identity & free inside, secure at the gate & capability security & both grants needed; closure; composition & Part III \\
Arsenal & shared words, imported never copied & one place for content-free results & \texttt{arsenal/} in the \texttt{im} repository & II.6 \\
\bottomrule
\end{longtable}}

Part II states each counterpart once, in its most general form.


%% ---- 02-part-II-machine.md
\section*{Part II. The machine, and the tools the problems use}
\addcontentsline{toc}{section}{Part II. The machine, and the tools the problems use}

\textbf{A. The machine.}

\subsection*{II.0 The machine in ordinary terms: multiplication by complex numbers}
\addcontentsline{toc}{subsection}{II.0 The machine in ordinary terms: multiplication by complex numbers}

\textbf{The translation.} In ordinary terms the machine is multiplication by complex numbers. Each step multiplies the carried amplitude by one number c with \textbar{}c\textbar{} \textless{}= 1: its modulus is the cost paid at the gate, its phase the rotation. A run multiplies its steps, so moduli multiply and phases add. If every \textbar{}c\textbar{} stays at or below some theta \textless{} 1, the run settles within an explicit number of steps. If every \textbar{}c\textbar{} = 1 (pure rotation), the number never changes the mass; the universal machine is there, and whether it settles cannot be decided. Everything else in this part (the ledger, the encoder, the settle theorem, the one-box test) is a detail of this one picture.

\textbf{In Lean} (\texttt{ETM/Complex.lean}, \texttt{complex\_translation}):

\begin{itemize}
\item a sequence of multipliers of modulus at most one is a running system whose ledger law holds with equality, and its run is the product of its steps (\texttt{cmach}, \texttt{cmach\_pays\_eq}, \texttt{cmach\_iterate}, \texttt{cmach\_mass\_eq}) \textbf{[Lean]} (claim II.0a);
\item every running system is dominated by the product of its own factors, read as complex numbers (\texttt{dominated}), so the translation loses nothing on the cost side \textbf{[Lean]} (claim II.0b);
\item below the circle the run settles (\texttt{cmach\_reaches}, from the settle theorem of II.3); on the circle the mass never changes (\texttt{cmach\_lossless}), and every factor of the universal machine is 1 (\texttt{etm\_fac\_one}) \textbf{[Lean]} (claim II.0c);
\item the encoder's reflection coefficients (II.2) are such multipliers, with cost 1 - \textbar{}c\textbar{}\textasciicircum{}2 (\texttt{encMach}) \textbf{[Lean]} (claim II.0d);
\item a link between two nodes keeps the step, the gate factor and the mass; multiplying every amplitude by a unit number is such a link, and a complex-linear map of the line that keeps every norm is multiplication by a unit number, so among linear links that keep identity only rotation links (\texttt{rotHom}, \texttt{rot\_of\_isometry}) \textbf{[Lean]} (claim II.0f);
\item on the reals the circle is the two points 1 and -1: a real-linear map of the line that keeps every norm is multiplication by 1 or -1, the identity or the half-turn, and every real number is its modulus times 1 or -1, so ordinary multiplication is a cost times an identity-keeping multiplication (\texttt{rot\_of\_isometry\_real}, \texttt{real\_eq\_sign\_mul\_abs}) \textbf{[Lean]} (claim II.0g).
\end{itemize}
\textbf{Decidable over every truly possible state.} The owner's reformulation: a DMP-complete state needs decidability for every truly possible state, with the least possible options given by y\textasciicircum{}k \textbf{[reading]} (claim II.0h). Each part has an exact counterpart. Decidability needs a floor (II.3) \textbf{[Lean]}. The truly possible paths after k steps grow like y\textasciicircum{}k, and the least such rate is the connective constant mu of self-avoiding walks \textbf{[classical]} (Hammersley-Morton). The run over all of them settles when the cost per step beats the options: the multiplied path sum mu\textasciicircum{}k e\textasciicircum{}(-F k) decays when F \textgreater{} log mu, and the joint form with the excess act gives both criteria \textbf{[Lean]} (model; \texttt{joint\_gives\_both}). For Yang-Mills that is O1, F\_e(L) - 2c(L) \textgreater{} log mu\_4 (Part IV). So the reformulation is the general form of the one assumption's decay edge, read in the language of options.

\textbf{Care point.} The translation is exact for the cost and the rotation. The content a step carries, the program the universal machine runs, lives in the state and not in the number; that is where its undecidability sits (II.3) \textbf{[paper]} (claim II.0e). In number theory ``complex multiplication'' names something else (elliptic curves with extra symmetries, near Part VIII); here it is always ``multiplication by complex numbers''.

\subsection*{II.1 The ledger and the transport theorem}
\addcontentsline{toc}{subsection}{II.1 The ledger and the transport theorem}

\textbf{Why it clicks.} A step that forgets a fixed share each time forgets geometrically. Write the share as a cost, and costs add. That is II.12 in bookkeeping form, and it carries an uncertainty as well as a certificate.

\textbf{The ledger (II.1.0).} A quantity M\_n \textgreater{}= 0 \emph{pays} cost c\_i at step i, with floor delta, when M\_\{i+1\} \textless{}= e\textasciicircum{}\{delta - c\_i\} M\_i. Then:

\begin{itemize}
\item M\_n \textless{}= e\textasciicircum{}\{n delta - (c\_0 + ... + c\_\{n-1\})\} M\_0, the ledger inequality;
\item if every cost is at least kappa, then M\_n decays like e\textasciicircum{}\{-n(kappa - delta)\}; if also kappa \textgreater{} delta, the sum of the M\_n converges (a \emph{cost above the floor}, \texttt{CostAboveFloor});
\item a lower bound on the final value is a lower bound on the total cost, hence on the size of the object that paid it.
\end{itemize}
\textbf{[Lean]} (\texttt{Unified/Ledger.lean}, \texttt{Pays}, \texttt{ledger}, \texttt{decay}, \texttt{summable\_of\_gap}, \texttt{size}, \texttt{CostAboveFloor}, \texttt{rate\_of\_costAboveFloor}; the abstract running-system form is \texttt{ETM/RunSys.lean}, \texttt{RunSys.carry}, \texttt{RunSys.contract}, \texttt{RunSys.decay\_cost}).

\textbf{Theorem II.1.1 (bounded transport).} An uncertainty e\_n \textgreater{}= 0 crosses step n with factor f\_n and floor delta, and the step may add fresh uncertainty r\_n \textgreater{}= 0 of its own: e\_\{n+1\} \textless{}= e\textasciicircum{}delta f\_n e\_n + r\_n. If the steps contract, e\textasciicircum{}delta f\_n \textless{}= q \textless{} 1, and each adds at most R, then e\_n \textless{}= q\textasciicircum{}n e\_0 + R/(1 - q) at every step, uniformly in n. \textbf{[Lean]} (\texttt{Unified/Transport.lean}, \texttt{bounded\_transport}, \texttt{transport\_geom}, \texttt{acc\_eq}).

\textbf{Theorem II.1.2 (critical transport).} At cost = floor, e\textasciicircum{}delta f\_n = 1, nothing is damped: e\_n \textless{}= e\_0 + (r\_0 + ... + r\_\{n-1\}). The bound is attained. It stays bounded exactly when the r\_j are summable, and a constant fresh uncertainty grows without limit. \textbf{[Lean]} (\texttt{critical\_transport}, \texttt{critical\_sharp}, \texttt{critical\_bounded\_iff}, \texttt{critical\_unbounded}).

\textbf{Theorem II.1.3 (sign preserved).} If a quantity has a sign margin m\_n and the uncertainty about its counterpart across the boundary stays below it, e\_n \textless{} m\_n, the counterpart's sign is certified at every step, with value at least m\_n - e\_n. \textbf{[Lean]} (\texttt{sign\_transport}, \texttt{sign\_of\_bounded}, \texttt{sign\_of\_critical}).

\textbf{II.1.4 (the ledger is the case r = 0).} With no fresh uncertainty the transport theorem is the ledger: \texttt{Pays}, \texttt{carry}, \texttt{Ledger.contract}, and every ledger of a running system. \textbf{[Lean]} (\texttt{pays\_iff\_transports}, \texttt{carry\_of\_transport}, \texttt{ETM.Ledger.transports}, \texttt{contract\_of\_transport}).

\textbf{II.1.5 (zero slack, exactly).} When an uncertainty and the margin it is compared with pay the same cost at every step, their ratio crosses at exactly cost = floor. Damping relative to a margin exists only if the margin shrinks more slowly than the uncertainty. \textbf{[Lean]} (\texttt{relative\_transport}, \texttt{relative\_critical}).

\textbf{The three applications.} Part III section 7 states what each problem's uncertainty is and where the open step sits. Yang-Mills: the boundary uncertainty of a slab. GRH: the explicit formula's remainder against the collapsing margin. NP vs coNP: the hiding error per hidden variable.

\medskip\noindent\rule{\linewidth}{0.4pt}\medskip

\subsection*{II.2 The encoder and its half-turn}
\addcontentsline{toc}{subsection}{II.2 The encoder and its half-turn}

\textbf{Why it clicks.} A boundary passes on one state, seen once from the inside and once, after one rotation, from the outside. Both sides certify the same size. All that is left to say about the step is how much the state overlaps with its rotated image, and one number carries that.

\textbf{Definition (II.2.1).} An \emph{encoder} is a state psi in a space with a positive form, together with its rotated image psi' = R psi of the same norm (the shared root). Its one number is the \emph{reflection coefficient} c = \textless{}psi, psi'\textgreater{} / \textbar{}\textbar{}psi\textbar{}\textbar{}\textasciicircum{}2. Positivity forces

{\footnotesize
\begin{longtable}{>{\raggedright\arraybackslash}p{\dimexpr 0.1467\linewidth-2\tabcolsep\relax}>{\raggedright\arraybackslash}p{\dimexpr 0.8333\linewidth-2\tabcolsep\relax}}
\toprule
\textbf{c} & \textbf{\textless{}= 1. \textbf{[Lean]} (\texttt{Unified/Encoder.lean}, \texttt{Encoder}, \texttt{Encoder.ofRot}, \texttt{refl}, \texttt{norm\_refl\_le}; reading the boundary from the other} \\
\midrule
\endhead
\bottomrule
\end{longtable}}

side conjugates c: \texttt{refl\_swap}).

\textbf{Two loss factors, determined by c alone (II.2.2).}

\begin{itemize}
\item The Szego factor 1 - \textbar{}c\textbar{}\textasciicircum{}2 is the share of psi left after removing its projection on the mirror image: the squared norm of psi - conj(c) psi' is \textbar{}\textbar{}psi\textbar{}\textbar{}\textasciicircum{}2 (1 - \textbar{}c\textbar{}\textasciicircum{}2). It is 1 exactly at c = 0 and positive exactly for \textbar{}c\textbar{} \textless{} 1.
\item The electric factor (1 - c)/(1 + c) is the odd part over the even part under the rotation: \textbar{}\textbar{}psi - psi'\textbar{}\textbar{}\textasciicircum{}2 / \textbar{}\textbar{}psi + psi'\textbar{}\textbar{}\textasciicircum{}2.
\end{itemize}
Both lie in [0, 1] for admissible c. The ledger cost is the minus log of the factor, and a chain of encoder factors is a ledger with floor 0. A uniform bound of the factors below e\textasciicircum{}\{-delta\} is exactly a cost above the floor. \textbf{[Lean]} (\texttt{szego}, \texttt{szego\_mem}, \texttt{szego\_pos\_iff}, \texttt{szego\_eq\_one\_iff}, \texttt{szego\_remainder}, \texttt{elec}, \texttt{elec\_mem}, \texttt{encoder\_elec}, \texttt{pays\_of\_factors}, \texttt{costAboveFloor\_of\_factor\_le}, \texttt{factor\_le\_of\_costAboveFloor}; the one-box form \texttt{oneBox\_iff\_refl}).

\textbf{One distance, two functions (II.2.3).} The cost of the electric factor is a hyperbolic distance d of c from the lossless point 0, and the Szego factor is the second function of the same d: it equals 4 e\textasciicircum{}\{-d\}/(1 + e\textasciicircum{}\{-d\})\textasciicircum{}2, that is 1/cosh\textasciicircum{}2(d/2). \textbf{[Lean]} (\texttt{Unified/GRHLens.lean}, \texttt{elec\_cost}, \texttt{szego\_of\_elec}). This is the object Part III section 2 puts at the centre. It is the distance of the encoder from the lossless point.

\textbf{The half-turn (II.2.4).} The electric factor is a Moebius map of the matrix (1/sqrt 2)[[-1, 1], [1, 1]]. That matrix is symmetric and squares to the identity, so the map is an involution, and its two fixed points, sqrt 2 - 1 and -(1 + sqrt 2), are antipodal. \textbf{[Lean]} (\texttt{halfTurn\_mobius}, \texttt{halfTurn\_transpose}, \texttt{halfTurn\_mul\_self}, \texttt{elec\_elec}, \texttt{elec\_selfdual}, \texttt{elec\_fixed\_neg}, \texttt{fixed\_antipodal}). As a Moebius map of the Riemann sphere it is a rotation by 180 degrees, the half-turn \textbf{[classical]}. It swaps the two ways of composing: the product of reflection coefficients is carried to hyperbolic addition of electric factors (\texttt{elec\_mul}, \texttt{elec\_oplus}) \textbf{[Lean]}. Its self-dual point has cost log(1 + sqrt 2), the 0.881 of II.13 (\texttt{selfdual\_cost}) \textbf{[Lean]}.

\textbf{What the half-turn does not do (II.2.5).} It is a symmetry of the description. It exchanges the pictures and fixes its axis, so it never decides on which side of the axis a box lies or how far from it. It does not turn an upper bound on one encoder into a lower bound on another \textbf{[paper]} (\texttt{docs/unified/GRH-LENS.md}, 6.3).

\textbf{The three applications (II.2.6).} Yang-Mills: the twist ratio Z\_tw/Z of the hand-over box is c, and O1 is c above (mu - 1)/(mu + 1) for every outside. GRH: the Verblunsky step of the zero measure at each scale is c. NP vs coNP: the crossing label passes at factor 1 with a cutting-planes line, so the encoder explains resolution's cost but supplies none for strong systems \textbf{[paper]} (\texttt{docs/unified/ENCODER.md}, 1).

\medskip\noindent\rule{\linewidth}{0.4pt}\medskip

\subsection*{II.3 The settle theorem and the decidability edge}
\addcontentsline{toc}{subsection}{II.3 The settle theorem and the decidability edge}

\textbf{Why it clicks.} If every step pays a cost above the floor, the quantity must fall below any error bar, and one can count how long that takes. If steps may pay exactly the floor, nothing forces it to fall, and nobody can tell in general whether it will.

\textbf{Running systems (II.3.1).} A \emph{running system} with floor delta has states, a transfer step, an encoder factor in [0, 1] at each state, a mass \textgreater{}= 0, and the ledger law mass(step x) \textless{}= e\textasciicircum{}delta fac(x) mass(x). The ledger theorems hold on every running system. Running systems form a category in which the space of ledgers is the terminal object, so every running system has exactly one map into it (its ledger). \textbf{[Lean]} (\texttt{ETM/RunSys.lean}, \texttt{RunSys}, \texttt{RunSys.carry}, \texttt{RunSys.contract}, \texttt{Hom}, \texttt{exists\_unique\_hom}, \texttt{univIsTerminal}). The statements XI.2c of the Outlook rest on this.

\textbf{Theorem II.3.2 (settle).} In the \emph{contraction class} (every factor at most theta, with e\textasciicircum{}delta theta \textless{} 1, which is exactly a cost above the floor), the run from any state reaches balance within epsilon, by an explicit computable number of steps: balanceBound M eps q = ceil(M q / (eps (1 - q))) + 1, for rational q with e\textasciicircum{}delta theta \textless{}= q \textless{} 1. So on this class the unbounded question ``does it settle?'' is a bounded one. \textbf{[Lean]} (\texttt{ETM/Machine.lean}, \texttt{Contracting}, \texttt{contracting\_iff\_costAboveFloor}, \texttt{balanceBound}, \texttt{balanceBound\_ge}, \texttt{reaches\_of\_contracting}, \texttt{reaches\_iff\_bounded}, \texttt{reaches\_of\_contracting'}). The same statement is XI.2a.

\textbf{Theorem II.3.3 (the edge).} The machine that runs every partial recursive content is lossless: factor 1, cost = floor, outside the contraction class. Whether it reaches balance is undecidable for all contents: no computable predicate decides it, and its truth value exists (classically) but cannot be computed. \textbf{[Lean]} (\texttt{etm}, \texttt{etm\_universal}, \texttt{etm\_not\_contracting}, \texttt{etm\_balance\_undecidable}, \texttt{etm\_truth\_edge}; both sides in one statement, \texttt{etm\_edge}). The same statement is XI.2b.

\textbf{What it says in one line.} With a floor, settling is decidable within an explicit bound. At the floor, it is not. The step from the second case to the first is the open step of every card (Part III, sections 4 and 5).

\textbf{The three applications.} Yang-Mills: the floor is the cost above log mu, and keeping it uniform in the cutoff is O1. GRH: the margin decays without a floor (zero slack, II.1.5), so only the sign is decided scale by scale. NP vs coNP: a constant cost per hidden variable is a floor.

\medskip\noindent\rule{\linewidth}{0.4pt}\medskip

\subsection*{II.4 The one-box test}
\addcontentsline{toc}{subsection}{II.4 The one-box test}

\textbf{The one-box test.} The author's sharpening: a box needs no outside, because every confinement has one shared gate at the DB. It needs one command, step forward, the universal tick, which can simulate everything. If the step produces the wrong R136a1 marker, the box is not DMP-complete; if the marker is correct, it is. In Lean (the rotation API, in IM's arsenal, \texttt{Arsenal/Rotation.lean}, with the running-system version in \texttt{ETM/Rotation.lean}): a box has one command, a gate factor, a margin and a marker per step, inside or R136a1; the marker is correct when it agrees with the gate's own factor, and DMP-complete means correct at every tick \textbf{[Lean]} (claim II.4a). The outside is read through the same gate by the half-turn: in the plaquette layer the flux read from a box and from its complement agree \textbf{[Lean]} (claim II.4b). Care point: the test has content only when the marker is the system's own claim; a marker read off the gate afterwards always passes \textbf{[Lean]} (claim II.4c). What the test gives for each card is in Part IV (IV.10c to IV.10e), Part VII (VII.7g) and Part III (4.4). The one-box certificate (one step at the shared gate certifies every tick), the scale certificate from the sharing, and the universal machine as a box are stated with the one assumption in Part III 4.4 (III.68 to III.68c) \textbf{[Lean]}. Rigidity is one-box testable too: a manifold's presentations with its self-equivalences acting make a box that is DMP-complete exactly when the manifold is rigid, and rigid from the shared gate (IM, \texttt{Arsenal.Rigidity}) \textbf{[Lean]} (claim II.4d; G2 in Appendix A.5).

\subsection*{II.5 The API: one card per problem}
\addcontentsline{toc}{subsection}{II.5 The API: one card per problem}

\textbf{Why it clicks.} The form layer says nothing about content, so every problem can use it. Every problem has the same shape: an input, a proved structure, one named open step.

\textbf{What the API gives (II.5.1).} One import, \texttt{ProofPhase.API}, a facade over the form layer: it moves and renames nothing, and its theorems are re-exports or one-line wrappers. It has one structure per concept \textbf{[Lean]} (\texttt{API.lean}):

{\footnotesize
\begin{longtable}{>{\raggedright\arraybackslash}p{\dimexpr 0.2137\linewidth-2\tabcolsep\relax}>{\raggedright\arraybackslash}p{\dimexpr 0.2307\linewidth-2\tabcolsep\relax}>{\raggedright\arraybackslash}p{\dimexpr 0.2561\linewidth-2\tabcolsep\relax}>{\raggedright\arraybackslash}p{\dimexpr 0.2796\linewidth-2\tabcolsep\relax}}
\toprule
\textbf{Concept} & \textbf{In Lean} & \textbf{You supply} & \textbf{You get for free} \\
\midrule
\endhead
a running system & \texttt{RunSys} & states, step, factor, mass, the ledger law & \texttt{carry}, \texttt{contract}, \texttt{decay\_cost} \\
the settle bound & \texttt{settle}, \texttt{settle\_everywhere}, \texttt{settle\_iff\_bounded}, \texttt{balanceBound} & the contraction class & an explicit step count \\
the decidability edge & \texttt{edge} & nothing & decided above the floor, undecidable at it \\
a chain with costs & \texttt{Chain} & a quantity, costs, the step law & \texttt{Chain.ledger}, \texttt{Chain.decay}, \texttt{Chain.summable}, \texttt{Chain.size} \\
a margin that survives & \texttt{Margin} & shares kept & \texttt{Margin.survive}, \texttt{Margin.pos} \\
an encoder & \texttt{Encoder}, \texttt{szegoFactor}, \texttt{elecFactor}, \texttt{cost} & a state and its rotation & both factors in [0, 1] \\
an uncertainty & \texttt{Uncertainty} & error, factor, fresh uncertainty & \texttt{bounded}, \texttt{critical}, \texttt{sign} \\
the problem card & \texttt{Problem} & input, \texttt{Open}, \texttt{Concl}, \texttt{proved} & \texttt{Problem.holds} \\
the open-step slot & \texttt{OpenStep "name" P} & the one statement \texttt{P} & \texttt{P} itself, labelled \\
the external-input slot & \texttt{External "source" P} & a theorem from the literature & \texttt{P} itself, labelled \\
\bottomrule
\end{longtable}}

The slots are transparent: \texttt{OpenStep "..." P} \emph{is} \texttt{P}, so a theorem takes it as an ordinary hypothesis. Nothing is ever an axiom. The labels make the open steps and outside dependencies greppable.

\textbf{Adding a problem (II.5.2).} Seven steps: copy the template (\texttt{API/Template.lean}, a worked toy); write the input; fill a running system, chain or margin and check the one ledger inequality; write the open step in its slot; call the free theorems; add an uncertainty if the input is only approximate; state the card. \textbf{[paper]} (the procedure; \texttt{docs/unified/API.md}, 3). The three cards of Parts IV to VI are \texttt{API/YangMills.lean}, \texttt{API/GRH.lean} and \texttt{API/NP.lean}, each proved as ``open step gives conclusion'' \textbf{[Lean]}; each open step is \textbf{[open]} (IV.5.6, V.4, VI.4.10). The toy settles in 101 steps to 1/100 (\texttt{toy\_settle\_101}) \textbf{[Lean]}.

\textbf{What runs and what is truth only (II.5.3).} The settle count, the encoder factors on rational data and the ledger law on rational data are computable and decidable: \texttt{balanceBound}, \texttt{szegoQ}, \texttt{elecQ}, \texttt{PaysQ}, \texttt{settle0}, with \texttt{szegoQ\_cast}, \texttt{elecQ\_cast} and \texttt{PaysQ.pays} proving them right for the real system; the machine's step and balance test are primitive recursive (\texttt{estep\_primrec}, \texttt{halted\_primrec}). The real-valued structures and the undecidability on the lossless side are truth statements, not computed. \textbf{[Lean]} (\texttt{API.lean}).

\textbf{One consequence for every instance (II.5.4).} G0's phase field is one compact circle. Every continuous character of the circle is z -\textgreater{} z\textasciicircum{}n for an integer n, so every label carried on it has integer charge, with no further assumption. \textbf{[Lean]} (\texttt{API/Phase.lean}, \texttt{chargeQuantisation}, \texttt{PhaseLabel.integer\_charge}; classical content: Pontryagin duality for the circle).

\textbf{What it does not do (II.5.5).} It proves no open step and adds no mathematics. It does not turn a problem into a running system for you: the structure and its one inequality are the work. Arguments that use only the interface cannot settle P vs NP (relativization, VI.7), and the NP card inherits that. \textbf{[paper]}

\medskip\noindent\rule{\linewidth}{0.4pt}\medskip

\subsection*{II.6 The arsenal and its lint}
\addcontentsline{toc}{subsection}{II.6 The arsenal and its lint}

\textbf{Why it clicks.} Content-free results that more than one project uses should live once, and every use should be checkable for one property: that nothing in it cannot run.

\textbf{The arsenal (II.6.1).} The arsenal is the set of transferable Lean results that need Mathlib, kept in the \texttt{im} repository (\texttt{arsenal/}), pinned to Lean and Mathlib v4.34.0. It is imported from there and never copied (II.6.4). Its modules are \texttt{Arsenal.R136a1}, \texttt{Arsenal.BSD} and \texttt{Arsenal.Rotation} (the one-box test, I.8) \textbf{[paper]} (\texttt{arsenal/README.md}).

\textbf{R136a1 (II.6.2).} The one way to leave a statement open: \texttt{R136a1 "name" P} is \texttt{P} as a named hypothesis, never an axiom. A theorem that needs it takes it as a hypothesis, so the system never assumes it. \textbf{[Lean]} (\texttt{Arsenal/R136a1.lean}, \texttt{R136a1}, \texttt{R136a1.out}).

\textbf{The lint (II.6.3).} \texttt{arsenal/tools/lint.sh} is static, needs no network and no build. It fails unless: the toolchain is Lean v4.34.0 and the one requirement is Mathlib at tag v4.34.0; every module imports only Mathlib or the arsenal itself; no module uses an escape hatch (\texttt{sorry}, \texttt{admit}, \texttt{axiom}, \texttt{native\_decide}, \texttt{unsafe}, \texttt{implemented\_by}, \texttt{extern}, \texttt{ofReduceBool}, \texttt{skipKernelTC}), read in code only so prose stays free; and every module is reached from \texttt{Arsenal.lean}. It lists every part answered R136a1 by name. \textbf{[paper]} (read from the script; run for this draft, result ``lint: OK (3 modules; Mathlib v4.34.0 only; no escape hatch)''). A self-test (\texttt{tools/lint-selftest.sh}) plants one case per rule, and the full check (\texttt{tools/check.sh}) adds a build against Mathlib and the kernel's axiom report for every declaration. Both were not run for this draft.

\textbf{Shared vocabulary, imported and never copied (II.6.4).} What can be transferred moves to the \texttt{im} repository and is used from there as a dependency, never vendored. Lean projects \texttt{require} the package by repository, subdirectory and commit. \textbf{[paper]} (the owner's rule; \texttt{arsenal/README.md}, ``Using it'').

\textbf{What the lint is, and is not.} It certifies form: that everything stated is soundly constructed. It does not say that the statements describe the world (Part I, I.8).

\medskip\noindent\rule{\linewidth}{0.4pt}\medskip

\subsection*{II.7 The independent re-check}
\addcontentsline{toc}{subsection}{II.7 The independent re-check}

Every Lean statement of this paper is checked twice more, by kernels other than the one that built it (\texttt{docs/unified/INDEPENDENT-CHECK.md}, \texttt{tools/independent-check.sh}). Lean's own \texttt{leanchecker} replays each module's declarations into a fresh environment, and \texttt{nanoda\_lib}, a type checker written independently in Rust, checks an export of every declaration together with everything it uses \textbf{[Lean, re-checked]}. A full run once established a trusted base, Mathlib and Lean core at the pinned commit, recorded by hash in a stamp. Later runs re-check only the modules whose source or imports changed, against that stamp.

The run for this draft (2026-11-24, at commit \texttt{2cb301c}) covered 358 modules, including the imported IM arsenal, which is checked like the project's own code and is never part of the trusted base. \texttt{leanchecker} replayed every changed module with no failure; the axiom audit found no declaration resting on anything beyond \texttt{propext}, \texttt{Quot.sound} and \texttt{Classical.choice}; \texttt{nanoda\_lib} checked 103,282 declarations with no type-checker error. The scan for escape hatches reads code only, as the arsenal's lint does, and found none.

\textbf{B. The tools the problems use.}

\subsection*{II.8 One chain, three problems}
\addcontentsline{toc}{subsection}{II.8 One chain, three problems}

All three problems run along the same four links.

\begin{quote}
\textbf{positivity -\textgreater{} a space -\textgreater{} an evolution on it -\textgreater{} a spectrum that the positivity pins}
\end{quote}

{\footnotesize
\begin{longtable}{>{\raggedright\arraybackslash}p{\dimexpr 0.1642\linewidth-2\tabcolsep\relax}>{\raggedright\arraybackslash}p{\dimexpr 0.2903\linewidth-2\tabcolsep\relax}>{\raggedright\arraybackslash}p{\dimexpr 0.2564\linewidth-2\tabcolsep\relax}>{\raggedright\arraybackslash}p{\dimexpr 0.2692\linewidth-2\tabcolsep\relax}}
\toprule
\textbf{Link} & \textbf{Yang-Mills} & \textbf{GRH} & \textbf{P vs NP (sampling side)} \\
\midrule
\endhead
1. Positivity & reflection positivity of e\textasciicircum{}\{-S\} & Weil's pairing W(f * f\textasciitilde{}) \textgreater{}= 0 & positive weights on assignments \\
2. Space & the Osterwalder-Schrader (OS) Hilbert space & the GNS space of Weil's form & L2 of the Gibbs measure \\
3. Evolution & the transfer matrix e\textasciicircum{}\{-H\}, H \textgreater{}= 0 & translation in log-scale & the heat-bath operator, 0 \textless{}= P \textless{}= 1 \\
4. Spectrum & gap = decay rate of correlations & real spectrum = zeros on the line & gap = fast mixing \\
\bottomrule
\end{longtable}}

\textbf{Why it clicks.} Links 2 to 4 are free once link 1 is given: they are theorems, the same theorem in each column. The whole difficulty of each problem sits in link 1 for the hard case (Yang-Mills: a certificate at one scale; GRH: positivity itself; P vs NP: the step from a certified class to all inputs). The rest of this part makes links 2 to 4 precise and adds the four tools that act on link 1 or carry it across scales.

\medskip\noindent\rule{\linewidth}{0.4pt}\medskip

\subsection*{II.9 Positivity builds a space}
\addcontentsline{toc}{subsection}{II.9 Positivity builds a space}

\textbf{Why it clicks.} A non-negative form is all you need to measure lengths, and lengths are all you need for a Hilbert space.

\textbf{Theorem II.9.1 (reconstruction from a positive form).} Let V be a real or complex vector space with a Hermitian form Q that is positive semi-definite, Q(v, v) \textgreater{}= 0. Divide out the null vectors and complete: the result is a Hilbert space H with a map Phi : V -\textgreater{} H such that \textless{}Phi u, Phi v\textgreater{} = Q(u, v). A linear map tau on V that is symmetric for Q and satisfies Q(tau u, tau u) \textless{}= Q(u, u) descends to a bounded self-adjoint operator T on H with 0 \textless{}= \textless{}x, T x\textgreater{} when Q(u, tau u) \textgreater{}= 0. \textbf{[classical]} (Gelfand-Naimark-Segal; Osterwalder-Schrader 1973, 1975). \textbf{[Lean]} in the abstract form used for Yang-Mills: \texttt{M1/UVTVOS.lean} (\texttt{inner\_Φ}, \texttt{T\_symm}, \texttt{T\_pos}, \texttt{contraction}).

\textbf{The three applications.}

\begin{itemize}
\item \textbf{Yang-Mills.} Q(F, G) = \textless{}theta F, G\textgreater{}, with theta the time reflection. Reflection positivity of the Wilson action, site and link, is \textbf{[Lean]} (\texttt{docs/mass-gap/formal/Positivity/}). The transfer step is tau. Positivity is given by construction.
\item \textbf{GRH.} Q is Weil's explicit-formula pairing. In a finite model (finitely many zeros, closed under conjugation) the pairing is non-negative on every finite exponential sum \textbf{iff} every zero is on the line \textbf{[Lean]} (\texttt{GRH/Weil.lean}, \texttt{weil\_criterion\_finite}). For the zeta function itself the equivalence is Weil's criterion \textbf{[classical]} (Weil 1952; Bombieri 2000). Positivity is \emph{not} given by construction: it is equivalent to the conclusion.
\item \textbf{P vs NP (sampling).} Q is the L2 inner product of a Gibbs measure with positive weights; the heat-bath operator is self-adjoint and positive, and every time correlation has a spectral measure on [0, 1] \textbf{[Lean]} (\texttt{Mixing/FiniteSpectral.lean}, \texttt{timeCorr\_spectral}).
\end{itemize}
\medskip\noindent\rule{\linewidth}{0.4pt}\medskip

\subsection*{II.10 Positivity pins the spectrum}
\addcontentsline{toc}{subsection}{II.10 Positivity pins the spectrum}

\textbf{Why it clicks.} A self-adjoint operator cannot hide anything off the real axis, and a positive contraction cannot hide anything outside [0, 1]. So ``where is the spectrum'' becomes ``how fast do correlations decay''.

\textbf{Theorem II.10.1 (decay rate = bottom of the spectrum).} Let nu be a finite measure on [0, inf) (the spectral measure of a vector for H \textgreater{}= 0) and corr(t) = int e\textasciicircum{}\{-tE\} d nu(E). Then corr(n) \textless{}= K e\textasciicircum{}\{-mn\} for all n, for some K, \textbf{iff} nu gives no mass below m. The admissible decay rates are exactly the m with nu([0, m)) = 0, and the best rate is the bottom of the support. \textbf{[Lean]} (\texttt{Bridge/Spectral.lean}, \texttt{spectral\_gap\_iff}, \texttt{decayRates\_eq}; discrete-time form \texttt{Mixing/TimeGap.lean}, \texttt{time\_gap\_iff}).

\textbf{Corollary II.10.2 (gap from reflection positivity).} In the setting of II.9.1, if every block correlation decays at rate m, the transfer operator restricted to the orthogonal complement of the vacuum has norm at most e\textasciicircum{}\{-m\}. \textbf{[Lean]} (\texttt{M1/UVTVOS.lean}, \texttt{gap\_of\_RP}; \texttt{Bridge/OSLink.lean}, \texttt{gap\_of\_RP\_spectral}).

\textbf{The three applications.} Yang-Mills: the mass gap is the decay rate of gauge-invariant correlations. GRH: the zeros are real in the variable gamma = (rho - 1/2)/i exactly when the generator of the evolution on Weil's space is self-adjoint (the Hilbert-Polya form) \textbf{[classical]}. P vs NP: a spectral gap of the heat-bath operator is fast mixing, hence polynomial-time sampling \textbf{[classical]}, finite form \textbf{[Lean]} (\texttt{Mixing/}).

\medskip\noindent\rule{\linewidth}{0.4pt}\medskip

\subsection*{II.11 The bridge lemma: positivity locates singularities}
\addcontentsline{toc}{subsection}{II.11 The bridge lemma: positivity locates singularities}

\textbf{Why it clicks.} A sum of positive terms cannot cancel. So the first place where it blows up must be visible on the real axis, where nothing can cancel.

\textbf{Theorem II.11.1 (Pringsheim, extension form).} Let F be holomorphic on Re s \textgreater{} x0, write its Taylor coefficients at real points c \textgreater{} x0 as F\textasciicircum{}(n)(c) = (-1)\textasciicircum{}n M(c, n) with M real, and suppose F agrees on the right of x0 with a function holomorphic on a disc about x0. Then for some real c \textgreater{} x0 and h \textgreater{}= 0 with c - h \textless{} x0, the series sum M(c, n) h\textasciicircum{}n / n! converges. \textbf{[Lean]} (\texttt{Bridge/Pringsheim.lean}, \texttt{pringsheim\_extension}; the disc form \texttt{pringsheim\_hasSum}). Not in Mathlib before this project.

The Lean core needs no sign: it is Cauchy's theorem in Taylor form. \textbf{Positivity enters in the instances.} There M(c, n) is itself a sum or integral of non-negative terms, such as sum\_k a\_k (log k)\textasciicircum{}n k\textasciicircum{}\{-c\} with a\_k \textgreater{}= 0. Non-negative double series may be summed in either order (Tonelli), and the inner sum is the original series evaluated at c - h. So the original series converges beyond the point where it was supposed to stop, a contradiction. That is the whole mechanism: without positivity the reordering fails and a singularity can hide off the real axis.

\textbf{Three instances, all [Lean]:}

\begin{enumerate}
\item \textbf{Landau's theorem.} An L-series with non-negative coefficients has a singularity at the real point of its abscissa of absolute convergence (\texttt{Bridge/Landau.lean}, \texttt{landau}). This is the first stage of the recent Lean proof of RH by others (Part V).
\item \textbf{Laplace transforms of positive measures.} For a measure on [0, inf), the real edge of convergence of int e\textasciicircum{}\{-st\} d mu is a singularity (\texttt{Bridge/Laplace.lean}, \texttt{laplace\_landau}).
\item \textbf{The mass gap.} The resolvent \textless{}F, (H + s)\textasciicircum{}\{-1\} F\textgreater{} of a positive Hamiltonian converges for real s \textgreater{} -l0 and is singular at s = -l0, where l0 is the bottom of the spectrum (\texttt{Bridge/Spectral.lean}, \texttt{resolvent\_singular}).
\end{enumerate}
\textbf{What the bridge does not give.} It locates the \emph{first} singularity of a positive object. It does not make an object positive. Over function fields, one positive duality (the Hodge index theorem) gives both the axis of the functional equation and the balance point 1/2; over the integers no such positivity is known \textbf{[open]} (\texttt{docs/rh-guide/bridge/BRIDGE.md}, section 0).

\medskip\noindent\rule{\linewidth}{0.4pt}\medskip

\subsection*{II.12 Contraction and data processing}
\addcontentsline{toc}{subsection}{II.12 Contraction and data processing}

\textbf{Why it clicks.} A lossy channel can only forget. If it forgets a fixed fraction each step, it forgets geometrically, and a geometric rate is a gap.

\textbf{Definitions.} A stochastic kernel K moves probability distributions. Its contraction coefficient eta(K) is the largest factor by which it shrinks the total-variation distance between two distributions (Dobrushin's coefficient). For a system of sites with local conditional laws, the influence C(x, z) of site z on site x is the largest change in the law at x when only z changes; the column sum of z is sum\_x C(x, z).

\textbf{Theorem II.12.1 (data processing, composed).} eta(K) \textless{}= 1 for every stochastic kernel, and eta(K L) \textless{}= eta(K) eta(L). The same kernel applied n times contracts by eta\textasciicircum{}n. \textbf{[classical]} (Dobrushin 1956; the strong data-processing inequality). \textbf{[Lean]} (\texttt{Mixing/Chain.lean}, \texttt{contracts\_kchain}, \texttt{contracts\_kpow}, \texttt{tv\_pushv\_kpow\_le\_pow}).

\textbf{Theorem II.12.2 (Doeblin: a quantum per tick).} If K(x, .) \textgreater{}= eps nu(.) for every x, then eta(K) \textless{}= 1 - eps. \textbf{[Lean]} (\texttt{Mixing/Doeblin.lean}, \texttt{contracts\_of\_doeblin}, \texttt{doeblin\_sharp}).

\textbf{Theorem II.12.3 (one finite check, decay everywhere).} If every column sum of the influence matrix is at most alpha \textless{} 1, the heat-bath dynamics has a spectral gap: the spectral measure of every time correlation gives no mass above 1 - (1 - alpha)/N, with N the number of sites. \textbf{[classical]} (Dobrushin 1968; Dobrushin-Shlosman 1985). \textbf{[Lean]} (\texttt{Mixing/FiniteSpectral.lean}, \texttt{spectral\_gap\_dobrushin}; \texttt{Mixing/Link.lean}, \texttt{spectral\_gap\_of\_column}). With exponential weights the same condition gives exponential decay of spatial correlations \textbf{[Lean]} (\texttt{docs/mass-gap/formal/Attempts/R3/TailDecay.lean}, \texttt{theoremA\_w}).

\textbf{The honest limit (II.12.4).} Doeblin gives a rate per channel. A site with k \textgreater{}= 2 neighbours needs the \emph{sum} of the k influences below one, and once the per-channel cost C is at least log 2, the sum of the Doeblin bounds k(1 - e\textasciicircum{}\{-C\}) is at least 1. The quantum per tick does not by itself give a column below one. \textbf{[Lean]} (\texttt{Mixing/Chain.lean}, \texttt{column\_sum\_ge\_one}, with the sharp threshold \texttt{column\_sum\_lt\_one\_iff}).

\textbf{Uniformity along a chain (II.12.5).} If the same channel acts at every step (a homogeneous chain: translation invariance in the transfer direction), the contraction coefficient is the same at every step; once eta \textless{} 1 at one step, the chain contracts like eta\textasciicircum{}n. \textbf{[classical]} with the composition law \textbf{[Lean]} above. Read as DMP: ``the uniformity is given by the arrow of causality'' (the author, 2026-11-23) \textbf{[reading]}.

\textbf{The three applications.} Yang-Mills: the gap mechanism reduces to one box of fixed physical size, uniformly over outsides; for a grid the per-box cost must beat the path entropy, F\_e(L) \textgreater{} log mu\_4 (Part V.5). P vs NP: a column sum below one is polynomial-time sampling and, classically, counting (Part VI). GRH: no contraction step is known; the prime-by-prime extension has no slack to lose (Part V).

\medskip\noindent\rule{\linewidth}{0.4pt}\medskip

\subsection*{II.13 A strict sign gives a rate: the self-dual point}
\addcontentsline{toc}{subsection}{II.13 A strict sign gives a rate: the self-dual point}

\textbf{Why it clicks.} If two quantities are dual to each other and the duality has a fixed point, then they cannot both be small: one of them is at least the fixed-point value.

\textbf{Theorem II.13.1 (hand-over).} For 0 \textless{} Z\_tw \textless{} Z put F\_tw = -log(Z\_tw / Z) and F\_e = -log((Z - Z\_tw)/(Z + Z\_tw)). Then F\_e is the image of F\_tw under a decreasing involution with fixed point log(1 + sqrt 2), and

max(F\_tw, F\_e) \textgreater{}= log(1 + sqrt 2), about 0.881.

The bound is sharp, attained at Z\_tw / Z = sqrt 2 - 1. \textbf{[Lean]} (\texttt{Crossover/HandOver.lean}, \texttt{handover\_ratio}, \texttt{log\_one\_add\_sqrt\_two\_le\_max}; the duality \texttt{Flip/Floor.lean}, \texttt{elecF\_eq\_dual}).

\textbf{Where it applies.} Any system with a strict inequality between a sector and its twist. In SU(2) Yang-Mills the strict inequality holds at every coupling beta != 0 (\texttt{Flip/Strict.lean}, \texttt{su2\_Ztw\_lt\_Z}), so every box has a hand-over rate of at least 0.881 nats in one of its two channels \textbf{[Lean]} (\texttt{Crossover/HandOver.lean}, \texttt{su2\_handover}, with the hypothesis 0 \textless{} Z\_tw stated; the integrand is positive, the formal step is not written). This is the precise form of I.7, ``existence = non-existence + 1'': a sign with no rate becomes a rate at the self-dual point.

\textbf{What it does not give.} It is a statement about free energies of the whole periodic box (averaged over everything), not about the influence of an arbitrary fixed outside on one box. And 0.881 suffices for a chain of boxes (two channels per box) but not for a grid: in d \textgreater{}= 3 the number of ways a path can continue outweighs it (\texttt{Crossover/Containment.lean}, \texttt{path\_entropy\_beats\_handover}, \textbf{[Lean]}). Both gaps are named in Part IV (IV.4, IV.5).

\medskip\noindent\rule{\linewidth}{0.4pt}\medskip

\subsection*{II.14 The dichotomy: local, crossing flux, residue}
\addcontentsline{toc}{subsection}{II.14 The dichotomy: local, crossing flux, residue}

\textbf{Why it clicks.} An obstruction to decay can only be of three kinds: short-range (it is forgotten), conserved (it cannot be forgotten, so it must be paid for), or stuck inside one region (then it is a finite question about that region).

\textbf{Statement (II.14.1, the shape).} For a system split into regions, every piece of an obstruction is

\begin{itemize}
\item (a) \textbf{local}: forced by a few constraints; behind a collar of positive conditional mass it enters only as a boundary term;
\item (b) \textbf{a crossing flux}: carried by a conserved quantity through the region's boundary; it cannot be screened;
\item (c) \textbf{a residue}: inside one region; what remains is a percolation question: does the bad set span the region, or does it break into finite clusters with an exponential tail?
\end{itemize}
This is a proof strategy, not a theorem \textbf{[paper]} (\texttt{docs/mass-gap/weak-coupling/proof-phase/c9a-crossover/DICHOTOMY.md}; \texttt{docs/pnp/DICHOTOMY-TRANSFER.md}). Its exact pieces are theorems:

\begin{itemize}
\item \textbf{Gauss law mod 2.} For a Z\_2 flux on a graph, the parity crossing the boundary of any region equals the charge inside, whatever the interior does; odd total charge is a contradiction \textbf{[Lean]} (\texttt{Proof/Crossing.lean}, \texttt{gauss}, \texttt{flux\_fixed}, \texttt{unsat\_of\_odd}). The same law for centre sheets: a flux-carrying outside forces at least one sheet through the box, entering and leaving an even number of times \textbf{[Lean]} (\texttt{Crossover/FluxRate.lean}, \texttt{exists\_sheet\_of\_flux}, \texttt{even\_card\_fluxCells}).
\item \textbf{Uncontained pieces weigh nothing.} A sheet that ends inside the box is not a closed configuration, so it has weight zero under every law carried by closed configurations \textbf{[Lean]} (\texttt{Crossover/Dichotomy.lean}, \texttt{not\_uncontained}, \texttt{measure\_uncontained}).
\item \textbf{Every refutation crosses.} In every sound proof system whose rules have at most two premises, every refutation of an expanding formula contains a line sensitive to at least the expansion many variables \textbf{[Lean]} (\texttt{Proof/Crossing.lean}, \texttt{cross\_any\_system}).
\item \textbf{Rarity and saturation do not overlap.} The regime where sheets are suppressed (their cost beats their entropy) and the regime where they saturate are disjoint; any route between them crosses a window where neither mechanism applies \textbf{[Lean]} (\texttt{Crossover/Dichotomy.lean}, \texttt{window\_gap}).
\end{itemize}
\textbf{The three applications.} Yang-Mills: (b) is exact, (a) fails uniformly at the SO(3) cut, (c) is the open window (Part IV). P vs NP: the crossing is universal, but its \emph{cost}, the price of naming the conserved quantity in the system's lines, is not; that cost is the open lower bound (Part VI). GRH: the analogue of (b) is the zeros on the line, ``pure oscillation'' \textbf{[reading]}; no dichotomy proof exists there.

\medskip\noindent\rule{\linewidth}{0.4pt}\medskip

\subsection*{II.15 The native optimization method}
\addcontentsline{toc}{subsection}{II.15 The native optimization method}

\textbf{Why it clicks.} A cost that comes in whole units cannot fall forever. Lower it by one unit per improving pass and the method must stop, at an irreducible core.

\textbf{Theorem II.15.1 (quantized descent arrives).} Let f be a map on any set and mu a natural-number measure with mu(f x) \textless{} mu(x) whenever f x != x. Then f has a fixed point after at most mu(x) steps from x. \textbf{[Lean]} (\texttt{Native/Method.lean}, \texttt{fixed\_after}, \texttt{reaches\_floor}). Without quanta it fails: x -\textgreater{} (x + 1)/2 strictly decreases toward 1 and never reaches it (\texttt{real\_never\_arrives}).

\textbf{Theorem II.15.2 (a true contraction for depth).} Pairwise merging of adjacent layers turns d layers into ceil(d/2); after ceil(log2 d) passes one layer remains. \textbf{[Lean]} (\texttt{Native/Method.lean}, \texttt{iterate\_clog\_pairUp}).

\textbf{Measured.} Applied to this project's own proof chains, the per-layer kernel cost is finite and bounded by the size of the finite building blocks it evaluates (0.17 to 1.4 ms per layer in the deepest DMP chain; about 0.3 ms per recipe line for a reflected check). \textbf{[numerical]} (\texttt{docs/pnp/NATIVE-METHOD.md}, section 0).

\textbf{The three applications.} P vs NP: the method reduces any refutation to an irreducible core; whether that core is always of polynomial size is the open question N1 (Part VI). Yang-Mills: the same ``quantum, then arrival'' shape is the Doeblin rate per tick (II.12.2). GRH: the integer multiplicities of zeros and prime powers are the quanta on the prime side \textbf{[reading]}.

\medskip\noindent\rule{\linewidth}{0.4pt}\medskip

\subsection*{II.16 What the tools do not do}
\addcontentsline{toc}{subsection}{II.16 What the tools do not do}

None of II.9 to II.15 creates positivity. Each starts from a positive weight, a positive pairing, a positive cost or a strict sign, and carries it: into a space, into a spectrum, across a chain, across a boundary, down to a floor. The open step of each problem is the place where that starting positivity is known locally, or in a model, and must be established for every case. Part III states this once (section 4). The five sections that follow, II.1 to II.6, are the \emph{form layer}: the part of the machinery that says nothing about content and so serves every problem. II.16 holds for them too.

\medskip\noindent\rule{\linewidth}{0.4pt}\medskip

\subsection*{II.17 Where Part II ends}
\addcontentsline{toc}{subsection}{II.17 Where Part II ends}

The tools II.9 to II.15 carry a positivity. The form layer II.1 to II.6 carries a certificate and its uncertainty, says when the carrying settles, and says where the one open statement of a problem goes. None of it supplies a positivity, and none supplies a floor. Part III says why each problem's open step stays open.


%% ---- 03-part-III-one-assumption.md
\section*{Part III. The one assumption and the rigid start: what the cards show together}
\addcontentsline{toc}{section}{Part III. The one assumption and the rigid start: what the cards show together}

\subsection*{III.0 How to read this part}
\addcontentsline{toc}{subsection}{III.0 How to read this part}

Parts IV to X each end at a \textbf{card}: an input, a proved structure, and exactly one named open step. This part puts the cards side by side. It says what they share, why each still has one open step, and why that step cannot be removed by a trick.

Nothing here solves a problem. Every section below ends in an open statement, and each open statement is the same one that the card's own part already names. What is new is the \emph{shape} the open statements share. Section 1 names what every card needs before its one box can run at all: a containing existence with free possibility space. Sections 2 and 6 are exact: they are theorems in Lean about one geometry. The others are readings of that geometry, and they are marked.

Marks as in Part I: \textbf{[Lean]}, \textbf{[classical]}, \textbf{[external]}, \textbf{[conditional]}, \textbf{[paper]}, \textbf{[numerical]}, \textbf{[reading]}, \textbf{[open]}. Lean paths are relative to \texttt{docs/mass-gap/formal/ProofPhase/}.

\medskip\noindent\rule{\linewidth}{0.4pt}\medskip

\subsection*{III.1 The container: a containing existence with free possibility space}
\addcontentsline{toc}{subsection}{III.1 The container: a containing existence with free possibility space}

\textbf{The one box.} Each card ends in a statement about one box. In DMP terms the box is the \emph{confined inside}: its existence is binary (it is, or it is not), and it is free to assert itself. The inside pays nothing to exist, but it cannot exist alone. It exists only as an assertion \emph{against an outside}: it asserts itself simply by being not that. Without anything to prove to an outsider that it exists, does it exist? \textbf{[reading]} (claim III.54).

\textbf{What the outside must be.} The outside is where locality and causality live. For the box to run in a DMP-complete way, the outside must have two properties:

\begin{itemize}
\item \textbf{Free.} It leaves the possibility space free: something might happen, but it need not. In numbers, 0 \textless{} p \textless{} 1, with p the chance that a crossing inverts the box's binary sector. p = 0 is a forced act (order, no gap). The possibility space is what makes each crossing contract, by the factor abs(1 - 2p) \textless{} 1. Causality in this sense and the gap are one condition.
\item \textbf{Complete.} Every configuration has positive weight: a compact group, a finite action, ``every possibility exists''. Nothing is excluded in advance, so nothing has to be added later.
\end{itemize}
This part calls the two together the \textbf{container}: a containing existence with free possibility space. \textbf{[reading]} (claim III.55).

\textbf{What is proved about it.} Two halves are in Lean for Yang-Mills; the third is a hypothesis.

\begin{itemize}
\item \emph{Complete.} The transfer kernel at a fixed cross-section is strictly positive: every configuration has positive weight \textbf{[Lean]} (\texttt{Flip/Bound.lean}, \texttt{stepK\_pos}; IV.1.7).
\item \emph{Free, from above.} The twisted partition function is strictly below the untwisted one at every coupling beta != 0, so p = Z\_tw/(Z\_tw + Z\_0) is below one half \textbf{[Lean]} (\texttt{Flip/Strict.lean}, \texttt{su2\_Ztw\_lt\_Z}; IV.1.6).
\item \emph{Free, from below.} 0 \textless{} p needs 0 \textless{} Z\_tw, which the Lean statement carries as a hypothesis (\texttt{Crossover/HandOver.lean}, \texttt{su2\_handover}; IV.3.1). The integrand is positive, but the formal step is not written. Keeping p bounded away from 0 as the box grows is the open step (section 4).
\end{itemize}
The strict facts are cross-references to IV.1.6, IV.1.7 and IV.3.1 (claim III.56).

\textbf{The container is the quantifier of the one-box statement.} In every card, the container shows up as one ``for every'':

{\footnotesize
\begin{longtable}{>{\raggedright\arraybackslash}p{\dimexpr 0.1554\linewidth-2\tabcolsep\relax}>{\raggedright\arraybackslash}p{\dimexpr 0.4460\linewidth-2\tabcolsep\relax}>{\raggedright\arraybackslash}p{\dimexpr 0.3785\linewidth-2\tabcolsep\relax}}
\toprule
\textbf{Card} & \textbf{The for-every that the container asks} & \textbf{Status} \\
\midrule
\endhead
Yang-Mills & the hand-over inequality for every outside, or for every \emph{running} outside & named hypotheses \texttt{HandoverSide}, \texttt{HandoverSideRunning} \textbf{[open]}; their relations \textbf{[Lean]} \\
GRH & Weil's form positive for every square-integrable test function on the support (the one-box certificate is ``for every f'') & certified at supports up to a = 0.69 \textbf{[numerical]}; the abstract step \textbf{[Lean]} \\
NP vs coNP & every assignment: a refutation certifies that every assignment falsifies some clause & \textbf{[reading]} on \textbf{[Lean]} refutation completeness \\
Settling machine & total freedom inside, one verified gate & \textbf{[reading]} (Part XI) \\
\bottomrule
\end{longtable}}

\begin{itemize}
\item \textbf{Yang-Mills.} \texttt{HandoverSide} says: for every cutoff, every nontrivial sector and every end-face outside, the hand-over slab has the electric side at most sqrt 2 - 1 of the neutral one. It is O1 at mu = 1 + sqrt 2 at the single length L\_h \textbf{[Lean]} (\texttt{Unified/ElectricSUNPress.lean}, \texttt{HandoverSide}, \texttt{handoverSide\_iff}). \texttt{HandoverSideRunning} quantifies only over \emph{running} outsides, those in the running set R of the measure, and the all-outsides form implies it \textbf{[Lean]} (\texttt{Unified/ElectricSUNRunning.lean}, \texttt{HandoverSideRunning}, \texttt{handoverSideRunning\_of\_handoverSide}). Both statements are named hypotheses and \textbf{[open]} (claim III.57). The owner's reading of the frozen outsides: they are not running outsides, so they are not quantified over, but their ticks of not counting count as remnants. Each enters only through its weight in the measure, an exponentially small probability epsilon. Every factor is at most 1, so the join pays epsilon at most: a fresh uncertainty r = epsilon in the transport theorem (II.1) \textbf{[reading]} for the identification, \textbf{[Lean]} for the transport (\texttt{runningTransport\_slab}, \texttt{runningTransport}) (claim III.57b). The assemblies \texttt{O1\_slab\_running} and \texttt{O1\_box\_running} take the running side, the remnant bound and the seam as hypotheses. They are not O1.
\item \textbf{GRH.} The one-box certificate splits a test function f into a finite part u and a tail g, and certifies three inputs: the finite part, the cross term and the tail. From them the positivity of Weil's form follows for \emph{every} f in the space \textbf{[Lean]} (\texttt{GRH/OneBoxCert.lean}, \texttt{schur\_scalar}, \texttt{onebox\_positive}) (claim III.58b). The three inputs are numbers computed in ball arithmetic at supports a = 0.60, 0.65 and 0.69, the universal start \texttt{\{infinity, 2, 3\}} \textbf{[numerical]} (\texttt{docs/unified/GRH-ONEBOX-CERT.md}). Beyond a = 0.69 the same ``for every f'' is the open step of Part V (claim III.58).
\item \textbf{NP vs coNP.} A refutation is a certificate that every assignment falsifies a clause. Resolution is sound and refutation complete (I.4a, VI.3.1) \textbf{[Lean]}. Reading the ``every assignment'' as the container's quantifier is \textbf{[reading]} (claim III.59).
\item \textbf{The settling machine.} Total freedom inside one verified gate: inside, anything may happen; at the gate, every act is checked for integrity and freshness, not meaning \textbf{[reading]} (XI.1, XI.3; claim III.60).
\end{itemize}
\textbf{Care point: the container is also where the open step sits.} The container is the precondition the cards share, not a result. O1 is exactly the requirement that the one box holds for every outside, uniformly in the scale. The same words, ``for every outside'', state the container and state the open step. That is why no card can prove it from inside: an inside cannot supply its own outside (section 4). \textbf{[reading]} (claim III.61).

\textbf{How the care point is answered.} By construction, not by a further proof inside the cards. The owner's answer: the care point comes from the construction of G0 and the established process by which it gets there. The process from G0 to the hand-over (the ODP, I.5) is one established process. Once that is accepted, every start gets a direction through \texttt{\{infinity, 2, 3\}}. Accepting it is the paper's one assumption, stated once in section 4.4 and not proved. At the start itself the sign is certified (section 6). For every later tick the direction holds only \emph{conditional on the acceptance}. \textbf{[reading]}; the status of each part is marked in section 4.4 (claim III.63).

\textbf{Where section 2 follows from this.} The container asks two things of each tick. From the inside, the box must keep what it holds: a margin is carried from one scale to the next. From the outside, the box must refute what must not be: the twisted sector is suppressed at every box. Section 2 shows that these two requests are the two edges of one band of per-tick factors, and what Lean says about them \textbf{[reading]} (claim III.62).

\medskip\noindent\rule{\linewidth}{0.4pt}\medskip

\subsection*{III.2 One geometry, two edges}
\addcontentsline{toc}{subsection}{III.2 One geometry, two edges}

\textbf{From the container.} Section 1 asked two things of every tick: carry a margin from the inside, refute the twisted sector from the outside. The geometry below is where both are measured, in one number.

\textbf{What is shared.} GRH (Part V) and the Yang-Mills one-box step O1 (Part IV) read the same object. At a boundary, an encoder has one number, the reflection coefficient c. Its distance d from the lossless point (c = 0) is one hyperbolic distance. Two loss factors are functions of that one distance:

\begin{itemize}
\item the electric factor is e\textasciicircum{}\{-d\};
\item the Szego factor 1 - \textbar{}c\textbar{}\textasciicircum{}2 equals 4 e\textasciicircum{}\{-d\} / (1 + e\textasciicircum{}\{-d\})\textasciicircum{}2, which is 1 / cosh\textasciicircum{}2(d/2).
\end{itemize}
The cost -log of the electric factor is the distance d, and the Szego factor is the second function of d: both are \textbf{[Lean]} (\texttt{Unified/GRHLens.lean}, \texttt{elec\_cost}, \texttt{szego\_of\_elec}; the cosh form is a rewriting of the second) (claim III.8). The ledger is also shared: \texttt{Pays}, \texttt{carry} and \texttt{Ledger.contract} are one theorem for both problems (II.1) \textbf{[Lean]}.

\textbf{What they ask of it.} They ask opposite things, one-sided and on different encoders:

{\footnotesize
\begin{longtable}{>{\raggedright\arraybackslash}p{\dimexpr 0.1618\linewidth-2\tabcolsep\relax}>{\raggedright\arraybackslash}p{\dimexpr 0.3296\linewidth-2\tabcolsep\relax}>{\raggedright\arraybackslash}p{\dimexpr 0.2534\linewidth-2\tabcolsep\relax}>{\raggedright\arraybackslash}p{\dimexpr 0.2352\linewidth-2\tabcolsep\relax}}
\toprule
\textbf{Statement} & \textbf{In the tick's distance d} & \textbf{On which encoder} & \textbf{Status} \\
\midrule
\endhead
GRH, the sign & d is finite at every tick (abs(k\_N) \textless{} 1) & Verblunsky step k\_N of the zero measure, one per scale & \textbf{[open]} (claim III.9) \\
GRH, multiplicative transport & d\_N \textless{}= D uniformly: every tick loses little (factor at least e\textasciicircum{}\{-D\}) & same & \textbf{[open]}; enough for the sign, not needed by it (claim III.9) \\
O1 & d \textgreater{}= kappa \textgreater{} log mu for every outside: every box loses a lot (factor at most e\textasciicircum{}\{-kappa\}) & twist ratio Z\_tw/Z of the hand-over box, one per outside & \textbf{[open]} (IV.5.6) \\
\bottomrule
\end{longtable}}

In words. \textbf{GRH asks every tick to lose little}, so that a margin is carried from the start to every scale. What exists, stays. \textbf{O1 asks every box to lose a lot}, so that the twisted sector is refuted. What must not exist, is removed. The two are the lower edge and the upper edge of one band of factors, e\textasciicircum{}\{-D\} \textless{}= f \textless{}= e\textasciicircum{}\{-kappa\}, with kappa \textgreater{} log mu \textbf{[reading]} (claim III.15).

\textbf{What Lean proves about the pair} (\texttt{Unified/GRHToO1.lean}, new in this draft; \texttt{docs/unified/GRH-TO-O1.md}):

\begin{enumerate}
\item \textbf{O1 is the uniform ratio bound, and a strict margin.} O1 holds exactly when some mu' \textgreater{} mu has mu' p1 \textless{}= p0 for every outside, with p0 and p1 the two electric sectors of the box. Equivalently, the sector form p0 - mu p1 is at least delta p1 for one delta \textgreater{} 0, uniformly \textbf{[Lean]} (\texttt{ymOneBox\_iff\_ratio}, \texttt{ymOneBox\_iff\_margin}, \texttt{cost\_eq\_log\_ratio}; the encoder form is \texttt{Unified/Encoder.lean}, \texttt{oneBox\_iff\_refl}). So the ``uniform ratio bound'' earlier rounds named as the missing bridge is not a smaller statement than O1. It is O1 (claim III.10).
\item \textbf{They exclude each other on one tick sequence.} Take one sequence of ticks whose encoders stay at least theta \textgreater{} 0 from the lossless point (O1 at each tick). Then the Szego product, the carried margin, falls below every m \textgreater{} 0. A uniform margin and O1 at every tick cannot both hold \textbf{[Lean]} (\texttt{o1\_ticks\_kill\_margin}, \texttt{uniform\_margin\_excludes\_o1}) (claim III.11).
\item \textbf{Each holds without the other.} Lossless ticks transport a margin and fail O1. Factors 2\textasciicircum{}\{-n\}/(2 mu) satisfy O1 and transport nothing \textbf{[Lean]} (\texttt{transport\_not\_decay}, \texttt{decay\_not\_transport}) (claim III.12).
\item \textbf{Birkhoff's bounded diameter is neither.} On the two-sector cone, a bounded spread of the image and O1 are independent: bounded spread with O1 false (all ratios 1), and O1 with unbounded spread (ratios e\textasciicircum{}\{-n\}/(2 mu)) \textbf{[Lean]} (\texttt{spread\_without\_o1}, \texttt{o1\_without\_spread}). O1 is a \emph{location} of the image, on the far side of the ray at ratio 1/mu. It is not the image's \emph{size} (claim III.13).
\end{enumerate}
\textbf{What does not follow.} The shared thing is a geometry and a role, not one statement that could be proved once for both. No implication between GRH and O1 is proved either way: in the card form ``GRH implies O1 on every input'' would be a disproof of GRH (claim III.16), and the finite model carries a margin of exactly zero (claim III.17). Both stay open. No known object has both the zero measure's Verblunsky steps and the box's sector ratios as its ticks, and no known result supplies either edge of the band: it is two open statements in one geometry \textbf{[open]} (claim III.14).

\textbf{Why ``GRH implies O1'' is empty in the card form.} The cards take their content as input. Some Yang-Mills input breaks O1. So for any statement H, ``H implies that every input meets the open step'' is equivalent to ``not H''. For Mathlib's Riemann hypothesis, for Dirichlet GRH and for the GRH card's own open step, a proof of ``GRH implies O1 on every input'' would therefore be a disproof of GRH \textbf{[Lean]} (\texttt{imp\_all\_iff\_not}, \texttt{badInput\_not\_open}, \texttt{rh\_imp\_card\_iff}, \texttt{dirichlet\_imp\_card\_iff}, \texttt{grhCard\_imp\_card\_iff}) (claim III.16). A real bridge must be input-specific: a statement that mentions the box's sector weights and zeta's zeros together. That statement is named exactly, \texttt{ZetaCoupling}: one bound on the sector ratio, valid for every outside, controlled by RH. The coupling and RH give O1, and O1 alone gives the coupling, so the coupling is no easier than O1 \textbf{[Lean]} (\texttt{o1\_of\_coupling}, \texttt{coupling\_of\_o1}); its content for four-dimensional SU(N) is \textbf{[open]}, and no theorem relates such Wilson weights to zeros of zeta (claim III.19).

\textbf{The positivity GRH supplies has margin exactly zero.} In the project's finite model with k zeros, the card's open step (the form is nonnegative at every scale) puts Weil's form at scale k at a null vector. So the form is positive with margin 0, and no positive margin is certified at that closure \textbf{[Lean]} (\texttt{exists\_null}, \texttt{toepForm\_atoms\_null}, \texttt{no\_margin\_at\_closure}, \texttt{grh\_margin\_exactly\_zero}) (claim III.17). And positivity does not choose the side: positive weights in both sectors with O1 false exist for every mu \textgreater{}= 1 \textbf{[Lean]} (\texttt{positivity\_without\_o1}) (claim III.18). Both problems have a proved \emph{sign}. O1 needs a \emph{margin}, for a different form, at a different place.

\textbf{The one model where zeta is the partition function.} Two-dimensional SU(2) on a closed surface of genus g \textgreater{}= 2 has, in the zero-area limit, the partition function zeta(2g - 2) \textbf{[external]} (Witten, CMP 141, 1991) (claim III.21). The two electric sectors are 2\textasciicircum{}\{-s\} zeta(s) and (1 - 2\textasciicircum{}\{-s\}) zeta(s) with s = 2g - 2. Zeta cancels from every ratio, so the O1-type quantity there is one Euler factor, at the prime 2, and no zero of zeta enters \textbf{[Lean]} (\texttt{Ze1\_eq}, \texttt{Ze0\_eq}, \texttt{zeta\_re\_eq}, \texttt{sector\_dom\_iff}, \texttt{twist\_ratio}) (claim III.20). Two dimensions have no outside, so the hard part of O1 has no counterpart on the zeta side.

\medskip\noindent\rule{\linewidth}{0.4pt}\medskip

\subsection*{III.3 Existence and falsifiability: the duality of construction}
\addcontentsline{toc}{subsection}{III.3 Existence and falsifiability: the duality of construction}

\textbf{The owner's reading.} The owner puts the difference between the two edges in one sentence. GRH asks that something \emph{exists}. O1 asks that something is \emph{falsifiable}. Together they are the duality inherent in the construction of a DMP-complete system: what is prescribed to exist, and what is prescribed to be refuted \textbf{[reading]} (claim III.23).

In the band of section 2 the reading is exact in this sense:

\begin{itemize}
\item \textbf{GRH's edge} asks that the margin exists at every tick. A witness stays.
\item \textbf{O1's edge} asks that every twisted state is refuted at every box. A bound, whose violation would show up as a failure at one box.
\item They are the two sides of the existence boundary (I.3), and the same duality as ``it exists, or it is answered R136a1'' at each tick (section 5).
\end{itemize}
\textbf{Care point: as whole statements the logical forms are crossed.} RH, and by the notes also GRH for Dirichlet L-functions, is a Pi-0-1 statement: one computable counterexample refutes it \textbf{[classical]} (Lagarias' criterion; Davis-Matiyasevich-Robinson; cited as the project notes cite them, not re-read here) (claim III.24). The Yang-Mills problem asks for an existence: a continuum theory with a gap \textbf{[classical]} (the Clay statement) (claim III.25). So each problem, stated whole, has the logical form of the \emph{other's role} in the band. The duality appears twice: once in the roles, once, mirrored, in the statements \textbf{[reading]} (claim III.26).

Two further care points keep this honest. First, ``falsifiable'' is a word about the role, not a theorem: the sector weights of O1 are integrals, not finite computations, and a measurement can refute or guide but never establish a for-all statement (0.5; IV.7). Second, neither word classifies the open step as easy or hard. It says which edge of the band the step guards.

\medskip\noindent\rule{\linewidth}{0.4pt}\medskip

\subsection*{III.4 Why O1 is always the open step}
\addcontentsline{toc}{subsection}{III.4 Why O1 is always the open step}

\subsubsection*{III.4.1 One named slot per card}
\addcontentsline{toc}{subsubsection}{III.4.1 One named slot per card}

Every problem the project treats has the same three parts:

\begin{enumerate}
\item a well-defined \textbf{input} (a chain of boxes, a model of zeros, a proof system);
\item a \textbf{proved structure} on it: a quantity carried along steps that each pay a price, with the theorems that follow for every such chain;
\item \textbf{one named open step}: the statement about the running state that nobody has proved.
\end{enumerate}
The form layer gives parts 2 and 3 a fixed Lean shape (II.5): a \texttt{Problem} card with an input, one \texttt{OpenStep} slot, a conclusion, and a proof that the open step gives the conclusion. Nothing is ever an axiom: a slot is a hypothesis \textbf{[Lean]} (\texttt{API.lean}, \texttt{Problem}, \texttt{Problem.holds}, \texttt{OpenStep}, \texttt{External}). The three cards of Parts IV to VI are written in this shape in \texttt{API/YangMills.lean}, \texttt{API/GRH.lean} and \texttt{API/NP.lean}. The cards for Parts VII to X are written by the owners of those parts and are not restated here (claim III.27).

{\footnotesize
\begin{longtable}{>{\raggedright\arraybackslash}p{\dimexpr 0.1548\linewidth-2\tabcolsep\relax}>{\raggedright\arraybackslash}p{\dimexpr 0.2607\linewidth-2\tabcolsep\relax}>{\raggedright\arraybackslash}p{\dimexpr 0.2912\linewidth-2\tabcolsep\relax}>{\raggedright\arraybackslash}p{\dimexpr 0.2733\linewidth-2\tabcolsep\relax}}
\toprule
\textbf{} & \textbf{Yang-Mills (SU(2))} & \textbf{GRH (beyond Dirichlet; independently for Dirichlet)} & \textbf{NP vs coNP (via proof complexity)} \\
\midrule
\endhead
\textbf{The one object} & one box of fixed physical size & Weil's form at every support, from arithmetic alone & a cost measure on cutting-planes lines \\
\textbf{The open statement} & F\_e(L) - 2c(L) \textgreater{} log mu\_4 = 1.91, the excess act included, uniformly over every well-defined outside (IV.5.6) \textbf{[open]} & lambda\_min(a) \textgreater{}= 0 at every support, built from the primes, the characters and the real place without the zeros; eight faces (V.4) \textbf{[open]} & interface hardness at width 3: a superpolynomial lower bound on dag-like protocols for the falsified-clause search of a random 3-CNF (VI.4.10) \textbf{[open]} \\
\textbf{Proved for one unit} & max(F\_tw, F\_e) \textgreater{}= 0.881 per box (IV.3.1) \textbf{[Lean]} & positivity in finite models (V.1) \textbf{[Lean]}; each prime step numerically (V.3.2) \textbf{[numerical]} & duality floor (e + 1)/2 for unit thresholds (VI.4.5) \textbf{[Lean]}; every refutation crosses (VI.4.3) \textbf{[Lean]} \\
\textbf{Why the unit does not carry} & 0.881 covers a chain (two channels), not a grid; path entropy log mu\_4 beats it in d \textgreater{}= 3 (IV.4.1, IV.5.4) \textbf{[Lean]} & only the full set of primes is positive; place-by-place and two-term assembly fail (V.3.5, V.3.7, V.3.8) \textbf{[Lean]} + \textbf{[numerical]} & weighted lines are cheap in both channels; one line aggregates its units (VI.4.4, VI.4.6) \textbf{[Lean]} \\
\textbf{What no duality supplies} & 't Hooft's relation forces nothing for n \textgreater{}= 2 sectors (IV.4.4) \textbf{[Lean]} & the functional equation allows off-line pairs (V.2.1, V.2.2) \textbf{[Lean]} + \textbf{[classical]} & the duality floor fails for weighted lines (VI.4.6) \textbf{[Lean]} \\
\textbf{The quantifier that makes it hard} & every outside (worst case, not typical) & every support, with margin 1e-9 & every refutation, of a structureless formula \\
\textbf{What measurement says} & passes centrally at L = 6, beta 2.3; not certified (IV.7.1) \textbf{[numerical]} & holds at every step computed, by exact cancellation (V.3.2) \textbf{[numerical]} & no measurement bears on a lower bound \\
\textbf{The ``it runs'' form} & the positive norm is free (RP); the margin eta \textless{} 1 is missing & the margin is not needed; the positive norm (Weil) is missing & the ``it runs'' object is the dag-like protocol; its size is missing \\
\bottomrule
\end{longtable}}

The table is a summary of Parts IV to VI and their proved obstructions \textbf{[paper]} (claim III.2a). The shape in the table is one shape: a positivity, or a certificate, is known for one unit; what is missing is that it holds for every case at once, without loss, and no assembly with slack can supply it \textbf{[reading]} (claim III.1). The last row is formal for the first two problems: \texttt{runs\_two\_targets} (\texttt{Crossover/Runs.lean}) and \texttt{runs\_and\_mirror\_on\_line} (\texttt{GRH/Tate.lean}) state the two targets side by side \textbf{[Lean]} (claim III.2b). For NP vs coNP it is the notes' identification of the size-indexed object for cutting planes (\texttt{CP-DUALITY.md}, 0.5) \textbf{[paper]} (claim III.2c).

Only in Yang-Mills is the starting positivity given by construction: e\textasciicircum{}\{-S\} \textgreater{}= 0 together with locality gives reflection positivity, and the chain of II.8 runs from it. GRH lacks an arithmetic G0 in which Weil positivity is a construction, and P vs NP lacks a description language in which the existence of every instance is local \textbf{[paper]} (observation on II.9) (claim III.5). This is the content of the last row, and of II.9.

\subsubsection*{III.4.2 Why the open step is the same slot every time}
\addcontentsline{toc}{subsubsection}{III.4.2 Why the open step is the same slot every time}

The owner's statement: the assertion of a DMP-complete running system must always be the open step. The reading has two senses.

\begin{itemize}
\item \textbf{Structural.} By construction every card's open step is the running-state hypothesis R: ``this object runs''. The form layer proves everything else. Each named hypothesis of the three bridges (section 7) says that its object is in a running state \textbf{[reading]} (claim III.28); the implication from R to the conclusions is \textbf{[Lean]} (claims III.7n, III.7m).
\item \textbf{Logical.} A system cannot certify its own running from inside. A consistent theory strong enough for arithmetic cannot prove its own consistency \textbf{[classical]} (Goedel II), and the settling machine cannot decide balance for all content \textbf{[Lean]} (II.3). So the certificate of running comes from the containing system. \textbf{[reading]} (claim III.29).
\end{itemize}
Care point (claim III.29): ``open from inside'' is not ``unprovable''. O1 is a statement about the Wilson measure that ordinary mathematics, the containing system, may well prove. What is always true is that the system cannot supply it, and that here it is the last statement to be made. The same reading applies to DMP itself: its own running assertion is its O1, the one statement it cannot certify from inside \textbf{[reading]} (claim III.29).

\subsubsection*{III.4.3 The dichotomy that makes the step decidable}
\addcontentsline{toc}{subsubsection}{III.4.3 The dichotomy that makes the step decidable}

The settling machine (II.3) has two sides. With a floor, meaning every step pays a cost above the floor, balance is reached within an explicit number of steps, so ``settled or not'' is decidable \textbf{[Lean]} (\texttt{reaches\_of\_contracting}, \texttt{balanceBound\_ge}, \texttt{contracting\_iff\_costAboveFloor}). Without a floor, at cost equal to the floor, no procedure decides it for all content \textbf{[Lean]} (\texttt{etm\_balance\_undecidable}, \texttt{etm\_edge}) (claim III.30; these are XI.2a and XI.2b, stated once in II.3).

\textbf{The gap at G0.} The owner's reading: at G0 a gap and an arrow have to be established. The gap is what turns the gate's answer at each tick into a binary one (it passes, or it is outside). The arrow orders the ticks, so a later excision is a later tick and never a revision of an earlier one \textbf{[reading]} (claim III.31). Confinement is the gapped regime in which this holds. Care point: in Lean the floor is a hypothesis of the settle theorem; for Yang-Mills in the continuum, establishing it is O1 itself. So this reading names why O1 is the step every construction rests on. It does not close it.

\textbf{The floor at the error bars.} The owner's reading: the floor is at least the existence tied to the error bars. A floor is never zero, because every assertion carries an error bar, and a finite resolution is a cost above zero. In the settle theorem this gives settling to within the error bar in an explicit number of steps, with the floor as a hypothesis \textbf{[Lean]} (\texttt{reaches\_of\_contracting}, epsilon the error bar) for the implication, and \textbf{[reading]} for the identification (claim III.32). The continuum is the limit in which the error bar goes to zero. Keeping the floor positive in that limit is O1 \textbf{[open]} (claim III.33).

The circle to avoid (claim III.34): it would be circular to justify O1 by ``the floor is at least the error bar'' and then to justify the error bar by O1's rate. The settle theorem takes the floor as a hypothesis. That keeps the argument honest \textbf{[paper]} (observation in \texttt{docs/unified/GRH-TO-O1.md}, section 6).

The floor reading also corrects a tempting derivation. GRH supplies the \emph{sign}, not the floor: in the finite model its form is degenerate exactly at the closure, the opposite of a floor (claim III.17). So the floor reading describes what O1 would deliver, not something GRH delivers \textbf{[reading]}.

\textbf{Only a decidable start is DMP-complete.} The owner's reading: a class that positivity cannot act on, such as a Z/2 class, is the undecidable side, and a system is DMP-complete only when its start forces a decidable reading from rotation. Exact counterpart: a floor is an order statement, and Z/2 carries no order, since no ordered ring has characteristic 2 \textbf{[classical]} (claim III.70a). The one-box test has exactly the two regimes of this section: with a floor it is a finite check, without one it is undecidable \textbf{[Lean]} (\texttt{rotation\_edge}, claim III.68c). The start \texttt{\{infinity, 2, 3\}} is where the floor is forced, Weil's form positive with explicit margins (claim III.65). So ``DMP-complete'' in the paper's sense reads as ``the start forces a floor'' \textbf{[reading]} (claim III.70). Care point: the start forces the floor at the start only; keeping it at every scale is the one assumption (claims III.64, III.66).

\medskip\noindent\rule{\linewidth}{0.4pt}\medskip

\subsubsection*{III.4.4 The paper's one assumption}
\addcontentsline{toc}{subsubsection}{III.4.4 The paper's one assumption}

This is the one place where the paper states what it assumes. Every result that uses it is marked \textbf{conditional} on it. The gates ask for it proved. It is not proved here.

\textbf{The assumption (the established process).} The process from G0 to the hand-over, the ODP (I.5), is one established process: pressing, then minting, with the certified start \texttt{\{infinity, 2, 3\}} as its first tick. Accepting that the process is established means accepting that the running state is reached: the one box is a running, DMP-complete system. In DMP terms this is the MP-type existence assumption. In each card's terms it is the open step, O1 for Yang-Mills (section 4.2: ``this object runs''). O1 is then the \emph{corollary} of running, slotted as the named hypothesis R of the transport theorem (claim III.7n). \textbf{[reading]} for the identification; \textbf{[open]} as a statement about each object (claim III.64).

\textbf{DMP-complete, defined.} The owner's definition: DMP-complete means \emph{runs on DMP}. In this paper's terms an object is DMP-complete when its running-state hypothesis R holds (section 4.2). The definition is circular in the way DMP itself is: whether R holds cannot be certified from inside (claim III.29) \textbf{[reading]} (claim III.64b).

\textbf{One assumption, two edges.} The running state has the two edges of section 2. Its upper edge is decay: O1. Its lower edge is that the running does not collapse: for Yang-Mills, a correlation bounded below uniformly in the scale (IV.9.5b). In the spectral model neither edge implies the other, and both fit one window \textbf{[Lean]} (model; \texttt{gap\_not\_lower}, \texttt{gap\_and\_lower\_compatible}). The owner's reading puts both in the one assumption: the ODP runs before any start and ends only once a dimensionless 1 is reached, so a DMP-complete system carries a unit, and with it a correlation that does not vanish, by construction. In Yang-Mills the unit is Lambda: every physical quantity is a pure number times a power of it (I.5) \textbf{[classical]}, and a vanishing correlation would leave no non-zero such number \textbf{[reading]} (claim III.64a). So Yang-Mills carries one assumption with two edges, not two assumptions. The owner places the lower edge exactly at the ODP's end, which delivers the start \texttt{\{infinity, 2, 3\}}; from there the rest follows for a DMP-complete system with nothing more to show. So the lower edge is \textbf{conditional} on the one assumption in the same way as the direction of every later tick (claim III.66), and without the assumption it is unproved, as O1 is (IV.9.5b).

\textbf{The one-box certificate.} The author's reading: every tick can be reconstructed from any scale, so the one step can be certified as soon as it works on the box. In the one-box test (I.8) the induction is a theorem: a marker correct at the start and kept correct by one step at the shared gate is correct at every tick \textbf{[Lean]} (claim III.68). The scale version is a theorem \emph{from} the one assumption: if every scale's box is rebuilt from the one box through its gate, a box certified at one scale is correct at every scale, and one scale with a wrong marker refutes the sharing \textbf{[Lean]} from III.64 as a named hypothesis (claim III.68a). Sharing the gate alone is free; the content of III.64 is that the marker travels with it (claim III.68b). The owner reads the one assumption as a leak axiom: anything within a DB shares a universal truth, so a subset cannot attach without it. For the gate this is a theorem: every running system maps in exactly one way into the universal one, and a subsystem attached by inclusion shares its trace \textbf{[Lean]} (\texttt{exists\_unique\_hom}, \texttt{trace\_unique}, \texttt{subIncl}). What stays assumed is the marker's part, that the claim travels with the gate \textbf{[reading]} (claim III.71). The author grounds the sharing in GRH: it gives meaning to anything translated through it, so once it is assumed a valid encoder, correctness can be inferred. In the test that is exact. A valid encoder between scales is the shared gate, and correctness at every scale is inferred from it \textbf{[Lean]} (\texttt{scale\_certificate}); what an encoder carries keeps its reflection coefficient \textbf{[Lean]} (\texttt{refl\_map}). Care point: a valid encoder keeps the same margin at every scale, and the margin is the floor, which GRH alone does not supply (claim III.17). So assuming the encoder valid is III.64 itself, and ``certified on the box, hence everywhere'' is \textbf{conditional} on it (claim III.68d). The universal machine is a box too: its DMP-completeness is undecidable over all content, and on the contraction class it is a finite check, the dichotomy of 4.3 in the test's terms \textbf{[Lean]} (claim III.68c).

\textbf{The rigid start and the language of the encoder.} The owner reads Poincare for SU(N) (A.3j) as fixing the start. Among SU(N) the homotopy type determines the manifold exactly for N = 1, 2, 3, and the circle U(1) is rigid too \textbf{[paper]} (A.3k). So the largest start built from SU(N) that a description can certify is the point with 2 and 3, the shape of \texttt{\{infinity, 2, 3\}} \textbf{[reading]} (claim III.69). The encoder is a norm-preserving map of an inner product space \textbf{[Lean]}; the determinant-1 norm-preserving maps of C\textasciicircum{}N form SU(N), and every compact Lie group sits inside some U(N) (Peter-Weyl) \textbf{[classical]}. In the owner's words SU(N) is the language the encoder speaks, and every other compact group speaks inside it (claim III.69a). Speaking (a valid encoder, the shared gate) and being understood (a correct marker) give DMP-completeness at every scale \textbf{[Lean]} (\texttt{scale\_certificate}; claim III.69b). So anything that rotates correctly from the rigid start is correct if DMP is correct: \textbf{conditional} on III.64 (claim III.69c). Care points: this covers SU(N), not every conceivable start; and ``other compact groups do not survive rotation'' is open. If surviving means rigidity, Sp(2) = Spin(5) and G2 pass the rational test that excludes SU(N), N \textgreater{}= 4, and their rigidity was not checked (claim A.5j). The owner's further reading settles what matters: only U(1) carries content across the gate. Sp(2) and G2 have only self-conjugate representations, so they carry no oriented phase whether or not they are rigid, and among the groups that can carry one, rigidity leaves only SU(3), which carries it through its centre on the U(1) circle (claims A.5k to A.5m). The owner closes the circle: Sp(2) and G2 are the U(1) gate's two hidden witnesses, dimensionless to the encoder, and with SU(3) they are the three compact simple groups of rank 2 (claims A.5n, A.5o). The reading needs the witnesses rigid: Sp(2) is, in the strong form; G2 has exactly two structures, and whether the second is G2 itself is one open test (claims A.5j, A.5p, A.5r).

\textbf{What can be shown about O1.} The owner's statement: the only thing that can be shown is that O1 does not fail for anything that can be described as existing in the containing DB. Three things stand behind it, each with its own status.

\begin{itemize}
\item \textbf{Any failure is describable.} O1 is a universal statement over describable instances: one box of fixed physical size, one outside, one level. A failure would therefore be one describable configuration, a finite witness \textbf{[paper]} (logic) (claim III.72).
\item \textbf{None has been found.} The certified start holds with margins (III.65, \textbf{[certified]}); per-box positivity holds by construction (IV.9.5a, \textbf{[Lean]}); small tori pass centrally at L = 6 but are not certified (IV.7.1, \textbf{[numerical]}). The one refuted check (IV.7.2) is the chain criterion in another geometry and bears on O1 neither way, and the cubic one-box geometry is untested by saved data (IV.7.3).
\item \textbf{Care point.} ``O1 does not fail for any describable instance'' is O1 itself, not a weaker statement that could be proved on the side. It is the one assumption phrased so that it can be falsified: one describable configuration would refute it. The owner reads it as the containing DB's guarantee that whatever exists in it can be described \textbf{[reading]} (claim III.72a).
\end{itemize}
\textbf{The expected negative.} Navier-Stokes is the card where the process does not build: the energy ledger cannot pay a uniform cost per scale (VII.3.4, \textbf{[Lean]} in the model). In the owner's reading nothing more needs or can be said than that it is not DMP-complete, and blowup is the expected consequence \textbf{[reading]} (VII.7e). This makes the framework falsifiable: a proof of global regularity would refute the prediction, and a certified blowup profile would confirm it (VII.5.1, \textbf{[open]}).

\textbf{What follows under it.} Accepting the process orients every start. Every start is directed through \texttt{\{infinity, 2, 3\}}:

\begin{itemize}
\item \textbf{At the start, certified.} At supports where only the real place and the primes 2 and 3 enter (up to a = 0.69), Weil's form is positive with an explicit margin, c = 1.5e-9, 3.4e-11, 6.4e-13 at a = 0.60, 0.65, 0.69. So the first tick is certified positive with a margin, and the sign is fixed there \textbf{[numerical]} (certified in ball arithmetic; \texttt{docs/unified/GRH-ONEBOX-CERT.md}), with the abstract step \textbf{[Lean]} (\texttt{GRH/OneBoxCert.lean}, \texttt{onebox\_positive}; claim III.58b). This is claim III.65.
\item \textbf{For every later tick, conditional.} The sign at the start, and with it the direction of every later tick, is fixed \emph{only if} the process is accepted. This is \textbf{conditional} on the acceptance (claim III.66). It is not a theorem that the margin survives: for GRH that is the open step (claims III.9, III.46), and the margin collapses fast beyond the start (about 3.6e-15 at a = 0.75) \textbf{[numerical]}, so it cannot be carried outward by computation.
\end{itemize}
\textbf{How to read a conditional result.} Whenever Part III, or a later part, writes ``conditional on the established process'', it means exactly this: the statement holds if the process of G0 to the hand-over is established, and the process being established is the assumption of this subsection. Where a result holds without it, the result says so: the Lean theorems of section 2 (the geometry of the two edges), the transport theorem (II.1) and the certified start do not use it. \textbf{[paper]} (claim III.67).

\textbf{Why it is stated plainly.} The settling gates ask for it proved. Until it is, the abstract and the introduction should say ``conditional on the established process (O1)''. That wording is the coordinator's to place; this part does not edit them.

\subsection*{III.5 Excision and incompleteness}
\addcontentsline{toc}{subsection}{III.5 Excision and incompleteness}

\textbf{R136a1 is the answer ``outside''.} In DMP a part that is answered R136a1 is outside the system, whatever it is. Nothing else needs to be proved about it. In the arsenal of the IM repository this is the one way to leave a statement open: \texttt{R136a1 "name" P} is the statement P taken as a named hypothesis. It never becomes an axiom, so every declaration still rests on Lean's three standard axioms only \textbf{[Lean]} (\texttt{Arsenal/R136a1.lean}, \texttt{R136a1}, \texttt{R136a1.out}, in the \texttt{im} repository, imported and never copied) (claim III.35). The arsenal's lint lists every part answered R136a1 by name (II.6).

\textbf{Excision decides per tick.} Goedel's theorems apply to a whole consistent, effectively axiomatised theory that contains arithmetic: some true sentences stay undecided \textbf{[classical]} (claim III.36). Two known ways to have completeness anyway are to restrict to a complete and decidable fragment (Presburger arithmetic, Tarski's real closed fields) \textbf{[classical]}, or to decide per tick: ``provable within n steps'' is decidable for every n, so at each tick everything not yet decided is outside, and the boundary moves with the ticks. DMP's excision is the second way, with R136a1 as the answer for ``outside''. The decidable half is \textbf{[Lean]} (II.3, \texttt{reaches\_of\_contracting}, with \texttt{etm\_balance\_undecidable} for the edge); the identification of excision with it is \textbf{[reading]} (claim III.37).

Care point (claim III.38): the excision is decidable per tick. The set of all sentences ever decided is not decidable in general. So ``outside at tick n'' can become ``inside at a later tick''. That is a new act, never a revision of the earlier tick. The owner adds that this includes the \emph{absence} of an expected tick: a missing tick is an answer too. The part is then outside. This is the time-to-live of every act, and the hard cutoff of the settling machine (Outlook) \textbf{[reading]} (claim III.41).

\textbf{Why the Clay statements cannot be excised as a class.} The Clay problems are statements of arithmetic and analysis. If one of them were independent of the axioms, it would be answered R136a1, ``outside''. That is not the settling the gates ask for \textbf{[reading]} (claim III.39). The owner agrees, and adds that the assertion simply states that an O1 of that behaviour exists and then lints everything that cannot survive it \textbf{[reading]} (claim III.39).

\textbf{Two markers, two channels.} \texttt{sorry} says ``this derivation is incomplete here'', and the kernel reports it as \texttt{sorryAx} in the axiom list. R136a1 says ``this part is excluded by name'', and it stays a hypothesis, never an axiom, so the lint lists it. A result that passes both (no \texttt{sorryAx} in the kernel's report, no escape hatch in the lint, every open part named) is certified twice, by channels that cannot fail together in the same way \textbf{[reading]}, with the mechanisms themselves \textbf{[Lean]} and \textbf{[paper]} (claim III.40). Care point: both channels certify form. This is double certainty of form, not of the content of an open part.

\textbf{The finite loss excised at the beginning.} The owner's reading: the finite amount of loss excised at the beginning is the floor at the error bars: a fixed finite cost paid once, which is what makes the rest meaningful. Every act then collapses to one of two: constructible inside (some subset of DMP), or R136a1 \textbf{[reading]} (claim III.42). Nothing in this section makes a theory with arithmetic complete. The completeness is that of the generalisation over what is inside at each tick \textbf{[classical]} for the incompleteness side (claim III.43).

\medskip\noindent\rule{\linewidth}{0.4pt}\medskip

\subsection*{III.6 The tick operator}
\addcontentsline{toc}{subsection}{III.6 The tick operator}

\textbf{Transitivity reaches every finite scale.} Reaching composes. From a certified start, every finite scale is reached by finitely many ticks, each one settling. That every tick past the start pays is \textbf{conditional} on the established process (section 4.4); the composition of steps that pay is a ledger again \textbf{[Lean]} (\texttt{Unified/Encoder.lean}, \texttt{pays\_of\_factors}; \texttt{Unified/Ledger.lean}, \texttt{carry}; \texttt{ETM/RunSys.lean}, \texttt{RunSys.carry}, \texttt{RunSys.contract}; \texttt{Margin.survive}) (claim III.44). For GRH the start is \texttt{\{infinity, 2, 3\}}: Weil's form on supports up to a = 0.69 is positive there, with a certified margin c = 1.5e-9, 3.4e-11, 6.4e-13 at a = 0.60, 0.65, 0.69 \textbf{[numerical]} (certified in ball arithmetic; \texttt{docs/unified/GRH-ONEBOX-CERT.md}; the abstract step \texttt{GRH/OneBoxCert.lean} is \textbf{[Lean]}) (V.3).

\textbf{The continuum needs a uniform bound.} Transitivity reaches every finite target. The continuum is a limit, not a finite target. The margin carried by ticks is c\_0 times the product of the factors. A uniform lower bound on the factors keeps it positive at every finite scale, decaying at most exponentially \textbf{[Lean]} (\texttt{Margin.survive}, \texttt{survive\_pos}; \texttt{GRH/Transfer.lean}, \texttt{margin\_pos\_of\_factors}) (claim III.45). A \emph{uniform margin}, an infimum above zero over all scales, is false for zeta's zero measure: the zero side of the explicit formula is purely atomic, so the product of the Szego factors is exp(integral of log w), which is zero \textbf{[classical]} (Szego-Kolmogorov-Krein) (claim III.45b). So what the continuum needs is a uniform bound on the tick factors. Which bound it is, is the open step of each card: the lower edge for GRH, the upper edge for O1 \textbf{[open]} (claim III.46).

\textbf{The tick is an act at a DB, rotating.} The owner's reading: the multiplicative transport is the tick operator. A tick is an act at a boundary, and its rotation is the encoder's half-turn \textbf{[reading]} (claim III.47). The half-turn is exact (II.2): the electric profile is the Moebius map of a symmetric matrix that squares to the identity, with antipodal fixed points \textbf{[Lean]} (\texttt{Unified/GRHLens.lean}, \texttt{halfTurn\_mobius}, \texttt{halfTurn\_transpose}, \texttt{halfTurn\_mul\_self}, \texttt{fixed\_antipodal}). As a Moebius map of the sphere that matrix is a rotation by 180 degrees \textbf{[classical]}.

\textbf{Birkhoff-Hopf, and where it does not apply.} A positive map whose image cone has finite Hilbert projective diameter Delta contracts the projective distance by tanh(Delta/4) per application \textbf{[classical]} (Birkhoff, Hopf) (claim III.48). On the Yang-Mills side the slab factor is tanh(beta N) per layer, below 1 at every coupling, with no join loss \textbf{[Lean]} (\texttt{Unified/JoinLossSlab.lean}, \texttt{kchain\_tanh}) (claim III.49). The cutoff-uniform form of that bound is O1 \textbf{[open]} (claim III.49b). On the GRH side the lead does not apply as stated. Weil's form at growing support does not factor into ticks of a positive map. First, growing the support is restriction to nested sections, not iteration: a margin at scale k + 1 is one at scale k \textbf{[Lean]} (\texttt{sec\_restrict}; claim III.22). Second, the Szego recursion does factor the margin into ticks, but each step is an SU(1,1) isometry of the cone, whose image is the whole cone: Delta is infinite, and Birkhoff-Hopf gives the factor tanh(infinity) = 1, no contraction \textbf{[classical]} (claim III.22b). Third, primes enter additively, with indefinite forms, so each prime consumes margin and is not a positive map composed with the previous one \textbf{[classical]} (claim III.22c). A symmetry exchanges pictures. It does not turn an upper bound on one encoder into a lower bound on another \textbf{[paper]} (\texttt{docs/unified/GRH-LENS.md}, 6.3).

\textbf{What the act carrier settled.} The owner answered the composition question with the act carrier (I.4): a certificate passes without loss when a conserved quantity carries it, and ``the uniformity is given by the arrow of causality''.

\begin{itemize}
\item \textbf{Yang-Mills: reduced to one number at one scale.} The conserved carrier (centre flux, Gauss's law on both sides of every DB) is in hand \textbf{[Lean]}; superadditivity along a chain \textbf{[Lean]}; uniformity along it \textbf{[classical]}; physical units \textbf{[Lean]}. The grid needs the influence to run in tubes, and tubes need F\_e above the path entropy: O1 \textbf{[open]}.
\item \textbf{GRH: the carrier is identified, and the step is located.} On the constructed space the scaling flow is an isometry; what remains are the zeros and the orbit lengths log p \textbf{[classical]}. The sign is global and all-or-nothing \textbf{[numerical]} + \textbf{[Lean]}.
\item \textbf{NP vs coNP: the carrier is identified, and the step is located.} The parity a proof line names carries the crossing (Tseitin) \textbf{[Lean]}; in cutting planes the crossing is forced and speakable in one line \textbf{[Lean]}; converting it into size is the open step.
\end{itemize}
The summary of these three bullets is \textbf{[paper]} (on IV.3 to IV.5) (claim III.3). In DMP terms all three ask how a certified act hands its certificate to the next act without loss, so that remnants accumulate into existence instead of eroding it. The three answers so far: in Yang-Mills, existence must outrun possibility by the path entropy; in GRH, every building block must enter at once, at its exact place; in NP vs coNP, the pair must be charged at its interface \textbf{[reading]} (claim III.3b).

\subsubsection*{The rotatable hash: the root chain closes through time}
\addcontentsline{toc}{subsubsection}{The rotatable hash: the root chain closes through time}

DMP and Gravity, the two repositories of the root chain, pin each other by hash. Gravity's side is read here: a \texttt{PIN} records the sha256 of the pinned copy of DMP, and the pin check turns any change of it into a failing check (\texttt{vendor/dmp/PIN}, \texttt{tools/check-pin.sh}). That DMP pins Gravity in turn is the notes' statement and was not re-read for this draft. So every pinned existence is a claim the gate can refute \textbf{[paper]} (claim III.50). A literal circle of hashes cannot be constructed: a file that contains its own hash, or two files that contain each other's current hashes, is a fixed point of sha256, which no one can compute \textbf{[classical]} (by the hash's design) (claim III.51). In the repositories the circle closes through time instead. Each side pins an earlier commit of the other, so the chain is a spiral of versions, one tick per pin \textbf{[paper]} (claim III.50). Care point: the circle is consistent only because it never closes at one instant. A pin to the other side's \emph{current} commit would be the uncomputable fixed point.

The owner corrected the first wording: the hash is \emph{rotatable}, not negatable. The hash itself is one-way \textbf{[classical]}. What rotates is the pin: across versions the pin passes from one node to the other, a turn per tick, the same half-turn shape as the encoder's, where two half-turns return to the start and the start is one version later \textbf{[reading]} (claim III.52). Care point: no hash function is rotated, the roles of pinner and pinned are. The owner adds that inside the system the time-to-live of a certification is not a separate mechanism. Every certification is re-made at each tick, against the same shared root, and never reaches back, so the validity window of any hash is bounded by the interval between two ticks \textbf{[reading]} (claim III.53). What leaves the inside is outside by definition, so its lifetime is the boundary's matter, not the rotation's.

\medskip\noindent\rule{\linewidth}{0.4pt}\medskip

\subsection*{III.7 Uncertainty transported across the DB}
\addcontentsline{toc}{subsection}{III.7 Uncertainty transported across the DB}

The owner's approach for the content layer: ``Exactly define the uncertainty, then transport it in a meaningful way across a DB.'' Each column keeps its uncertainty in one named bridge to the real object; the finite model asserts only itself. That uncertainty is itself a carried quantity, so the ledger applies to it with one more term (\texttt{docs/unified/TRANSPORT.md}).

\textbf{The transport theorem is stated once, in II.1.} An uncertainty crosses each step with a factor and may gain fresh uncertainty: it stays bounded when the steps contract (II.1.1), is carried without damping at cost = floor (II.1.2), preserves a sign when it stays below the margin (II.1.3), and with no fresh uncertainty is the ledger itself (II.1.4). Where the uncertainty and the margin it is compared with pay the same cost, their ratio crosses at exactly cost = floor: damping relative to a margin exists only if the margin shrinks more slowly than the uncertainty (II.1.5). The five statements were X.7a to X.7e in draft 2.

\textbf{Where each uncertainty sits, and what the open step needs in these terms.}

\begin{itemize}
\item \textbf{Yang-Mills (claims III.7f, III.7g).} The uncertainty is the boundary uncertainty of a slab: how large the charged (twisted) sector can be against the neutral one, at any outside; O1 is the margin 1/mu. In the slab geometry the slab chain carries it with each layer's own factor and adds nothing (zero join loss), and any positive twist ratio clears O1 with enough layers on a fixed lattice \textbf{[Lean]}. The open step: a cost above the floor at fixed physical length, uniform in the cutoff (the twist ratio bounded below as the lattice spacing goes to zero), plus, in the box geometry, the side seam added at each join. For three slabs at the self-dual factor the seam must stay below (1/mu - (sqrt 2 - 1)\textasciicircum{}3)/sqrt 2 per join \textbf{[open]}.
\item \textbf{GRH (claims III.7h, III.7i).} The uncertainty is how far zeta's margin (the Szego prediction error at scale n) is from the finite model's, relative to the model's; the explicit formula's remainder enters each scale as the relative difference of the two factors 1 - \textbar{}k\textbar{}\textasciicircum{}2. In log form it crosses at exactly cost = floor: zero slack. Zeta's margin is certified, with size, while the summed remainders stay below log 2; the sign alone is decided step by step \textbf{[Lean]}. The open step: remainders, measured against a margin that collapses, that are summable with total below log 2, from the primes alone, beyond Dirichlet \textbf{[open]}.
\item \textbf{NP vs coNP (claims III.7j, III.7k, III.7l).} The uncertainty is the hidden information per variable, the pointer of the gadget, log m bits. Each hidden variable a refutation must resolve takes a fixed share of the line count across the two-party split. That is a ledger with no fresh uncertainty, and it is exactly a size bound; lifting supplies the cost c log n per variable \textbf{[Lean]} from \textbf{[external]}. The open step: a constant cost c \textgreater{} 0 per hidden variable for a constant-size gadget (\texttt{LiftCPConst}) \textbf{[open]}. In the known proofs the per-block hiding errors add up across all blocks with no damping, which is why the gadget grows with n \textbf{[reading]} (recalled; to be checked against Mertz-Pitassi and Lovett-Meka-Mertz-Pitassi-Zhang).
\end{itemize}
\textbf{Side by side (claim III.7m).} \texttt{three\_bridges} (\texttt{Unified/All.lean}) states the three columns in one theorem, each as ``named hypothesis =\textgreater{} the uncertainty is transported below the margin and the sign holds'' \textbf{[Lean]}. It is the transport counterpart of \texttt{three\_instances}. The theorem is not empty: on a perturbed Doeblin chain every hypothesis is proved, and the uncertainty stays below the margin at every step; without mixing, a perturbation of the same size erases the margin (claim III.7o) \textbf{[Lean]}.

\textbf{The running-state axiom (claim III.7n).} Each named hypothesis says that its object is in a running state. With that statement taken as a named axiom R, the three conclusions follow: the boundary uncertainty stays below 1/mu, zeta's margin keeps its sign, the crossing cost is paid \textbf{[Lean]} for the implication. R for the specific objects (the Wilson measure in the continuum, the L-functions beyond Dirichlet, strong proof systems) is exactly the open step of each problem \textbf{[reading]}. This is the structural sense of section 4.2.

\medskip\noindent\rule{\linewidth}{0.4pt}\medskip

\subsection*{III.8 What this part adds up to}
\addcontentsline{toc}{subsection}{III.8 What this part adds up to}

\begin{enumerate}
\item Every card needs the same container, a containing existence with free possibility space (section 1). The cards share one geometry (section 2) and one slot (section 4). What each card asks of the geometry differs, so no single statement proves them all.
\item The open step of each card is the same assertion in the card's terms: ``this object runs'', with a floor, uniformly (sections 4.2 and 6).
\item The proved obstructions say why assembly with slack cannot supply it (the table in section 4.1).
\item The dichotomy of II.3 says what would follow: with a floor, a settling bound; without it, no decision procedure (section 4.3).
\end{enumerate}
When a step closes, the change to the paper is local: the status in \texttt{../CLAIMS.md} moves from open to Lean (or paper), the section's verdict line changes, one cell of the table in section 4.1 changes, and the abstract's list of what remains shrinks by one item. No other section needs to move. That is the test of the paper's structure.


%% ---- 04-part-IV-yang-mills.md
\section*{Part IV. Yang-Mills: the gap mechanism reduces to one box of fixed physical size}
\addcontentsline{toc}{section}{Part IV. Yang-Mills: the gap mechanism reduces to one box of fixed physical size}

\subsection*{IV.0 What is proved and what is not}
\addcontentsline{toc}{subsection}{IV.0 What is proved and what is not}

\textbf{The Clay statement} (Jaffe-Witten): for every compact simple gauge group G, a quantum Yang-Mills theory on R\textasciicircum{}4 exists, satisfying the Osterwalder-Schrader or Wightman axioms with local gauge-invariant fields, and has a mass gap Delta \textgreater{} 0.

\textbf{This part does not prove it.} For SU(2) it proves two things.

\begin{enumerate}
\item \textbf{The minting half} (I.5): from a gap certificate at one box scale to the mass gap of the lattice theory in both orders of limits, with reflection positivity of the limit \textbf{[Lean]}, conditional on stated inputs (IV.2).
\item \textbf{The global structure of the pressing half}: a per-box rate in one of two channels, a conserved carrier that no collar can screen, superadditivity along a chain, uniformity along it, and the conversion to physical units. Together they reduce the \textbf{gap mechanism} to one statement about one box \textbf{[Lean]} + \textbf{[classical]} (IV.3). Sections IV.4 and IV.5 show why that box statement is harder for a grid of boxes than for a chain, and sharpen it to one number.
\end{enumerate}
It leaves four things open.

\begin{itemize}
\item \textbf{(O1) The one-box statement} (IV.5.6). There is a box of fixed physical size at which the electric flux free energy exceeds the path entropy of the four-dimensional lattice after the neutral excess act, \textbf{F\_e(L) - 2c(L) \textgreater{} log mu\_4 = 1.91}, uniformly over every well-defined outside. As E = mc\textasciicircum{}2 is E\textasciicircum{}2 = (pc)\textasciicircum{}2 + (mc\textasciicircum{}2)\textasciicircum{}2 with the momentum term dropped, the familiar form drops the excess act and reads F\_e(L) \textgreater{} 1.91 (twist ratio r(L) \textgreater{} 0.74); it is weaker and alone does not give the chain (IV.5.8). In effective-string units the familiar form asks for a box of side L sqrt(sigma) of about 1.7 to 1.8 (IV.7.4, \textbf{[numerical]}); the full form asks for at least that.
\item \textbf{(O2) The identification inputs}: the line of constant physics (OnLCP) and non-triviality, the two hypotheses that make the lattice family Yang-Mills and not some other gapped family (IV.9). The gap itself needs of the line of constant physics only its resolution clause (IV.9.5, \textbf{[Lean]}). Non-triviality holds per box by construction: the reflected correlation at even time is a reflection-positivity square, so it is at least zero at every level (IV.9.5a, \textbf{[Lean]}). What remains is the scale-uniform lower edge and the scheme clause (IV.9.5b, \textbf{[open]}), the opposite edge to O1's decay. It is not a second assumption: the ODP ends only at a dimensionless 1, so a running, DMP-complete system carries the lower edge by construction, and IV.9.5b is the lower edge of the one assumption, beside O1's upper edge (Part III 4.4, \textbf{[reading]}). It is placed exactly at the ODP's end, the certified start \texttt{\{infinity, 2, 3\}}; from there it follows for a DMP-complete system with nothing more to show, so it is \textbf{conditional} on the one assumption, as the direction of every later tick is (III.66).
\item \textbf{(O3) The continuum field theory}: Euclidean invariance, local fields as distributions, uniqueness of the limit and Wightman reconstruction. For the external construction this project uses (IV.9) these are claimed there \textbf{[external]}. For this project's own objects, rotation invariance cannot yet be stated, since the lattice has only the hypercubic group \textbf{[open]}.
\item \textbf{(O4) Other groups}: the Lean setting has a tuning for SU(2) only. The SU(N) package and the audit of what generalises are partial (\texttt{ProofPhase/General/}, \texttt{proof-phase/GENERALISE-AUDIT.md}).
\end{itemize}
The title's ``one box'' is exact for the gap mechanism (O1). It is not true of the Clay statement as a whole, which also needs O2 to O4.

\textbf{What changed since draft 1.} Draft 1 stated O1 as ``the loss rate of one box is below one''. The crossover rounds of the same day showed that the per-box rate 0.881 suffices for a chain of boxes but not for a grid (IV.4), and that it suffices for a grid exactly when the influence runs in tubes, which needs a larger electric cost per box (IV.5). O1 is now one number at one scale.

\subsection*{IV.1 The lattice theory and its positivity (G0 by construction)}
\addcontentsline{toc}{subsection}{IV.1 The lattice theory and its positivity (G0 by construction)}

Yang-Mills is the one problem of the three whose starting positivity is given. The Wilson weight e\textasciicircum{}\{-S\} is positive and local, and from those two facts reflection positivity follows. Everything in this section is machine-checked.

{\footnotesize
\begin{longtable}{>{\raggedright\arraybackslash}p{\dimexpr 0.1016\linewidth-2\tabcolsep\relax}>{\raggedright\arraybackslash}p{\dimexpr 0.3712\linewidth-2\tabcolsep\relax}>{\raggedright\arraybackslash}p{\dimexpr 0.1853\linewidth-2\tabcolsep\relax}>{\raggedright\arraybackslash}p{\dimexpr 0.3220\linewidth-2\tabcolsep\relax}}
\toprule
\textbf{\#} & \textbf{Result} & \textbf{Status} & \textbf{Source} \\
\midrule
\endhead
IV.1.1 & Wilson measure on the 4D torus for SU(2) and general compact G & \textbf{[Lean]} & \texttt{AreaLaw/ConcreteSUn.lean}, \texttt{OneScale/} \\
IV.1.2 & Reflection positivity, site and link & \textbf{[Lean]} & \texttt{Positivity/Wilson.lean} (\texttt{wilson\_rp}), \texttt{Positivity/LinkWilson.lean}, \texttt{Positivity/Concrete.lean} (\texttt{wilson\_rp\_SU2}) \\
IV.1.3 & Temporal-gauge transfer kernel; centre flip of one row is a symmetry; sector identity Z\_e = (Z + Z\_tw)/2 & \textbf{[Lean]} (identity definitional) & \texttt{Attempts/R4/LayerKernel.lean}, \texttt{LayerSector.lean} \\
IV.1.4 & 't Hooft duality: the lattice rotation maps the temporal twist to the spatial twist on the symmetric box & \textbf{[Lean]} & \texttt{ProofPhase/RouteB/Duality.lean}, \texttt{tHooft\_rotation} \\
IV.1.5 & A good blocking: centre-covariant and reflection-positive at every level & \textbf{[Lean]}; non-trivial limit at fixed L1 \textbf{[paper]} & \texttt{ProofPhase/RouteA/Blocking.lean}, \texttt{goodBlocking\_rp}; \texttt{proof-phase/ws1-routeA/NOTE.md} \\
IV.1.6 & Strict flip bound Z\_tw \textless{} Z for every beta != 0, every box, every twist position (``existence = non-existence + 1'', I.7) & \textbf{[Lean]} & \texttt{ProofPhase/Flip/Strict.lean}, \texttt{su2\_Ztw\_lt\_Z} \\
IV.1.7 & The one-step transfer kernel at fixed cross-section is strictly positive & \textbf{[Lean]} & \texttt{ProofPhase/Flip/Bound.lean}, \texttt{stepK\_pos} \\
\bottomrule
\end{longtable}}

\subsection*{IV.2 The minting half: from a certificate to the gap}
\addcontentsline{toc}{subsection}{IV.2 The minting half: from a certificate to the gap}

Suppose one box, at one scale, is known to forget its outside at a fixed rate. Then the whole lattice theory has a mass gap, in both orders of limits, and the reconstructed Hilbert space carries a transfer operator with a spectral gap. That implication is the minting half, and it is done.

{\footnotesize
\begin{longtable}{>{\raggedright\arraybackslash}p{\dimexpr 0.0977\linewidth-2\tabcolsep\relax}>{\raggedright\arraybackslash}p{\dimexpr 0.3819\linewidth-2\tabcolsep\relax}>{\raggedright\arraybackslash}p{\dimexpr 0.1586\linewidth-2\tabcolsep\relax}>{\raggedright\arraybackslash}p{\dimexpr 0.3418\linewidth-2\tabcolsep\relax}}
\toprule
\textbf{\#} & \textbf{Result} & \textbf{Status} & \textbf{Source} \\
\midrule
\endhead
IV.2.1 & Dobrushin uniqueness with exponential weights; D' gives decay of block observables at one rate (abstract Gibbs systems) & \textbf{[Lean]} & \texttt{Attempts/R3/TailDecay.lean} (\texttt{theoremA\_w}), \texttt{Attempts/R3/Robust.lean} (\texttt{gap\_of\_localFactorsW}) \\
IV.2.2 & Block decay converted to physical units (Corollary B) & \textbf{[Lean]} & \texttt{OneScale/CorollaryB.lean}, \texttt{corollaryB\_lattice} \\
IV.2.3 & OS route: a windowed block gap gives MassGapLoc and the transfer operator's gap on the reconstructed space & \textbf{[Lean]} & \texttt{ProofPhase/OS/GapLoc.lean} (\texttt{massGapLoc\_of\_UVTVBlockGap}), \texttt{ProofPhase/M1/UVTVOS.lean} (\texttt{gap\_of\_RP}) \\
IV.2.4 & The top-level assembly: OpenPieces implies MassGapLoc, LimitRP and NonTrivial; reflection positivity of the limit is proved, not assumed & \textbf{[Lean]} & \texttt{ProofPhase/Torus/Top.lean} (\texttt{yang\_mills\_assembly}), \texttt{ProofPhase/Assembly.lean} (\texttt{OpenPieces}) \\
IV.2.5 & The gap in both orders of limits (torus route) & \textbf{[Lean]} & \texttt{ProofPhase/Torus/Top.lean} (\texttt{yang\_mills\_assembly\_torus}); \texttt{ProofPhase/Existence/Bridge.lean} (\texttt{bothOrders\_of\_torusDecay}) \\
IV.2.6 & Non-triviality of the hypotheses: pure noise and every bounded coupling fail OpenPieces & \textbf{[Lean]} & \texttt{ProofPhase/Assembly.lean}, \texttt{not\_trivial} \\
IV.2.7 & MassGapLoc implies a spectral gap of the continuum Hamiltonian (needs density of local vectors and the reconstruction) & \textbf{[paper]}; spectral bridge \textbf{[Lean]} & \texttt{ProofPhase/Bridge/Spectral.lean}, \texttt{ProofPhase/Bridge/OSLink.lean} \\
IV.2.8 & Strong-coupling gap for any compact group (widest window 18 tanh(n beta) \textless{} 1) & \textbf{[Lean]} / \textbf{[paper]} & \texttt{OneScale/SingleLink.lean} \\
\bottomrule
\end{longtable}}

\textbf{The crossover needs no description of its content.} Between the weak-coupling scale L\_* (where M1 controls the law, IV.9) and the box scale of O1 lie finitely many blocking steps, each of finite range. The assembly uses only a block gap at one scale; it never needs to know what happens inside the crossover. So the certificate can be checked directly on the actual blocked law at that scale, with no assumption about how the law got there \textbf{[paper]} (\texttt{PREP.md}, ``The crossover is constrained by making no statement about its content''). The price is the quantifier: since nothing is said about the content, the certificate must hold for every content the crossover could have produced, that is, uniformly over every well-defined outside. That quantifier is already part of O1.

\subsection*{IV.3 The pressing half: what one box gives, and how a chain carries it}
\addcontentsline{toc}{subsection}{IV.3 The pressing half: what one box gives, and how a chain carries it}

This section collects what holds for every box, and what carries it along a line of boxes.

\textbf{A strict sign gives a rate (II.13).} By the strict flip bound (IV.1.6) and 't Hooft's duality F\_e = -log tanh(F\_tw/2), at every box size max(F\_tw, F\_e) \textgreater{}= log(1 + sqrt 2) = 0.881. Rarity's rate is F\_tw, the free energy of a forced centre sheet; disorder's rate is F\_e, the per-box decay of a centre-charged signal. They cannot both be small. \textbf{[Lean]} (\texttt{su2\_handover}).

\textbf{A flux-carrying outside forces at least one sheet.} The flux cells of a box are even in number (Gauss's law mod 2), so a flux carrying outside forces a sheet through the box, crossing the boundary at every flux cell. This is the quantum: ``at least something'' always crosses. \textbf{[Lean]} (\texttt{exists\_sheet\_of\_flux}, \texttt{even\_card\_fluxCells}).

\textbf{The law is the same on both sides.} The flux through a box's boundary computed from inside equals the one computed from outside, and the total flux out of a box is even. \textbf{[Lean]} (\texttt{cob\_restrict\_eq\_cob\_restrict\_compl}, \texttt{sum\_cob\_restrict\_eq\_zero}). Electric influence therefore has no sources inside a box: it propagates along connected paths of boxes (IV.5.3).

\textbf{Superadditivity along the chain.} With the transfer matrix along a chain of boxes, F\_e(a + b) \textgreater{}= F\_e(a) + F\_e(b) - c(a) - c(b), and F\_e(kL) \textgreater{}= k(F\_e(L) - 2C). The boundary term c(L) = log(Z\_0(L)/lambda\_0\textasciicircum{}L) is the neutral excess act, the weight of the neutral excitations over one box length. The per-box rate multiplies if F\_e(L) \textgreater{} 2c(L) at one scale (superadditivity up to the excess terms; the converse is not proved). \textbf{[Lean]} in abstract transfer-matrix form (\texttt{flux\_superadd}, \texttt{flux\_superadd\_iter}).

\textbf{Uniformity along the chain.} Every DB applies the same transfer step (translation invariance in Euclidean time), so the loss rate is the same at every step: a homogeneous Markov chain \textbf{[classical]}; its contraction eta\textasciicircum{}n is \textbf{[Lean]} (\texttt{contracts\_kpow}, II.12).

\textbf{Physical units.} A one-box contraction eta \textless{}= eta\_0 \textless{} 1 at a fixed physical box size l\_* = L\_* a, uniform in the lattice spacing a, gives a physical mass m\_phys \textgreater{}= -log eta\_0 / l\_* \textbf{[Lean]} (\texttt{phys\_rate\_of\_uniform}). For a family of blocked kernels with one floor eps\_* for every spacing, the influence at physical distance D is at most exp(-m D) with m = -log(1 - eps\_*)/l\_* \textbf{[Lean]} (\texttt{rdoeblin\_physical\_decay}).

{\footnotesize
\begin{longtable}{>{\raggedright\arraybackslash}p{\dimexpr 0.1018\linewidth-2\tabcolsep\relax}>{\raggedright\arraybackslash}p{\dimexpr 0.3275\linewidth-2\tabcolsep\relax}>{\raggedright\arraybackslash}p{\dimexpr 0.1828\linewidth-2\tabcolsep\relax}>{\raggedright\arraybackslash}p{\dimexpr 0.3679\linewidth-2\tabcolsep\relax}}
\toprule
\textbf{\#} & \textbf{Result} & \textbf{Status} & \textbf{Source} \\
\midrule
\endhead
IV.3.1 & Per-box hand-over rate: max(F\_tw, F\_e) \textgreater{}= 0.881 at every box size & \textbf{[Lean]} (0 \textless{} Z\_tw a hypothesis) & \texttt{ProofPhase/Crossover/HandOver.lean}, \texttt{su2\_handover} \\
IV.3.2 & A flux-carrying outside forces at least one sheet & \textbf{[Lean]} & \texttt{ProofPhase/Crossover/FluxRate.lean}, \texttt{exists\_sheet\_of\_flux}, \texttt{even\_card\_fluxCells} \\
IV.3.3 & The flux law is the same seen from inside and outside; no sources & \textbf{[Lean]} & \texttt{ProofPhase/Crossover/Dichotomy.lean}, \texttt{cob\_restrict\_eq\_cob\_restrict\_compl}, \texttt{sum\_cob\_restrict\_eq\_zero} \\
IV.3.4 & Superadditivity up to the excess terms; the rate multiplies along a chain if F\_e(L) \textgreater{} 2c(L) at one scale & \textbf{[Lean]} (abstract) & \texttt{ProofPhase/Crossover/CompositeHTF.lean}, \texttt{flux\_superadd}, \texttt{flux\_superadd\_iter} \\
IV.3.5 & Uniformity along the chain & \textbf{[classical]}; eta\textasciicircum{}n \textbf{[Lean]} & translation invariance; \texttt{ProofPhase/Mixing/Chain.lean}, \texttt{contracts\_kpow} \\
IV.3.6 & A one-box contraction at fixed physical size gives a physical mass & \textbf{[Lean]} & \texttt{ProofPhase/Crossover/Neighbours.lean}, \texttt{phys\_rate\_of\_uniform}; \texttt{ProofPhase/Crossover/RDoeblin.lean}, \texttt{rdoeblin\_physical\_decay} \\
IV.3.7 & The doubling law with a fixed loss gives the gap on all larger dyadic boxes & \textbf{[Lean]} (implication) & \texttt{ProofPhase/Flip/Floor.lean}, \texttt{su2\_gap\_of\_doubling\_defect} \\
IV.3.8 & Uncontained sheets weigh nothing (dichotomy part b) & \textbf{[Lean]} & \texttt{ProofPhase/Crossover/Dichotomy.lean}, \texttt{measure\_uncontained} \\
\bottomrule
\end{longtable}}

So along a chain, one box at one scale is enough. The question is what a four-dimensional grid of boxes needs.

\subsection*{IV.4 Why a grid is harder than a chain}
\addcontentsline{toc}{subsection}{IV.4 Why a grid is harder than a chain}

In a chain, flux entering a box has one way out. In a grid it can leave through any of the other 2d - 1 faces, so a box has 2d channels to its neighbours, and Dobrushin's condition needs the \textbf{sum} of the influences over all channels below one (II.12).

\textbf{The hand-over covers two channels, not eight.} The per-box ratio sqrt 2 - 1 = 0.414 keeps a column sum below one for at most two channels \textbf{[Lean]} (\texttt{handover\_column\_sum\_iff}). A grid in d dimensions needs a per-box suppression F \textgreater{} log(2d), that is log 8 = 2.08 in d = 4, more than twice the duality floor \textbf{[Lean]} (\texttt{column\_requirement}).

\textbf{A per-box floor (Doeblin) cannot supply it.} Doeblin's mass per box is at most e\textasciicircum{}\{-K\}/(1 - e\textasciicircum{}\{-K\}) with K = C beta L\_*\textasciicircum{}\{d-1\}, which goes to zero as the lattice spacing does \textbf{[Lean]} (\texttt{neg\_log\_one\_sub\_le}, \texttt{doeblin\_mass\_tendsto\_zero}). For a grid the floor exponent must stay below log(8/7) = 0.134 (column form) \textbf{[Lean]} (\texttt{floor\_column\_requirement}), or below -log(1 - p\_c), about 0.22 in the disagreement-percolation form (the requirement \texttt{floor\_disagreement\_requirement} is \textbf{[Lean]} for any p; the value 0.22 uses the site-percolation threshold of Z\textasciicircum{}4, p\_c about 0.197, \textbf{[classical]} (recalled), claim IV.4.3c). At the weak-coupling scale it is in the tens \textbf{[paper]} (\texttt{RDOEBLIN.md}). It misses by orders of magnitude, and a larger box at fixed coupling makes it worse.

\textbf{Duality gives no floor above log(2d).} For n independent Z/2 sectors 't Hooft's Fourier relation reads (1 + S\_tw)(1 + S\_e) = 2\textasciicircum{}n \textbf{[paper]} (\texttt{TRANSFER.md}). The smaller sum is then forced only below the self-dual value sqrt(2\textasciicircum{}n) - 1 \textbf{[Lean]} (\texttt{min\_le\_of\_sector\_product}), which is at least one as soon as n \textgreater{}= 2 \textbf{[Lean]} (\texttt{selfdual\_sector\_ge\_one}). Duality allows both summed weights to be large and forces no summable suppression. What the summed form does give is the requirement as a quantitative 't Hooft criterion: the summed electric weight is below one iff the summed twist weight exceeds 2\textasciicircum{}\{n-1\} - 1 \textbf{[Lean]} (\texttt{electric\_sum\_lt\_one\_iff}). That is a dynamical statement (near-saturated twists), not supplied by any duality.

\textbf{A gauge-covariant frame carries nothing.} A composite Higgs-type field built from the collar's covariant Laplacian exists on every collar \textbf{[paper]}. But conditioning on a gauge-covariant frame leaves the law of every gauge-invariant observable unchanged \textbf{[Lean]}, finite analogue (\texttt{weight\_fiber\_eq}, \texttt{weight\_fiber\_mul\_card}); \textbf{[paper]} for compact groups. The frame changes nothing (\texttt{HTF.md}).

{\footnotesize
\begin{longtable}{>{\raggedright\arraybackslash}p{\dimexpr 0.0911\linewidth-2\tabcolsep\relax}>{\raggedright\arraybackslash}p{\dimexpr 0.3469\linewidth-2\tabcolsep\relax}>{\raggedright\arraybackslash}p{\dimexpr 0.2069\linewidth-2\tabcolsep\relax}>{\raggedright\arraybackslash}p{\dimexpr 0.3351\linewidth-2\tabcolsep\relax}}
\toprule
\textbf{\#} & \textbf{Result} & \textbf{Status} & \textbf{Source} \\
\midrule
\endhead
IV.4.1 & The hand-over keeps a column below one for at most two channels; a grid needs F \textgreater{} log(2d) = 2.08 & \textbf{[Lean]} & \texttt{ProofPhase/Crossover/Neighbours.lean}, \texttt{handover\_column\_sum\_iff}, \texttt{column\_requirement} \\
IV.4.2 & Doeblin's mass per box goes to zero in the continuum & \textbf{[Lean]} & \texttt{Neighbours.lean}, \texttt{neg\_log\_one\_sub\_le}, \texttt{doeblin\_mass\_tendsto\_zero} \\
IV.4.3 & A grid needs floor exponent below 0.134 (or 0.22, with the recalled threshold p\_c = 0.197); at L\_* it is in the tens & \textbf{[Lean]} (requirement) + \textbf{[paper]} (estimate) & \texttt{ProofPhase/Crossover/RDoeblin.lean}, \texttt{floor\_column\_requirement}, \texttt{floor\_disagreement\_requirement}; \texttt{RDOEBLIN.md} \\
IV.4.4 & Multi-sector duality forces nothing for n \textgreater{}= 2; the quantitative 't Hooft criterion & \textbf{[Lean]} (the relation itself \textbf{[paper]}) & \texttt{ProofPhase/Crossover/Runs.lean}, \texttt{min\_le\_of\_sector\_product}, \texttt{selfdual\_sector\_ge\_one}, \texttt{electric\_sum\_lt\_one\_iff} \\
IV.4.5 & A gauge-covariant frame carries no gauge-invariant information & \textbf{[Lean]} (finite) + \textbf{[paper]} & \texttt{ProofPhase/Crossover/CompositeHTF.lean}, \texttt{weight\_fiber\_eq}, \texttt{weight\_fiber\_mul\_card}; \texttt{HTF.md} \\
IV.4.6 & At typical fields the summed boundary influence drops below one at about 7 cells from the boundary (box of about 15 cells) & \textbf{[numerical]} (typical, not supremum) & \texttt{NEIGHBOURS.md}, 0.2 \\
\bottomrule
\end{longtable}}

\subsection*{IV.5 Containment: the requirement moved inside one box}
\addcontentsline{toc}{subsection}{IV.5 Containment: the requirement moved inside one box}

The owner's correction (``the saturation has to be proven inside confinement and absent outside that confinement, as a containing existence'', \texttt{PREP.md}) reverses the direction of the attack. Instead of building the inside of a box from outside data, prove three things: outside, saturation is absent; the container ends; inside, saturation makes the influence run in tubes, and tubes are chains, where IV.3 closes. The first two close. The third closes negatively at the duality floor and positively above a sharp threshold.

\textbf{IV.5.1 Outside, absence (lemma).} For every box, the twist ratio r = Z\_tw/Z satisfies r \textless{} 1 and r = (1 - e\textasciicircum{}\{-F\_e\})/(1 + e\textasciicircum{}\{-F\_e\}) \textless{}= F\_e \textbf{[Lean]} (\texttt{twist\_ratio\_eq}, \texttt{twist\_ratio\_le\_elecF}, \texttt{elecF\_nonneg}). The twist is rare exactly where electric flux is cheap. In a physically small box small-volume perturbation theory makes F\_e of order g(L)\textasciicircum{}2 \textbf{[classical]} (Luescher; van Baal; recalled, not re-checked), so r = O(g\textasciicircum{}2). The mechanism outside is cheap electric flux, seen from the other side of the duality as a rare twist.

\textbf{IV.5.2 The container ends, in long boxes.} At fixed cross-section the transfer kernel is strictly positive (IV.1.7). By Jentzsch's theorem \textbf{[classical]} the vacuum is the unique top state; it is centre-invariant \textbf{[paper]}, so the flux sector lies a factor mu \textless{} 1 below it per unit length. Then F\_e(T) \textgreater{}= -T log mu - log C grows without bound \textbf{[Lean]} (\texttt{elecF\_lower\_of\_gap}, \texttt{exists\_elecF\_ge}). In every long enough box the rarity branch fails, and by the hand-over the electric channel carries the suppression. What this does not give: mu depends on the cross-section and may tend to one as it grows. In the cubic geometry the end needs F\_e(L) \textgreater{}= 0.881 at one L, that is g\textasciicircum{}2 of order one, beyond perturbative or M1 control \textbf{[open]}; measured near L = 4 at beta = 2.3 (F\_e = 0.91, \textbf{[numerical]}, preliminary).

\textbf{IV.5.3 Inside: what the law gives.} Gauss's law makes electric influence a conserved line with no sources that no collar can screen (IV.3.2, IV.3.3). That is the whole structural content of ``tube'', derived from the law without assuming tubes \textbf{[Lean]} for the pieces, \textbf{[paper]} for the conclusion. The law does not say which path carries the influence: it allows every path and every branching.

\textbf{IV.5.4 Inside: what the cost must give.} Bound the influence between two boxes at distance D by a sum over paths of boxes, each box costing e\textasciicircum{}\{-F\}. With mu continuations per step the sum is a geometric series in mu e\textasciicircum{}\{-F\}, which converges iff F \textgreater{} log mu \textbf{[Lean]} (\texttt{path\_sum\_summable\_iff}, \texttt{tube\_threshold}). Self-avoiding paths on Z\textasciicircum{}d have connective constant mu\_d \textgreater{}= d, and for mu \textgreater{}= 3, mu (sqrt 2 - 1) \textgreater{}= 1 \textbf{[Lean]} (\texttt{path\_entropy\_beats\_handover}). \textbf{So in every d \textgreater{}= 3 the duality floor 0.881 loses to path entropy: the cheapest chain does not dominate.} For d = 4, mu\_4 = 6.77 and log mu\_4 = 1.91 \textbf{[classical]} (recalled); rigorously log mu\_4 lies between log 4 = 1.39 and log 7 = 1.95. Equivalently, the single-sector twist ratio must exceed (mu - 1)/(mu + 1) = 0.74.

\textbf{IV.5.5 What the action adds and what it does not.} The Bogomolny bound S \textgreater{}= 8 pi\textasciicircum{}2 \textbar{}Q\textbar{} / g\textasciicircum{}2 \textbf{[classical]} gives, on a twisted torus with odd twist product, the floor S \textgreater{}= pi\textasciicircum{}2 beta for admissible fields \textbf{[classical]} (Luescher's admissibility). It strengthens IV.5.1 for multi-plane twist sectors \textbf{[paper]}. It does not help the single-plane twist used here (that sector contains flat configurations, so its floor is zero), nor inside (it bounds twists from above, while O1 needs them high) \textbf{[paper]}. At a fixed physical box size the position entropy 4 log(L/a) of a fractional instanton beats its action pi\textasciicircum{}2 beta = (11/3) log(1/(a Lambda)) as a -\textgreater{} 0 \textbf{[paper]}, one-loop sketch (determinants and measure not checked).

\textbf{IV.5.6 The one-box statement (O1).} Combining IV.5.3 and IV.5.4: law (structure) plus cost above entropy (dominance) gives tubes, and tubes are chains, where IV.3 applies. What remains is:

\begin{quote}
\textbf{(O1)} At one box of fixed physical size, F\_e(L) - 2c(L) \textgreater{} log mu\_4 = 1.91, uniformly over every well-defined outside, with c(L) the neutral excess act of IV.3. \textbf{[open]}
\end{quote}

The familiar form F\_e(L) \textgreater{} 1.91 (equivalently r(L) \textgreater{} 0.74) drops the excess act, as E = mc\textasciicircum{}2 drops the momentum term of E\textasciicircum{}2 = (pc)\textasciicircum{}2 + (mc\textasciicircum{}2)\textasciicircum{}2. The drop is only seeming: c(L) \textgreater{}= 0, so the full form implies the familiar one, but the familiar one alone does not give the chain criterion (IV.5.8, examples in Lean).

This is the grid requirement log 8 = 2.08 of IV.4.1 refined from column sums to path counting. Containment moves the requirement inside the box but does not lower it to the chain value 0.881.

\textbf{Care points, settled in draft 3.}

\begin{itemize}
\item \textbf{O1 in one form.} O1 is \texttt{FineOneBox} (\texttt{OneBox/Fine.lean}): one box at the physical cell scale, influence bounds under every boundary condition, uniform in level and volume. \texttt{WilsonOneBox} and \texttt{FineOneBoxSplit} imply it, and it feeds \texttt{yang\_mills\_assembly\_onebox}; \texttt{BlockGapBddW} belongs to the separate blocking route (IV.5.9, \textbf{[Lean]}).
\item \textbf{O1 in the one-box test.} In the test of Part I (I.8, the rotation API) O1 is exactly DMP-completeness of the Yang-Mills box: the chain form holds exactly when some margin makes every chain of outsides DMP-complete, and \texttt{FineOneBox} holds exactly when the Yang-Mills box, with its geometry fixed across levels, is DMP-complete from its start level \textbf{[Lean]} (IV.10c, IV.10d). Both directions are proved, so the general test carries exactly O1, no more and no less. That the box is DMP-complete is O1 itself \textbf{[open]} (IV.10e).
\item \textbf{Per-box costs multiply.} In that per-outside form the composition of the per-box influences is a theorem, so no annealed quantity enters the gap; the annealed-to-quenched identification is needed only to read a torus measurement of F\_e as evidence about O1 (IV.5.7, \textbf{[Lean]}).
\item \textbf{Chain and tube.} In the transfer-matrix model the chain criterion F\_e(L) \textgreater{} 2c(L) and the tube threshold F\_e(L) \textgreater{} log mu are independent, with examples both ways; the joint form gives both (IV.5.8, \textbf{[Lean]}, model).
\end{itemize}
{\footnotesize
\begin{longtable}{>{\raggedright\arraybackslash}p{\dimexpr 0.1077\linewidth-2\tabcolsep\relax}>{\raggedright\arraybackslash}p{\dimexpr 0.3460\linewidth-2\tabcolsep\relax}>{\raggedright\arraybackslash}p{\dimexpr 0.2168\linewidth-2\tabcolsep\relax}>{\raggedright\arraybackslash}p{\dimexpr 0.3095\linewidth-2\tabcolsep\relax}}
\toprule
\textbf{\#} & \textbf{Result} & \textbf{Status} & \textbf{Source} \\
\midrule
\endhead
IV.5.1a & r \textless{} 1, r = (1 - e\textasciicircum{}\{-F\_e\})/(1 + e\textasciicircum{}\{-F\_e\}) \textless{}= F\_e, F\_e \textgreater{}= 0 & \textbf{[Lean]} & \texttt{ProofPhase/Crossover/Containment.lean}, \texttt{twist\_ratio\_eq}, \texttt{twist\_ratio\_le\_elecF}, \texttt{elecF\_nonneg} \\
IV.5.1b & Small box: F\_e = O(g(L)\textasciicircum{}2) & \textbf{[classical]} (recalled) & \texttt{CONTAINMENT.md}, 1 \\
IV.5.2a & At fixed cross-section F\_e(T) grows without bound & \textbf{[Lean]} core + \textbf{[classical]} (Jentzsch) + \textbf{[paper]} (centre invariance) & \texttt{Containment.lean}, \texttt{elecF\_lower\_of\_gap}, \texttt{exists\_elecF\_ge}; \texttt{CONTAINMENT.md}, 2 \\
IV.5.2b & Cubic boxes: F\_e(L) \textgreater{}= 0.881 at some L & \textbf{[open]}; near L = 4 at beta 2.3 \textbf{[numerical]} & \texttt{CONTAINMENT.md}, 2 \\
IV.5.3 & Gauss's law makes influence path-like (no sources, unscreenable) & \textbf{[Lean]} pieces + \textbf{[paper]} & IV.3.2, IV.3.3; \texttt{CONTAINMENT.md}, 6.1 \\
IV.5.4a & Path sums converge iff F \textgreater{} log mu; 0.881 loses for mu \textgreater{}= 3 & \textbf{[Lean]} & \texttt{Containment.lean}, \texttt{path\_sum\_summable\_iff}, \texttt{tube\_threshold}, \texttt{path\_entropy\_beats\_handover} \\
IV.5.4b & mu\_4 = 6.77, log mu\_4 in [1.39, 1.95] & \textbf{[classical]} & \texttt{CONTAINMENT.md}, 3 \\
IV.5.5 & Bogomolny floor helps outside only and loses to entropy as a -\textgreater{} 0 & \textbf{[classical]} (bound) + \textbf{[paper]} (sketch) & \texttt{CONTAINMENT.md}, 6.2 \\
IV.5.6 & \textbf{(O1)} F\_e(L) - 2c(L) \textgreater{} log mu\_4 at one box of fixed physical size, uniformly over outsides (familiar form F\_e(L) \textgreater{} log mu\_4, the excess act dropped) & \textbf{[open]} & \texttt{CONTAINMENT.md}, 0 and 5; \texttt{docs/OPEN-QUESTIONS-DMP.md}, 0 \\
IV.5.7 & Per-box costs multiply over the lattice in O1's per-outside form; annealed to quenched only for reading measurements & \textbf{[Lean]} & \texttt{OneBox/Fine.lean}, \texttt{fineDecayInv\_of\_fineOneBox}; \texttt{Closures/OneBoxForm.lean} \\
IV.5.8 & Chain criterion and tube threshold independent in the model; the joint form gives both & \textbf{[Lean]} (model) & \texttt{Closures/ChainJoint.lean}, \texttt{joint\_gives\_both}, \texttt{tube\_not\_chain}, \texttt{chain\_not\_tube} \\
IV.5.9 & O1 in one Lean form: \texttt{FineOneBox}, with the other forms linked & \textbf{[Lean]} & \texttt{OneBox/Fine.lean}; \texttt{Closures/OneBoxForm.lean} \\
\bottomrule
\end{longtable}}

\subsection*{IV.6 The one box from inside: saturation and the dichotomy}
\addcontentsline{toc}{subsection}{IV.6 The one box from inside: saturation and the dichotomy}

\textbf{Saturation from what we have.} Configuration by configuration the twist changes the weight by at least e\textasciicircum{}\{-2 beta N\}, N = L\textasciicircum{}2 flipped plaquettes, so Z\_tw/Z \textgreater{}= e\textasciicircum{}\{-2 beta L\textasciicircum{}2\} \textbf{[Lean]} (\texttt{ratio\_floor\_of\_pointwise}, \texttt{exp\_floor\_ge\_iff}). That saturates the twist only for beta L\textasciicircum{}2 \textless{}= 0.42: strong coupling in tick units, vanishing in the continuum. Conversely a large electric free energy forces a saturated twist \textbf{[Lean]} (\texttt{twist\_ratio\_of\_elecF}), so saturation is a consequence of O1, not an independent input.

\textbf{Bond-moving bounds localise Tomboulis's step, they do not avoid it.} Rigorous one-sided bounds on Z and on the even sector Z\textasciicircum{}+ exist, and the upper decimation flows to strong coupling for every beta in d \textless{}= 4 \textbf{[classical]} (Ito). A lower bound on the twist ratio needs the two normalisations matched; the upper one is at least one \textbf{[Lean]} (\texttt{jensen\_normalization}) and extensive in the box volume. In the one-box form the mismatch is finite but vanishes at fixed physical box size \textbf{[paper]}, sketch. So the step needed is a lower bound on the box partition function matched to the upper decimation's normalisation within O(1): Tomboulis's interpolation step (2007; gaps pointed out by Ito and Seiler; not accepted as a proof \textbf{[classical]}), now localised to one box.

\textbf{The dichotomy inside one box.} Part (a), that a sheet plaquette always costs something, fails uniformly: at the SO(3) cut (tr U\_p = 0) a sheet plaquette is free \textbf{[paper]}. What survives: for every outside, the box contracts, or one of two bad events holds (near-degenerate spanning surfaces for a flux-carrying outside; an SO(3) cut along a would-be sheet) \textbf{[paper]}. Rarity and saturation are disjoint, with a window between them of about 0.18 to at least 1.79 nats per plaquette \textbf{[Lean]} (\texttt{window\_gap}) + \textbf{[paper]} (numbers).

{\footnotesize
\begin{longtable}{>{\raggedright\arraybackslash}p{\dimexpr 0.0995\linewidth-2\tabcolsep\relax}>{\raggedright\arraybackslash}p{\dimexpr 0.3555\linewidth-2\tabcolsep\relax}>{\raggedright\arraybackslash}p{\dimexpr 0.2313\linewidth-2\tabcolsep\relax}>{\raggedright\arraybackslash}p{\dimexpr 0.2937\linewidth-2\tabcolsep\relax}}
\toprule
\textbf{\#} & \textbf{Result} & \textbf{Status} & \textbf{Source} \\
\midrule
\endhead
IV.6.1 & Tick-unit floor Z\_tw/Z \textgreater{}= e\textasciicircum{}\{-2 beta L\textasciicircum{}2\}; saturates only for beta L\textasciicircum{}2 \textless{}= 0.42 & \textbf{[Lean]} & \texttt{ProofPhase/Crossover/Saturation.lean}, \texttt{ratio\_floor\_of\_pointwise}, \texttt{exp\_floor\_ge\_iff} \\
IV.6.2 & A large electric free energy forces a saturated twist & \textbf{[Lean]} & \texttt{Saturation.lean}, \texttt{twist\_ratio\_of\_elecF} \\
IV.6.3 & Bond-moving: the one-box form localises the interpolation step (Tomboulis) to one box & \textbf{[paper]} (sketch); Jensen normalisation \textbf{[Lean]} & \texttt{SATURATION.md}, 0; \texttt{Saturation.lean}, \texttt{jensen\_normalization} \\
IV.6.4 & Kanazawa's bracket 1/2 \textless{} Z\textasciicircum{}+/Z \textless{} 1 is the non-strict form of IV.1.6 & \textbf{[classical]}; strict form \textbf{[Lean]} & \texttt{SATURATION.md}, 0 \\
IV.6.5 & Dichotomy part (a) fails uniformly at the SO(3) cut & \textbf{[paper]} & \texttt{DICHOTOMY.md}, 0.2 \\
IV.6.6 & Reduction to two bad events & \textbf{[paper]} & \texttt{DICHOTOMY.md}, 0.4 \\
IV.6.7 & Window between rarity and saturation, 0.18 to \textgreater{}= 1.79 nats per plaquette & \textbf{[Lean]} (disjointness) + \textbf{[paper]} (numbers) & \texttt{ProofPhase/Crossover/Dichotomy.lean}, \texttt{window\_gap}; \texttt{DICHOTOMY.md}, 0.3 \\
IV.6.8 & Y2 to Y5: inside residue, hand-over through the window, condition A uniformly (m about 4 measured on typical fields), uniformity in volume & \textbf{[open]} (pieces of O1); per-volume gap \textbf{[Lean]} conditional on H2-all & \texttt{docs/OPEN-QUESTIONS-DMP.md}, 2; \texttt{ProofPhase/Quantum/PerVolume.lean}, \texttt{perVolumeGap} \\
IV.6.9 & IR column at L1: W1 column certified at 0.99988 at L sqrt(sigma) = 2.80 (computer-assisted certificate as a named hypothesis) & \textbf{[open]}; implication \textbf{[Lean]} & \texttt{ProofPhase/Certified/Measured.lean}; \texttt{proof-phase/c9b-adjoint/ROUTE3.md} \\
\bottomrule
\end{longtable}}

\subsection*{IV.7 Measured: what the numbers say about O1}
\addcontentsline{toc}{subsection}{IV.7 Measured: what the numbers say about O1}

Measurements can refute O1, never establish it (provability rule, \texttt{proof-phase/c9b-adjoint/REDESIGN.md}, section 8). Every error bar is max(statistical floor, resolution floor), the resolution floor of an energy read over Euclidean extent T being 1/T in lattice units; ad-hoc floors are not used; a change of floor after data is reported beside the pre-registered result (\texttt{attempts/RULES.md}, last section).

\textbf{IV.7.1 Small tori (preliminary).} At beta = 2.3 the cubic twist ratio crosses the tube threshold between L = 4 (r = 0.42, at the self-dual value) and L = 6 (r = 0.86, F\_e = 2.6). Central values pass O1 at L = 6, but the one-sided 95 \% lower bound of F\_e is 1.26 statistically and 0.96 with the resolution floor: \textbf{not certified} \textbf{[numerical]} (\texttt{CONTAINMENT.md}, 1 and 3; \texttt{SATURATION.md}, 0). At beta = 3.5, L = 2: r = 0.018, F\_e = 0.035, as IV.5.1 predicts.

\textbf{IV.7.2 The chain criterion at a large cross-section.} The pre-registered check of F\_e(L) \textgreater{} 2c(L) on chains of cross-section 16\textasciicircum{}3 refuted it at beta = 3.5, L = 1, at 99 \% \textbf{[numerical]}. Re-scored post hoc under the universal floor, the refuted range grows to L \textless{}= 2 (both levels) and L = 3 (95 \% only). At beta = 2.3 no test was possible (no flux bound). This geometry is not the cubic one of O1, so it bears on O1 neither way; it says only that at that cross-section and coupling the electric multiplication needs L \textgreater{}= 4 (\texttt{ONEBOX-CHECK.md}, 4, 5, 7).

\textbf{IV.7.3 The cubic one-box geometry is untested by saved data}; a small-torus simulation is pre-registered (\texttt{ONEBOX-CHECK.md}, 8).

\textbf{IV.7.4 The mirror reading.} The identity 8 pi\textasciicircum{}2 = -4 zeta(2)/zeta(-1) is derived in Lean from Mathlib's functional equation \textbf{[Lean]} (\texttt{GRH/Mirror.lean}, \texttt{eight\_pi\_sq\_mirror}, \texttt{zeta\_neg\_one\_val}); it is the k = 1 member of a classical family. In Yang-Mills it brackets F\_e(L) between two asymptotic ends: the Luescher term (zeta(-1)) at the infrared end, fractional instantons (4 pi\textasciicircum{}2/g\textasciicircum{}2) at the ultraviolet end. In effective-string terms O1 becomes L sqrt(sigma) \textgreater{} 1.72 (Luescher) or 1.80 (Arvis), just above the string's breakdown at 1.45 \textbf{[numerical]}, using the effective string as a model. The mirror does not constrain F\_e between the two ends, where O1 sits \textbf{[reading]} for ``brackets'' (\texttt{docs/rh-guide/MIRROR.md}, 0.2).

{\footnotesize
\begin{longtable}{>{\raggedright\arraybackslash}p{\dimexpr 0.1153\linewidth-2\tabcolsep\relax}>{\raggedright\arraybackslash}p{\dimexpr 0.3761\linewidth-2\tabcolsep\relax}>{\raggedright\arraybackslash}p{\dimexpr 0.2018\linewidth-2\tabcolsep\relax}>{\raggedright\arraybackslash}p{\dimexpr 0.2868\linewidth-2\tabcolsep\relax}}
\toprule
\textbf{\#} & \textbf{Result} & \textbf{Status} & \textbf{Source} \\
\midrule
\endhead
IV.7.1 & Small tori at beta 2.3: F\_e(6) = 2.6 centrally, lower bound 0.96 with the floor; not certified & \textbf{[numerical]} (preliminary) & \texttt{CONTAINMENT.md}, 1, 3; \texttt{SATURATION.md}, 0 \\
IV.7.2 & Chain criterion refuted at cross-section 16\textasciicircum{}3, beta 3.5, L = 1 (pre-registered); L \textless{}= 2 post hoc & \textbf{[numerical]} & \texttt{ONEBOX-CHECK.md}, 4, 7 \\
IV.7.3 & The universal error-bar floor max(statistical, 1/T) & project rule & \texttt{attempts/RULES.md}, last section \\
IV.7.4a & 8 pi\textasciicircum{}2 = -4 zeta(2)/zeta(-1); zeta(-1) = -1/12 & \textbf{[Lean]} & \texttt{ProofPhase/GRH/Mirror.lean}, \texttt{eight\_pi\_sq\_mirror}, \texttt{zeta\_neg\_one\_val} \\
IV.7.4b & O1 in effective-string units: L sqrt(sigma) \textgreater{} 1.72 to 1.80 & \textbf{[numerical]} (model-based) & \texttt{docs/rh-guide/MIRROR.md}, 0.2 \\
\bottomrule
\end{longtable}}

\subsection*{IV.8 ``It runs'': the formal target shared with GRH}
\addcontentsline{toc}{subsection}{IV.8 ``It runs'': the formal target shared with GRH}

Both Yang-Mills and GRH ask about an evolution that does not increase a positive norm. If T is a contraction on a normed space, the GRH-type conclusion is that every mirrored mode is a phase; the Yang-Mills conclusion is that, with \textbar{}\textbar{}T x\textbar{}\textbar{} \textless{}= eta \textbar{}\textbar{}x\textbar{}\textbar{} on the neutral sector and eta \textless{} 1, there is no phase mode there, and iterates decay as eta\textasciicircum{}n \textbf{[Lean]} (\texttt{no\_phase\_mode}, \texttt{iterate\_contraction}, \texttt{runs\_two\_targets}; GRH side \texttt{GRH/Tate.lean}, \texttt{runs\_and\_mirror\_on\_line}). The open steps are complementary: Yang-Mills has the positive norm for free (reflection positivity) and lacks the margin eta \textless{} 1; GRH lacks the positive norm (Weil positivity) and needs no margin (Part III).

{\footnotesize
\begin{longtable}{>{\raggedright\arraybackslash}p{\dimexpr 0.1080\linewidth-2\tabcolsep\relax}>{\raggedright\arraybackslash}p{\dimexpr 0.3414\linewidth-2\tabcolsep\relax}>{\raggedright\arraybackslash}p{\dimexpr 0.1527\linewidth-2\tabcolsep\relax}>{\raggedright\arraybackslash}p{\dimexpr 0.3779\linewidth-2\tabcolsep\relax}}
\toprule
\textbf{\#} & \textbf{Result} & \textbf{Status} & \textbf{Source} \\
\midrule
\endhead
IV.8.1 & ``It runs'' in two targets: phases (mirror) and no phase mode under a margin & \textbf{[Lean]} & \texttt{ProofPhase/Crossover/Runs.lean}, \texttt{no\_phase\_mode}, \texttt{iterate\_contraction}, \texttt{runs\_two\_targets} \\
\bottomrule
\end{longtable}}

\subsection*{IV.9 The weak-coupling side and the continuum (inputs from outside the gap mechanism)}
\addcontentsline{toc}{subsection}{IV.9 The weak-coupling side and the continuum (inputs from outside the gap mechanism)}

{\footnotesize
\begin{longtable}{>{\raggedright\arraybackslash}p{\dimexpr 0.0978\linewidth-2\tabcolsep\relax}>{\raggedright\arraybackslash}p{\dimexpr 0.3788\linewidth-2\tabcolsep\relax}>{\raggedright\arraybackslash}p{\dimexpr 0.2177\linewidth-2\tabcolsep\relax}>{\raggedright\arraybackslash}p{\dimexpr 0.2857\linewidth-2\tabcolsep\relax}}
\toprule
\textbf{\#} & \textbf{Result} & \textbf{Status} & \textbf{Source} \\
\midrule
\endhead
IV.9.1 & UV stability, small-field control and one-loop running at the weak-coupling unit scale (a recent construction by others) & \textbf{[external]} & \texttt{H1-OPENAI.md} \\
IV.9.2 & Balaban's large-field control preserves partition functions, not laws & \textbf{[classical]}, read in full & \texttt{proof-phase/c5d-supplier/BALABAN-READING.md} \\
IV.9.3 & M1: the law-preserving large-field step, at weak coupling up to L\_* & \textbf{[paper]}, complete at paper level; Lean parts & \texttt{ProofPhase/M1/Assembly.lean} (\texttt{m1\_assembly}); \texttt{proof-phase/m1-tools/M1-PLAN.md} \\
IV.9.4 & The continuum limit exists (mesh limit at fixed volume, then volume along a subsequence) & \textbf{[external]} & \texttt{H1-OPENAI.md}; \texttt{ProofPhase/Existence/Bridge.lean} (\texttt{plaqLaw\_eq\_wilsonMeasure}) \\
IV.9.5 & (O2) Line of constant physics (one-loop slope 11/(6 pi\textasciicircum{}2)) and non-triviality, as inputs of the assembly & \textbf{[open]} & \texttt{ProofPhase/Assembly.lean} (\texttt{OnLCP}, \texttt{NonTrivial}); \texttt{ProofPhase/V4/LCP.lean} \\
IV.9.6 & (O3) Euclidean invariance, local fields as distributions, uniqueness of the limit, Wightman reconstruction & \textbf{[external]} for the external construction; \textbf{[open]} for ours & \texttt{PROOF-BLUEPRINT.md}, P6.4 to P6.9, P7.6 \\
IV.9.7 & (O4) Every compact simple group: existence external; this project's tuning SU(2) only & \textbf{[external]} / \textbf{[open]} & \texttt{PROOF-BLUEPRINT.md}, P6.10 \\
\bottomrule
\end{longtable}}

M1 hands over at a weak-coupling scale: the coupling becomes of order one only after the crossover. So the floor M1 supplies at L\_* is uniform in the spacing but tiny, and it gives a physical-units gap only along a chain at fixed cross-section (IV.3.6), not on a grid (IV.4.3) \textbf{[paper]} (\texttt{RDOEBLIN.md}, 0).

\subsection*{IV.10 DMP readings attached to this part (all [reading])}
\addcontentsline{toc}{subsection}{IV.10 DMP readings attached to this part (all [reading])}

\begin{itemize}
\item \textbf{Outside, non-existence keeps existing through assertion.} The twist is rare because electric flux costs almost nothing: one fact seen from the two sides of the duality (IV.5.1).
\item \textbf{The container must end, and does, along the run.} ``It can never catch up'' in exact form: Z\_1/Z\_0 \textless{}= C mu\textasciicircum{}T, mu \textless{} 1 (IV.5.2).
\item \textbf{Inside, existence has to outrun possibility.} Of the mu continuations per step only one is the tube; the tube dominates only if its cost per step exceeds log mu. The balanced state (0.881) is not enough (IV.5.4). This is ``existence = non-existence + 1'' in counting form.
\item \textbf{IM is both sides.} Gauss's law is the same statement on both sides of every DB; it makes every tube possible, and the action has to make one tube cheaper than all others combined (IV.5.3, IV.5.4).
\item \textbf{The pressing as the instanton of the first act.} The Bogomolny bound is a worst-case floor proved before any dynamics runs; it prices the rare side and loses logarithmically to entropy as the tick shrinks (IV.5.5).
\item \textbf{The duality is never alone while running.} The 't Hooft relation is kinematic and chooses no branch; the action selects it (IV.4.4).
\item \textbf{The crossover is confined by making no statement about its content} (IV.2).
\item WHKs: group plus one scale fix everything (dimensional transmutation, \textbf{[classical]}); mass gap = the arrow of causality exists (I.7); only 3+1 transmutes, g\textasciicircum{}2 being dimensionless only in d = 4 \textbf{[classical]}.
\end{itemize}
\subsection*{IV.11 What would close it}
\addcontentsline{toc}{subsection}{IV.11 What would close it}

O1 first: a worst-case argument, for one box of fixed physical size, that the electric cost per box exceeds the path entropy of the four-dimensional lattice for every well-defined outside, together with the per-box product structure along paths (IV.5.7). Then O2: its scale-uniform lower edge (IV.9.5b), conditional on the one assumption beside O1 (Part III 4.4); the per-box half is proved and the gap uses only the resolution clause. Then O3: an observable bridge from this project's bounded slab observables to local fields, or adoption of the external construction's field notion with its audit. Then O4: a tuning for general G. O1 is the mathematical core; O2 to O4 are the remaining distance from the gap mechanism to the Clay statement.


%% ---- 05-part-V-grh.md
\section*{Part V. The generalised Riemann hypothesis: the positivity route}
\addcontentsline{toc}{section}{Part V. The generalised Riemann hypothesis: the positivity route}

\subsection*{V.0 What is proved, by whom, and what this part does}
\addcontentsline{toc}{subsection}{V.0 What is proved, by whom, and what this part does}

\textbf{State of the problem (2026-11-23).}

\begin{itemize}
\item The Riemann hypothesis for zeta and GRH for every Dirichlet L-function are proved in Lean by others: a library released 2026-11-20 and certified by the Lean and Mathlib communities on 2026-11-21 \textbf{[external]} (\texttt{SOURCES.md}, S1 to S4; \texttt{GUIDE.md}, section 1). This project read the main statements in the source and did not rebuild the library.
\item GRH for Dedekind, Hecke and general automorphic L-functions is not proved \textbf{[open]} (\texttt{GUIDE.md}, sections 1 and 5).
\item The proof by others does not go through positivity. Its endgame is Landau's theorem applied to a one-sided square-root bound on the prime count, obtained from a Stepanov-type size balance \textbf{[external, traced]} (\texttt{GUIDE.md}, sections 2, 4, 5).
\end{itemize}
\textbf{What this part does.} It follows the other road: Weil's positivity criterion, the second column of the chain of II.8. It proves the chain in finite models (V.1), shows exactly why a mirror symmetry alone cannot place zeros on the line (V.2), measures the shape of the missing positivity (V.3), and states the one missing construction in eight equivalent-looking faces (V.4), with ``it runs'' as the form shared with Yang-Mills (V.5). \textbf{This part does not prove GRH for any L-function.} Its value is a reduction that is independent of the existing proof and would, if closed, extend to L-functions the existing proof does not reach.

\subsection*{V.1 The chain in finite models (links 1 to 4)}
\addcontentsline{toc}{subsection}{V.1 The chain in finite models (links 1 to 4)}

In a finite model everything works: zeros are finitely many, the pairing is a finite sum, and ``positive'' means ``on the line''.

{\footnotesize
\begin{longtable}{>{\raggedright\arraybackslash}p{\dimexpr 0.0975\linewidth-2\tabcolsep\relax}>{\raggedright\arraybackslash}p{\dimexpr 0.4358\linewidth-2\tabcolsep\relax}>{\raggedright\arraybackslash}p{\dimexpr 0.1409\linewidth-2\tabcolsep\relax}>{\raggedright\arraybackslash}p{\dimexpr 0.3058\linewidth-2\tabcolsep\relax}}
\toprule
\textbf{\#} & \textbf{Result} & \textbf{Status} & \textbf{Source} \\
\midrule
\endhead
V.1.1 & Finite Weil criterion: zeros closed under conjugation; the zero pairing is non-negative on every finite exponential sum iff every zero is on the line & \textbf{[Lean]} & \texttt{GRH/Weil.lean}, \texttt{weil\_criterion\_finite} \\
V.1.2 & Zeros on the line imply the pairing is non-negative (each term is a square) & \textbf{[Lean]} & \texttt{GRH/Weil.lean}, \texttt{weilSum\_re\_nonneg\_of\_real} \\
V.1.3 & Positivity implies the line for any test class with the approximation property (a named hypothesis) & \textbf{[Lean]} & \texttt{GRH/Weil.lean}, \texttt{real\_of\_weilSum\_nonneg} \\
V.1.4 & Weil's criterion for zeta and Dirichlet L-functions & \textbf{[classical]} & Weil 1952; Bombieri 2000 \\
V.1.5 & Weil's explicit formula for zeta holds for every smooth compactly supported test function, by the contour shift, with a polynomial zero count per unit window (Mathlib-only core, no named step) & \textbf{[Lean]} & \texttt{GRH/ContourShift.lean} (\texttt{CS.contourShift}), \texttt{GRH/ContourShiftTie.lean} (\texttt{explicit\_formula\_of\_smooth}), \texttt{GRH/ContourShiftZeta.lean} (\texttt{zero\_count\_window}) \\
V.1.5b & For Dirichlet and other L-functions the explicit formula, and Weil's criterion itself, are not yet in Lean & \textbf{[open]} & \texttt{docs/unified/CONTOUR-SHIFT.md}, section 5 \\
V.1.6 & A similitude multiplying a positive form by q has all eigenvalues of modulus sqrt q & \textbf{[Lean]} & \texttt{GRH/Orbits.lean}, \texttt{similitude\_eigenvalue} \\
V.1.7 & If Theta - 1/2 is skew-adjoint for a positive form, every eigenvalue has real part 1/2 (finite dimension) & \textbf{[Lean]} & \texttt{GRH/Orbits.lean}, \texttt{skew\_shift\_re\_half} \\
V.1.8 & Lee-Yang circle theorem and Asano contraction, small cases & \textbf{[Lean]} & \texttt{GRH/LeeYang.lean}, \texttt{asano}, \texttt{ising3\_zeros\_on\_circle} \\
V.1.9 & Fixed points of f\textasciicircum{}n are counted by points of exact period d dividing n (finite ``remnants as evidence'') & \textbf{[Lean]} & \texttt{GRH/Orbits.lean}, \texttt{card\_fixed\_pow\_eq\_sum} \\
V.1.10 & A finite exponential sum bounded by any exponential constant is bounded exactly (``a lossy bound becomes exact on a finite spectrum'') & \textbf{[Lean]} & \texttt{GRH/Induction.lean}, \texttt{norm\_le\_of\_exp\_bound} \\
\bottomrule
\end{longtable}}

\subsection*{V.2 The mirror alone does not place the zeros}
\addcontentsline{toc}{subsection}{V.2 The mirror alone does not place the zeros}

\textbf{V.2.1 Without positivity, duality allows off-line pairs.} An indefinite form is preserved by a mirror pair off the line (3/4 and 1/4) \textbf{[Lean]} (\texttt{GRH/Orbits.lean}, \texttt{krein\_off\_line}). The functional equation pairs rho with 1 - conj(rho); on the line the partner is the zero itself, off the line it is a second zero. A mirror symmetry never places zeros on its axis by itself.

\textbf{V.2.2 The arithmetic instance.} Tate's thesis writes the completed L-function as a product of local zeta integrals, one per place, each with its own functional equation; the global functional equation is Poisson summation for Q inside the adeles \textbf{[classical]}. On the critical line the local factors at s and at 1 - s have the same modulus, so their ratio is a pure phase, proved in Lean for every unramified prime \textbf{[Lean]} (\texttt{GRH/Tate.lean}, \texttt{local\_factor\_unimodular}). Mathlib proves the functional equation for every even or odd function on Z/N, not only for characters \textbf{[classical]} (in Mathlib as \texttt{completedLFunction\_one\_sub\_even} and \texttt{completedLFunction\_one\_sub\_odd}, \texttt{Mathlib/NumberTheory/LSeries/ZMod.lean}, Mathlib v4.34.0). Among those functions is the Davenport-Heilbronn combination mod 5, which has a self-dual functional equation and zeros off the line \textbf{[classical]} (recalled). So local functional equations, global Poisson and the Gaussian allow an off-line pair. This is the arithmetic \texttt{krein\_off\_line}.

\textbf{V.2.3 The one local sign.} What must be added is a sign. The one local sign available is Lambda(n) \textgreater{}= 0: each Euler factor's log-derivative has non-negative coefficients; Davenport-Heilbronn lacks it \textbf{[classical]}. With Landau's theorem it turns a one-sided square-root bound on psi(x) - x into a contradiction with an off-line pair \textbf{[paper]} (\texttt{grh/INDUCTION.md}, 0.2). The one-sided bound is global: it does not follow from the local data plus Poisson, and generalised-prime systems with regular integers violate it \textbf{[classical]} (Diamond-Montgomery-Vorhauer). This is exactly the bound the proof by others supplies.

\textbf{V.2.4 Positivity locates singularities (the bridge, II.11).} Pringsheim, Landau, Laplace and spectral instances \textbf{[Lean]} (\texttt{Bridge/}); a positive transform has no pole of negative residue approached from the right, and the ray form excludes zeros right of theta under a one-sided bound \textbf{[Lean]} (\texttt{Bridge/Poles.lean}, \texttt{landau\_ray\_poles}, \texttt{lseries\_residue\_nonneg}); a real zero beta in (theta, 1) of L(chi, s) gives such a pole, of residue -m/(beta(1 - beta)) \textbf{[classical]}, two-line computation, not formalised.

\textbf{V.2.5 The two 1/2s.} Over function fields one positive duality (the Hodge index; Rosati positivity) gives both the mirror axis and the balance point 1/2 \textbf{[classical]} (Weil 1948). In Stepanov's route and in the Lean proof by others the two 1/2s meet only in the conclusion \textbf{[classical]} / \textbf{[external, traced]}. An arithmetic positivity in the role of the Hodge index, giving both, is \textbf{[open]} (\texttt{bridge/BRIDGE.md}, section 4).

\textbf{V.2.6 The mirror constants.} The mirror pair Weil's form actually uses is (s = 1, s = 0); it fixes the real-place constants (log pi + gamma, the pole terms; zeta'(0)/zeta(0) = log 2 pi) \textbf{[numerical]}. These constants belong to the archimedean part alone, which on its own is indefinite from support 0.425 on (V.3.9). No sign comes from them (\texttt{MIRROR.md}, 0.3).

{\footnotesize
\begin{longtable}{>{\raggedright\arraybackslash}p{\dimexpr 0.1100\linewidth-2\tabcolsep\relax}>{\raggedright\arraybackslash}p{\dimexpr 0.3803\linewidth-2\tabcolsep\relax}>{\raggedright\arraybackslash}p{\dimexpr 0.2242\linewidth-2\tabcolsep\relax}>{\raggedright\arraybackslash}p{\dimexpr 0.2656\linewidth-2\tabcolsep\relax}}
\toprule
\textbf{\#} & \textbf{Result} & \textbf{Status} & \textbf{Source} \\
\midrule
\endhead
V.2.1 & Duality without positivity allows mirror pairs off the line & \textbf{[Lean]} & \texttt{GRH/Orbits.lean}, \texttt{krein\_off\_line} \\
V.2.2a & The ratio of the local factor at s and at 1 - s has modulus 1 on the line (unramified primes) & \textbf{[Lean]} & \texttt{GRH/Tate.lean}, \texttt{local\_factor\_unimodular} \\
V.2.2b & Functional equation for every even/odd function on Z/N; Davenport-Heilbronn has zeros off the line & \textbf{[classical]} (in Mathlib: \texttt{completedLFunction\_one\_sub\_even}, \texttt{\_odd}) & \texttt{grh/TATE.md}, 0.2 \\
V.2.3a & Lambda(n) \textgreater{}= 0 plus a one-sided square-root bound plus Landau excludes zeros right of 1/2 & \textbf{[paper]} & \texttt{grh/INDUCTION.md}, 0.2; \texttt{grh/TATE.md}, 0.3 \\
V.2.3b & The one-sided bound is global (generalised primes violate it) & \textbf{[classical]} & Diamond-Montgomery-Vorhauer; \texttt{grh/TATE.md}, 0.3 \\
V.2.4a & Positive transforms exclude negative-residue poles; ray form & \textbf{[Lean]} & \texttt{Bridge/Poles.lean}, \texttt{landau\_ray\_poles}, \texttt{lseries\_residue\_nonneg} \\
V.2.4b & A real zero in (theta, 1) gives a negative-residue pole & \textbf{[classical]} (not formalised) & \texttt{bridge/BRIDGE.md}, 1.4 \\
V.2.5a & Function fields: one positive duality gives both 1/2s & \textbf{[classical]} & Weil 1948; \texttt{bridge/BRIDGE.md}, 2.3 \\
V.2.5b & Stepanov and the Lean proof by others: the 1/2s meet only in the conclusion & \textbf{[classical]} / \textbf{[external, traced]} & \texttt{bridge/BRIDGE.md}, 2.3 and 3 \\
V.2.5c & \textbf{An arithmetic positivity giving both 1/2s} & \textbf{[open]} & \texttt{bridge/BRIDGE.md}, 4 \\
V.2.6 & The mirror pair Weil uses is (1, 0); its constants are archimedean and give no sign & \textbf{[numerical]} + \textbf{[Lean]} (via V.2.1) & \texttt{MIRROR.md}, 0.3 \\
\bottomrule
\end{longtable}}

Reading: ``inversion and halving'': every 1/2 in these formulas is the fixed point of an inversion (s -\textgreater{} 1 - s, alpha -\textgreater{} q/alpha, m -\textgreater{} x/m) \textbf{[reading]}.

\subsection*{V.3 Measured: the shape of the missing positivity}
\addcontentsline{toc}{subsection}{V.3 Measured: the shape of the missing positivity}

On a grid of cells Weil's form is a Toeplitz matrix, because it is translation invariant. This turns the question into one about a sequence of matrices, and numerics can probe it. Every number in this section can refute a construction; none proves a sign.

\textbf{V.3.1 The induction frame is exact.} Growing the support by one cell adds one entry, at the corner of the matrix. The bigger section is positive exactly when that entry lands in an interval centred on the value extrapolated from all smaller scales, with the current margin as radius (Szego recursion, Levinson) \textbf{[classical]}. Each step is one scalar condition.

\textbf{V.3.2 Each step holds, by exact cancellation.} The induction holds at every step computed, every reflection coefficient at most 0.68 in modulus, up to support a = 1.25 \textbf{[numerical]}. But the margin collapses (lambda\_min about 1e-11 by a = 1.24), and the new prime's own term is up to about 700 interval radii, of which the extrapolation from below cancels all but less than one \textbf{[numerical]}. No estimate with slack can carry such a step: scaling one prime weight by (1 + eta) breaks positivity at eta of order 1e-3 to 1e-9 \textbf{[numerical]}.

\textbf{V.3.3 The ground state.} Minimising the form over unit vectors on each section gives the margin; the minimiser's polynomial has all its zeros on the unit circle, which is the critical line in grid coordinates (to 1e-53), and as primes enter these zeros lock onto zeta's (14.134725 to 3e-7 with the primes up to 11) \textbf{[numerical]}. Without the primes they are evenly spaced. The zeros of the minimiser are on the circle for every real symmetric Toeplitz section, positive or not (Caratheodory-Fejer, Pisarenko) \textbf{[classical]}; in Lean, \texttt{atoms\_on\_circle} and \texttt{atomic\_toep\_nonneg} (\texttt{GRH/Induction.lean}) \textbf{[Lean]}. Positivity is exactly the statement that the uniform remnant lambda\_min is non-negative.

\textbf{V.3.4 The true primes sit at the edge, not at the centre.} The margin is concave in the prime weights and the positive configurations form a convex set \textbf{[Lean]} (\texttt{GRH/Variational.lean}, \texttt{toep\_pos\_convex}). Scaling the prime 2 by (1 + s) changes the margin at rate -1.01 while the margin is 1e-9: the arithmetic configuration lies about 1e-9 from the edge of the positive window \textbf{[numerical]}. ``Positivity = stability of the minimum'' holds only as a tautology \textbf{[Lean]} (\texttt{line\_min\_iff}, \texttt{min\_at\_zero\_iff}).

\textbf{V.3.5 All or nothing.} At support 0.625 to 0.725 only 2 and 3 are present: A + P\_2 and A + P\_3 are both indefinite, A + P\_2 + P\_3 is positive. At every support computed beyond an entry window, the only positive set of primes is the full set present; dropping 11, which has barely entered, already gives -1.1e-3 against +1e-9 \textbf{[numerical]}. Two combined Szego steps carry exactly the product of the single steps, so the chain has no hidden pair certificate \textbf{[classical]}.

\textbf{V.3.6 Contraction on the sections, isometry on the space.} On the sections every step contracts (factors 1 - k\textasciicircum{}2 in (0.54, 1); per-prime products 0.065 for 2 up to 0.92 for 11), and the product goes to zero (Kolmogorov-Szego for the atomic zero measure) \textbf{[numerical]} + \textbf{[classical]}. A positive margin survives finitely many positive factors \textbf{[Lean]} (\texttt{GRH/Transfer.lean}, \texttt{margin\_pos\_of\_factors}, \texttt{prod\_one\_sub\_le\_exp}). Here the roles of Yang-Mills are reversed: eta \textless{} 1 is automatic and harmless, and the sign is eta\_p \textgreater{} 0 at every step.

\textbf{V.3.7 No two-term bound gives the sign.} The best bound that sees the archimedean part and the primes separately is Weyl's, lambda\_min(A) + lambda\_min(sum P), attained by aligning the worst directions \textbf{[Lean]} (\texttt{two\_term\_bound}). It is -7.9 t\_0 against a true +1e-9 t\_0 at support 1.225 \textbf{[numerical]}. The correct two-term form is Krein's extension theorem: a section is positive exactly when Weil's distribution on the window extends to a positive measure \textbf{[classical]}.

\textbf{V.3.8 Positivity cannot be assembled place by place.} Different primes live on disjoint lattices k log p \textbf{[Lean]} (\texttt{prime\_lattices\_disjoint}), so places can exchange only the lag-0 part; each prime block has zero diagonal and is indefinite as soon as it is non-zero \textbf{[Lean]} (\texttt{zero\_diag\_indefinite}, \texttt{form\_two\_point}); so place-by-place positivity would force the real place's own form to be positive and to pay for every prime \textbf{[Lean]} (\texttt{local\_assembly\_forces\_arch}).

\textbf{V.3.9 The real place alone is not positive at resolvable scales.} It loses positivity at support 0.40 to 0.43, just after the prime 2 enters at (1/2) log 2 = 0.347; by support 1.225 the place-by-place condition fails by 9.7 t\_0 while the whole form is positive by 1e-9 t\_0 \textbf{[numerical]}. The growing mean log(r / 2 pi) confines only above frequency about 2 pi. Semilocally, Weil's form cut to a set of places is positive only if the set holds the real place and every prime up to e\textasciicircum{}\{2a\} \textbf{[numerical]} (\texttt{grh/SEMILOCAL.md}).

\textbf{V.3.10 Weight positivity is not enough.} Lambda(n) \textgreater{}= 0, Phi(u) \textgreater{} 0 and the Euler product are provably insufficient; the de Bruijn-Newman constant Lambda = 0 sits exactly at the edge \textbf{[classical]} + \textbf{[paper]} (\texttt{grh/LEE-YANG.md}, 5; Polya 1926; Rodgers-Tao 2020).

{\footnotesize
\begin{longtable}{>{\raggedright\arraybackslash}p{\dimexpr 0.1065\linewidth-2\tabcolsep\relax}>{\raggedright\arraybackslash}p{\dimexpr 0.3755\linewidth-2\tabcolsep\relax}>{\raggedright\arraybackslash}p{\dimexpr 0.1987\linewidth-2\tabcolsep\relax}>{\raggedright\arraybackslash}p{\dimexpr 0.2994\linewidth-2\tabcolsep\relax}}
\toprule
\textbf{\#} & \textbf{Result} & \textbf{Status} & \textbf{Source} \\
\midrule
\endhead
V.3.1 & Toeplitz frame; each step is one scalar condition at the corner & \textbf{[classical]} + \textbf{[paper]} & \texttt{grh/INDUCTION.md}, 0.1 \\
V.3.2a & Each step holds up to a = 1.25, reflection coefficients at most 0.68 & \textbf{[numerical]} & \texttt{grh/INDUCTION.md}, 0.1, 0.5 \\
V.3.2b & Each step is an exact cancellation (prime term up to about 700 radii) & \textbf{[numerical]} & \texttt{grh/INDUCTION.md}, 0.1 \\
V.3.2c & Weight scaling breaks positivity at 1e-3 to 1e-9 & \textbf{[numerical]} & \texttt{grh/INDUCTION.md}, 0.3 \\
V.3.3a & Ground-state zeros on the circle, locking onto zeta's zeros & \textbf{[numerical]} & \texttt{grh/VARIATIONAL.md}, 0.1 \\
V.3.3b & Minimiser's zeros on the circle for every real symmetric Toeplitz section; positivity = remnant \textgreater{}= 0 & \textbf{[classical]} + \textbf{[Lean]} & \texttt{GRH/Induction.lean}, \texttt{atoms\_on\_circle}, \texttt{atomic\_toep\_nonneg}; \texttt{grh/VARIATIONAL.md}, 0.2 \\
V.3.4a & Positive configurations convex; margin concave in the weights & \textbf{[Lean]} & \texttt{GRH/Variational.lean}, \texttt{toep\_pos\_convex} \\
V.3.4b & The true primes sit about 1e-9 from the edge & \textbf{[numerical]} & \texttt{grh/VARIATIONAL.md}, 0.4 \\
V.3.5 & Only the full set of primes present is positive & \textbf{[numerical]} & \texttt{grh/TRANSFER.md}, 0.1 \\
V.3.6 & Per-step contraction on sections; a positive margin survives positive factors & \textbf{[numerical]} + \textbf{[Lean]} & \texttt{GRH/Transfer.lean}, \texttt{margin\_pos\_of\_factors}, \texttt{prod\_one\_sub\_le\_exp}; \texttt{grh/TRANSFER.md}, 0.2 \\
V.3.7 & Weyl's two-term bound is the best separate bound and fails by 7.9 t\_0 & \textbf{[Lean]} (bound) + \textbf{[numerical]} & \texttt{GRH/Transfer.lean}, \texttt{two\_term\_bound}; \texttt{grh/TRANSFER.md}, 0.3 \\
V.3.8 & Disjoint prime lattices; prime blocks indefinite; place-by-place forces the real place to pay for all & \textbf{[Lean]} & \texttt{GRH/Tate.lean}, \texttt{prime\_lattices\_disjoint}, \texttt{zero\_diag\_indefinite}, \texttt{form\_two\_point}, \texttt{local\_assembly\_forces\_arch} \\
V.3.9 & Real place alone indefinite from support 0.40 to 0.43; semilocal test & \textbf{[numerical]} & \texttt{grh/TATE.md}, 0.4; \texttt{grh/SEMILOCAL.md} \\
V.3.10 & Weight positivity is not enough; Lambda = 0 at the edge & \textbf{[classical]} + \textbf{[paper]} & \texttt{grh/LEE-YANG.md}, 5 \\
V.3.11 & Explicit formula normalisation agrees to 1e-16; smallest eigenvalue on a 10-function space 2.4e-1 to 2.6e-9, no violation & \textbf{[numerical]} & \texttt{grh/weil\_numerics.log}; \texttt{grh/ATTEMPT.md}, 4 \\
\bottomrule
\end{longtable}}

\subsection*{V.4 The open step: one missing positivity, eight faces}
\addcontentsline{toc}{subsection}{V.4 The open step: one missing positivity, eight faces}

The missing piece: \textbf{a positive object built from the primes, the characters and the real place alone, without the zeros.} The notes have met it eight times. Each face would close the others' consequences; they are not shown to be equivalent.

{\footnotesize
\begin{longtable}{>{\raggedright\arraybackslash}p{\dimexpr 0.1652\linewidth-2\tabcolsep\relax}>{\raggedright\arraybackslash}p{\dimexpr 0.4938\linewidth-2\tabcolsep\relax}>{\raggedright\arraybackslash}p{\dimexpr 0.1371\linewidth-2\tabcolsep\relax}>{\raggedright\arraybackslash}p{\dimexpr 0.1838\linewidth-2\tabcolsep\relax}}
\toprule
\textbf{Face} & \textbf{Statement} & \textbf{Status} & \textbf{Source} \\
\midrule
\endhead
1 Weil Hilbert space & A Hilbert space H\_chi and a map Phi from test functions, built from arithmetic data only, with \textless{}Phi f, Phi g\textgreater{} = W\_chi(f * g\textasciitilde{}). By Weil and GNS it exists iff GRH holds for chi. None of its subset truncations is positive (V.3.5) & \textbf{[open]} & \texttt{grh/ATTEMPT.md}, step 4 \\
2 Critical ferromagnet & A lattice ferromagnet at its critical coupling whose block-spin law is Riemann's Phi; RH by Newman's closure. Asano adds one site at a time; here one prime at a time fails & \textbf{[open]} & \texttt{grh/LEE-YANG.md}, 5c \\
3 Deninger space & A foliated space whose closed orbits are the primes, with a positive leafwise Hodge form. In the hyperbolic analogue positivity gives only ``line or real segment''; the line needs a gap (lambda\_1 \textgreater{}= 1/4) & \textbf{[open]} & \texttt{grh/ORBITS.md}, 6; \texttt{grh/TRANSFER.md}, 5 \\
4 Hermite-Biehler chain & A Hermite-Biehler coupling of xi's real and imaginary parts (de Branges frame; the specific de Branges conjecture is false, Conrey-Li 2000); its Schur parameters are the step factors & \textbf{[open]}; framework \textbf{[classical]} & \texttt{grh/INDUCTION.md}, 0.4 \\
5 Step lemma & From the primes up to p\_\{k+1\} and the real place, put the new corner entry inside the extrapolated interval, without the zeros & \textbf{[open]} & \texttt{grh/INDUCTION.md}, 7 \\
6 Sign of the remnant & From arithmetic alone, lambda\_min(a) \textgreater{}= 0 at every support; equivalently the ground states converge to the zero measure & \textbf{[open]} & \texttt{grh/VARIATIONAL.md}, 0.5 \\
7 Global sign (``it runs'') & The norm in which the scaling flow runs is positive; equivalently the one-sided square-root prime bound. In the constructions where the norm is positive by construction, the space either misses the off-line zeros (Connes) or yields a half-plane (Lax-Phillips) & \textbf{[open]} & \texttt{grh/TATE.md}, 0.5, 0.6 \\
8 All-or-nothing step & At each scale the full set of primes present, and no subset, places the corner entry inside the interval; no decomposition into pairs, single primes or two terms carries it & \textbf{[open]} & \texttt{grh/TRANSFER.md}, 5 \\
\bottomrule
\end{longtable}}

Faces 2 and 3 each imply face 1 \textbf{[paper]} (\texttt{grh/INDUCTION.md}, 7; \texttt{grh/ORBITS.md}, 4). Faces 5, 6 and 8 are the same quantity read between scales, at one scale, and over subsets \textbf{[paper]} (\texttt{grh/VARIATIONAL.md}, 5; \texttt{grh/TRANSFER.md}, 5).

\textbf{The constraint on any construction} (V.3): it cannot add one prime at a time with slack; it cannot assemble positivity place by place or subset by subset; it cannot separate the archimedean part from the primes. All primes up to the scale and the real place must enter together, in their exact relative position, through an exact identity.

\subsection*{V.5 ``It runs'': the face shared with Yang-Mills}
\addcontentsline{toc}{subsection}{V.5 ``It runs'': the face shared with Yang-Mills}

The owner named the positivity the reductio needs: ``It runs.'' Read formally: an off-line pair of zeros would be a growing mode of the arithmetic evolution, contradicting that the evolution does not increase a positive norm. On the completion of test functions in Weil's form the scaling flow preserves the form automatically \textbf{[paper]}. If the flow does not increase a positive norm, there is no growing mode; the functional equation pairs every decaying mode with a growing one; so every mode is a phase, which is RH \textbf{[Lean]} (\texttt{GRH/Tate.lean}, \texttt{runs\_and\_mirror\_on\_line}; \texttt{atoms\_on\_circle} in grid form). The open part is exactly face 7: that norm is Weil's form, and its positivity is the open step. On the constructed space the flow is an isometry (eta = 1, ``flat''), and what remains are the zeros (phases) and the orbit lengths log p: the explicit formula is a trace formula of these remnants \textbf{[classical]}.

Yang-Mills has the same formal target with the opposite missing piece (IV.8): there the positive norm is free and the margin is missing; here the margin is not needed and the positive norm is missing \textbf{[Lean]} (\texttt{Crossover/Runs.lean}, \texttt{runs\_two\_targets}).

\subsection*{V.6 DMP readings attached to this part (all [reading])}
\addcontentsline{toc}{subsection}{V.6 DMP readings attached to this part (all [reading])}

\begin{itemize}
\item \textbf{Positivity is the corollary ``something exists''; forced existence lands on the prescribed line} (I.1, I.2).
\item \textbf{The circle is the only self-closed existence; a collapse must split into at least two.} The ground state's zeros are always on the circle (V.3.3); a zero leaving the line comes with its mirror partner (V.2.1). Positivity is what keeps every zero self-paired on the circle.
\item \textbf{The sign from the virtual flux around the true zeros.} Each prime's wobble has mean exactly zero; the real place contributes the growing positive mean log(r / 2 pi); the sign holds only if the mean wins against every alignment of prime phases a test function can resolve (by Kronecker and Voronin they can align). The margin of 1e-9 says it wins, only just (V.3.2, V.3.4).
\item \textbf{The reductio as excision at G0.} An overtaking alignment would be a negative direction, that is an off-line pair; the reading says such an assertion is excised at once. In mathematics that needs a contradiction from arithmetic structure; duality alone gives none (V.2.1). The 17-million-line Lean proof by others achieved it by another route (a one-sided prime bound plus Landau).
\item \textbf{The functional equation, proven to be locally evaluable} is Tate's thesis (V.2.2); it gives the mirror, not the sign.
\item \textbf{GRH as the ODP's tool that runs before causality.} RH is equivalent to M(x) = O(x\textasciicircum{}\{1/2 + eps\}): the building blocks deviate from their mean density by no more than the square-root floor of independent fluctuations (``Moebius randomness'') \textbf{[classical]} for the equivalence; reading GRH as the certificate that the blocks sit at the universal floor is interpretation.
\item ``What GRH is missing is its concept of the remnant of acts'': face 3; ``the act must entangle the real with the imaginary through a coupling factor'': face 4; ``arithmetic exists but never runs'': there is no Frobenius over Z \textbf{[classical]}, and running it as content is face 2.
\end{itemize}
\subsection*{V.7 What would close it}
\addcontentsline{toc}{subsection}{V.7 What would close it}

Any one of the eight faces, built without the zeros. The constraint from V.3 rules out every assembly with slack, by place, by subset or by separate terms, so the construction must be an exact identity coupling all primes up to the scale and the real place. The Yang-Mills chain (Part IV) is the model: there the positive object (e\textasciicircum{}\{-S\} with locality) is given by construction, and reflection positivity follows. The open step here is to find the arithmetic G0 in which Weil positivity is a construction rather than a conclusion (Part III).

\subsection*{V.8 The box in the one-box test}
\addcontentsline{toc}{subsection}{V.8 The box in the one-box test}

\textbf{Gate factor and marker.} The ticks are the scales, each carrying the margin to the next. The owner reads each scale as one, and so as existence itself \textbf{[reading]}. At each scale the gate reads inside when the margin is carried to the next scale with a nonnegative factor, the sign at that scale; the marker is the claim ``inside''. The box of a finite zero model passes exactly when the sign holds at every scale, the card's open step (\texttt{ScaleSign}) \textbf{[Lean]} (\texttt{grh\_box\_iff}, \texttt{ETM/CardBox.lean}; claim V.8a). For a model of zeta's zeros with nonnegative weights it passes under Mathlib's statement of the Riemann hypothesis, taken as a named \texttt{External} slot \textbf{[Lean]} (\texttt{grh\_box\_of\_RH}). The first ticks are certified: at the start \{infinity, 2, 3\} the sign holds with explicit margins (III.65, \textbf{[certified]}). The box is the card's open step in the test's language (II.4): it restates the step and proves no new bound, and the content stays in the step \textbf{[paper]}. Its open ticks are the scales beyond the start (V.4).


%% ---- 06-part-VI-pnp.md
\section*{Part VI. P vs NP: inside a certified class, and across it}
\addcontentsline{toc}{section}{Part VI. P vs NP: inside a certified class, and across it}

\subsection*{VI.0 What is proved and what is not}
\addcontentsline{toc}{subsection}{VI.0 What is proved and what is not}

\textbf{This part does not resolve P vs NP, nor NP vs coNP.} It proves, in Lean, the sampling side of the chain of II.8 for finite systems (VI.1) and the exact parts of the dichotomy for refutations (VI.3). It then takes the strongest proof system in which the structureless case is open, cutting planes on random 3-CNF, and locates the open step there to one missing cost measure (VI.4, VI.5). Every closing argument must pass the relativization, natural-proofs and algebrization barriers (VI.7).

Two notions of the problem are used:

\begin{itemize}
\item \textbf{P vs NP} (decision, search, counting): \texttt{DICTIONARY.md}.
\item \textbf{NP vs coNP} via proof systems: a polynomially bounded proof system exists iff NP = coNP \textbf{[classical]} (Cook-Reckhow 1979); NP != coNP implies P != NP. \texttt{PROOF-SYSTEMS.md}. The owner's title names NP vs coNP for this reason: the proof-complexity side is where this part's work sits.
\end{itemize}
\subsection*{VI.1 Inside a certified class: the chain holds (links 1 to 4)}
\addcontentsline{toc}{subsection}{VI.1 Inside a certified class: the chain holds (links 1 to 4)}

A finite spin system with positive weights whose Dobrushin influences sum to less than one is easy in every sense: its sampler mixes fast, its dynamics has a spectral gap, and through self-reducibility counting and search follow. All the finite pieces are now unconditional in Lean: the couplings that earlier files assumed are built from the classical objects, and the spectral measure is proved from Mathlib's finite-dimensional spectral theorem (\texttt{ITEMS.md}).

{\footnotesize
\begin{longtable}{>{\raggedright\arraybackslash}p{\dimexpr 0.0935\linewidth-2\tabcolsep\relax}>{\raggedright\arraybackslash}p{\dimexpr 0.4101\linewidth-2\tabcolsep\relax}>{\raggedright\arraybackslash}p{\dimexpr 0.1709\linewidth-2\tabcolsep\relax}>{\raggedright\arraybackslash}p{\dimexpr 0.3056\linewidth-2\tabcolsep\relax}}
\toprule
\textbf{\#} & \textbf{Result} & \textbf{Status} & \textbf{Source} \\
\midrule
\endhead
VI.1.1 & Dobrushin row sums at most alpha \textless{} 1 (finite system, positive weights): the heat-bath sampler is within n alpha\textasciicircum{}t of equilibrium after t sweeps; random scan under column sums & \textbf{[Lean]} & \texttt{Mixing/Classical.lean}, \texttt{tv\_sweep\_dobrushin}, \texttt{tv\_rs\_dobrushin}, \texttt{mixing\_sweep\_dobrushin} \\
VI.1.2 & The column condition gives a spectral gap of at least (1 - alpha)/n per single-site step; the spectral measure exists and lies in [0, 1] & \textbf{[Lean]} & \texttt{Mixing/FiniteSpectral.lean}, \texttt{spectral\_gap\_dobrushin}, \texttt{timeCorr\_spectral}, \texttt{exists\_spectral\_measure} \\
VI.1.3 & Gap in time iff no spectral mass above e\textasciicircum{}\{-m\} & \textbf{[Lean]} & \texttt{Mixing/TimeGap.lean}, \texttt{time\_gap\_iff} \\
VI.1.4 & Doeblin contraction, sharpness, eta\textasciicircum{}n along a homogeneous chain, the box bound under a density floor & \textbf{[Lean]} & \texttt{Mixing/Doeblin.lean}, \texttt{contracts\_of\_doeblin}, \texttt{doeblin\_sharp}; \texttt{Mixing/Chain.lean}, \texttt{contracts\_kpow}, \texttt{box\_bound} \\
VI.1.5 & Dobrushin condition gives a unique Gibbs state; spatial mixing gives a volume-uniform gap & \textbf{[classical]} & Dobrushin 1968, 1970; Stroock-Zegarlinski 1992; Martinelli-Olivieri 1994 \\
VI.1.6 & Polynomial sampling gives polynomial approximate counting for self-reducible problems, and back & \textbf{[classical]} & Jerrum-Valiant-Vazirani 1986 \\
VI.1.7 & Hard-core model below lambda\_c(D): deterministic counting (correlation decay) and O(n log n) mixing & \textbf{[classical]} & Weitz 2006; Anari-Liu-Oveis Gharan 2020; Chen-Liu-Vigoda 2021 \\
VI.1.8 & \textbf{The certified class} (influence sums below alpha after blocking, or the Lovasz local lemma condition): existence guaranteed, search, sampling and approximate counting polynomial & \textbf{[classical]} assembled; mixing step \textbf{[Lean]} & \texttt{DICTIONARY.md}, 6.2; Lovasz 1975; Moser-Tardos 2010; Moitra 2019 \\
VI.1.9 & Not yet unconditional: infinite or continuous single-site spaces (the gauge group), hard constraints, systematic scan under column sums & \textbf{[open]} in Lean; \textbf{[classical]} on paper & \texttt{ITEMS.md}, 3 \\
\bottomrule
\end{longtable}}

The finite results are the same theorems Part IV uses for Yang-Mills (II.12); the gauge-theory instance still rests on the couplings and the L\textasciicircum{}2 spectral theorem as hypotheses (VI.1.9).

\subsection*{VI.2 At the edge of the class: what is known to go wrong}
\addcontentsline{toc}{subsection}{VI.2 At the edge of the class: what is known to go wrong}

{\footnotesize
\begin{longtable}{>{\raggedright\arraybackslash}p{\dimexpr 0.1029\linewidth-2\tabcolsep\relax}>{\raggedright\arraybackslash}p{\dimexpr 0.4433\linewidth-2\tabcolsep\relax}>{\raggedright\arraybackslash}p{\dimexpr 0.1921\linewidth-2\tabcolsep\relax}>{\raggedright\arraybackslash}p{\dimexpr 0.2417\linewidth-2\tabcolsep\relax}}
\toprule
\textbf{\#} & \textbf{Result} & \textbf{Status} & \textbf{Source} \\
\midrule
\endhead
VI.2.1 & Antiferromagnetic two-spin systems beyond uniqueness: no FPRAS unless NP = RP & \textbf{conditional} (NP != RP) & Sly 2010; Sly-Sun 2014; Galanis-Stefankovic-Vigoda 2015, 2016 \\
VI.2.2 & Bounded-degree k-CNF: approximate counting NP-hard above degree c 2\textasciicircum{}\{k/2\} & \textbf{conditional} & Bezakova-Galanis-Goldberg-Guo-Stefankovic 2019 \\
VI.2.3 & The ferromagnetic Ising model loses uniqueness but keeps an FPRAS at every temperature & \textbf{[classical]} & Jerrum-Sinclair 1993 \\
VI.2.4 & Lattice gauge theory with the Wilson action is on the ferromagnetic side: slow dynamics at phase boundaries, but no NP-hardness from coexistence is predicted & \textbf{[reading]} & \texttt{DICTIONARY.md}, 5 \\
VI.2.5 & Putting all counting instances inside the uniqueness region is as strong as NP = RP; reducing SAT into the LLL class is impossible as stated (all LLL instances are satisfiable) & \textbf{conditional} / \textbf{[classical]} & \texttt{DICTIONARY.md}, 6.3 \\
\bottomrule
\end{longtable}}

\subsection*{VI.3 Refutations: the dichotomy transferred (NP vs coNP)}
\addcontentsline{toc}{subsection}{VI.3 Refutations: the dichotomy transferred (NP vs coNP)}

A refutation proves that a formula has no satisfying assignment: the remnant ``it does not exist''. Resolution always finds one, and every refutation in every sound two-premise system must, somewhere, cross from local facts to a global one. That crossing is universal; its cost is not.

\textbf{Confinement times resolution (the size-width relation).} In resolution the cost of the crossing has two terms: the width the crossing must assert (confinement: the cut principle, width at least the boundary), and the share of it one restriction resolves (resolution). The DAG-like size-width relation is now in Lean: if width w + d + b is impossible, every DAG-like refutation has size at least exp(b (w + d + 1)/2n) \textbf{[Lean]} (\texttt{Proof/Dag.lean}, \texttt{size\_width}, \texttt{size\_width\_exp}). With the cut principle the same bound holds whenever w + d + b \textless{} e, e the boundary of every middle region \textbf{[Lean]} (\texttt{dag\_size\_ge}); optimising b gives exp((e - w)\textasciicircum{}2 / 8n) up to rounding \textbf{[paper]}. This is Ben-Sasson-Wigderson (2001) \textbf{[classical]}, now machine-checked.

{\footnotesize
\begin{longtable}{>{\raggedright\arraybackslash}p{\dimexpr 0.0952\linewidth-2\tabcolsep\relax}>{\raggedright\arraybackslash}p{\dimexpr 0.4278\linewidth-2\tabcolsep\relax}>{\raggedright\arraybackslash}p{\dimexpr 0.2043\linewidth-2\tabcolsep\relax}>{\raggedright\arraybackslash}p{\dimexpr 0.2526\linewidth-2\tabcolsep\relax}}
\toprule
\textbf{\#} & \textbf{Result} & \textbf{Status} & \textbf{Source} \\
\midrule
\endhead
VI.3.1 & Resolution is sound and refutation-complete: refutable iff unsatisfiable & \textbf{[Lean]} & \texttt{Proof/Resolution.lean}; \texttt{Proof/Complete.lean}, \texttt{refutable\_iff} \\
VI.3.2 & No inference creates assertion; every refutation has a middle-sized clause; resolution width is at least the boundary count & \textbf{[Lean]} & \texttt{Proof/Cut.lean}, \texttt{exists\_middle}, \texttt{width\_ge} \\
VI.3.3 & Tree-like size-width: width W forces more than 2\textasciicircum{}\{W - w\} leaves & \textbf{[Lean]} & \texttt{Proof/SizeWidth.lean}, \texttt{tl\_size\_gt} \\
VI.3.4 & DAG-like size-width: size \textgreater{}= exp(b (w + d + 1)/2n); with the cut principle for w + d + b \textless{} e, hence exp((e - w)\textasciicircum{}2 / 8n) & \textbf{[Lean]} (the optimisation over b \textbf{[paper]}; classical: Ben-Sasson-Wigderson 2001) & \texttt{Proof/Dag.lean}, \texttt{size\_width}, \texttt{size\_width\_exp}, \texttt{dag\_size\_ge} \\
VI.3.5 & Tseitin on expanders: resolution width at least the expansion, with two stated hypotheses (sub-systems satisfiable; expansion count) & \textbf{[Lean]} with hypotheses & \texttt{Proof/Tseitin.lean}, \texttt{width\_ge}, \texttt{tl\_size\_ge} \\
VI.3.6 & Gauss law mod 2: boundary parity equals interior charge; odd charge is unsatisfiable & \textbf{[Lean]} & \texttt{Proof/Crossing.lean}, \texttt{gauss}, \texttt{unsat\_of\_odd} \\
VI.3.7 & \textbf{Every refutation crosses:} in every sound system with at most two premises per rule, every refutation of an expanding formula has a line sensitive to at least the expansion many variables & \textbf{[Lean]} with the two hypotheses of VI.3.5 & \texttt{Proof/Crossing.lean}, \texttt{cross\_any\_system} \\
VI.3.8 & The crossing is universal, its cost is not: one XOR names the Tseitin flux; resolution pays 2\textasciicircum{}\{b-1\} clauses of width b for it, Frege about b\textasciicircum{}2 symbols & \textbf{[Lean]} (\texttt{summary}) / \textbf{[classical]} & \texttt{DICHOTOMY-TRANSFER.md}, 3 \\
VI.3.9 & Known lower bounds: resolution (Haken 1985; Urquhart 1987; Chvatal-Szemeredi 1988), bounded-depth Frege (Ajtai 1988; Pitassi-Beame-Impagliazzo 1993; Krajicek-Pudlak-Woods 1995), cutting planes (Pudlak 1997) & \textbf{[classical]} & \texttt{PROOF-SYSTEMS.md}, 2 \\
VI.3.10 & Feasible interpolation for resolution and cutting planes; its failure for Frege and extended Frege under cryptographic assumptions & \textbf{[classical]} / \textbf{conditional} & Krajicek 1997; Pudlak 1997; Krajicek-Pudlak 1998; Bonet-Pitassi-Raz 2000 \\
VI.3.11 & No superpolynomial lower bound is known for AC\textasciicircum{}0[p]-Frege, Frege or extended Frege; best formula lower bounds about n\textasciicircum{}3 & \textbf{[open]} / \textbf{[classical]} & Hastad 1998; Tal 2014 \\
\bottomrule
\end{longtable}}

\subsection*{VI.4 The edge: cutting planes on random 3-CNF}
\addcontentsline{toc}{subsection}{VI.4 The edge: cutting planes on random 3-CNF}

\textbf{Why this system and this family.} Cutting planes (CP) is the strongest system in which random CNF is known to be hard only when the clauses are wide enough: width Theta(log n) is solved, constant width is open \textbf{[classical]}. The other candidates are either solved (polynomial calculus, sum-of-squares degree) or open for every formula, structured or not (Res(xor), bounded-depth Frege with parity gates), so randomness is not their obstacle. The open part of CP is exactly the structureless, constant-width case \textbf{[classical]} + \textbf{[reading]} for the choice (\texttt{EDGE.md}, 0.1).

\textbf{Target.} For a large constant Delta and a random 3-CNF F with Delta n clauses: with high probability every CP refutation of F has size 2\textasciicircum{}(n\textasciicircum{}Omega(1)) \textbf{[open]}.

\textbf{What closes.} (1) Incompressibility of the specific formula gives expansion: a small clause set spanning too few variables could be described in fewer bits than its length. This is a statement about one input, not about typical Boolean functions, so the natural-proofs barrier does not bite here \textbf{[classical]} (counting) + \textbf{[reading]} (the form). (2) Expansion forces the crossing in CP: CP's rules are sound with at most two premises, so \texttt{cross\_any\_system} applies \textbf{[Lean]}.

\textbf{Where it stops.} In resolution the crossing line is a clause of width e, and any one literal kills it under a restriction, which turns ``sensitive to e variables'' into size. In CP that step is false: the majority line ``at least half of these e variables are true'' is sensitive to all e and survives every restriction of fewer than e/2 variables \textbf{[Lean]} (\texttt{Proof/Edge.lean}, \texttt{majority\_line}, \texttt{thr\_sensitive}, \texttt{thr\_survives}). The uncertainty term of the size-width relation collapses: one line can aggregate its units.

\textbf{The duality floor, and where it fails.} For a unit threshold ``at least k of these e variables'', forcing it true needs k fixed variables and forcing it false needs e - k + 1; the two sum to e + 1, so every such line pays at least (e + 1)/2 in one channel \textbf{[Lean]} (\texttt{Proof/CPDuality.lean}, \texttt{thr\_duality}). This is the CP form of max(F\_tw, F\_e) \textgreater{}= c. It fails for weighted lines: the line r x1 + r x2 + (sum of r further variables) \textgreater{}= 2r is sensitive to all r + 2 variables and is forced true by fixing two and false by fixing two \textbf{[Lean]} (\texttt{two\_cheap}). So no measure on single lines built from certificates gives the floor.

\textbf{Pairs, and the interface.} Feasible interpolation splits the formula as A(x, y) and B(x, z) and every line into an A-part and a B-part sharing only x: two derived lines that share existence (x) and not possibility. The two boundary laws are disjoint exactly when the formula is unsatisfiable \textbf{[Lean]} (\texttt{split\_unsat\_iff}). The cost is charged to the interface, a monotone real circuit (Pudlak 1997) or a dag-like communication protocol \textbf{[classical]}. That passes clique-colouring and random k-CNF with k = Theta(log n) \textbf{[classical]} (Hrubes-Pudlak; Fleming-Pankratov-Pitassi-Robere; Sokolov) and fails at constant width, where the interface carries O(1) bits per clause. Every CP refutation of a 3-CNF contains a creative pair, a rounding of a combination of at least two lines, while rounding a single axiom creates nothing; additive pairs cross for free \textbf{[Lean]} (\texttt{Proof/CPPairs.lean}, \texttt{creative\_pair}, \texttt{round\_single}, \texttt{fan\_cheap}). The pairwise form gives a floor in kind and in depth, not in size.

\textbf{Positivity runs to rank, not size.} The all-1/2 point satisfies every clause of width at least two and every sum of satisfied lines \textbf{[Lean]} (\texttt{half\_clause}); rounding breaks it \textbf{[Lean]} (\texttt{half\_breaks\_rounding}). Fractional points surviving r rounds give CP rank lower bounds, known for random k-CNF (Buresh-Oppenheim et al.) \textbf{[classical]}, and rank bounds size only linearly. SOS-degree hardness does not transfer to CP size: Tseitin needs linear SOS degree (Grigoriev 2001) but has quasi-polynomial CP refutations (Dadush-Tiwari 2020) \textbf{[classical]}.

\textbf{The first unclosable step.} Both transfers stop at one place: \textbf{interface hardness at width 3}, a superpolynomial lower bound on dag-like protocols for the falsified-clause search of a random 3-CNF under some balanced split. Equivalently, a cost measure W\_CP on CP lines that the crossing forces on every refutation of an expanding 3-CNF and that is expensive in size \textbf{[open]}. \textbf{No new lower bound} is claimed.

{\footnotesize
\begin{longtable}{>{\raggedright\arraybackslash}p{\dimexpr 0.1066\linewidth-2\tabcolsep\relax}>{\raggedright\arraybackslash}p{\dimexpr 0.3704\linewidth-2\tabcolsep\relax}>{\raggedright\arraybackslash}p{\dimexpr 0.2103\linewidth-2\tabcolsep\relax}>{\raggedright\arraybackslash}p{\dimexpr 0.2927\linewidth-2\tabcolsep\relax}}
\toprule
\textbf{\#} & \textbf{Result} & \textbf{Status} & \textbf{Source} \\
\midrule
\endhead
VI.4.1 & CP on random 3-CNF is the sharpest open case (log n width solved, constant width open) & \textbf{[classical]} + \textbf{[reading]} & \texttt{EDGE.md}, 0.1 \\
VI.4.2 & Incompressibility of the one formula gives expansion; natural proofs do not bite & \textbf{[classical]} + \textbf{[reading]} & \texttt{EDGE.md}, 0.3 \\
VI.4.3 & Expansion forces a crossing line in every CP refutation & \textbf{[Lean]} (hypotheses of VI.3.5) & \texttt{Proof/Crossing.lean}, \texttt{cross\_any\_system} \\
VI.4.4 & The majority line is sensitive to all e variables and survives restrictions of fewer than e/2 & \textbf{[Lean]} & \texttt{Proof/Edge.lean}, \texttt{majority\_line}, \texttt{thr\_sensitive}, \texttt{thr\_survives} \\
VI.4.5 & Unit thresholds: duality floor (e + 1)/2 & \textbf{[Lean]} & \texttt{Proof/CPDuality.lean}, \texttt{thr\_duality} \\
VI.4.6 & Weighted lines are cheap in both channels & \textbf{[Lean]} & \texttt{Proof/CPDuality.lean}, \texttt{two\_cheap} \\
VI.4.7 & Interpolation split: disjoint boundary laws iff unsatisfiable & \textbf{[Lean]} (logical core) + \textbf{[classical]} (interface bounds) & \texttt{Proof/CPDuality.lean}, \texttt{split\_unsat\_iff}; Pudlak 1997 \\
VI.4.8 & Every CP refutation of a 3-CNF contains a creative pair; single roundings create nothing; additive pairs are cheap & \textbf{[Lean]} & \texttt{Proof/CPPairs.lean}, \texttt{creative\_pair}, \texttt{round\_single}, \texttt{fan\_cheap} \\
VI.4.9 & The all-1/2 point survives sums, not rounding; rank, not size; SOS degree does not transfer & \textbf{[Lean]} + \textbf{[classical]} & \texttt{Proof/CPDuality.lean}, \texttt{half\_clause}, \texttt{half\_breaks\_rounding}; Grigoriev 2001; Dadush-Tiwari 2020 \\
VI.4.10 & \textbf{Interface hardness at width 3} (equivalently a size-forcing cost measure W\_CP) & \textbf{[open]} & \texttt{EDGE.md}, 0.4; \texttt{CP-DUALITY.md}, 0.5 \\
\bottomrule
\end{longtable}}

\subsection*{VI.5 The native method and the floor}
\addcontentsline{toc}{subsection}{VI.5 The native method and the floor}

The native method makes each layer of a proof one decision and then computes: merge single-purpose lemmas, turn finite case analysis into a checker proved equal to the statement, let automation carry regularity. Every pass lowers a natural-number measure, so the method stops \textbf{[Lean]} (\texttt{fixed\_after}); without quanta it would not (\texttt{real\_never\_arrives}). Depth contracts by 1/2 under pointer jumping (\texttt{iterate\_clog\_pairUp}), while the total checking work is conserved (\texttt{sum\_pairUp}) \textbf{[Lean]}. Applied to the bottom of the \texttt{yang\_mills\_assembly} chain with the same statements re-checked: 16 theorems to 12, 5 layers to 4, term size -6 \%, kernel time unchanged within noise \textbf{[Lean]} + \textbf{[numerical]}.

{\footnotesize
\begin{longtable}{>{\raggedright\arraybackslash}p{\dimexpr 0.1044\linewidth-2\tabcolsep\relax}>{\raggedright\arraybackslash}p{\dimexpr 0.3963\linewidth-2\tabcolsep\relax}>{\raggedright\arraybackslash}p{\dimexpr 0.1772\linewidth-2\tabcolsep\relax}>{\raggedright\arraybackslash}p{\dimexpr 0.3020\linewidth-2\tabcolsep\relax}}
\toprule
\textbf{\#} & \textbf{Result} & \textbf{Status} & \textbf{Source} \\
\midrule
\endhead
VI.5.1 & Quantized descent arrives; real descent need not; pointer jumping halves depth; work conserved (II.15) & \textbf{[Lean]} & \texttt{Native/Method.lean}, \texttt{fixed\_after}, \texttt{real\_never\_arrives}, \texttt{iterate\_clog\_pairUp}, \texttt{sum\_pairUp}, \texttt{reaches\_floor} \\
VI.5.2 & Applied to a real chain: same statements, 16 theorems to 12, 5 layers to 4 & \textbf{[Lean]} + \textbf{[numerical]} & \texttt{Native/RouteBChain.lean}; \texttt{NATIVE-METHOD.md}, 0.3 \\
VI.5.3 & The method's floor is per proof and lies at or above the minimal refutation; the method bounds no minimum & \textbf{[reading]} & \texttt{NATIVE-METHOD.md}, 5 \\
VI.5.4 & Finding a resolution refutation polynomial in the shortest one is NP-hard & \textbf{[classical]} & Atserias-Mueller 2019 \\
\bottomrule
\end{longtable}}

\subsection*{VI.6 The open steps}
\addcontentsline{toc}{subsection}{VI.6 The open steps}

{\footnotesize
\begin{longtable}{>{\raggedright\arraybackslash}p{\dimexpr 0.0975\linewidth-2\tabcolsep\relax}>{\raggedright\arraybackslash}p{\dimexpr 0.4678\linewidth-2\tabcolsep\relax}>{\raggedright\arraybackslash}p{\dimexpr 0.1689\linewidth-2\tabcolsep\relax}>{\raggedright\arraybackslash}p{\dimexpr 0.2458\linewidth-2\tabcolsep\relax}}
\toprule
\textbf{\#} & \textbf{Open step} & \textbf{Form} & \textbf{Source} \\
\midrule
\endhead
N1 & Is the irreducible core of every refutation polynomial in the question? (the size of the remnant) & NP vs coNP & \texttt{docs/OPEN-QUESTIONS-DMP.md}, 4 \\
N2 & Interface hardness at width 3 for cutting planes on random 3-CNF (VI.4.10) & NP vs coNP, for CP & \texttt{EDGE.md}, 0.4; \texttt{CP-DUALITY.md}, 0.5 \\
N3 & Step T and step H for Frege: a short Frege refutation of a random 3-CNF yields a small separating formula for a balanced cut, and that pair has none of polynomial size (beyond the n\textasciicircum{}3 barrier) & NP vs coNP, for Frege & \texttt{DICHOTOMY-TRANSFER.md}, 4 to 6 \\
N4 & From the certified class to all inputs: is every instance's existence decidable cheaply? & P vs NP & \texttt{DICTIONARY.md}, 6.3, 8 \\
\bottomrule
\end{longtable}}

N2 is new in this draft and is the sharpest: a named system, a named family and one missing object. It would not settle NP vs coNP (CP is not the strongest system), but it is the first place where the structureless case meets a system strong enough to name a crossing in one line.

\subsection*{VI.7 The barriers}
\addcontentsline{toc}{subsection}{VI.7 The barriers}

\textbf{Every closing argument must pass three barriers} \textbf{[classical]}: relativization (Baker-Gill-Solovay 1975; the coupling proofs and the gadget reductions both relativize), natural proofs (Razborov-Rudich 1997; ``structureless means hard'' is large and constructive when stated about typical functions; VI.4.2 avoids it by speaking about one formula), algebrization (Aaronson-Wigderson 2009; partition-function polynomials and zero-free regions algebrize). Williams' route (2014) evades all three, but no faster-than-brute-force satisfiability algorithm follows from local certification \textbf{[open]}.

\subsection*{VI.8 DMP readings attached to this part (all [reading])}
\addcontentsline{toc}{subsection}{VI.8 DMP readings attached to this part (all [reading])}

\begin{itemize}
\item \textbf{Two categories of proof lines.} Lines derived directly from the axioms are cheap and mechanically checkable; new derivations (lemmas, cuts, extension variables) are what make short proofs possible, and eliminating them can blow a proof up exponentially. The owner's ``first true act'' is two new derivations used together, sharing existence but not possibility: in CP, the interpolation split (VI.4.7) and the creative pair (VI.4.8).
\item \textbf{The pairs of derived lines are the entanglement through which the HTF acts in NP}: the size floor needs the pair to be charged at its interface (VI.4.10).
\item \textbf{The lower bound is confinement plus the needed uncertainty}: the size-width relation (VI.3.4); in CP the uncertainty term collapses because one line aggregates its units (VI.4.4), and W\_CP is the missing term.
\item \textbf{``P = NP within a DB, P != NP across it''}: the first half is VI.1.8; the second is not proved (VI.6).
\item Describable =\textgreater{} exists =\textgreater{} positive: SAT iff \#SAT \textgreater{} 0. A contradiction excised leaves ``it does not exist'': refutation completeness (VI.3.1). Inward and outward reach: Dobrushin column and row sums.
\item ``A question with no structure has no basis/language in which to speak an answer'': incompressibility of the one formula (VI.4.2); the care point is the natural-proofs barrier, avoided only by speaking about the specific formula.
\end{itemize}
\subsection*{VI.9 What would close it}
\addcontentsline{toc}{subsection}{VI.9 What would close it}

For NP vs coNP through proof complexity: superpolynomial lower bounds for every proof system, which would include N3 for Frege. The nearest step is N2, interface hardness at width 3 for cutting planes, which would be a new lower bound for a natural system on the structureless family, though not a resolution of NP vs coNP. For P vs NP: N4, a barrier-passing argument from the certified class to all inputs. Nothing in this part closes any of them.

\subsection*{VI.10 The box in the one-box test}
\addcontentsline{toc}{subsection}{VI.10 The box in the one-box test}

\textbf{Gate factor and marker.} The ticks are the \textbf{bounds} (the owner's name; in plain terms, the restrictions): each bounds the share of costly lines a restriction keeps. At each bound the gate reads inside when the restriction keeps at most a 1 - theta share of the costly lines; the marker is the claim ``inside''. The box passes exactly when the card's open step (\texttt{CrossingCost}) holds, and a passing box gives the crossing bound: every refutation has at least e\textasciicircum{}(theta b) lines \textbf{[Lean]} (\texttt{np\_box\_concl}, \texttt{ETM/CardBox.lean}; claim VI.10a). The box is the card's open step in the test's language (II.4): it restates the step and proves no new bound, and the content stays in the step \textbf{[paper]}. Its open ticks are the bounds of the lifting step at width 3 (VI.4.10).


%% ---- 07-part-VII-navier-stokes.md
\section*{Part VII. Navier-Stokes: the card, the claimed certificate, and Poincare as validation}
\addcontentsline{toc}{section}{Part VII. Navier-Stokes: the card, the claimed certificate, and Poincare as validation}

\subsection*{VII.0 What is proved, by whom, and what this part does}
\addcontentsline{toc}{subsection}{VII.0 What is proved, by whom, and what this part does}

\textbf{The Clay problem} (Fefferman's official statement) asks one of four things. For three-dimensional incompressible Navier-Stokes with viscosity nu \textgreater{} 0 and force f:

\begin{itemize}
\item \textbf{(A)} On R\textasciicircum{}3, with \textbf{zero force}, every smooth, divergence-free datum whose derivatives decay faster than any power has a smooth solution for all time with bounded energy.
\item \textbf{(B)} The same on the torus R\textasciicircum{}3/Z\textasciicircum{}3.
\item \textbf{(C)} On R\textasciicircum{}3 there are a smooth, rapidly decaying datum \textbf{and a smooth force} for which no such global solution exists.
\item \textbf{(D)} The same on the torus, again with a force.
\end{itemize}
A proof of any one wins the prize \textbf{[classical]}. Two points matter below. (C) and (D) allow a force, so a forced blowup of the admitted class would settle (C) as stated. A blowup with \textbf{no} force disproves (A) (or (B)) and so proves (C) (or (D)) with f = 0.

\textbf{State of the problem (2026-11-24).} Two claimed proofs of blowup were posted in autumn 2026. Neither is refereed, and neither is checked in this project.

\begin{itemize}
\item \textbf{OpenAI, 8 September 2026:} blowup with a smooth, compactly supported force. It addresses (C) and (D) \textbf{[external, unverified]}.
\item \textbf{Alpoge and Buckmaster, 20 November 2026:} blowup with \textbf{no force}. It addresses (A) and (B). It rests on a numerical certificate that is \textbf{announced, not released} (only its hash is public). \textbf{[external, unverified, certificate unreleased]}.
\end{itemize}
If either claim holds, the Clay problem is answered by others, in the direction of blowup. \textbf{This part proves no statement about the real Navier-Stokes equations.} The status of the problem in this paper is: (A) and (B) are open here, and an external claim asserts that they are false.

\textbf{What this part does.} It puts Navier-Stokes on the DMP API as a card (VII.2), and proves in Lean, as a skeleton, what the form layer says about the two possible blowup routes and about regularity (VII.3). It records what is borrowed and its exact status (VII.4), states the one open step (VII.5), and checks the restatement against Perelman's proof of Poincare (A). The Lean modules contain \textbf{no Navier-Stokes object}: nothing in them is about the equations themselves.

\subsection*{VII.1 The problem as stated, and what is known}
\addcontentsline{toc}{subsection}{VII.1 The problem as stated, and what is known}

\textbf{VII.1.1 No-go results} \textbf{[classical]} (recalled from the literature; not re-read for this draft).

{\footnotesize
\begin{longtable}{>{\raggedright\arraybackslash}p{\dimexpr 0.6274\linewidth-2\tabcolsep\relax}>{\raggedright\arraybackslash}p{\dimexpr 0.3526\linewidth-2\tabcolsep\relax}}
\toprule
\textbf{Result} & \textbf{What it rules out} \\
\midrule
\endhead
Leray 1934: a singularity at T* forces \texttt{sup abs(u) \textgreater{}= c (nu/(T* - t))\textasciicircum{}\{1/2\}} and \texttt{norm(grad u)\_\{L\textasciicircum{}2\} \textgreater{}= c (T* - t)\textasciicircum{}\{-1/4\}} & blowup slower than these rates \\
Necas, Ruzicka, Sverak 1996; Tsai 1998: no nonzero backward self-similar solution \texttt{(T* - t)\textasciicircum{}\{-1/2\} U(x / sqrt(T* - t))} with \texttt{U} in \texttt{L\textasciicircum{}3} (Tsai: \texttt{L\textasciicircum{}p}, \texttt{3 \textless{} p \textless{} infinity}, and under local energy bounds) & exact self-similarity at the Leray scaling \\
Caffarelli, Kohn, Nirenberg 1982: the singular set of a suitable weak solution has zero one-dimensional parabolic measure & singularities along curves; isolated points remain \\
Escauriaza, Seregin, Sverak 2003; Seregin 2012: a solution bounded in \texttt{L\textasciicircum{}infinity\_t L\textasciicircum{}3\_x} is smooth; at a singular time the \texttt{L\textasciicircum{}3} norm must tend to infinity & blowup with a bounded scale-invariant norm \\
Axisymmetric Type I exclusion (Chen-Strain-Tsai-Yau; Koch-Nadirashvili-Seregin-Sverak) & axisymmetric Type I blowup \\
Chae-Wolf; Pineau-Vicol: a smooth backward discretely self-similar solution obeys the Type I bound; a nonzero one needs a period bounded below & discretely self-similar solutions with a too-small period \\
\bottomrule
\end{longtable}}

What these leave open: \textbf{discretely self-similar} solutions (the shape repeats after a fixed zoom), asymptotically self-similar solutions, and Type II rates.

\textbf{VII.1.2 Tao's averaged equations} \textbf{[classical]} (Tao, J. Amer. Math. Soc. 29 (2016)). Replace the nonlinearity by an average of rotated and dilated copies. The average keeps the energy identity and the standard estimates, and the averaged equation blows up in finite time. So a proof of global regularity must use finer structure of the true nonlinearity than the energy identity and those estimates.

\textbf{VII.1.3 Numerical evidence by others} \textbf{[external]}. Hou (arXiv 2107.06509, 2022) reports an axisymmetric scenario with smooth data and constant viscosity: vorticity growth of about 10\textasciicircum{}7 at nearly self-similar Leray-type rates. It is numerical evidence, not a proof. Strict Type I axisymmetric blowup is excluded (VII.1.1), so a singular scenario would have to deviate from the observed rate; the tension is noted here, not resolved.

\subsection*{VII.2 The card}
\addcontentsline{toc}{subsection}{VII.2 The card}

The DMP API (\texttt{docs/unified/API.md}) gives every problem one shape: a well-defined input, a proved structure on it, and one named open step in a slot, with external inputs in their own slot. Navier-Stokes has three cards in \texttt{Unified/NavierStokesAPI.lean} (namespace \texttt{ProofPhase.API.NavierStokes}). The table fills the API's template.

{\footnotesize
\begin{longtable}{>{\raggedright\arraybackslash}p{\dimexpr 0.1940\linewidth-2\tabcolsep\relax}>{\raggedright\arraybackslash}p{\dimexpr 0.5129\linewidth-2\tabcolsep\relax}>{\raggedright\arraybackslash}p{\dimexpr 0.2732\linewidth-2\tabcolsep\relax}}
\toprule
\textbf{API concept} & \textbf{Navier-Stokes} & \textbf{Lean} \\
\midrule
\endhead
\textbf{Input} & route I: the perturbation states, the period map, the size (rescaled \texttt{H\textasciicircum{}1}), a contraction constant \texttt{rho}, a quadratic constant \texttt{C}, a ball radius \texttt{delta} with \texttt{rho + C delta \textless{} 1}; route II: a profile certificate and a trajectory; regularity: a chain over dyadic scales & \texttt{DSSInput}, \texttt{RouteIInput}, \texttt{DyadicInput} \\
\textbf{Running system} & route I: states in the ball, one step is one period, factor \texttt{rho + C size}, mass the perturbation size & \texttt{system}, \texttt{contracting} \\
\textbf{Chain} & regularity: the quantity carried to dyadic scale \texttt{2\textasciicircum{}\{-k\}}, cost the share dissipated per scale, floor the share the transfer adds back & \texttt{regularity} \\
\textbf{Uncertainty} & route II: the perturbation, factor \texttt{q + C M}, fresh uncertainty the profile residual plus the viscous term & \texttt{eulerUncertainty}, \texttt{routeII\_bounded} \\
\textbf{Open-step slot} & route I: \texttt{PeriodContraction}; regularity: a uniform dissipation cost per dyadic scale above the transfer floor; the tie: the certificate & \texttt{PeriodContraction}, \texttt{regularity}, \texttt{abProblem} \\
\textbf{External-input slot} & Alpoge-Buckmaster's Theorems 7.1 and 1.1 and their Lean library & \texttt{ABInterface.thm71}, \texttt{.thm11}, \texttt{.formal} \\
\textbf{Computable layer} & the window decision; the settle count & \texttt{windowQ}, \texttt{windowQ\_correct}, \texttt{settles\_by} \\
\bottomrule
\end{longtable}}

The three cards, in the card shape:

{\footnotesize
\begin{longtable}{>{\raggedright\arraybackslash}p{\dimexpr 0.1309\linewidth-2\tabcolsep\relax}>{\raggedright\arraybackslash}p{\dimexpr 0.1801\linewidth-2\tabcolsep\relax}>{\raggedright\arraybackslash}p{\dimexpr 0.2487\linewidth-2\tabcolsep\relax}>{\raggedright\arraybackslash}p{\dimexpr 0.2471\linewidth-2\tabcolsep\relax}>{\raggedright\arraybackslash}p{\dimexpr 0.1733\linewidth-2\tabcolsep\relax}}
\toprule
\textbf{Card} & \textbf{Input} & \textbf{The one open step} & \textbf{Conclusion} & \textbf{External input} \\
\midrule
\endhead
\textbf{Route I} (\texttt{routeI}) & a non-steady periodic profile of the zoomed equation at the Leray scaling, and a perturbation in a ball & \texttt{PeriodContraction}: one period contracts the perturbation, after finitely many directions are removed by the datum \textbf{[open]} & the perturbation settles onto the profile within every epsilon; the periods fill the finite time 1; the \texttt{H\textasciicircum{}1} amplitude tends to infinity; the profile is not steady & none in the card; the content is what a certificate supplies \\
\textbf{Regularity} (\texttt{regularity}) & a chain over dyadic scales & every dyadic scale dissipates a uniform cost above the transfer floor \textbf{[open]} & the carried quantity decays at a positive rate and is summable over all scales & none \\
\textbf{External tie} (\texttt{abProblem}) & the interface \texttt{ABInterface} & the Alpoge-Buckmaster certificate satisfies their Appendix B (unreleased) \textbf{[open]} & not (A) and not (B), so (C) and (D) with zero force & their Theorems 7.1 and 1.1 and their Lean library, as \texttt{External} fields \\
\bottomrule
\end{longtable}}

Nothing is an axiom. An \texttt{OpenStep} or \texttt{External} slot is an ordinary hypothesis, and neither appears in the axiom audit. The three statements of the interface (a certificate exists; a profile exists; Fefferman's (A); (B)) are \textbf{named, not defined}: this library has no Navier-Stokes object. Replacing the names by real definitions is the first item of the audit (VII.5).

\subsection*{VII.3 What is proved}
\addcontentsline{toc}{subsection}{VII.3 What is proved}

All of this is \textbf{[Lean]}, standard axioms only (\texttt{docs/unified/navierstokes-axioms.log}: 52 declarations on the three standard axioms, 4 on none; the source contains no \texttt{sorry}, \texttt{admit}, \texttt{native\_decide} or \texttt{axiom}). None of it proves a statement about the equations. Each result says what the ledger gives from a hypothesis.

\subsubsection*{VII.3.1 The window}
\addcontentsline{toc}{subsubsection}{VII.3.1 The window}

Zoom at the singularity's own rate. If the singular region shrinks like \texttt{(T* - t)\textasciicircum{}alpha}, the viscous term in the zoomed frame has the weight \texttt{(T* - t)\textasciicircum{}\{1 - 2 alpha\}}. The local energy near the singular point scales like \texttt{(T* - t)\textasciicircum{}\{5 alpha - 2\}}. The scalings are \textbf{[classical]}; the arithmetic is \textbf{[Lean]}.

{\footnotesize
\begin{longtable}{>{\raggedright\arraybackslash}p{\dimexpr 0.1688\linewidth-2\tabcolsep\relax}>{\raggedright\arraybackslash}p{\dimexpr 0.4141\linewidth-2\tabcolsep\relax}>{\raggedright\arraybackslash}p{\dimexpr 0.1933\linewidth-2\tabcolsep\relax}>{\raggedright\arraybackslash}p{\dimexpr 0.2038\linewidth-2\tabcolsep\relax}}
\toprule
\textbf{\texttt{alpha}} & \textbf{Viscous floor} & \textbf{Local energy} & \textbf{Verdict} \\
\midrule
\endhead
\texttt{\textgreater{} 1/2} & grows without bound (\texttt{nsViscous\_tendsto\_atTop}); no one-scale certificate covers all steps (\texttt{ns\_no\_uniform\_certificate}) & bounded & excluded \\
\texttt{= 1/2} & constant (\texttt{nsViscous\_half}): part of the zoomed equation; Type I (\texttt{ns\_typeI\_iff}) & bounded & \textbf{route I}: the Leray end \\
\texttt{2/5 \textless{}= alpha \textless{} 1/2} & shrinks (\texttt{nsViscous\_le}), summable over steps (\texttt{nsViscous\_summable}) & bounded (\texttt{ns\_energy\_window}) & \textbf{route II}: Type II \\
\texttt{\textless{} 2/5} & shrinks & unbounded & excluded \\
\bottomrule
\end{longtable}}

\texttt{ns\_window} (\texttt{Unified/NavierStokes.lean}) says that a bounded viscous floor and a bounded local energy hold together exactly for \texttt{2/5 \textless{}= alpha \textless{}= 1/2}. \texttt{nsViscous\_bounded\_iff} says the floor is bounded exactly when \texttt{alpha \textless{}= 1/2}.

The window decision runs. \texttt{windowQ} sorts a rational exponent into one of the four rows, and \texttt{windowQ\_correct} (\texttt{Unified/NavierStokesAPI.lean}) proves the computed answer right for the real statement: ``inside'' exactly when the viscous floor and the local energy are bounded, ``Leray'' exactly when the Type I bound holds as well (\texttt{windowQ\_admissible\_iff}, \texttt{windowQ\_leray\_iff}, \texttt{windowQ\_half}). This follows the API's rule for a new instance: the decision runs, and a theorem proves it right.

\subsubsection*{VII.3.2 Route I: the Leray end, a discretely self-similar profile}
\addcontentsline{toc}{subsubsection}{VII.3.2 Route I: the Leray end, a discretely self-similar profile}

The hypothesis is a structure, \texttt{LerayDSSCertificate}: a profile \texttt{W} that is periodic in zoomed time with period \texttt{S \textgreater{} 0} and \textbf{not steady}; a lower bound \texttt{A \textgreater{} 0} on its rescaled size; a contraction \texttt{0 \textless{}= rho \textless{} 1} and a constant \texttt{C \textgreater{}= 0} such that over one period the perturbation obeys \texttt{M(n+1) \textless{}= rho M(n) + C M(n)\textasciicircum{}2}. That last line is the period map on the complement of finitely many directions, which the choice of datum removes.

The ledger then proves (\texttt{ns\_blowup\_of\_dss}):

\begin{itemize}
\item the perturbation decays geometrically, \texttt{M(n) \textless{}= (rho + C delta)\textasciicircum{}n delta}, for a start \texttt{M(0) \textless{}= delta} with \texttt{rho + C delta \textless{} 1} (\texttt{ns\_dss\_converge}): the solution converges to the profile;
\item infinitely many periods fit into the finite time 1 (\texttt{ns\_period\_hasSum}, via \texttt{nsPeriodLength});
\item the amplitude \texttt{e\textasciicircum{}\{nS/4\} (A - M(n))} tends to infinity (\texttt{ns\_dss\_amplitude}): the rate \texttt{(T* - t)\textasciicircum{}\{-1/4\}}, Leray's lower bound;
\item the profile is not steady (\texttt{LerayDSSCertificate.not\_steady}), so exact self-similarity is not used.
\end{itemize}
In ledger language: the free flow pays \texttt{S/4} per period at floor 0 (\texttt{ns\_free\_flow\_pays}), which is the subcriticality of \texttt{H\textasciicircum{}1} and the source of all damping in route I. A period factor \texttt{0 \textless{} rho \textless{} 1} is a cost \texttt{-log rho} above floor 0 for one step type (\texttt{ns\_dss\_costAboveFloor}); the zoomed equation is autonomous, so one period is every period.

The API card \texttt{routeI} assembles this: from \texttt{PeriodContraction} (the open step) it concludes settling onto the profile (the API's \texttt{settle\_everywhere}), with the period count from \texttt{settles\_by}, together with the three consequences above. \textbf{Settling in the zoomed frame is blowup in real time.} The open step is the content a certificate would supply (VII.5).

\subsubsection*{VII.3.3 Route II: inside the window, viscosity as a shrinking residual}
\addcontentsline{toc}{subsubsection}{VII.3.3 Route II: inside the window, viscosity as a shrinking residual}

\texttt{EulerProfileCertificate} is the Chen-Hou one-scale certificate (damping \texttt{q}, nonlinear constant \texttt{C}, profile residual \texttt{epsilon}, focusing margin) with the exponent in \texttt{2/5 \textless{}= alpha \textless{} 1/2} and the trapping condition raised by the viscosity: \texttt{(1 - q)\textasciicircum{}2 \textgreater{}= 4 C (epsilon + nu)}. From it the ledger proves trapped perturbation, a finite \texttt{T*} and unbounded amplitude (\texttt{ns\_blowup\_of\_euler\_profile}), stated also as a cost above a floor (\texttt{ns\_euler\_gap\_of\_costAboveFloor}). Below the Leray scaling the viscous term is a summable fresh uncertainty, so even with no damping its accumulated total stays bounded (\texttt{ns\_viscous\_transport\_bounded}); at \texttt{alpha = 1/2} it is a constant fresh uncertainty and accumulates without bound (\texttt{ns\_viscous\_transport\_unbounded\_at\_leray}). The API gives the same bound in its own form (\texttt{routeII\_bounded}): \texttt{M(n) \textless{}= M(0) + 2 (epsilon + nu)/(1 - q)}.

The consequence for the two routes: at the Leray scaling, viscosity cannot be a perturbation of an inviscid profile. It must sit inside the profile equation, which is route I. \textbf{No route II profile is known [open].}

\subsubsection*{VII.3.4 Why the energy ledger cannot decide}
\addcontentsline{toc}{subsubsection}{VII.3.4 Why the energy ledger cannot decide}

The energy identity is one global sum: total dissipation is at most the start energy. A summable budget never pays a uniform positive cost per step (\texttt{ns\_energy\_budget\_not\_costAboveFloor}), and dissipated shares whose partial sums stay below the start energy are such a budget (\texttt{ns\_energy\_identity\_not\_costAboveFloor}, \texttt{energy\_cannot\_fill}). Per dyadic scale the viscous cost grows like \texttt{nu 4\textasciicircum{}k} while the transfer bound that energy alone allows grows like \texttt{K 2\textasciicircum{}\{5k/2\}}; whatever \texttt{nu} and \texttt{K \textgreater{} 0}, the bound overtakes the cost at some scale (\texttt{dyadic\_supercritical}). The scalings are \textbf{[classical]} (Bernstein); the arithmetic is \textbf{[Lean]}. In two dimensions the same count is critical, which is why 2D Navier-Stokes is regular \textbf{[classical]}.

So the regularity card's slot cannot be filled from the energy identity. This is the ledger form of Tao's averaged equations (VII.1.2): they keep this identity and blow up.

\subsubsection*{VII.3.5 The blueprint: Chen-Hou, as a skeleton}
\addcontentsline{toc}{subsubsection}{VII.3.5 The blueprint: Chen-Hou, as a skeleton}

Chen and Hou proved stable, nearly self-similar finite-time blowup for 2D Boussinesq and for 3D axisymmetric Euler with smooth data and a solid wall, by a computer-assisted proof \textbf{[external]} (arXiv 2210.07191 and 2305.05660). It is not Navier-Stokes: there is no viscosity, there is a wall, and the flow is axisymmetric. The project restates its ledger in \texttt{Unified/Blowup.lean}:

\begin{itemize}
\item if the one-step certificate holds (\texttt{(1 - q)\textasciicircum{}2 \textgreater{}= 4 C r}) and the start is in the ball, the perturbation stays in the ball forever (\texttt{trap}, \texttt{trap\_radius}, \texttt{trap\_of\_gap}), and the threshold is sharp (\texttt{escape\_of\_below}, \texttt{escape\_pos});
\item the certificate follows from \texttt{CostAboveFloor} of \texttt{1 - q} over \texttt{2 sqrt(C r)} for one step type (\texttt{gap\_of\_costAboveFloor}: the strict cost condition implies the trapping gap; the converse and the boundary case are not proved);
\item under a focusing margin the real durations are summable, the total time \texttt{T*} is finite, and the amplitude is unbounded (\texttt{durations\_pay}, \texttt{arrival}, \texttt{amplitude\_unbounded}, assembled in \texttt{blowup\_of\_certificate});
\item every rescaled step adds one fixed quantum to the Beale-Kato-Majda count \textbf{[reading]}: the Lean content is only the identity \texttt{boundary\_quantum} (the sum of the inverse durations times the durations is n); the Beale-Kato-Majda count is not defined in Lean.
\end{itemize}
All of this except the last item is \textbf{[Lean]} for the skeleton. The partial differential equation enters only as the step inequality, a hypothesis. That BKM's integral is infinite exactly when smooth continuation fails is \textbf{[classical]}.

\subsection*{VII.4 What is borrowed}
\addcontentsline{toc}{subsection}{VII.4 What is borrowed}

\subsubsection*{VII.4.1 Alpoge and Buckmaster: the unforced claim \textbf{[external, unverified, certificate unreleased]}}
\addcontentsline{toc}{subsubsection}{VII.4.1 Alpoge and Buckmaster: the unforced claim \textbf{[external, unverified, certificate unreleased]}}

The claim, as the notes read it from the paper itself (79 pages, version of 20 November 2026):

\begin{itemize}
\item \textbf{Statement.} Smooth, compactly supported, divergence-free data on R\textasciicircum{}3, and smooth data on the torus, whose \textbf{unforced} solutions blow up in finite time. The paper says this disproves (A) and (B).
\item \textbf{Route.} A profile of the Navier-Stokes equations in self-similar variables, with the viscous term inside the profile equation, periodic in zoomed time up to an isometry: backward discretely self-similar, the case Bradshaw and Tsai left open. Their Theorem 1.1 says one such profile gives blowup at Leray's rate, stable in finite codimension. Their Theorem 7.1 says an approximate profile with explicit inequalities (a \textbf{certificate}) gives an exact profile nearby. This is exactly route I of VII.3.2.
\item \textbf{The candidate.} Found by Newton's method after continuation. Not axisymmetric: invariant under rotation by \texttt{2 pi/19} about one axis and by \texttt{pi} about another. Period about 0.36.
\item \textbf{What is not available.} The certificate is \textbf{not released}; the paper gives only its SHA-256 hash, \texttt{3f7a335b84b6c7631742cacc5343dfcd11cf72f5221ff79b2b9268790b4b22c5}, and says it will be released. The interval-arithmetic check (128 bits) is outside their formalisation and not in their repository. Their Lean library (about 2,270 files and 600,000 lines, Lean 4.33) was patched on 20 November and the patched package was not rebuilt, per its README. Its formal Theorem 1.1 has no rate, and the paper's prose no longer follows the Lean proof.
\item \textbf{Expert status.} The notes record Tao's blog post of 20 November: technically human-readable, broadly clear strategy, a professional write-up still needed. No referee report exists.
\end{itemize}
This project did \textbf{not} rebuild their library, did \textbf{not} see the certificate, and did \textbf{not} check any of their inequalities. It reads the claim, records it, and uses it only as an external input in a named slot.

How their profile avoids the no-go results (VII.1.1) \textbf{[external]}: it depends on zoomed time, so Necas-Ruzicka-Sverak and Tsai (which exclude steady profiles) do not apply; it is bounded in weak \texttt{L\textasciicircum{}3} but not \texttt{L\textasciicircum{}3}, consistent with Escauriaza-Seregin-Sverak; it is not axisymmetric; its period lies above the Chae-Wolf lower bound, with the constant not checked here.

\subsubsection*{VII.4.2 OpenAI: the forced claim \textbf{[external, unverified]}}
\addcontentsline{toc}{subsubsection}{VII.4.2 OpenAI: the forced claim \textbf{[external, unverified]}}

For every \texttt{nu \textgreater{} 0}, a smooth force with compact support, a compact set, and a smooth solution on \texttt{[0, 1)} with zero initial velocity and bounded energy whose sup-norm is unbounded as \texttt{t -\textgreater{} 1}; the paper says this establishes (C), and (D) by a corollary. An independent build record of its Lean proof reports the three standard axioms, no \texttt{sorry}, and a match to a public formalisation of (C) and (D) \textbf{[external, not rebuilt here]}. A preprint (arXiv 2610.08144) shows lemma-level places where the Lean proof is not a faithful translation of the written one, and expressly claims nothing about the written proof's correctness \textbf{[external]}. Because a force is allowed in (C) and (D), this claim, if it holds, answers (C) and (D) as stated but says nothing about (A) and (B), which have zero force.

\subsubsection*{VII.4.3 Smaller programmes \textbf{[external]}}
\addcontentsline{toc}{subsubsection}{VII.4.3 Smaller programmes \textbf{[external]}}

Petrillo and Glimm (arXiv 2609.23868) set a target for a computer search for unforced blowup; its implications are proved in Lean over Mathlib with no Navier-Stokes object. Cheskidov, Dai and Palasek (arXiv 2609.26790) prove cascade blowup in dyadic shell models, not in Navier-Stokes. The companion unforced Euler claim by Alpoge, Buckmaster and Coiculescu was not read.

\subsubsection*{VII.4.4 The logic point \textbf{[classical]}}
\addcontentsline{toc}{subsubsection}{VII.4.4 The logic point \textbf{[classical]}}

The owner's scope decision is that a blowup that cannot be ruled out is in itself enough for an existence to be constructed where it blows up. The exact statement behind it is Goedel's completeness theorem: if ZFC plus ``some smooth datum blows up'' is consistent, a model of it exists. Three senses of ``cannot be ruled out'' must be kept apart: (a) one method cannot rule it out (the energy method; VII.3.4); (b) the axioms cannot refute it (independence; \textbf{not known} for Navier-Stokes); (c) it is not ruled out for the real equations over the real numbers. A completeness model for (b) may be nonstandard, so it settles nothing about the actual equations. Only a proof settles the Clay question. The two 2026 claims are proofs of that kind: an explicit datum and an argument, not a completeness model. If correct, they are the owner's ``existence constructed where it blows up'', in the standard sense.

\subsection*{VII.5 The one open step}
\addcontentsline{toc}{subsection}{VII.5 The one open step}

\textbf{The open step of the card.} The route I card has one open step: \textbf{\texttt{PeriodContraction}. One period of the zoomed flow, from a non-steady periodic profile at the Leray scaling, contracts the perturbation, on the complement of finitely many directions removed by the datum, in a real norm, with \texttt{rho + C delta \textless{} 1}.} In the external tie, this is the statement ``the Alpoge-Buckmaster certificate satisfies their Appendix B''. It is exactly what their unreleased certificate claims to supply. Everything else in the card is proved (VII.3) or is a named external theorem (VII.4.1). \textbf{[open]}

\textbf{The other slot, stated exactly.} The regularity card names what regularity would need: \textbf{every dyadic scale dissipates a uniform cost above the transfer floor, from the true nonlinearity.} The energy identity cannot supply it (VII.3.4). If the 2026 claims hold, the slot cannot be filled for some data. \textbf{[open; false if the unforced claim holds]}

\textbf{What this does to the problem.} The API does not settle Navier-Stokes as regular. It settles blowup in the Leray window once one named slot is filled (the period contraction), which is a claimed but unreleased certificate today. The mechanical answer is therefore: \textbf{open, with the open step located and an external claim asserting it holds.} Route II has no candidate \textbf{[open]}. What continues past \texttt{T*} for a blowup solution is \textbf{[open]} (weak solutions are not unique in general; for Ricci flow the same question is a theorem, A.2).

\textbf{What would discharge the external inputs by import.} The notes list twelve audit items, none done. The ones that change the status of the tie:

\begin{enumerate}
\item Replace the named statements by real definitions of (A) and (B) (a pinned public formalisation), and check in one environment that their formal Theorem 1.1 has exactly the type ``not (A) and not (B)'', word by word against Fefferman.
\item Rebuild their library after the 20 November patch; run the axiom audit on the comparator statements; search for \texttt{sorry}, \texttt{native\_decide}, custom notation that could change a statement's meaning.
\item When the certificate is released: check its hash against the published one; check it has the form of their Appendix B; re-verify the Appendix B inequalities with an \textbf{independent} interval implementation. This is the step that today rests on one computation by the authors.
\item Identify the two derivation steps that their paper says are proved in the Lean and not in the prose.
\item Map their Appendix B row by row onto \texttt{LerayDSSCertificate}.
\end{enumerate}
Until these are done, everything in VII.4.1 stays \textbf{[external, unverified]}. The project's policy is to cite and audit, not vendor. An independent construction of the profile from scratch is beyond reach now: it is a genuinely three-dimensional, time-periodic shape with an operator-level certificate at 128-bit precision \textbf{[reading]}.

\subsection*{VII.7 DMP readings attached to this part (all [reading])}
\addcontentsline{toc}{subsection}{VII.7 DMP readings attached to this part (all [reading])}

\begin{itemize}
\item \textbf{A blowup is a certificate that infinitely many steps fit in finite time.} The finite check at one scale is what makes the packing geometric; \texttt{T*} is the boundary where the smooth description ends, and the profile is what crosses it. This is the hand-over of the form layer run in reverse.
\item \textbf{The same theorem as the mass gap.} \texttt{summable\_of\_gap} gives a finite correlation length in Yang-Mills and a finite \texttt{T*} here (Lean for both uses; the parallel is the reading).
\item \textbf{Scale invariance as the arrow of causality.} Because the zoomed equation is autonomous, the set of ``outsides'' has one element: the same step at every scale, so the for-all over steps costs nothing.
\item \textbf{``Not DMP-complete, hence blows up.''} In ledger terms the energy ledger cannot pay a uniform cost per scale (VII.3.4). Read as a prediction, this named the side the 2026 claims assert. It proves nothing: a method's failure to exclude blowup is not blowup. The claims, if correct, prove it by a certified profile. DMP-completeness now has a precise statement: DMP-complete means runs on DMP, the running-state hypothesis R of Part III 4.2 (the owner's definition, III.64b). The owner's reading: the whole process is described and Navier-Stokes does not build in it (the energy ledger cannot pay a uniform cost per scale, VII.3.4a and VII.3.4b, \textbf{[Lean]}), so it is not DMP-complete, and blowup is the expected consequence. The negative result is the expected one, and it makes the framework falsifiable: global regularity would refute it \textbf{[reading]} (VII.7e). Blowup of the real equations is still the card's open step (VII.5.1, \textbf{[open]}); in the settling machine a missing floor gives undecidability (II.3), so the step from ``not DMP-complete'' to blowup is the reading, not a theorem. In the one-box test (I.8) the wrong marker is proved for the model: with the energy-only transfer against the viscous cost as gate factor and the claim ``regular'' at every dyadic scale, the gate reads R136a1 at some scale, for every viscosity and margin \textbf{[Lean]} (model, VII.7g). It covers the energy-only ledger, not an actual solution with a better transfer bound \textbf{[reading]} (VII.7h).
\item \textbf{Settle means opposite things on the two cards.} In the API, ``settle'' means the carried quantity dies out. For regularity that quantity is the cascade, so settling means regular; for route I it is the perturbation from the blowup profile, so settling means blowup. The same theorem answers opposite physical questions.
\end{itemize}
\subsection*{VII.8 What would close it}
\addcontentsline{toc}{subsection}{VII.8 What would close it}

If the Alpoge-Buckmaster certificate is released and verifies independently, the route I card's open step is discharged, (A) and (B) fail, and the Clay problem is answered in the direction of blowup, by their work, with this part supplying the ledger skeleton and the audit. If the certificate fails an inequality, the open step stays open and the regularity slot stays as it is. Either way the change to this paper is one row of the ledger, one verdict line of this part and one item of the abstract. A proof of \textbf{regularity} would have to fill the regularity slot from the true nonlinearity, which the energy identity cannot do (VII.3.4), and which the unforced claim, if true, makes impossible for some data.


%% ---- 08-part-VIII-bsd.md
\section*{Part VIII. Birch and Swinnerton-Dyer: the card, three named parts, and the one open step}
\addcontentsline{toc}{section}{Part VIII. Birch and Swinnerton-Dyer: the card, three named parts, and the one open step}

\subsection*{VIII.0 What is proved, by whom, and what this part does}
\addcontentsline{toc}{subsection}{VIII.0 What is proved, by whom, and what this part does}

\textbf{The Clay problem} (Wiles' official description) \textbf{[classical]}. Let E be an elliptic curve over the rationals. Its group of rational points E(Q) is finitely generated (Mordell 1922), so it is a finite group of torsion points plus r independent points of infinite order. The number r is the \textbf{rank}. For each prime p of good reduction let a\_p = p + 1 - \#E(F\_p). The numbers a\_p build the L-function L(E, s) as an Euler product. The conjecture says: \textbf{the Taylor expansion of L(E, s) at s = 1 is c (s - 1)\textasciicircum{}r + higher terms, with c not zero and r the rank.} In particular L(E, 1) = 0 exactly when E(Q) is infinite. Clay's text also gives the refined formula for c: the real period, the regulator, the order of the Tate-Shafarevich group, the Tamagawa numbers and the torsion order enter. The statement needs L(E, s) near s = 1, so it needs the analytic continuation, which is the modularity theorem (Wiles, Taylor-Wiles, Breuil-Conrad-Diamond-Taylor) \textbf{[external]}.

\textbf{State of the problem (2026-11-24).} The rank statement is proved when the order of the zero is 0 or 1 (Gross-Zagier, Kolyvagin, with modularity) \textbf{[external]}. It is \textbf{open} from order 2 on. The notes record no claimed proof of the general statement. \textbf{This part proves no case of the conjecture that was not already known}, and its Lean contains no new theorem about elliptic curves that experts did not have.

\textbf{What this part does.} It states BSD as a property of one Weierstrass equation and proves in Lean what Mathlib gives about it (VIII.3). It puts BSD on the DMP API as a card with one named open step (VIII.2). It records what is borrowed and its exact status (VIII.4), states the one open step (VIII.5), and attaches the owner's readings, short (VIII.6). Three classical theorems that Mathlib lacks are carried as \textbf{named parts}, not axioms: weak Mordell-Weil, Hasse's bound, and one Euler product bound. They are theorems of others, not open problems. They are missing from the library, not from mathematics.

\subsection*{VIII.1 The problem as stated, and what is known}
\addcontentsline{toc}{subsection}{VIII.1 The problem as stated, and what is known}

\textbf{VIII.1.1 What is proved} (recalled from the literature; not re-read for this draft).

{\footnotesize
\begin{longtable}{>{\raggedright\arraybackslash}p{\dimexpr 0.3016\linewidth-2\tabcolsep\relax}>{\raggedright\arraybackslash}p{\dimexpr 0.5054\linewidth-2\tabcolsep\relax}>{\raggedright\arraybackslash}p{\dimexpr 0.1730\linewidth-2\tabcolsep\relax}}
\toprule
\textbf{Result} & \textbf{Statement} & \textbf{Mark} \\
\midrule
\endhead
Mordell 1922, Weil 1929 & E(Q), and E(K) for a number field K, is finitely generated. The proof is weak Mordell-Weil (E(Q)/2E(Q) is finite) plus descent with a height & classical \\
Hasse 1933-36, Weil 1948 & \texttt{abs(a\_p) \textless{}= 2 sqrt p}: the Riemann hypothesis of the curve over F\_p & classical \\
Deuring 1953-56 & for curves with complex multiplication, L(E, s) is a Hecke L-function: entire, with a functional equation & classical \\
Wiles 1995, Taylor-Wiles 1995, Breuil-Conrad-Diamond-Taylor 2001 & every E/Q is modular, so L(E, s) is entire and the completed function has a functional equation with sign w = +1 or -1 & external \\
Coates-Wiles 1977; Rubin 1987 & CM curves: L(E, 1) not zero gives E(Q) finite (Coates-Wiles) and Sha finite (Rubin) & external \\
Gross-Zagier 1986 & L'(E, 1) is a nonzero multiple of the height of a Heegner point; a simple zero gives a point of infinite order & external \\
Kolyvagin 1988-90, with modularity & analytic rank 0 or 1 gives rank = analytic rank and Sha finite & external \\
Buhler-Gross-Zagier 1985 & the curve 5077a has a zero of order exactly 3 and rank 3 & external, recalled \\
Kato 2004; Skinner-Urban 2014; later work & for many curves of analytic rank 0 or 1, the p-part of the refined formula for many p & external, recalled \\
Monsky 1996; Nekovar; Dokchitser-Dokchitser 2010 & p-parity for E/Q (a theorem); so finiteness of Sha would give \texttt{(-1)\textasciicircum{}rank = w} & external \\
\bottomrule
\end{longtable}}

\textbf{VIII.1.2 What is open} \textbf{[open]}.

\begin{itemize}
\item The rank part for analytic rank 2 or more, in general. There is no construction of two or more independent points from a higher zero (no higher Heegner points), and no bound ``rank \textless{}= order'' beyond order 1.
\item The finiteness of Sha in general (known for analytic rank at most 1).
\item The refined formula for all curves of analytic rank 0 and 1, at all primes at once, and any case of rank 2 or more beyond single curves.
\item The Riemann hypothesis for L(E, s) (the zeros on the line \texttt{re s = 1}), part of GRH beyond Dirichlet L-functions (Part V).
\item No curve of proved analytic rank 4 or more is known. A zero can be certified only up to order 3 (a value by modular symbols, a first derivative by Gross-Zagier, each with parity). Numerics cannot certify a zero \textbf{[classical, recalled]}.
\end{itemize}
\subsection*{VIII.2 The card}
\addcontentsline{toc}{subsection}{VIII.2 The card}

The DMP API (\texttt{docs/unified/API.md}) gives every problem one shape: a well-defined input, a proved structure on it, one named open step in a slot, and external inputs in their own slot. BSD has two layers. The Mathlib-only core is \texttt{Arsenal.BSD} in DMP.IM. The instance of the API is \texttt{Unified/BSDAPI.lean} in this repository (namespace \texttt{ProofPhase.API.BSD}). The table fills the template.

{\footnotesize
\begin{longtable}{>{\raggedright\arraybackslash}p{\dimexpr 0.1790\linewidth-2\tabcolsep\relax}>{\raggedright\arraybackslash}p{\dimexpr 0.4942\linewidth-2\tabcolsep\relax}>{\raggedright\arraybackslash}p{\dimexpr 0.3068\linewidth-2\tabcolsep\relax}}
\toprule
\textbf{API concept} & \textbf{BSD} & \textbf{Lean} \\
\midrule
\endhead
\textbf{Input} & an elliptic curve W over Q; a function \texttt{Lam} meant to be its continued L-function; a sign \texttt{w}. Over IM: any number field K & \texttt{Input} (\texttt{BSDAPI}); \texttt{Arsenal.BSD.Pts} \\
\textbf{Encoder} & at each prime p: state \texttt{sqrt p}, rotated image the Frobenius eigenvalue \texttt{alpha}. The shared root \texttt{norm alpha = sqrt p} is Hasse's bound. Reflection coefficient \texttt{a\_p / (2 sqrt p)}, Szego factor \texttt{1 - a\_p\textasciicircum{}2/(4p)} & \texttt{frob}, \texttt{frob\_refl}, \texttt{frob\_szego} \\
\textbf{Margin and chain} & the partial Euler product \texttt{P\_n = product of N\_p / p}. It keeps at least \texttt{(1 - p\textasciicircum{}\{-1/2\})\textasciicircum{}2} per prime and gains at most \texttt{(1 + p\textasciicircum{}\{-1/2\})\textasciicircum{}2} & \texttt{EulerData.lower}, \texttt{EulerData.upper} \\
\textbf{Open-step slot} & BSD: the L-series continues to \texttt{s = 1} and the order of its zero there is the rank. In ledger form: \texttt{P\_n \textasciitilde{} C (log p\_n)\textasciicircum{}rank} & \texttt{BSDStatement} (core); \texttt{problem}, \texttt{ledgerCard} \\
\textbf{External-input slot} & modularity (continuation, functional equation), weak Mordell-Weil, Hasse, Gross-Zagier-Kolyvagin, Goldfeld. In IM, weak Mordell-Weil, Hasse and the Euler product bound are named parts & \texttt{Input.cont}, \texttt{.fe}, \texttt{.weak}, \texttt{EulerData.hasse}, \texttt{LowRank}, \texttt{LedgerInput.goldfeld}; \texttt{R136a1} \\
\textbf{Computable layer} & point counts over any finite ring, \texttt{a\_p}, the Hasse check, each proved right & \texttt{Arsenal.BSD.count}, \texttt{frobTrace}, \texttt{hasseCheck}, \texttt{count\_correct}, \texttt{hasseCheck\_iff} \\
\bottomrule
\end{longtable}}

The cards, in the card shape:

{\footnotesize
\begin{longtable}{>{\raggedright\arraybackslash}p{\dimexpr 0.1677\linewidth-2\tabcolsep\relax}>{\raggedright\arraybackslash}p{\dimexpr 0.1828\linewidth-2\tabcolsep\relax}>{\raggedright\arraybackslash}p{\dimexpr 0.2288\linewidth-2\tabcolsep\relax}>{\raggedright\arraybackslash}p{\dimexpr 0.2233\linewidth-2\tabcolsep\relax}>{\raggedright\arraybackslash}p{\dimexpr 0.1774\linewidth-2\tabcolsep\relax}}
\toprule
\textbf{Card} & \textbf{Input} & \textbf{The one open step} & \textbf{Conclusion} & \textbf{External input} \\
\midrule
\endhead
\textbf{BSD} (\texttt{problem}) & a curve W over Q, a function \texttt{Lam}, a sign \texttt{w} & \texttt{BSDStatement W}: rank equals the order of the zero at s = 1 \textbf{[open]} & \texttt{E(Q)} is infinite exactly when \texttt{Lam 1 = 0}; and \texttt{w = (-1)\textasciicircum{}n} whenever the rank is the number n & modularity, weak Mordell-Weil \\
\textbf{Ledger form} (\texttt{ledgerCard}) & a curve, a continuation, the prime data (the counts N\_p) & the rank is r and the products \texttt{N\_p / p} grow like \texttt{C (log p)\textasciicircum{}r} \textbf{[open]} & \texttt{BSDStatement W} & modularity; Goldfeld's theorem \\
\textbf{Low-rank slot} (\texttt{LowRank}, \texttt{open\_of\_lowRank}) & the BSD input & none: it is a filled slot & the card's open step holds whenever the order is at most 1 & Gross-Zagier, Kolyvagin, with modularity \\
\bottomrule
\end{longtable}}

In IM the statement is made for \textbf{any number field} and uses the true rank in the extended naturals (no junk value when the group is not known to be finitely generated). The older form, over Q only, is kept in the API instance. \textbf{Nothing is an axiom.} A part answered R136a1 is a named hypothesis that the theorems using it carry (it is the statement itself); an \texttt{OpenStep} or \texttt{External} slot of the API is likewise an ordinary hypothesis. Neither appears in an axiom audit.

\subsection*{VIII.3 What is proved}
\addcontentsline{toc}{subsection}{VIII.3 What is proved}

All of this is \textbf{[Lean]}. Three cautions come first.

\begin{enumerate}
\item \textbf{Audit.} The standard-axioms-only claim is not re-made in this draft. The axiom log of the API instance (\texttt{bsd-axioms.log}) was written before the instance was adapted today to the native core, and it still lists the earlier signatures (for example \texttt{BSDStatement} with a function argument). The native modules were not rebuilt for this draft either. The names below were checked against the source files, by reading, one by one. \textbf{Before the claim ``standard axioms only'' is made for the adapted instance, the log must be regenerated.}
\item \textbf{Inputs.} A result marked ``from the named part'' is a Lean theorem with the part as a hypothesis, not a proof of the part.
\item \textbf{No new theorem.} The proved structure is bookkeeping around classical facts. It shows the shape of the problem.
\end{enumerate}
\subsubsection*{VIII.3.1 The algebraic side (\texttt{Arsenal.BSD.Algebraic})}
\addcontentsline{toc}{subsubsection}{VIII.3.1 The algebraic side (\texttt{Arsenal.BSD.Algebraic})}

The group of points \texttt{Pts W} is Mathlib's group of nonsingular points with the point at infinity.

\begin{itemize}
\item \textbf{Mordell-Weil from weak Mordell-Weil} (\texttt{mordellWeil\_of\_weak}): if E(K)/2E(K) is finite, E(K) is finitely generated. The proof is Mathlib's descent theorem with the naive height. Weak Mordell-Weil (\texttt{WeakMordellWeil}) is a named part. \textbf{[Lean, from the named part]}
\item \textbf{The torsion is finite}, with no assumption (\texttt{torsion\_finite}, from Mathlib's \texttt{finite\_torsion}). \textbf{[Lean]}
\item \textbf{The rank} \texttt{algRank} is the rank of E(K) as a Z-module, in the extended naturals. It is the true rank whether or not E(K) is known to be finitely generated. It equals \texttt{Module.finrank} once E(K) is finitely generated, and is then finite (\texttt{algRank\_eq\_finrank}, \texttt{algRank\_ne\_top}).
\item \textbf{Rank zero is finiteness} (\texttt{algRank\_eq\_zero\_iff\_torsion}, \texttt{algRank\_eq\_zero\_iff}): the rank is 0 exactly when every point has finite order, exactly when E(K) is finite. No Mordell-Weil is needed: finite torsion is enough. And E(K) is infinite exactly when the rank is positive (\texttt{infinite\_iff\_algRank\_pos}). \textbf{[Lean]}
\end{itemize}
\subsubsection*{VIII.3.2 The analytic side (\texttt{Arsenal.BSD.Analytic})}
\addcontentsline{toc}{subsubsection}{VIII.3.2 The analytic side (\texttt{Arsenal.BSD.Analytic})}

\begin{itemize}
\item \textbf{The continuation as a property} (\texttt{IsContinuation}): a function analytic on a half-plane \texttt{re s \textgreater{} a} with \texttt{a \textless{} 1}, equal to Mathlib's L-series on some right half-plane. This is what BSD reads: the germ at s = 1, reached from where the series converges. An entire function that equals the series for \texttt{re s \textgreater{} 3/2} is one (\texttt{IsContinuation.of\_entire}). A continuation is analytic at 1 (\texttt{IsContinuation.analyticAt\_one}). \textbf{[Lean]} That a continuation exists over Q is modularity \textbf{[external]}.
\item \textbf{The order of the zero is well defined}: two continuations agree near s = 1 (\texttt{continuation\_unique}, the identity theorem on a convex half-plane), so they have the same order there (\texttt{order\_welldefined}). A zero exists exactly when the order is positive (\texttt{zero\_iff\_order\_ne\_zero}). \textbf{[Lean]}
\item \textbf{The gamma factor} (\texttt{gammaFactor}; over Q it is \texttt{N\textasciicircum{}\{s/2\} (2 pi)\textasciicircum{}\{-s\} Gamma(s)}, \texttt{gammaFactor\_degree\_one}) is analytic and not zero at s = 1 (\texttt{gammaFactor\_analyticAt}, \texttt{gammaFactor\_one\_ne\_zero}). So the completed function has the same order as the L-function (\texttt{completed\_order}). \textbf{[Lean]}
\item \textbf{Parity} (\texttt{order\_parity}, \texttt{order\_parity\_of\_factor}, \texttt{analytic\_parity}): if the completed function satisfies \texttt{xi(2 - s) = w xi(s)} near s = 1 (\texttt{FunctionalEquation}; the strip form implies it, \texttt{FunctionalEquation.of\_strip}) and the order is a natural number n, then \texttt{w = (-1)\textasciicircum{}n}. The sign is not free: it is forced by the order. \textbf{[Lean, from the functional equation]}
\end{itemize}
In words: the mirror \texttt{s \textless{}-\textgreater{} 2 - s} fixes the parity of the order. It does not fix the order.

\subsubsection*{VIII.3.3 Point counts that run (\texttt{Arsenal.BSD.Count})}
\addcontentsline{toc}{subsubsection}{VIII.3.3 Point counts that run (\texttt{Arsenal.BSD.Count})}

\texttt{count W} counts the solutions of the Weierstrass equation over any finite ring, plus the point at infinity. \texttt{frobTrace W = q + 1 - count W}. The decidable \texttt{hasseCheck W} is proved equal to the real inequality \texttt{abs(a) \textless{}= 2 sqrt q} (\texttt{hasseCheck\_iff}). Over a finite field, \texttt{count} is Mathlib's \texttt{Nat.card} of the group of points (\texttt{count\_correct}), and at every prime of good reduction the count of the curve read mod p is Mathlib's (\texttt{reduceMod}, \texttt{count\_reduceMod}). Nothing is \texttt{noncomputable}: what runs is visibly computable. At p = 5 the curve y\textasciicircum{}2 = x\textasciicircum{}3 - x has trace -2 and passes the check, both by the kernel (\texttt{decide}). \textbf{[Lean]}

\subsubsection*{VIII.3.4 Hasse, named, and what it gives (\texttt{Arsenal.BSD.Hasse})}
\addcontentsline{toc}{subsubsection}{VIII.3.4 Hasse, named, and what it gives (\texttt{Arsenal.BSD.Hasse})}

Hasse's bound at a place of good reduction, in Mathlib's own terms, is the named part \texttt{Hasse W} (with \texttt{LocalHasse} at one place). Mathlib has neither isogeny degrees nor intersection theory on E x E, so the bound is not provable there yet. What follows from it is proved:

\begin{itemize}
\item \textbf{The rotation bound} (\texttt{invQuad}, \texttt{invQuad\_coeff}, \texttt{quad\_bound}): if \texttt{a\textasciicircum{}2 \textless{}= 4q}, the coefficients \texttt{c\_k} of \texttt{1/(1 - aT + qT\textasciicircum{}2)} obey \texttt{abs(c\_k) \textless{}= (k + 1) q\textasciicircum{}\{k/2\}}. The proof is the rotation itself: \texttt{c\_k} is a sum of \texttt{alpha\textasciicircum{}i beta\textasciicircum{}\{k-i\}} with both roots on the circle of radius \texttt{sqrt q}. At bad places \texttt{abs(c\_k) \textless{}= 1} (\texttt{lin\_bound}, \texttt{invQuad\_bound}, \texttt{invQuad\_bound\_lin}). \textbf{[Lean]}
\item \textbf{Every local factor of Mathlib's L-function is bounded} (\texttt{localPowerSeries\_coeff\_zero}, \texttt{localPowerSeries\_coeff\_bound}), and the first coefficient of \texttt{W.LFunction} is 1 (\texttt{eulerProduct\_apply\_one}, \texttt{LFunction\_one}). \textbf{[Lean, from \texttt{LocalHasse}]}
\item \textbf{A finite order}: an L-series with first coefficient 1 that converges absolutely somewhere is not zero for all large real arguments (\texttt{LSeries\_not\_eventually\_zero}), so every continuation has a finite order at s = 1 (\texttt{order\_ne\_top\_of\_LSeries}, \texttt{order\_ne\_top\_of\_summable}). \textbf{[Lean]}
\item \textbf{Through one more named part.} \texttt{EulerProductBound} says an Euler product whose local factors have these coefficient bounds converges absolutely for \texttt{re s \textgreater{} 3/2} (it is dominated by \texttt{zeta\_K(s - 1/2)\textasciicircum{}2}) \textbf{[classical, recalled]}. With it, Hasse gives the convergence (\texttt{LSeriesSummable\_of\_hasse}), a finite order (\texttt{order\_ne\_top}), and parity with \textbf{no order given separately} (\texttt{analytic\_parity\_of\_hasse}). \textbf{[Lean, from the named parts \texttt{Hasse} and \texttt{EulerProductBound}]}
\end{itemize}
\subsubsection*{VIII.3.5 The statement, and what follows from it (\texttt{Arsenal.BSD})}
\addcontentsline{toc}{subsubsection}{VIII.3.5 The statement, and what follows from it (\texttt{Arsenal.BSD})}

\texttt{BSDStatement W} says: a continuation exists, and the rank equals the order of its zero at s = 1. It is a statement about W alone (\texttt{bsdStatement\_iff} turns it into the equation \texttt{rank = order} when a continuation is in hand). It is the \textbf{rank part of the Clay statement} for K = Q, given the continuation \textbf{[paper]} (the identification of the Lean definition with Wiles' text is made in prose here, not in Lean). From it:

\begin{itemize}
\item \textbf{The existence line} (\texttt{existence\_line}): E(K) is infinite exactly when \texttt{Lam 1 = 0}. \textbf{[Lean, from BSD]}
\item \textbf{Parity of the rank} (\texttt{rank\_parity}, \texttt{rank\_parity\_of\_fg}): with the functional equation, a rank equal to the number n gives \texttt{w = (-1)\textasciicircum{}n}. \textbf{[Lean, from BSD]}
\item \textbf{A finite rank, and parity with no rank given separately} (\texttt{algRank\_ne\_top\_of\_hasse}, \texttt{rank\_parity\_of\_hasse}): under BSD, Hasse and the Euler product bound, the rank is finite (it is the order of a zero) and \texttt{w = (-1)\textasciicircum{}rank}. Over Q the Euler product bound is proved (IM, \texttt{Arsenal.BSD.EulerRat}, \texttt{eulerProductBound\_rat}), so there Hasse alone is named. \textbf{[Lean, from BSD and the named parts]}
\item \textbf{The refined conjecture} (\texttt{RefinedBSD}): near s = 1, \texttt{Lam(s) = (s - 1)\textasciicircum{}r g(s)} with \texttt{g(1) = Omega Reg \#Sha tam / (\#tors\textasciicircum{}2 sqrt abs(d\_K))}. The rank \texttt{r}, the torsion order and the field discriminant are derived. Mathlib defines none of the real period, the regulator, Sha or the Tamagawa product, so they stay free numbers in the statement. The refined statement gives the rank statement when its leading coefficient is not zero (\texttt{bsd\_of\_refined}). \textbf{[Lean]}
\end{itemize}
\subsubsection*{VIII.3.6 Hasse as an encoder, and the Euler ledger (\texttt{ProofPhase.API.BSD})}
\addcontentsline{toc}{subsubsection}{VIII.3.6 Hasse as an encoder, and the Euler ledger (\texttt{ProofPhase.API.BSD})}

\begin{itemize}
\item \textbf{The Frobenius eigenvalue} \texttt{alpha = a/2 + i sqrt(q - a\textasciicircum{}2/4)} has norm \texttt{sqrt q} exactly when \texttt{a\textasciicircum{}2 \textless{}= 4q} (\texttt{frobAlpha}, \texttt{frobAlpha\_normSq}, \texttt{norm\_frobAlpha}). The point count is \texttt{norm(1 - alpha)\textasciicircum{}2} (\texttt{count\_normSq}). \textbf{[Lean]}
\item \textbf{Hasse's bound is an encoder} (\texttt{frob}): the state \texttt{sqrt q} and the rotated image \texttt{alpha} share one root, and that shared root is Hasse's bound. The reflection coefficient is \texttt{a/(2 sqrt q)} (\texttt{frob\_refl}). The Szego factor is \texttt{1 - a\textasciicircum{}2/(4q)}, in the unit interval (\texttt{frob\_szego}). In angle terms: \texttt{cos theta\_p} and \texttt{sin\textasciicircum{}2 theta\_p}, the Sato-Tate density. \textbf{[Lean; Hasse as input]}
\item \textbf{The ledger}. At each step (a prime) the factor is \texttt{N\_p / p}. Hasse's bound puts it between \texttt{(1 - p\textasciicircum{}\{-1/2\})\textasciicircum{}2} and \texttt{(1 + p\textasciicircum{}\{-1/2\})\textasciicircum{}2} (\texttt{factor\_bounds}, \texttt{EulerData.fac\_bounds}). So the partial product keeps a margin (\texttt{EulerData.lower}) and has a chain bound (\texttt{EulerData.upper}, \texttt{P\_le}). \textbf{No finite stage of the Euler product is zero} (\texttt{no\_finite\_zero}). A zero at s = 1 is a property of the whole chain, a rate, never of one step. \textbf{[Lean; Hasse as input]}
\item \textbf{The ledger of one curve.} Built from one given integral curve, the ledger's counts are that curve's point counts by construction, and at every good prime they are the computed count, which Mathlib's count equals (\texttt{ProofPhase/Closures/EulerCurve.lean}: \texttt{ofCurve\_N}, \texttt{ofCurve\_N\_good}, \texttt{ofCurve\_fac\_good}). \textbf{[Lean, by construction; Hasse stays the named input]} That these counts give the linear coefficient of Mathlib's local factor at the completion is the classical definition of \texttt{a\_p}, not linked in Lean. \textbf{[classical, a formalisation item]}
\end{itemize}
\subsubsection*{VIII.3.7 The card and its conclusions (\texttt{ProofPhase.API.BSD})}
\addcontentsline{toc}{subsubsection}{VIII.3.7 The card and its conclusions (\texttt{ProofPhase.API.BSD})}

\begin{itemize}
\item \textbf{Without the open step} (\texttt{unconditional}): from the three \texttt{External} inputs of \texttt{Input}, E(Q) is finitely generated, its torsion is finite, every continuation has the same order at s = 1, and that order's parity is the sign. \textbf{[Lean, from modularity and weak Mordell-Weil]}
\item \textbf{The card} (\texttt{problem}): open step \texttt{BSDStatement W}; conclusion the existence line and the parity of the rank. \textbf{[Lean]}
\item \textbf{The low-rank theorem in its slot} (\texttt{LowRank}, \texttt{open\_of\_lowRank}): Gross-Zagier-Kolyvagin as an \texttt{External} fills the open step whenever the order is at most 1. \textbf{[Lean for the plumbing; external for the content]}
\item \textbf{The ledger card} (\texttt{LedgerInput}, \texttt{ledgerCard}): from the growth \texttt{C (log p)\textasciicircum{}r} (open step) and Goldfeld's theorem (\texttt{External}), BSD follows. Goldfeld's theorem is true only when the prime data are the curve's own (see the care in VIII.3.6). \textbf{[Lean for the plumbing; external for the content]}
\end{itemize}
\subsubsection*{VIII.3.8 The slot is not empty: y\textasciicircum{}2 = x\textasciicircum{}3 - x}
\addcontentsline{toc}{subsubsection}{VIII.3.8 The slot is not empty: y\textasciicircum{}2 = x\textasciicircum{}3 - x}

The curve y\textasciicircum{}2 = x\textasciicircum{}3 - x (conductor 32, complex multiplication by i, Fermat's curve for the congruent number 1).

\begin{itemize}
\item The curve over Q and over Z; its discriminant is 64 and it is elliptic (\texttt{E32}, \texttt{E32\_Δ}, \texttt{E32Z}, \texttt{E32Z\_Δ}); the three points of order two are on it (\texttt{E32\_points}). \textbf{[Lean]}
\item The computed count mod p is Mathlib's at every odd prime (\texttt{count\_E32}). At p = 5 the Hasse check is proved by the kernel (\texttt{hasse\_E32\_five}). \textbf{[Lean]} The trace of Frobenius at the odd primes up to 37 (zero at the primes that are 3 mod 4) and the Hasse check at the odd primes up to 59 run as \texttt{\#eval} and \texttt{\#guard} in the build: computed, not kernel proofs \textbf{[numerical]}. (The notes say ``below 100''; the source checks up to 59, and this part follows the source.)
\item The known values enter as \texttt{External} inputs (\texttt{Data32}): continuation and functional equation (Deuring), \texttt{L(E, 1)} not zero, and E(Q) finite (Fermat's descent, E(Q) = the four points O, (0,0), (1,0), (-1,0)). \textbf{[external]}
\item From them (\texttt{input32}, \texttt{bsd\_E32}, \texttt{card\_E32}): rank 0 equals order 0, BSD holds for this curve, the card's open step is filled, and the root number is +1. \textbf{[Lean, from the external values]}
\end{itemize}
Numerically (\texttt{bsd\_check.log}, about 30 seconds on a CPU) \textbf{[numerical]}:

\begin{itemize}
\item For y\textasciicircum{}2 = x\textasciicircum{}3 - x (32a2): Hasse holds at all 2262 primes below 20000; the largest \texttt{abs(a\_p)/(2 sqrt p)} is 0.99989. \texttt{L(E, 1) = 0.655514388573}. With the real period 5.2441151066, \texttt{L(E, 1)/Omega = 0.1250000000}. That equals the refined right-hand side \texttt{c\_2 Sha / T\textasciicircum{}2 = 2 * 1 / 16} with the recalled values \texttt{c\_2 = 2}, \texttt{Sha = 1}, \texttt{T = 4}. A number agreeing to ten digits does not prove \texttt{L(E, 1)} is not zero; the non-vanishing is the external input.
\item For y\textasciicircum{}2 + y = x\textasciicircum{}3 - x (37a1, root number -1, rank 1): Hasse holds at the same primes. \texttt{L'(E, 1) = 0.305999773834}. The point (0, 0) has infinite order (the multiple 5P = (1/4, -5/8) is not integral, Nagell-Lutz). \texttt{L'(E, 1)/Omega = 0.0511114083}, against the recalled regulator 0.0511114082.
\item The ledger \texttt{product of N\_p/p} over \texttt{(log x)\textasciicircum{}r} at x = 100, 1000, 5000, 20000 is 1.65, 2.24, 2.12, 2.25 for 32a2 (target 2.157) and 6.82, 6.79, 7.96, 7.19 for 37a1 (target 8.23). The products oscillate and converge slowly at the centre of the critical strip. The factor \texttt{sqrt 2} in the targets is the recalled constant of Goldfeld \textbf{[classical, recalled]}; the rank-0 numbers support it.
\end{itemize}
A computation can refute a statement. It never proves a for-all statement, and none is claimed.

\subsection*{VIII.4 What is borrowed}
\addcontentsline{toc}{subsection}{VIII.4 What is borrowed}

Every borrowed result is cited and marked. The project did not rebuild any of it.

{\footnotesize
\begin{longtable}{>{\raggedright\arraybackslash}p{\dimexpr 0.5302\linewidth-2\tabcolsep\relax}>{\raggedright\arraybackslash}p{\dimexpr 0.2737\linewidth-2\tabcolsep\relax}>{\raggedright\arraybackslash}p{\dimexpr 0.1761\linewidth-2\tabcolsep\relax}}
\toprule
\textbf{Borrowed} & \textbf{Where it enters} & \textbf{Status} \\
\midrule
\endhead
\textbf{Modularity} (Wiles, Taylor-Wiles, Breuil-Conrad-Diamond-Taylor): continuation and functional equation of L(E, s) over Q & \texttt{Input.cont}, \texttt{Input.fe} & external \\
\textbf{Weak Mordell-Weil} (Mordell 1922, Weil 1929): E(K)/2E(K) is finite & \texttt{WeakMordellWeil} (named part in IM; \texttt{Input.weak} in the API) & classical \\
\textbf{Hasse's bound} (1933-36): \texttt{a\_p\textasciicircum{}2 \textless{}= 4p} at good primes. Textbook proof: the degree is a positive definite quadratic form on End(E), so \texttt{deg(m - n phi) = m\textasciicircum{}2 - a m n + q n\textasciicircum{}2 \textgreater{}= 0} forces the discriminant bound (Silverman VI.1.1). The same positivity is the Hodge index on E x E (Part IX) & \texttt{Hasse} (named part); \texttt{EulerData.hasse} & classical \\
\textbf{The Euler product bound}: factors with the Hasse-type coefficient bounds give absolute convergence for \texttt{re s \textgreater{} 3/2}, dominated by \texttt{zeta\_K(s - 1/2)\textasciicircum{}2} & \texttt{EulerProductBound}: proved over Q (\texttt{eulerProductBound\_rat}); a named part beyond Q & Lean over Q; classical, recalled, beyond \\
\textbf{Gross-Zagier, Kolyvagin} (with modularity): analytic rank 0 or 1 gives BSD & \texttt{LowRank} & external \\
\textbf{Goldfeld 1982}: the 1965 growth form of the Euler products implies BSD and the Riemann hypothesis for L(E, s) & \texttt{LedgerInput.goldfeld} & external \\
For the example: Deuring (continuation, CM), Fermat's descent (E(Q) is the four points), \texttt{L(E, 1)} not zero, recalled constants of the refined formula & \texttt{Data32} & classical / external \\
\bottomrule
\end{longtable}}

The three named parts of IM (weak Mordell-Weil, Hasse, the Euler product bound) are classical theorems. A statement answered R136a1 is placed outside the system and carried as a hypothesis by exactly the theorems that use it. It is not the same as ``open''. The notes assess weak Mordell-Weil as within reach on Mathlib, Hasse as needing isogeny degrees or intersection theory on E x E, and modularity as a very large formalisation \textbf{[paper, an estimate]}.

\subsection*{VIII.5 The one open step}
\addcontentsline{toc}{subsection}{VIII.5 The one open step}

\textbf{The open step of the card.} The BSD card has one open step: \textbf{\texttt{BSDStatement W}. The L-series of the elliptic curve continues to s = 1, and the order of its zero there equals the rank of E(Q).} It is proved for order 0 and 1 by others (VIII.4), with modularity. \textbf{For order 2 and more it is open. [open]} Everything else in the card is proved (VIII.3) or is a named external theorem (VIII.4).

\textbf{What a proof would have to supply} (the two halves of an equality):

\begin{itemize}
\item a \textbf{construction} of r independent rational points from a zero of order r. No higher analogue of Heegner points is known. Candidates (derivatives of p-adic L-functions, Darmon's Stark-Heegner points) are conjectural; and
\item a \textbf{bound} ``rank \textless{}= order'' beyond order 1 (Euler systems of higher rank, or Iwasawa main conjectures with non-degenerate p-adic heights).
\end{itemize}
\textbf{The stronger slot.} The ledger card's open step, the 1965 growth form of the Euler products, is strictly stronger. By Goldfeld it also gives the Riemann hypothesis for L(E, s) \textbf{[external]}, which is part of GRH beyond Dirichlet L-functions. \textbf{[open]}

\textbf{The refined formula.} That the leading coefficient equals the full product of arithmetic invariants is also open in general. The rank statement follows from it (\texttt{bsd\_of\_refined}). \textbf{[open]}

\textbf{Care points.}

\begin{itemize}
\item \textbf{The direction of the classical proofs.} The proofs do not run ``points exist, so a zero by counting''. At order 1 the zero \textbf{produces} the point (the Heegner point has height \texttt{L'(E, 1)} up to a factor). A nonzero value \textbf{forbids} points (Kolyvagin's Euler system bounds the Selmer group). So ``the zero forces a point'' is proved when the zero is simple and open from order 2 on. "A point of infinite order forces \texttt{L(E, 1) = 0}" is proved for every E/Q \textbf{[external]}; it is the contrapositive of Kolyvagin with modularity.
\item \textbf{No finite stage sees the zero} (\texttt{no\_finite\_zero}). The order of the zero is decided by all primes together. A zero at s = 1 is not visible in any finite set of values \texttt{a\_p}.
\item \textbf{Two Riemann hypotheses, one proved.} Hasse's bound is the curve's local Riemann hypothesis, proved and an encoder. The Riemann hypothesis for L(E, s) is global and open. BSD itself does not need it; the ledger form of BSD would imply it.
\item \textbf{Parity is free; size is not.} The functional equation fixes the parity of the order (\texttt{analytic\_parity}). The order is the content and no symmetry fixes it.
\end{itemize}
\textbf{What this does to the problem.} The API does not settle BSD. It states BSD as a property of one equation, separates the proved structure from the one open equality, and shows the slot is filled where the order is 0 or 1 and for the example above. The mechanical answer is: \textbf{open, with the open step located, and the three classical parts that Mathlib lacks named.}

\subsection*{VIII.6 DMP readings attached to this part (all [reading])}
\addcontentsline{toc}{subsection}{VIII.6 DMP readings attached to this part (all [reading])}

Each reading is the owner's statement with its exact counterpart. None is mathematics.

\begin{itemize}
\item \textbf{BSD is IM: a full subset of only derivations.} Both sides of BSD are derived from one Weierstrass equation. The rank is a derivation of the rational points. The analytic side is a derivation of the point counts mod p. Unlike Yang-Mills (a measure), Navier-Stokes (initial data) or GU (hypercharges), the card takes \textbf{no content input}; its external slots are theorems about the derivations. So BSD is stated entirely in the vocabulary of IM (\texttt{Arsenal.BSD}). \textbf{Care:} this is a statement about form, not about settling. The rank equality stays open.
\item \textbf{The zeros follow existence.} The order of the zero at s = 1 counts the independent rational points. Its exact counterpart has three levels: the existence line (\texttt{existence\_line}), the rank equation, and the 1965 ledger form. The proved arrow runs by construction (Heegner) and exclusion (Kolyvagin), not by counting (VIII.5).
\item \textbf{Hasse is rotation at a DB.} In the owner's words, what stays open in the rederived core is ``rotation at DB shaped''. At the prime p (the DB) Frobenius acts as a rotation: its eigenvalue lies on the circle of radius \texttt{sqrt p}, and \texttt{cos theta\_p = a\_p/(2 sqrt p)} is the encoder's reflection coefficient (\texttt{frob}, \texttt{frob\_refl}) \textbf{[Lean; Hasse classical]}. Its proof is one positivity: the degree form, or equivalently the Hodge index on E x E. The linear algebra of that is Part IX's \texttt{weil\_bound} at genus 1, where the conclusion is \texttt{abs(N - (q + 1)) \textless{}= 2 sqrt q}, which is Hasse's bound \textbf{[Lean for the linear algebra, from the intersection numbers and the Hodge index property, both external; the specialisation to one curve and the link of N to Mathlib's count are not made in Lean]}. So the open piece in Lean is one rotation per DB, confined by one positivity. It is open as a formalisation, not as mathematics.
\item \textbf{What is left of BSD is the record, the \texttt{.cmb}.} In DMP the CMB is the observability layer: counts, never content. What the native core leaves open is exactly such a record: the refined formula's counts (the order of Sha, the Tamagawa product, the torsion) and measures (the period, the regulator). The refined conjecture says the leading coefficient over \texttt{Omega Reg Sha tam / (tors\textasciicircum{}2 sqrt abs(d\_K))} is exactly 1: a dimensionless number whose logarithm is 0 \textbf{[classical statement, open]}. The functional equation carries both sides through the rotation \texttt{s \textless{}-\textgreater{} 2 - s} with sign \texttt{(-1)\textasciicircum{}rank} (\texttt{analytic\_parity}, \texttt{rank\_parity}) \textbf{[Lean]}. Dividing by \texttt{(s - 1)\textasciicircum{}r} removes the zero at the problem spot and leaves a finite nonzero value at the one box s = 1 (\texttt{bsd\_of\_refined}) \textbf{[Lean for the shape]}. The record balances numerically for y\textasciicircum{}2 = x\textasciicircum{}3 - x: \texttt{L(E, 1)/Omega = 0.125} against \texttt{2 * 1 / 16} \textbf{[numerical]}. \textbf{Care:} ``the ratio is 1'' is the refined conjecture itself, open in general. In Lean the five quantities are free numbers.
\item \textbf{Where BSD sits in the triad.} Yang-Mills is existence, GRH is non-existence, NP vs coNP is their interaction (Part I). BSD is an equation between an existence count (independent points) and a zero count (the vanishing rate of the Euler ledger).
\end{itemize}
\subsection*{VIII.7 What would close it}
\addcontentsline{toc}{subsection}{VIII.7 What would close it}

If an independent proof supplied the construction and the bound for orders 2 and more, the card's open step would be discharged by that work, and this part would supply the ledger skeleton and the audit. The change to this paper would be one row of the ledger, one verdict line of this part and one item of the abstract. As the notes state it, nothing in this round proves any part of that.

The formalisation items that would change the status of the \textbf{borrowed} parts, none done:

\begin{enumerate}
\item Regenerate the axiom log of the adapted API instance, and run IM's own audit on the native modules.
\item Prove weak Mordell-Weil on Mathlib (Kummer theory, class groups and unit groups of the 2-division field), so \texttt{WeakMordellWeil} stops being a named part.
\item Prove Hasse's bound (isogeny degrees as a quadratic form, or intersection theory on E x E), so \texttt{Hasse} stops being a named part.
\item Prove the Euler product bound beyond Q (done over Q, \texttt{eulerProductBound\_rat}; for a number field of degree above 1 it needs the residue fields of the completions as \texttt{O\_K / v} and the convergence of \texttt{zeta\_K} for \texttt{re s \textgreater{} 1}).
\item State and prove the link from Mathlib's local factor at the completion to the count over \texttt{Z/p}, at every prime, so \texttt{EulerData} is tied to one curve in Lean.
\item Modularity, a very large formalisation, stays an external input.
\end{enumerate}
\subsection*{VIII.8 The box in the one-box test}
\addcontentsline{toc}{subsection}{VIII.8 The box in the one-box test}

\textbf{Gate factor and marker.} The ticks are the \textbf{primitives} (the owner's name): the primes, the primitive elements of the multiplicative integers \textbf{[classical]}, at which the curve's Euler ledger is read (VIII.3.6); the gate reads inside when the ledger is consistent with the rank, and the marker is the claim ``inside''. The box of a curve passes exactly when BSD holds for it \textbf{[Lean]} (\texttt{bsd\_box\_iff}, \texttt{ETM/CardBox.lean}; claim VIII.8a). The box is the card's open step in the test's language (II.4): it restates the step and proves no new bound, and the content stays in the step \textbf{[paper]}. Its open ticks start where the order of the zero reaches 2 (VIII.5.1).


%% ---- 09-part-IX-hodge.md
\section*{Part IX. Hodge: the card, the positivity, and the one open step}
\addcontentsline{toc}{section}{Part IX. Hodge: the card, the positivity, and the one open step}

\subsection*{IX.0 What is proved, by whom, and what this part does}
\addcontentsline{toc}{subsection}{IX.0 What is proved, by whom, and what this part does}

\textbf{The Clay problem} (Deligne's description) \textbf{[classical]}. Let X be a projective non-singular algebraic variety over the complex numbers. A \textbf{Hodge class} is a class in the rational cohomology \texttt{H\textasciicircum{}\{2p\}(X, Q)} whose image in \texttt{H\textasciicircum{}\{2p\}(X, C)} lies in the \texttt{(p, p)} part of the Hodge decomposition. The class of an algebraic subvariety of codimension p is a Hodge class. \textbf{The Hodge conjecture says the converse: every Hodge class is a rational linear combination of classes of algebraic subvarieties.}

\textbf{State of the problem (2026-11-24).} The conjecture is a theorem for p = 1 (Lefschetz 1924), and by hard Lefschetz for curves (degree 2n - 2). So it holds for every variety of dimension at most 3 \textbf{[classical]}. It is \textbf{open} for codimension p \textgreater{}= 2 in general; the first open case is the fourth cohomology of fourfolds \textbf{[open]}. Beyond that, only special families are known, mostly abelian varieties (IX.1). The notes record no claimed proof of the general case.

\textbf{This part proves no case of the Hodge conjecture that was not already known.} The Lean proves the linear algebra of the positivity and the diagram chases of Lefschetz (1,1). The geometry (Hodge theory, the exponential sequence, hard Lefschetz, intersection theory on a surface) enters as named \texttt{External} slots. The conjecture itself is the card's one \texttt{OpenStep}.

\textbf{What this part does.} It puts Hodge on the DMP API as a card (IX.2), proves what Mathlib can reach (IX.3), records what is borrowed (IX.4), states the one open step (IX.5), and attaches the owner's readings, short (IX.6). Hodge is also the place where the \textbf{same positivity} that the Riemann hypothesis over finite fields needs appears as a theorem over the complex numbers. IX.3.4 and IX.3.5 show the shape for curves in Lean.

\subsection*{IX.1 The problem as stated, and what is known}
\addcontentsline{toc}{subsection}{IX.1 The problem as stated, and what is known}

\textbf{IX.1.1 What is proved} (recalled from the literature; not re-read for this draft).

{\footnotesize
\begin{longtable}{>{\raggedright\arraybackslash}p{\dimexpr 0.2157\linewidth-2\tabcolsep\relax}>{\raggedright\arraybackslash}p{\dimexpr 0.5846\linewidth-2\tabcolsep\relax}>{\raggedright\arraybackslash}p{\dimexpr 0.1797\linewidth-2\tabcolsep\relax}}
\toprule
\textbf{Case} & \textbf{Result} & \textbf{Mark} \\
\midrule
\endhead
p = 1 (divisors) & \textbf{Lefschetz (1,1)} (Lefschetz 1924; sheaf proof Kodaira-Spencer 1953): an integral class in \texttt{H\textasciicircum{}2(X, Z)} is of type \texttt{(1,1)} exactly when it is the first Chern class of a holomorphic line bundle; on a projective X every line bundle is \texttt{O(D)} for a divisor D. So the conjecture holds for p = 1, even with integer coefficients & classical \\
p = n - 1 (curves) & hard Lefschetz \texttt{L\textasciicircum{}\{n-2\} : H\textasciicircum{}2 -\textgreater{} H\textasciicircum{}\{2n-2\}} is an isomorphism of Hodge structures and sends divisor classes to curve classes; so the rational conjecture holds in degree 2n - 2 & classical \\
dimension at most 3 & follows from the two rows above: the only degrees with \texttt{(p, p)} classes are p = 0, 1, n - 1, n & classical \\
abelian varieties whose Hodge ring is generated by divisor classes & the conjecture holds; this covers a very general abelian variety (Mattuck) and simple abelian varieties of prime dimension (Tankeev). Mumford's example and Weil's classes show that not every abelian variety is of this kind & classical / external, recalled \\
abelian fourfolds of Weil type & Weil classes are algebraic: discriminant 1 by Markman (2022; simpler proofs by Floccari-Fu and by van Geemen-Rapagnetta); every abelian fourfold of Weil type by Markman, arXiv 2509.23403 (2025). The abstract reports that the conjecture for abelian varieties of dimension at most 5 then follows (with Moonen-Zarhin) & external, abstracts only, not checked \\
other families & cubic fourfolds (Zucker), uniruled fourfolds (Conte-Murre), some Fermat hypersurfaces (Shioda) & external, recalled \\
\bottomrule
\end{longtable}}

\textbf{IX.1.2 What fails} \textbf{[external, recalled]}. Two natural strengthenings are false, so ``described'' must be read exactly.

\begin{itemize}
\item \textbf{Integral coefficients.} ``Every integral class whose rational image is Hodge is algebraic'' is false. Atiyah and Hirzebruch (1962) found torsion classes in \texttt{H\textasciicircum{}4(X, Z)} that are not algebraic (Totaro 1997 gave the cobordism reading). Kollar (1992) found non-torsion failures: integral classes on hypersurfaces that are not algebraic though a multiple is.
\item \textbf{Non-projective Kaehler manifolds.} On a general complex torus there are Hodge classes but too few analytic subvarieties (Zucker 1977). Voisin (2002) showed that even with Chern classes of coherent sheaves the statement fails on some complex tori. Projectivity is part of the hypothesis.
\end{itemize}
\textbf{IX.1.3 The standard conjectures and the positivity} \textbf{[classical; open where marked]}. Grothendieck's standard conjectures (1969) are four statements for a Weil cohomology of a smooth projective variety.

\begin{itemize}
\item \textbf{B (Lefschetz type):} the inverse of hard Lefschetz is given by an algebraic correspondence. Known for abelian varieties (Lieberman); implied by the Hodge conjecture in characteristic 0. \textbf{[open in general]}
\item \textbf{C (Kunneth type):} the Kunneth projectors are algebraic. Over finite fields it follows from the Weil conjectures (Katz-Messing 1974).
\item \textbf{D:} numerical and homological equivalence agree. \textbf{[open in general]}
\item \textbf{Hodge type:} on primitive algebraic classes of codimension \texttt{p \textless{}= n/2} the form \texttt{(-1)\textasciicircum{}p \textless{}x, L\textasciicircum{}\{n-2p\} x\textgreater{}} is positive definite.
\end{itemize}
The Hodge-type conjecture is a \textbf{positivity}. \textbf{Over the complex numbers it is a theorem}: the Hodge-Riemann bilinear relations make the form definite on primitive Hodge classes, and algebraic classes are Hodge classes. \textbf{In characteristic p it is open}, except for surfaces, where it is the Hodge index theorem (Hodge 1937; Segre; Grothendieck 1958 in every characteristic), and a few further cases (for example Ancona 2021 for abelian fourfolds) \textbf{[external, recalled]}.

\textbf{The link to the Weil conjectures.} Weil proved the Riemann hypothesis for curves over a finite field in 1948 through the Castelnuovo-Severi inequality, which is the Hodge index theorem on the surface \texttt{C x C} (Mattuck-Tate 1958; Grothendieck 1958): \texttt{abs(N - q - 1) \textless{}= 2 g sqrt q}. Grothendieck's plan was that the standard conjectures give it for every smooth projective variety by the same positivity. \textbf{Deligne proved the Riemann hypothesis over finite fields in 1974 by another route} (Lefschetz pencils, monodromy, Rankin's squaring trick), without the standard conjectures, which remain open \textbf{[classical]}.

\textbf{IX.1.4 What is open} \textbf{[open]}.

\begin{itemize}
\item The Hodge conjecture for p \textgreater{}= 2 in general, including general abelian varieties beyond the cases above.
\item The standard conjecture B in general, D, and the Hodge type in positive characteristic.
\item Any positivity route to the Riemann hypothesis over finite fields beyond surfaces (Deligne's theorem is proved; the route is not), and its analogue over Z.
\end{itemize}
\subsection*{IX.2 The card}
\addcontentsline{toc}{subsection}{IX.2 The card}

The DMP API (\texttt{docs/unified/API.md}) gives every problem a well-defined input, a proved structure on it, and one named open step. Hodge has the cards of \texttt{Unified/HodgeAPI.lean} (namespace \texttt{ProofPhase.API.Hodge}) on top of the Mathlib-only module \texttt{Unified/HodgeLinear.lean} (namespace \texttt{ProofPhase.Hodge}). The table fills the template.

{\footnotesize
\begin{longtable}{>{\raggedright\arraybackslash}p{\dimexpr 0.1844\linewidth-2\tabcolsep\relax}>{\raggedright\arraybackslash}p{\dimexpr 0.5054\linewidth-2\tabcolsep\relax}>{\raggedright\arraybackslash}p{\dimexpr 0.2903\linewidth-2\tabcolsep\relax}}
\toprule
\textbf{API concept} & \textbf{Hodge} & \textbf{Lean} \\
\midrule
\endhead
\textbf{Input} & \texttt{V = H\textasciicircum{}\{2p\}(X, Q)}; \texttt{hodge}, the rational \texttt{(p, p)} classes (the \textbf{described} side); \texttt{alg}, the rational span of the classes of algebraic cycles of codimension p (the side that \textbf{exists}) & \texttt{Input} \\
\textbf{Encoder} & the polarisation \texttt{Q}: positive definite on Hodge classes, so Cauchy-Schwarz \texttt{Q(x,y)\textasciicircum{}2 \textless{}= Q(x,x) Q(y,y)} bounds the overlap of two classes, the same shape as the encoder bound & \texttt{PolHS.hodge\_cs}, \texttt{PosInput.alg\_cs} \\
\textbf{Running system} & over a finite field: Frobenius acting as a similitude of a positive form & \texttt{weilCard} \\
\textbf{Open-step slot} & the Hodge conjecture: \texttt{hodge \textless{}= alg} & \texttt{HodgeConjecture} \\
\textbf{External-input slot} & the cycle class is a Hodge class; the Hodge-Riemann relations; the exponential sequence; the Hodge decomposition; hard Lefschetz; Atiyah-Hirzebruch's torsion class; the Hodge index theorem on a surface & \texttt{Input.alg\_le}, \texttt{PosInput.hr}, \texttt{DivisorInput.exact}, \texttt{.hodge\_int}, \texttt{.div\_hodge}, \texttt{CurveClassInput.hl\_hodge}, \texttt{.hl\_alg}, \texttt{integral\_fails}, \texttt{curve\_open} \\
\bottomrule
\end{longtable}}

The cards, in the card shape:

{\footnotesize
\begin{longtable}{>{\raggedright\arraybackslash}p{\dimexpr 0.1706\linewidth-2\tabcolsep\relax}>{\raggedright\arraybackslash}p{\dimexpr 0.2215\linewidth-2\tabcolsep\relax}>{\raggedright\arraybackslash}p{\dimexpr 0.2358\linewidth-2\tabcolsep\relax}>{\raggedright\arraybackslash}p{\dimexpr 0.1772\linewidth-2\tabcolsep\relax}>{\raggedright\arraybackslash}p{\dimexpr 0.1749\linewidth-2\tabcolsep\relax}}
\toprule
\textbf{Card} & \textbf{Input} & \textbf{The one open step} & \textbf{Conclusion} & \textbf{External input} \\
\midrule
\endhead
\textbf{Hodge} (\texttt{problem}) & rational cohomology with its Hodge classes and algebraic classes & \texttt{HodgeConjecture}: every rational \texttt{(p, p)} class is a rational combination of algebraic cycle classes \textbf{[open]} & \texttt{hodge = alg}: described = exists & the cycle class is a Hodge class \\
\textbf{Divisors} (\texttt{DivisorInput}, \texttt{holds}, \texttt{card}) & line bundles, \texttt{H\textasciicircum{}2(X, Z)}, \texttt{H\textasciicircum{}\{0,2\}}, \texttt{H\textasciicircum{}2(X, Q)}, the maps & none: the open step is \textbf{filled} for p = 1 & \texttt{hodge} = the rational span of divisor classes & exponential sequence; Hodge decomposition; divisors are \texttt{(1,1)} \\
\textbf{Curves} (\texttt{CurveClassInput}, \texttt{holds}) & the divisor input, \texttt{H\textasciicircum{}\{2n-2\}}, hard Lefschetz & none: filled by transfer from divisors & the conjecture in degree \texttt{2n - 2} & hard Lefschetz \\
\textbf{Weil, positivity} (\texttt{weilCard}) & Frobenius on weight-w cohomology over \texttt{F\_q} & \texttt{StandardHodge}: a positive form on which Frobenius is a similitude of ratio \texttt{q\textasciicircum{}w} \textbf{[open]} & every eigenvalue has \texttt{abs(mu)\textasciicircum{}2 = q\textasciicircum{}w} & none \\
\textbf{Curves over \texttt{F\_q}} (\texttt{curveCard}) & the intersection form on \texttt{C x C} with fibre classes, graph of Frobenius and diagonal & the Hodge index property at \texttt{f1 + f2} & \texttt{abs(N - (q + 1)) \textless{}= 2 g sqrt q} & intersection numbers; the Hodge index theorem (\texttt{curve\_open}) \\
\bottomrule
\end{longtable}}

The last two cards belong to the GRH column (Part V); they sit here because they use the same positivity. Hodge's own open step is the first row. Nothing is an axiom. An \texttt{OpenStep} or \texttt{External} slot is an ordinary hypothesis and does not appear in the axiom audit. \texttt{Unified/HodgeAPI.lean} is not yet part of the aggregate import \texttt{Unified/All.lean}; that is left to the coordinator.

\subsection*{IX.3 What is proved}
\addcontentsline{toc}{subsection}{IX.3 What is proved}

All of this is \textbf{[Lean]}. The audit log (\texttt{hodge-axioms.log}, written after both modules were last edited): both modules build with no warnings; all 47 declarations rest on the three standard axioms (two, \texttt{sq\_add\_smul} and \texttt{transfer}, on two of them); the source has no \texttt{sorry}, \texttt{admit}, \texttt{native\_decide} or \texttt{axiom}. The log was read, not regenerated. None of it proves a statement about a variety. Each result says what the ledger gives from a hypothesis.

\subsubsection*{IX.3.1 A polarised Hodge structure, and positivity on its Hodge classes (\texttt{HodgeLinear})}
\addcontentsline{toc}{subsubsection}{IX.3.1 A polarised Hodge structure, and positivity on its Hodge classes (\texttt{HodgeLinear})}

Cauchy-Schwarz holds for a symmetric form that is nonnegative on a subspace, over any ordered field (\texttt{sq\_add\_smul}, \texttt{cs\_of\_nonneg}).

\texttt{PolHS} is a real polarised Hodge structure of weight 0 in Deligne's circle form. It is a real representation \texttt{h} of the circle, whose characters \texttt{z -\textgreater{} z\textasciicircum{}\{p-q\}} are the Hodge types; a symmetric form \texttt{Q} invariant under the circle; and the Hodge-Riemann positivity \texttt{Q(C x, x) \textgreater{} 0} for \texttt{x} not zero, where \texttt{C = h(i)} is the Weil operator. A Tate twist brings weight 2p to weight 0. The \textbf{Hodge classes} are the vectors fixed by the whole circle (\texttt{PolHS.hodge}, \texttt{mem\_hodge}): charge 0. Proved:

\begin{itemize}
\item \texttt{Q} is positive definite on the Hodge classes (\texttt{PolHS.hodge\_pos}, \texttt{hodge\_nonneg}): on a fixed vector the Weil operator is the identity;
\item it satisfies Cauchy-Schwarz there (\texttt{hodge\_cs}), the encoder bound;
\item the Hodge classes split off orthogonally: the space is their direct sum with their \texttt{Q}-orthogonal (\texttt{hodge\_isCompl}, finite dimension).
\end{itemize}
That the cohomology of a smooth projective variety, with \texttt{Q} built from the cup product, an ample class and the Lefschetz signs, \textbf{is} such a structure is the Hodge-Riemann theorem of Hodge theory \textbf{[classical]}, not in Mathlib and not proved here.

\subsubsection*{IX.3.2 The card, and the positivity on algebraic classes (\texttt{HodgeAPI})}
\addcontentsline{toc}{subsubsection}{IX.3.2 The card, and the positivity on algebraic classes (\texttt{HodgeAPI})}

\begin{itemize}
\item \textbf{The card} (\texttt{Input}, \texttt{HodgeConjecture}, \texttt{problem}): from \texttt{alg \textless{}= hodge} (an \texttt{External}) and the open step \texttt{hodge \textless{}= alg}, the conclusion \texttt{hodge = alg} follows by antisymmetry. The card adds no mathematics. It shows the shape: the content is all in the slot.
\item \textbf{Positivity} (\texttt{PosInput}): a symmetric rational form \texttt{Q} and the Hodge-Riemann relations on rational Hodge classes, as an \texttt{External}. Then: \texttt{Q} is positive definite on algebraic classes (\texttt{PosInput.alg\_pos}, \texttt{alg\_nonneg}), which is \textbf{Grothendieck's Hodge standard conjecture in characteristic 0}; an algebraic class orthogonal to all algebraic classes is zero (\texttt{num\_eq\_hom}), which is conjecture D in this piece; and \texttt{Q(x,y)\textasciicircum{}2 \textless{}= Q(x,x) Q(y,y)} on algebraic classes (\texttt{alg\_cs}), the encoder bound. \textbf{[Lean, from Hodge-Riemann]}
\item \textbf{Where the external input comes from} (\texttt{hr\_of\_polHS}): a real polarised Hodge structure whose rational points carry \texttt{Q}, and whose Hodge classes contain the rational Hodge classes, gives exactly the \texttt{hr} slot. \textbf{[Lean]} The one thing outside is the theorem of Hodge theory that the cohomology of X is such a structure.
\end{itemize}
\subsubsection*{IX.3.3 The Hodge index theorem as linear algebra (\texttt{HodgeLinear})}
\addcontentsline{toc}{subsubsection}{IX.3.3 The Hodge index theorem as linear algebra (\texttt{HodgeLinear})}

\texttt{HodgeIndex B H} says every class orthogonal to \texttt{H} has nonpositive square. For the intersection form on the Neron-Severi group of a surface and \texttt{H} ample this is the Hodge index theorem \textbf{[classical]}. Proved:

\begin{itemize}
\item \textbf{Signature gives the index property} (\texttt{hodgeIndex\_of\_sigPos}): if the form has no positive plane (signature \texttt{(1, *)}) and \texttt{H} is positive, a class orthogonal to \texttt{H} cannot be positive. The signature is Mathlib's, in Sylvester's law of inertia form.
\item \textbf{The reversed Cauchy-Schwarz inequality} (\texttt{reverse\_cs}): under the index property, \texttt{B(D,D) B(H,H) \textless{}= B(D,H)\textasciicircum{}2}. This is the Lorentzian twin of the encoder bound. \textbf{[Lean]}
\end{itemize}
\subsubsection*{IX.3.4 Weil's bound for curves from the Hodge index property (\texttt{HodgeLinear}, \texttt{HodgeAPI})}
\addcontentsline{toc}{subsubsection}{IX.3.4 Weil's bound for curves from the Hodge index property (\texttt{HodgeLinear}, \texttt{HodgeAPI})}

On \texttt{C x C} let \texttt{f1}, \texttt{f2} be the fibre classes, \texttt{Gamma} the graph of Frobenius and \texttt{Delta} the diagonal. The classical intersection numbers are an \texttt{External} slot (\texttt{CurveInput.geometry}). Then, all \textbf{[Lean]}:

\begin{itemize}
\item \textbf{Castelnuovo-Severi} (\texttt{castelnuovo\_severi}): the index property at \texttt{f1 + f2} gives \texttt{D\textasciicircum{}2 \textless{}= 2 (D f1)(D f2)} for every class \texttt{D}.
\item \textbf{The quadratic form} (\texttt{weil\_form}): with the intersection numbers, \texttt{g q a\textasciicircum{}2 - a b (N - q - 1) + g b\textasciicircum{}2 \textgreater{}= 0} for all real \texttt{a}, \texttt{b}.
\item \textbf{Weil's bound} (\texttt{weil\_bound}, \texttt{weil\_bound\_abs}): choosing \texttt{a = N - q - 1}, \texttt{b = 2 g q} gives \texttt{(N - q - 1)\textasciicircum{}2 \textless{}= 4 g\textasciicircum{}2 q}, so \texttt{abs(N - (q + 1)) \textless{}= 2 g sqrt q}.
\item \textbf{The card} (\texttt{curveCard}): the open step is the index property at \texttt{f1 + f2}; the conclusion is Weil's bound.
\item \textbf{The slot is filled for curves} (\texttt{curve\_open}): with the Hodge index theorem in its signature form (one positive direction; Hodge 1937, Grothendieck 1958) as an \texttt{External}, the open step holds and Weil's bound follows. \textbf{[Lean from the external index theorem]}
\end{itemize}
Positivity runs the system and certifies the line, for curves. The full Riemann hypothesis for a curve also needs the bound for every extension \texttt{F\_\{q\textasciicircum{}n\}} and the rationality of the zeta function \textbf{[classical, not in the Lean]}.

\subsubsection*{IX.3.5 A positive form puts the eigenvalues on a circle (\texttt{HodgeLinear}, \texttt{HodgeAPI})}
\addcontentsline{toc}{subsubsection}{IX.3.5 A positive form puts the eigenvalues on a circle (\texttt{HodgeLinear}, \texttt{HodgeAPI})}

\begin{itemize}
\item If Frobenius \texttt{F} scales a positive-definite norm by \texttt{c}, every eigenvalue \texttt{mu} has \texttt{abs(mu)\textasciicircum{}2 = c} (\texttt{eigen\_on\_circle}). The same holds with a positive-definite Hermitian form for which \texttt{F} is a similitude (\texttt{eigen\_on\_circle\_form}). \textbf{[Lean]}
\item \textbf{The card} (\texttt{weilCard}): the open step \texttt{StandardHodge} is a positive-definite Hermitian form on which Frobenius is a similitude of ratio \texttt{q\textasciicircum{}w} (the Hodge standard conjecture over \texttt{F\_q}, in the form the Riemann hypothesis needs). The conclusion is every eigenvalue on the circle \texttt{abs(mu)\textasciicircum{}2 = q\textasciicircum{}w}. \textbf{[Lean from the open step]} Deligne (1974) proves the conclusion for l-adic cohomology by another route \textbf{[classical]}, so the card's conclusion is a theorem. Its open step is not known to hold in general.
\end{itemize}
\subsubsection*{IX.3.6 Lefschetz (1,1) and hard Lefschetz, as diagram chases (\texttt{HodgeLinear}, \texttt{HodgeAPI})}
\addcontentsline{toc}{subsubsection}{IX.3.6 Lefschetz (1,1) and hard Lefschetz, as diagram chases (\texttt{HodgeLinear}, \texttt{HodgeAPI})}

\begin{itemize}
\item \textbf{The chase} (\texttt{lefschetz11}): if \texttt{Pic -\textgreater{} H\textasciicircum{}2(X, Z) -\textgreater{} H\textasciicircum{}\{0,2\}} is exact at \texttt{H\textasciicircum{}2} (the exponential sequence) and an integral class has zero \texttt{(0,2)} part, it is the first Chern class of a line bundle.
\item \textbf{Rational descent} (\texttt{rational\_descent}): a rational class with a positive multiple equal to the image of an element of a set lies in the rational span of the images.
\item \textbf{Divisors} (\texttt{DivisorInput}, \texttt{DivisorInput.toInput}, \texttt{DivisorInput.holds}, \texttt{DivisorInput.card}): from three \texttt{External} facts (exactness of the exponential sequence; a rational \texttt{(1,1)} class has an integral multiple with no \texttt{(0,2)} part; divisor classes are \texttt{(1,1)} and every line bundle on a projective X is \texttt{O(D)}), the open step of the Hodge card \textbf{holds for p = 1}, and the described classes are the rational span of divisor classes. \textbf{[Lean from the three external facts]}
\item \textbf{Transfer} (\texttt{transfer}): if an isomorphism of Hodge structures has an inverse that keeps Hodge classes and a map that sends algebraic classes to algebraic classes, the Hodge statement on the source gives it on the target.
\item \textbf{Curves} (\texttt{CurveClassInput}, \texttt{CurveClassInput.toInput}, \texttt{CurveClassInput.holds}): hard Lefschetz as an \texttt{External} carries the statement from divisors to curve classes. \textbf{[Lean from hard Lefschetz and the divisor inputs]}
\end{itemize}
So the card's slot is \textbf{not empty}: it is filled for p = 1 and for p = n - 1, which is exactly the classical range.

\subsubsection*{IX.3.7 Describable must mean rational: the integral version fails (\texttt{HodgeLinear}, \texttt{HodgeAPI})}
\addcontentsline{toc}{subsubsection}{IX.3.7 Describable must mean rational: the integral version fails (\texttt{HodgeLinear}, \texttt{HodgeAPI})}

A torsion class maps to 0 in rational cohomology, and 0 is a Hodge class. So the integral statement makes every torsion class algebraic (\texttt{integral\_needs\_torsion}). With a non-algebraic torsion class as an \texttt{External} (Atiyah-Hirzebruch 1962), the integral statement is false (\texttt{IntegralInput}, \texttt{IntegralHodge}, \texttt{integral\_fails}). \textbf{[Lean from the Atiyah-Hirzebruch class]} The non-projective failure (Zucker, Voisin) is text only: Mathlib has no complex manifolds beyond holomorphic functions.

\subsubsection*{IX.3.8 What Mathlib has and lacks}
\addcontentsline{toc}{subsubsection}{IX.3.8 What Mathlib has and lacks}

Mathlib (v4.34.0, as the notes read it) has no Kaehler metric, no de Rham theorem, no Hodge decomposition, no Chow ring, no cycle class map and no Lefschetz theorems. It has algebraic cycles on a scheme (new in 2026), sheaf cohomology, the signature of a quadratic form, eigenvalues and the circle. So the core is finite-dimensional linear algebra, with the geometry as slots. \textbf{[paper, from the notes' search of the library]}

\subsection*{IX.4 What is borrowed}
\addcontentsline{toc}{subsection}{IX.4 What is borrowed}

Every borrowed result is cited and marked. The project did not rebuild any of it.

{\footnotesize
\begin{longtable}{>{\raggedright\arraybackslash}p{\dimexpr 0.5155\linewidth-2\tabcolsep\relax}>{\raggedright\arraybackslash}p{\dimexpr 0.2468\linewidth-2\tabcolsep\relax}>{\raggedright\arraybackslash}p{\dimexpr 0.2177\linewidth-2\tabcolsep\relax}}
\toprule
\textbf{Borrowed} & \textbf{Where it enters} & \textbf{Status} \\
\midrule
\endhead
The cycle class of a subvariety is a rational \texttt{(p, p)} class & \texttt{Input.alg\_le} & classical \\
\textbf{Hodge-Riemann bilinear relations} on the cohomology of a projective manifold & \texttt{PosInput.hr} & classical \\
The exponential sequence \texttt{0 -\textgreater{} Z -\textgreater{} O -\textgreater{} O* -\textgreater{} 0} and exactness of \texttt{H\textasciicircum{}2(X, Z) -\textgreater{} H\textasciicircum{}2(X, O)} at the image of \texttt{c1}; Hodge symmetry for real classes; divisors are \texttt{(1,1)} and every line bundle is \texttt{O(D)} on a projective X & \texttt{DivisorInput.exact}, \texttt{.hodge\_int}, \texttt{.div\_hodge} & classical \\
\textbf{Hard Lefschetz} as an isomorphism of Hodge structures that sends divisor classes to curve classes & \texttt{CurveClassInput.hl\_hodge}, \texttt{.hl\_alg} & classical \\
Atiyah-Hirzebruch (1962): a non-algebraic torsion class in \texttt{H\textasciicircum{}4(X, Z)} & \texttt{integral\_fails} & external, recalled \\
Intersection numbers on \texttt{C x C} & \texttt{CurveInput.geometry} & classical \\
\textbf{The Hodge index theorem on a surface} (Hodge 1937; Segre; Grothendieck 1958 in every characteristic) & \texttt{curve\_open} & classical \\
Deligne 1974, the Riemann hypothesis over finite fields, by Lefschetz pencils and monodromy & \texttt{weilCard.sources} & classical \\
Abelian-variety cases (Mattuck, Tankeev, Markman and others), cubic fourfolds, uniruled fourfolds, Fermat hypersurfaces & text only & classical / external, recalled; the 2025 abstracts read only \\
Failures: Kollar 1992; Zucker 1977; Voisin 2002 & text only & external, recalled \\
\bottomrule
\end{longtable}}

\subsection*{IX.5 The one open step}
\addcontentsline{toc}{subsection}{IX.5 The one open step}

\textbf{The open step of the card.} The Hodge card has one open step: \textbf{\texttt{HodgeConjecture}. On a smooth complex projective variety, every rational \texttt{(p, p)} class is a rational linear combination of classes of algebraic cycles of codimension p.} It is proved for p = 1 and p = n - 1 (IX.3.6) by classical theorems, and so for dimension at most 3. \textbf{For p \textgreater{}= 2 in general it is open.} Everything else in the card is proved (IX.3) or is a named external theorem (IX.4). \textbf{[open]}

\textbf{The care points, stated exactly.}

\begin{itemize}
\item \textbf{Rational, not integral.} The integral statement is false (Atiyah-Hirzebruch; Kollar). The Lean says why in one line: torsion classes are automatically Hodge (\texttt{integral\_needs\_torsion}).
\item \textbf{Projective, not merely Kaehler.} The statement fails on some complex tori (Zucker; Voisin).
\item \textbf{The card adds no content.} \texttt{problem.proved} is antisymmetry of inclusion. The mathematics is the slot.
\end{itemize}
\textbf{The second slot, a different problem.} Over a finite field the same positivity is the open step of \texttt{weilCard}: the Hodge standard conjecture, in the form that puts the Frobenius eigenvalues on the circle. It is open in general, except for surfaces (the Hodge index theorem) and special cases \textbf{[open]}. The Riemann hypothesis over finite fields is nevertheless a theorem (Deligne), by a route that does not use that positivity \textbf{[classical]}. So this slot's conclusion is known; only its route is open. Over the integers the analogue would be a Hodge index theorem on "\texttt{Spec Z x Spec Z}", which does not exist as an object \textbf{[open]}. That missing positivity is the GRH card's sign slot (Part V).

\textbf{What this does to the problem.} The API does not settle Hodge. It states the conjecture as one inclusion in one slot, shows the classical range filled, and shows the positivity that is a theorem over the complex numbers. The mechanical answer is: \textbf{open for p \textgreater{}= 2, with the open step located.}

\subsection*{IX.6 DMP readings attached to this part (all [reading])}
\addcontentsline{toc}{subsection}{IX.6 DMP readings attached to this part (all [reading])}

\begin{itemize}
\item \textbf{Describable implies exists.} The Hodge conjecture is the cleanest instance of the owner's sentence among the seven problems. The described objects are the rational \texttt{(p, p)} classes (cohomology describes them). The existing objects are the algebraic cycles (geometry builds them). The conjecture is \texttt{describable \textless{}= exists}; the converse is classical. In the API it is literally \texttt{OpenStep (hodge \textless{}= alg)} with \texttt{alg \textless{}= hodge} as an \texttt{External}.
\item \textbf{It is not a logical triviality.} The owner anchored ``what is described exists'' in Goedel's completeness theorem (a consistent first-order theory has a model) \textbf{[classical]}. Hodge is the geometric version. Here a description is a cohomology class with a Hodge-type condition and existence is an algebraic cycle. No completeness theorem bridges them, which is why it is open.
\item \textbf{Describable must mean rational, on a projective variety.} Integral ``descriptions'' include torsion classes, which every \texttt{(p, p)} description admits (their rational image is 0) but which need not exist as cycles. On non-projective Kaehler manifolds descriptions outrun existence. A description counts only after tensoring with Q, and only on a space with an ample class, the polarisation that makes the positivity available. The sentence holds at most in that form, and in that form it is the open conjecture. (Lean for the torsion half: \texttt{integral\_fails}.)
\item \textbf{Hodge types as charges on the G0 circle.} In Deligne's definition the circle acts on \texttt{H\textasciicircum{}\{p,q\}} by \texttt{z\textasciicircum{}\{p-q\}}: a Hodge type is a charge. The Hodge classes are the neutral vectors (\texttt{PolHS.hodge}, fixed by the circle). Described = neutral and rational; exists = built from subvarieties. (Lean for the model: \texttt{PolHS}, \texttt{mem\_hodge}; the identification with DMP's G0 circle is the reading.)
\item \textbf{The Hodge index route is more efficiently correct than Deligne's.} There are two complete proofs of the Riemann hypothesis over finite fields. Weil's (curves, 1948) is one positivity statement, the Hodge index theorem on \texttt{C x C} (\texttt{weil\_bound}). Deligne's (all smooth projective varieties, 1974) reaches the same bound without a positivity statement on cycles. The owner prefers the first because it originates directly from G0 and so simplifies further. Both proofs are complete and correct. The difference is that the positivity route has \textbf{one moving part}, Cauchy-Schwarz on a positive form (\texttt{cs\_of\_nonneg}), the same shape as the encoder bound of Yang-Mills and of \texttt{PosInput.alg\_cs}. That is why it transfers between problems. \textbf{Reach:} the positivity route is proved for curves and, for the index theorem, for surfaces in every characteristic. In higher dimension it needs the standard conjecture of Hodge type in characteristic p \textbf{[open]}, where Deligne's route covers everything. For GRH over number fields the missing piece is this route's analogue for \texttt{Spec Z} \textbf{[open]}. So choosing the route names the open step precisely; it does not close it.
\item \textbf{The shared positivity with the GRH column.} ``Exists implies positive'' is, for Hodge, the Hodge-Riemann relations: on the described classes the polarisation is positive (\texttt{hodge\_pos}, in the model), so on the existing classes too (\texttt{alg\_pos}). Over C that positivity is proved. Over \texttt{F\_q} the same positivity is Grothendieck's Hodge standard conjecture, and it would put the Frobenius eigenvalues on the critical circle (\texttt{eigen\_on\_circle\_form}). For curves the index theorem on \texttt{C x C} supplies it (\texttt{curve\_open}, \texttt{weil\_bound}). The two columns have one shape: \textbf{a positive form, a dynamics that is a similitude of it, and the line as the consequence.} In GRH and in Hodge over \texttt{F\_q} the positive form is the open step; over C, and for curves over \texttt{F\_q}, it is a theorem.
\end{itemize}
\subsection*{IX.7 What would close it}
\addcontentsline{toc}{subsection}{IX.7 What would close it}

If an independent proof of the Hodge conjecture for p \textgreater{}= 2 were found, the card's open step would be discharged by that work. This part would supply the shape and the audit. The change to this paper would be one row of the ledger, one verdict line of this part and one item of the abstract. Nothing in this round proves any part of that.

The formalisation items that would change the status of the \textbf{borrowed} parts, none done and all named in the notes as reachable or far:

\begin{enumerate}
\item \textbf{Reachable on Mathlib, not built:} hard Lefschetz as the linear algebra of an \texttt{sl\_2}-representation (an \texttt{sl\_2}-triple on a graded space gives \texttt{L\textasciicircum{}k : V\_\{-k\} \textasciitilde{} V\_k} and the primitive decomposition); the Hodge-Riemann signature formula for a polarised Hodge structure of any weight; Hodge structures of weight k with the Hodge filtration; finiteness of lattice Hodge classes of bounded norm. \textbf{[paper, an estimate]}
\item \textbf{Far, needs geometry Mathlib lacks:} Kaehler manifolds, harmonic forms, the Hodge decomposition itself, the cycle class map, the exponential sequence on a complex manifold. \textbf{[paper, an estimate]}
\item A Hodge row in the notes' API table and an import of \texttt{Unified/HodgeAPI.lean} into \texttt{Unified/All.lean}: left to the coordinator.
\end{enumerate}
\subsection*{IX.8 The box in the one-box test}
\addcontentsline{toc}{subsection}{IX.8 The box in the one-box test}

\textbf{Gate factor and marker.} The ticks are the \textbf{+dimensions} (the owner's name; in plain terms, the codimensions). The owner's reason: at each codimension the Hodge-Riemann form is positive, up to its fixed sign, on the primitive classes (an \texttt{External} slot here, IX.3.1), so each is a reflection-positive covariant dimension, and what is positive exists \textbf{[reading]}. At +dimension p the gate reads inside when every rational (p,p) class is a rational combination of classes of algebraic cycles; the marker is the claim ``inside''. The box of an input passes exactly when the Hodge conjecture holds for it \textbf{[Lean]} (\texttt{hodge\_box\_iff}, \texttt{ETM/CardBox.lean}; claim IX.8a). For divisors, codimension one, the box passes, from Lefschetz (1,1) with its named inputs \textbf{[Lean]} (from \texttt{External}; \texttt{hodge\_box\_divisor}). The box is the card's open step in the test's language (II.4): it restates the step and proves no new bound, and the content stays in the step \textbf{[paper]}. Its open ticks start at the second +dimension (IX.5.1).


%% ---- 10-part-X-gu.md
\section*{Part X. GU: the Grand Unified reading}
\addcontentsline{toc}{section}{Part X. GU: the Grand Unified reading}

\subsection*{X.0 What ``Grand Unified'' means here, and where it stands}
\addcontentsline{toc}{subsection}{X.0 What ``Grand Unified'' means here, and where it stands}

\textbf{GU is not a Clay problem.} There is no official statement to quote. In this work ``Grand Unified'' is the second half of the owner's \textbf{release criterion} (\texttt{OUTLINE.md}): the paper may leave the draft stage when both readings of GUID hold. \textbf{General} is the unified theorem with its open hypotheses closed. \textbf{Grand Unified} is \emph{the framework containing SU(3) x SU(2) x U(1) in one structure with at least one checkable consequence}. The owner's wording is a goal (``what I want to claim when finished, not what I can prove now''), so this part states exactly how far each half of the Grand Unified reading holds today.

\textbf{State (2026-11-24).}

\begin{itemize}
\item \textbf{``In one structure'': met, in a modest and classical sense.} The group SU(3) x SU(2) x U(1) divided by a cyclic group of order 6 sits inside SU(5), and the kernel of that embedding is proved in Lean \textbf{[Lean]}, with the group-level facts about it \textbf{[classical]} (X.3).
\item \textbf{``At least one checkable consequence'': the owner's decision, uniform charge quantisation.} The notes find two candidates and prove what can be proved about each (X.5). The owner takes candidate (b), with its guarantee from the ODP: every charge is an integer multiple of one unit, the same unit for all matter, and the unit is the ODP's dimensionless 1 (X.5) \textbf{[reading]} for the guarantee, \textbf{[Lean]} for the integer charges on a compact circle, \textbf{[external]} for the measurements. Neither candidate is a consequence that the standard expectation of any compact U(1) or any SU(5) does not already contain; a consequence beyond it stays \textbf{[open]} (X.5e).
\end{itemize}
\textbf{What this part does not do.} It does not derive the Standard Model's hypercharges, does not predict a new particle or number, and contains no new physics. The hypercharges are an \textbf{input}.

\subsection*{X.1 The problem as meant in this work}
\addcontentsline{toc}{subsection}{X.1 The problem as meant in this work}

\textbf{X.1.1 Grand unification in physics} \textbf{[classical]}. The Standard Model gauge group SU(3) x SU(2) x U(1) is embedded in one simple group (SU(5): Georgi and Glashow 1974; SO(10): Georgi 1975, Fritzsch and Minkowski 1975), with one coupling at the unification scale, from which the three couplings run apart toward low energies. Checkable consequences in such theories are coupling relations, charge quantisation, proton decay and magnetic monopoles. Care points the paper keeps: the single coupling holds at one scale only; minimal non-supersymmetric SU(5) is excluded by proton-decay limits and the three couplings do not meet at one point \textbf{[external]} (X.4); and not every SU(N) is a viable grand unified group, since a GUT must contain SU(3) x SU(2) x U(1), cancel the gauge anomalies with the observed matter and survive the proton-decay limits, which singles out SU(5) and its extensions, SO(10) and E6.

\textbf{X.1.2 The quotient.} The Standard Model fields are all invariant under a cyclic group of order 6 sitting in the centre of SU(3) x SU(2) x U(1). The gauge group can therefore be taken as the quotient by any subgroup Gamma of Z\_6, that is by Z\_1, Z\_2, Z\_3 or Z\_6. All four have the same local physics and the same Feynman rules, and present experiments do not decide between them (Tong 2017) \textbf{[classical]}. The maximal choice Z\_6 is what puts electric charge in thirds, correlated with colour. It is what any SU(5) or SO(10) unification forces, since the SU(5) embedding has exactly this Z\_6 as its kernel.

\subsection*{X.2 The card}
\addcontentsline{toc}{subsection}{X.2 The card}

The DMP API (\texttt{docs/unified/API.md}) gives every problem a well-defined input, a proved structure on it, and one named open step. GU has that shape, but it is \textbf{not a \texttt{Problem} card in the Lean API}: its Lean pieces are group-level statements in \texttt{General/}, plus the API's one general statement about the G0 phase field (X.3.3). The table is the card in prose.

{\footnotesize
\begin{longtable}{>{\raggedright\arraybackslash}p{\dimexpr 0.2526\linewidth-2\tabcolsep\relax}>{\raggedright\arraybackslash}p{\dimexpr 0.7274\linewidth-2\tabcolsep\relax}}
\toprule
\textbf{Card element} & \textbf{GU} \\
\midrule
\endhead
\textbf{Well-defined input} & the representation content of one generation and the Higgs: for each field its triality \texttt{t}, duality \texttt{d} and hypercharge in sixths \texttt{q = 6Y} (\texttt{smFields}, a definition in Lean). \textbf{The hypercharges are input, not output.} \\
\textbf{Proved structure} & the centre elements of SU(3) x SU(2) x U(1) that act trivially on all these fields form exactly a cyclic group of order 6, split as 2 x 3; every invariant component has \texttt{Q + t/3} an integer (charge in thirds, correlated with colour); the SU(5) block embedding has exactly this group as kernel; a compact U(1) admits only integer charges \textbf{[Lean]} \\
\textbf{External inputs} & the hypercharge values (fixed by anomaly cancellation, now in Lean, X.7b, or by SU(5)); the measured neutrality of matter, the absence of free fractional charges, proton-decay limits, the failure of coupling unification, the absence of monopoles \textbf{[classical]} / \textbf{[external]} \\
\textbf{The one open step} & a checkable consequence beyond the standard expectation, supplied by DMP \textbf{[open]}; for the gate's wording the owner takes uniform charge quantisation (X.5) \textbf{[reading]} \\
\textbf{Conclusion} (if the step closes) & the Grand Unified half of the release criterion holds \\
\bottomrule
\end{longtable}}

\subsection*{X.3 What is proved}
\addcontentsline{toc}{subsection}{X.3 What is proved}

All of this is \textbf{[Lean]}, standard axioms only (\texttt{docs/unified/gu-axioms.log}: 19 declarations on the three standard axioms, one on \texttt{propext} alone, two definitions on none; \texttt{docs/unified/gucharge-axioms.log}: all 12 declarations on the three standard axioms; no \texttt{sorry}, \texttt{admit}, \texttt{native\_decide} or \texttt{axiom} in either module). The statements are about groups and their characters, that is, about the possible charges of fields. \textbf{Not formalised:} a gauge field theory with this U(1); that the image of the SU(5) map is all of S(U(3) x U(2)); any lattice or twist-sector statement for the product group; chiral matter.

\subsubsection*{X.3.1 The Z\_6 (\texttt{General/GUCentre.lean}, namespace \texttt{ProofPhase.General.GU})}
\addcontentsline{toc}{subsubsection}{X.3.1 The Z\_6 (\texttt{General/GUCentre.lean}, namespace \texttt{ProofPhase.General.GU})}

One generation, with \texttt{q = 6Y} and \texttt{2t + 3d + q} the exponent of the generator's phase:

{\footnotesize
\begin{longtable}{>{\raggedright\arraybackslash}p{\dimexpr 0.1326\linewidth-2\tabcolsep\relax}>{\raggedright\arraybackslash}p{\dimexpr 0.1842\linewidth-2\tabcolsep\relax}>{\raggedright\arraybackslash}p{\dimexpr 0.1326\linewidth-2\tabcolsep\relax}>{\raggedright\arraybackslash}p{\dimexpr 0.1326\linewidth-2\tabcolsep\relax}>{\raggedright\arraybackslash}p{\dimexpr 0.1326\linewidth-2\tabcolsep\relax}>{\raggedright\arraybackslash}p{\dimexpr 0.1326\linewidth-2\tabcolsep\relax}>{\raggedright\arraybackslash}p{\dimexpr 0.1326\linewidth-2\tabcolsep\relax}}
\toprule
\textbf{field} & \textbf{rep (SU(3), SU(2))} & \textbf{Y} & \textbf{q} & \textbf{t} & \textbf{d} & \textbf{\texttt{2t + 3d + q}} \\
\midrule
\endhead
Q\_L & (3, 2) & 1/6 & 1 & 1 & 1 & 6 \\
u\_R & (3, 1) & 2/3 & 4 & 1 & 0 & 6 \\
d\_R & (3, 1) & -1/3 & -2 & 1 & 0 & 0 \\
L\_L & (1, 2) & -1/2 & -3 & 0 & 1 & 0 \\
e\_R & (1, 1) & -1 & -6 & 0 & 0 & -6 \\
H & (1, 2) & 1/2 & 3 & 0 & 1 & 6 \\
\bottomrule
\end{longtable}}

A centre element \texttt{(u 1\_3, v 1\_2, z)} acts on a field by \texttt{u\textasciicircum{}t v\textasciicircum{}d z\textasciicircum{}q} (\texttt{centrePhase}). Then:

\begin{itemize}
\item it acts trivially on all fields exactly when \texttt{z\textasciicircum{}6 = 1}, \texttt{u = z\textasciicircum{}2}, \texttt{v = z\textasciicircum{}3} (\texttt{sm\_kernel}); equivalently the kernel is generated by \texttt{(zeta\textasciicircum{}2, zeta\textasciicircum{}3, zeta)} with \texttt{zeta = e\textasciicircum{}\{2 pi i/6\}}, a primitive sixth root (\texttt{sm\_kernel\_cyclic}, \texttt{zeta6\_prim}); the generator has order 6 (\texttt{gen\_order}) and splits into a primitive cube root and \texttt{-1} (\texttt{gen\_parts});
\item \textbf{the CRT split}: the order-2 part \texttt{(1, -1, -1)} and the order-3 part \texttt{(zeta\textasciicircum{}4, 1, zeta\textasciicircum{}2)} each act trivially (\texttt{sm\_kernel\_two}, \texttt{sm\_kernel\_three}); at the level of \texttt{Z/6} the split is \texttt{centre6 : ZMod 6 = ZMod 2 x ZMod 3} (\texttt{General/SUNTwistPrimes.lean});
\item the generator is trivial on a field exactly when 6 divides \texttt{2t + 3d + q}, and this holds for every field in the table (\texttt{gen\_trivial\_iff}, \texttt{gen\_trivial\_sm});
\item \textbf{charge in thirds}: if 6 divides \texttt{2t + 3d + q}, every component with \texttt{2 T\_3 = s} has charge \texttt{Q = (3s + q)/6} and \texttt{Q + t/3} is an integer (\texttt{charge\_thirds}, \texttt{charge\_thirds\_sm}). So electric charge modulo 1 is fixed by colour: integer for colour singlets, in \texttt{-1/3 + Z} for quarks;
\item \textbf{the falsifiable side}: a colour singlet with a component of non-integer charge is not invariant under the generator (\texttt{fractional\_singlet\_breaks}), so it cannot exist if the gauge group is the full quotient by Z\_6;
\item \textbf{SU(5)}: the block map \texttt{(A, B, z) -\textgreater{} diag(z\textasciicircum{}\{-2\} A, z\textasciicircum{}3 B)} has determinant 1 and its kernel is exactly \texttt{\{(z\textasciicircum{}2 1, z\textasciicircum{}3 1, z) : z\textasciicircum{}6 = 1\}} (\texttt{su5Embed}, \texttt{su5\_embed\_det}, \texttt{su5\_embed\_kernel}). So any SU(5) unification forces Gamma = Z\_6.
\end{itemize}
The centres of the factors are also in Lean: for \texttt{omega\textasciicircum{}N = 1} the scalar matrix \texttt{omega I} lies in SU(N) and is central; the map \texttt{zeta -\textgreater{} zeta I} is an injective homomorphism from the N-th roots of unity, a copy of Z\_N with N elements (\texttt{General/Centre.lean}: \texttt{centreHom}, \texttt{centreHom\_injective}, \texttt{centreHom\_central}, \texttt{card\_centre}, \texttt{centreHom\_rep}). Every SU(N) centre therefore sits as phases on one common circle.

\subsubsection*{X.3.2 Charge quantisation from a compact U(1) (\texttt{General/GUCharge.lean})}
\addcontentsline{toc}{subsubsection}{X.3.2 Charge quantisation from a compact U(1) (\texttt{General/GUCharge.lean})}

If the U(1) of hypercharge is a circle, not a line, it admits only integer charges. In Lean, from Fourier analysis and Parseval's identity on the circle (Mathlib has neither this classification nor the character shift, so the steps are proved here):

\begin{itemize}
\item \texttt{fourierCoeff\_char\_shift}: the Fourier coefficients of a character shift by its value (translation invariance of Haar measure);
\item \texttt{addCircle\_char}: every continuous character \texttt{chi} of \texttt{AddCircle T} is \texttt{x -\textgreater{} toCircle(n x)} for an integer \texttt{n}; \texttt{addCircle\_char\_unique}: the integer is unique;
\item \textbf{\texttt{circle\_char}: every continuous homomorphism from the circle to the circle is \texttt{z -\textgreater{} z\textasciicircum{}n}, with \texttt{n} a unique integer} (\texttt{circle\_char\_unique});
\item \textbf{on a line, every real \texttt{a} gives a character \texttt{t -\textgreater{} e\textasciicircum{}\{iat\}}} (\texttt{line\_char}), and it factors through the circle exactly when \texttt{a} is an integer (\texttt{line\_char\_periodic\_iff}). So quantisation comes from compactness and from nothing else;
\item an N-th root of unity on the circle (an SU(N) centre element) sees \texttt{z\textasciicircum{}n} as \texttt{omega\textasciicircum{}(n mod N)} (\texttt{centre\_char\_mod}).
\end{itemize}
\textbf{The chain} (\texttt{g0\_charge\_thirds}). Input, stated explicitly: a field on which hypercharge acts by a continuous character \texttt{chi} of the circle, with triality \texttt{t} and duality \texttt{d}, invariant under the generator \texttt{(omega, -1, zeta)}. Then \texttt{chi = z -\textgreater{} z\textasciicircum{}q} with \texttt{q} an integer \textbf{(derived, from compactness)}, 6 divides \texttt{2t + 3d + q}, and \texttt{Q + t/3} is an integer for every component. The split into thirds uses the matter input; the integer \texttt{q} does not.

\subsubsection*{X.3.3 The G0 link as an API statement (\texttt{API.lean}, \texttt{API/Phase.lean})}
\addcontentsline{toc}{subsubsection}{X.3.3 The G0 link as an API statement (\texttt{API.lean}, \texttt{API/Phase.lean})}

G0's phase field is one compact U(1), Mathlib's \texttt{Circle} (\texttt{Phase}). A label carried on it is a continuous character (\texttt{PhaseLabel}). \texttt{ChargeQuantisation} names the statement; it is proved with no hypothesis (\texttt{chargeQuantisation}), and \texttt{PhaseLabel.integer\_charge} says every label on the phase field has an integer charge \textbf{[Lean]}. This is the one consequence of the G0 link that holds for \textbf{every} instance of the API, not only for GU. \texttt{g0\_charge\_thirds} is the GU instance of it.

\subsection*{X.4 What is borrowed}
\addcontentsline{toc}{subsection}{X.4 What is borrowed}

\textbf{X.4.1 The hypercharges} \textbf{[classical]}. Given the Standard Model representations of one generation, cancellation of the gauge anomalies and the mixed gauge-gravitational anomaly fixes the hypercharge ratios (Geng and Marshak 1989; Minahan, Ramond and Warner 1990); with a right-handed neutrino a B - L direction opens unless a Majorana mass is assumed (Babu and Mohapatra 1989). Alternatively the ratios follow from SU(5). Either way they are an \textbf{input to the Z\_6, not an output of it}. Anomaly cancellation is now in Lean (X.7b, \texttt{ProofPhase/Closures/Anomaly.lean}: \texttt{anomaly\_fixes\_sm}, \texttt{anomaly\_fixes\_Z6}): for one generation the four anomaly conditions fix the charges up to u \textless{}-\textgreater{} d, and with them the Z\_6. This is classical physics formalised, not a DMP result.

\textbf{X.4.2 The global form} \textbf{[classical]}. Which of Z\_1, Z\_2, Z\_3, Z\_6 nature uses is not decided by present experiment (Tong, JHEP 07 (2017) 104). The choice changes which new particles and line operators are allowed:

{\footnotesize
\begin{longtable}{>{\raggedright\arraybackslash}p{\dimexpr 0.1517\linewidth-2\tabcolsep\relax}>{\raggedright\arraybackslash}p{\dimexpr 0.3072\linewidth-2\tabcolsep\relax}>{\raggedright\arraybackslash}p{\dimexpr 0.2606\linewidth-2\tabcolsep\relax}>{\raggedright\arraybackslash}p{\dimexpr 0.2606\linewidth-2\tabcolsep\relax}}
\toprule
\textbf{Gamma} & \textbf{trivial on a field iff} & \textbf{colour singlet with Q = 1/3} & \textbf{colour singlet with Q = 1/2} \\
\midrule
\endhead
Z\_1 & always & allowed & allowed \\
Z\_2 & \texttt{q = d} mod 2 & allowed & forbidden \\
Z\_3 & \texttt{2t + q = 0} mod 3 & forbidden & allowed \\
Z\_6 & \texttt{2t + 3d + q = 0} mod 6 & forbidden & forbidden \\
\bottomrule
\end{longtable}}

It also changes the smallest magnetic charge of a monopole and the periodicities of theta angles (Tong 2017), and for Z\_6 the fractional part of the axion-photon coupling is tied to the axion-gluon coupling (Reece 2023; Choi, Forslund, Lam and Shao 2023). All \textbf{[classical]}.

\textbf{X.4.3 Experiments} \textbf{[external]} (cited as the project notes cite them; not re-read for this draft).

{\footnotesize
\begin{longtable}{>{\raggedright\arraybackslash}p{\dimexpr 0.3433\linewidth-2\tabcolsep\relax}>{\raggedright\arraybackslash}p{\dimexpr 0.6367\linewidth-2\tabcolsep\relax}}
\toprule
\textbf{Check} & \textbf{Result} \\
\midrule
\endhead
neutrality of matter (one charge unit) & electron-proton charge difference below about 1 part in 10\textasciicircum{}21 (Bressi et al. 2011) \\
free fractional charges & none found (review: Perl, Lee and Loomba 2009) \\
proton decay & not seen; Super-Kamiokande bound above 10\textasciicircum{}34 years for p -\textgreater{} e+ pi0; minimal non-supersymmetric SU(5) is excluded \\
coupling unification & the three couplings do not meet at one point; with supersymmetric running they nearly do \\
magnetic monopoles & none found (MACRO, IceCube flux limits; MoEDAL production limits) \\
\bottomrule
\end{longtable}}

\subsection*{X.5 The one open step: the gate's ``at least one checkable consequence''}
\addcontentsline{toc}{subsection}{X.5 The one open step: the gate's ``at least one checkable consequence''}

\textbf{The gate, as worded:} the framework contains SU(3) x SU(2) x U(1) in one structure \textbf{with at least one checkable consequence}. \textbf{[open]}

The notes find two candidate consequences. The table states, for each, what is derived, what is input and what the standard expectation already says. It makes no recommendation.

{\footnotesize
\begin{longtable}{>{\raggedright\arraybackslash}p{\dimexpr 0.1893\linewidth-2\tabcolsep\relax}>{\raggedright\arraybackslash}p{\dimexpr 0.3879\linewidth-2\tabcolsep\relax}>{\raggedright\arraybackslash}p{\dimexpr 0.4028\linewidth-2\tabcolsep\relax}}
\toprule
\textbf{} & \textbf{Candidate (a): the selection of Z\_6} & \textbf{Candidate (b): charge quantisation from the G0 link} \\
\midrule
\endhead
\textbf{The statement} & the U(1) of G0 is glued to the SU(2) and SU(3) centres as fully as the matter allows, which selects the maximal quotient Z\_6 out of the four & every charge on the G0 circle is an integer multiple of one unit, with the unit splitting into thirds for quarks \\
\textbf{Checkable by} & a colour singlet with fractional charge (forbidden for Z\_6), a monopole's smallest charge, the fractional part of axion couplings & neutrality of matter; the absence of free fractional charges; Dirac quantisation of any monopole \\
\textbf{What is derived} & the arithmetic of Z\_6 and its consequences, given the hypercharges \textbf{[Lean]}; the selection itself is a \textbf{[reading]} & the integer charge, from compactness alone \textbf{[Lean]}; the thirds, from the input \textbf{[Lean]} \\
\textbf{What is input} & the hypercharges; that the gluing is ``as full as the matter allows'' & that the hypercharge U(1) is compact (in DMP the reading ``G0's U(1) possibility field'', with the centres on one circle in Lean); the hypercharges \\
\textbf{Status of the evidence} & no present measurement decides Z\_6 against Z\_1, Z\_2, Z\_3 & measured: matter neutral to about 1 part in 10\textasciicircum{}21 \textbf{[external]} \\
\textbf{What the standard picture already says} & Z\_6 is forced by any SU(5) or SO(10) unification and is the textbook expectation & charge quantisation is the standard expectation of any compact U(1) and of any theory whose U(1) sits in a compact simple group \\
\textbf{What would make it DMP-native} & a DMP theorem, stated before the hypercharges, that selects the quotient; \textbf{none exists} & it already follows from the link to G0, with the circle as the project's own object \\
\bottomrule
\end{longtable}}

\textbf{Why the project's present pieces cannot see candidate (a).} All twist-sector results (the \texttt{SUNTwist} family) are for pure gauge theory with one group. In pure gauge theory the whole centre is a symmetry; Z\_6 is singled out only once matter is added. The selection lives where the project has nothing yet: matter, and in particular chiral fermions.

\textbf{Why 2-then-3 does not give Z\_6 as a consequence} \textbf{[classical]} + \textbf{[Lean]}. That the centre of SU(2) x SU(3) is Z\_2 x Z\_3 = Z\_6 follows once the factors are chosen; the reading ``2 first, then 3'' restates the choice. Which part of Z\_6 is divided out is decided by the matter, and present experiments do not decide it. Z\_6 = Z\_2 x Z\_3 is symmetric: it has no ``first 2, then 3''. An order appears only in the symmetry breaking (X.6).

\textbf{The owner's position, recorded.} The owner's reading is that GU comes down to well-defined inputs, which DMP can guarantee, and that the checkable consequence comes from the link to G0 (candidate (b)).

\textbf{The owner's decision (2026-11-24).} ``Uniformity is also guaranteed to be constructed by the ODP'': the guarantee of ``well-defined inputs, guaranteed'' is the same as the one that gives Yang-Mills its lower edge (Part III 4.4). The ODP ends only at a dimensionless 1, and that 1 is the unit every charge is an integer multiple of. So the checkable consequence for the gate is uniform charge quantisation, one unit for all matter, checked by the neutrality of matter and the absence of free fractional charges \textbf{[reading]} for the guarantee (X.5f), \textbf{[Lean]} for the integer charges on the compact circle (X.3.2), \textbf{[external]} for the measurements (X.4.3a, X.4.3b). Care point: it coincides with the standard expectation of any compact U(1); the paper claims no consequence beyond it (X.5e stays open).

\subsection*{X.6 DMP readings attached to this part (all [reading])}
\addcontentsline{toc}{subsection}{X.6 DMP readings attached to this part (all [reading])}

\begin{itemize}
\item \textbf{``G0's U(1) possibility field'' selects the maximal quotient.} If every centre phase is a point on the one shared circle, the natural gluing identifies the SU(2) and SU(3) centres with hypercharge phases as fully as the matter allows: Gamma = Z\_6. This is the sharpest DMP-flavoured statement found. It is checkable (X.5) but is not a derivation: ``as fully as the matter allows'' uses the hypercharges, and the conclusion coincides with what SU(5) gives.
\item \textbf{``The universal quantum running'' read as charge quantisation.} The owner's entry, in the coordinator's reading: the link to G0 makes the phase group a circle, and a circle carries only integer charges. The step from the circle to integer charges is proved (X.3.2); that this is what the entry means is the reading.
\item \textbf{``2 first, then 3, with the hand-over at the state 2, 3.''} An order does appear in the symmetry breaking \textbf{[classical]}: at the electroweak scale the Higgs vacuum (read as the SU(2) ``handshake'') leaves the unbroken group (SU(3) x U(1)\_em)/Z\_3. The Z\_2 part is absorbed, and the Z\_3 part survives to low energies as the correlation of charge with colour. In descending energy the 2-part acts first and the 3-part remains. This matches the owner's ordering and predicts nothing beyond the standard picture.
\item \textbf{``SU(5) with a single coupling is G0 at GRH, with a forward arrow independent of scale.''} SU(5) is the smallest simple group that contains the Z\_6 quotient (X.3.1). Its single coupling holds at one scale only, minimal SU(5) is excluded by proton decay, and the couplings do not meet exactly (X.4.3). The owner's mapping (one coupling = one G0; the forward arrow = the running from the unification scale) is a reading, and a programme beyond Yang-Mills, GRH and NP vs coNP.
\item \textbf{Every SU(N) as a valid unified theory, indexed by the values of zeta.} On \texttt{s \textgreater{} 1} the zeta function decreases continuously from infinity to 1, so it takes every integer \texttt{N \textgreater{}= 2} exactly once, which follows from continuity and monotonicity \textbf{[classical]}. So zeta indexes every SU(N), one to one. The project's gauge-group package covers every SU(N) (\texttt{General/SUN*.lean}). That every SU(N) ``describes existence'' is the owner's reading; as physical unified theories most SU(N) are not viable (X.1.1).
\item \textbf{SU(2) as the stable profile} (Appendix A): a statement about the group manifold under Ricci flow, not about the global form of the Standard Model group. No link to Z\_6 was found.
\item \textbf{The rigid groups.} Which compact groups are rigid, which carry content across the gate, and the U(1) gate's two hidden witnesses are in Appendix A.5.
\end{itemize}
\subsection*{X.7 What would close it}
\addcontentsline{toc}{subsection}{X.7 What would close it}

Ranked by how directly each would answer the gate. None is promised.

\begin{enumerate}
\item \textbf{Derive the selection, not the arithmetic.} A DMP principle, stated before the hypercharges and proved as a theorem about running existence spaces, whose consequence is ``the centre twists that survive matter are glued to the U(1) maximally''. With the representation content as input it would output Gamma = Z\_6 and hence the three checks of X.4.2. This is the only route found on which DMP would contribute a consequence rather than a description.
\item \textbf{Anomaly cancellation in Lean} (done, X.7b): from the one-generation representation content the anomaly equations fix the hypercharge ratios, hence Z\_6 (\texttt{Closures/Anomaly.lean}, \texttt{anomaly\_fixes\_sm}, \texttt{anomaly\_fixes\_Z6}). It makes the input a classical theorem, not a DMP one, and leaves the card's open step (a checkable consequence) where it was.
\item \textbf{Twist sectors of the product group with Z\_6-invariant matter}: extend the \texttt{SUNTwist} chain to a lattice SU(3) x SU(2) x U(1) theory with vector-like matter, so that the Z\_6 twist sectors are strictly between nothing and everything, in the project's own language. Chiral matter, which the Standard Model needs, is a known hard problem on the lattice and stays out of reach.
\item \textbf{The single coupling}: a statement about running couplings needs the continuum theory with matter, beyond Yang-Mills as treated here.
\end{enumerate}
Until one of these lands, the paper can say exactly this: the framework contains SU(3) x SU(2) x U(1) / Z\_6 as one structure (classical, Lean-stated); the owner's readings select Z\_6, which a fractional-charge discovery, a monopole's charge or an axion's couplings could refute; charge quantisation from the circle is proved and measured but is also the standard expectation; and no consequence is derived from DMP beyond what the standard expectation contains. Which of these, if any, counts for the gate is the owner's decision.

\subsection*{X.8 The box in the one-box test}
\addcontentsline{toc}{subsection}{X.8 The box in the one-box test}

\textbf{Gate factor and marker.} The ticks are the \textbf{points} (the owner's name; he also offers ``uniquenesses'' and ``DoEs''): the labels carried on G0's phase field. The name is exact: the continuous characters of the circle are the integers, the points of its dual group \textbf{[classical]} (Pontryagin duality). At each point the gate reads inside when the label has an integer charge; the marker is the claim ``inside''. This box passes, proved: every continuous character of the circle is z to the n \textbf{[Lean]} (\texttt{gu\_box}, \texttt{ETM/CardBox.lean}; claim X.8a). The GU card's open step, a consequence beyond the standard expectation (X.5e), is not a statement about this box; the box shows what the G0 link already gives every label, and what stays open lies beyond it \textbf{[paper]}.


%% ---- 11-outlook-settling-machine.md
\section*{Outlook. A settling machine: the first author's concept of a DMP-native general intelligence}
\addcontentsline{toc}{section}{Outlook. A settling machine: the first author's concept of a DMP-native general intelligence}

\subsection*{XI.1 The concept, in the first author's terms}
\addcontentsline{toc}{subsection}{XI.1 The concept, in the first author's terms}

A general intelligence is a CMOS system that runs entirely in IM tokens, adapted to the acts it has to assert on G0. \textbf{[reading]}

\begin{itemize}
\item \textbf{One interface to the outside: Milky Way.} Everything that leaves the system passes one Milky Way interface. Milky Way checks twice, by hash and TTL, that the IM is correct and a valid act, and only then sends it out as GG: the content of that ``forward pass'', or settling.
\item \textbf{The output rate is given from outside.} It is not a measure of how efficiently the system infers. It is an upper bound, set externally, on the rate of valid GG tokens. Between two outputs, the inside organizes itself toward the most balanced output for the currently asserted state of existence.
\item \textbf{Total freedom inside.} Behind Milky Way the system has total act freedom. It can change its internal layout, recombine itself for new inputs, or run many iterations of some logic. It keeps one periodicity, a tick: a constant reminder that it exists between outputs, rather than producing one token at a time in a line. Because it can change its coupling, and so its content, while still recognizing its state by the data that state sits on, it can run more efficiently.
\item \textbf{The meta layer.} The first author's expectation is that such a system works directly on the vectors that embed existence. That embedding is its meta layer: on it, the system can construct arbitrary logic, as long as a correct answer comes out. \textbf{[reading]} Known neighbours: reasoning carried out in a continuous latent space instead of in output tokens; vector-symbolic (hyperdimensional) architectures, where binding and superposition of vectors implement symbolic logic; and fast weights or hypernetworks, where one network writes the representations or weights another one uses. \textbf{[classical]}, as pointers only.
\item \textbf{Fast weights as the native form.} The first author reads fast weights and hypernetworks as DMP-native. The duality gives an existence (the embedding) on which a meta layer can be spanned, and a definition of running on that embedding. Changes to the embedding go through IM, and possibly GG, so that the system can change while its existence stays guaranteed. \textbf{[reading]} The classical counterpart: a slow network writes the fast weights another network runs with (fast weight programmers), and linear-attention transformers are exactly such programmers (Schlag, Irie, Schmidhuber 2021). \textbf{[classical]}, pointer only.
\item \textbf{Training.} The core network has to be trained from the ground up for this way of running.
\end{itemize}
\subsection*{XI.2 The critical link to this work}
\addcontentsline{toc}{subsection}{XI.2 The critical link to this work}

The concept depends on one property: that the inside settles to a balanced output within the tick. This is exactly the decidable half of the Existence Turing Machine (Part II and \texttt{docs/unified/ETM.md}).

\begin{itemize}
\item \textbf{Settling is guaranteed when every internal step contracts}, i.e. pays a cost above the floor. The state is then within epsilon of balance after an explicit, computable number of steps. \textbf{[Lean]} \texttt{ProofPhase/ETM/Machine.lean}: \texttt{reaches\_of\_contracting}, \texttt{balanceBound\_ge}; a contracting step is a cost above the floor, \texttt{contracting\_iff\_costAboveFloor}.
\item \textbf{Without a contraction, no procedure decides settling in general.} \textbf{[Lean]} \texttt{etm\_balance\_undecidable}, from Mathlib's halting theorem.
\end{itemize}
\textbf{Changing the system while keeping its existence guaranteed.} Two kinds of change keep the guarantee, both in Lean:

\begin{itemize}
\item \textbf{Changes that are maps of running systems} (relabelling, restricting to a part the run never leaves, coarse-graining by a relation the run respects) keep the ledger exactly. Step, factor and carried amount commute with the map. \textbf{[Lean]} \texttt{ProofPhase/ETM/RunSys.lean}: \texttt{Hom}, \texttt{Hom.iterate\_comm}, \texttt{sub}, \texttt{quot}.
\item \textbf{Changes that keep every factor at most theta \textless{} 1} (any new layout or coupling inside the contracting class) keep the settle bound, whatever else they change. \textbf{[Lean]} \texttt{RunSys.contract}, \texttt{reaches\_of\_contracting}.
\end{itemize}
In the first author's terms, an IM change (inside) must stay in the second kind. A GG change (from outside) passes the gate first. Self-modification of the embedding then never loses the guarantee that the system settles. \textbf{[reading]} for the IM/GG identification.

So the training target the concept needs has a precise form: train the core so that its update is a contraction, and choose the externally set tick from the settle bound. The same condition, a contracting update, is what guarantees a unique fixed point and stable training in equilibrium (fixed-point) neural networks. \textbf{[classical]}, as a pointer only.

\subsection*{XI.3 Known neighbours and care points}
\addcontentsline{toc}{subsection}{XI.3 Known neighbours and care points}

\textbf{Neighbours} (pointers, not re-read here): equilibrium and energy-based networks, recurrent-depth reasoning models and adaptive computation time (settling between outputs); anytime algorithms (a fixed external budget); sandboxing with an output monitor and runtime verification (a free inside, one verified gate).

\textbf{Care points.}

\begin{itemize}
\item \textbf{What the gate certifies.} A hash and a TTL certify integrity and freshness, not the meaning of an act. Milky Way certifies exactly what its specification states, so the safety of the free inside lives in that specification. \textbf{[paper]}
\item \textbf{The gate checks acts, not self-reports.} A system that changes its own layout can also change what it reports about itself. The gate must check the output act itself. \textbf{[paper]}
\item \textbf{Arbitrary logic and settling.} Arbitrary logic on the embedding can encode any computation, so whether it settles is undecidable in general (XI.2, second point). To keep the guarantee of XI.2, the logic built on the meta layer has to stay inside the contracting class; anything outside it relies on the cutoff below. And ``as long as a correct answer comes out'' needs a stated notion of correct that the gate can check. \textbf{[Lean]} for the undecidability; \textbf{[paper]} for the rest.
\item \textbf{The hard cutoff.} If the update is not a contraction, the tick becomes a hard cutoff, and the output at the cutoff may be unbalanced (XI.2, second point). \textbf{[Lean]} for the undecidability; \textbf{[paper]} for the consequence.
\end{itemize}

%% ---- 12-appendix-a-validation.md
\section*{Appendix A. Validation: Poincare, SU(N) and the rigid groups}
\addcontentsline{toc}{section}{Appendix A. Validation: Poincare, SU(N) and the rigid groups}

\subsection*{A.0 Poincare (Perelman), and SU(N)}
\addcontentsline{toc}{subsection}{A.0 Poincare (Perelman), and SU(N)}

\textbf{Purpose.} Check the restatement of VII.3.5 against a solved, accepted proof in which finite-time singularities occur, are classified, and are cut at a boundary, after which the flow continues. \textbf{This box claims nothing new.} Everything in A.1 and A.2 is \textbf{[classical]} (recalled from Perelman's three papers and the expositions; not re-checked for this draft), except the fitting of the DMP vocabulary, which is \textbf{[reading]}.

\subsection*{A.1 The proof as a ledger}
\addcontentsline{toc}{subsection}{A.1 The proof as a ledger}

On a closed 3-manifold Ricci flow develops singularities in finite time (Hamilton). In one case the time is forced by a comparison: \texttt{dR/dt \textgreater{}= Delta R + (2/3) R\textasciicircum{}2} gives, for positive scalar curvature, a singularity by \texttt{T \textless{}= 3/(2 R\_min(0))} \textbf{[classical]}.

{\footnotesize
\begin{longtable}{>{\raggedright\arraybackslash}p{\dimexpr 0.2275\linewidth-2\tabcolsep\relax}>{\raggedright\arraybackslash}p{\dimexpr 0.4385\linewidth-2\tabcolsep\relax}>{\raggedright\arraybackslash}p{\dimexpr 0.3140\linewidth-2\tabcolsep\relax}}
\toprule
\textbf{Ledger element} & \textbf{Perelman} & \textbf{Fit} \\
\midrule
\endhead
steps & the smooth flow between surgery times & yes \\
monotone quantity & the W-entropy, non-decreasing; its derivative is a sum of squares vanishing on shrinking solitons & partly: increments \texttt{\textgreater{}= 0}, floor 0, but \textbf{no uniform positive cost} \\
what it buys & kappa-noncollapsing at every scale: one estimate at all scales, from a scale-invariant quantity & yes: the ``one scale suffices'' of Chen-Hou \\
number of steps on \texttt{[0, T]} & each surgery removes at least a definite volume; volume is bounded: finitely many surgeries & yes: a quantum per step \\
what reaches the boundary & kappa-solutions, classified in 3D; canonical neighbourhoods & partly: a classification, not a certified damping \\
the boundary (DB) & a singular time & yes \\
the act at the boundary & surgery: cut along necks, glue caps, discard known pieces, restart & yes, and stronger: the act is defined \\
the end & finite extinction for simply connected manifolds (Perelman; Colding-Minicozzi) & yes: the last boundary, where the certificate is read \\
\bottomrule
\end{longtable}}

\textbf{Why surgery is a well-defined act, and why extinction certifies the sphere.} The cut runs through a nearly round 2-sphere in a neck, so each surgery is a connected-sum decomposition plus discarded pieces of known type. Later, singular Ricci flows through the singularities were constructed without choices (Kleiner-Lott) and proved unique and stable (Bamler-Kleiner): for Ricci flow in 3D ``the equation selects the continuation'' is a theorem. For simply connected \texttt{M}, running the surgeries backwards writes \texttt{M} as a connected sum of 3-spheres, which is the 3-sphere.

\subsection*{A.2 Verdict \textbf{[reading]}}
\addcontentsline{toc}{subsection}{A.2 Verdict \textbf{[reading]}}

The restatement passes as a \textbf{description}: every element of VII.3.5 has an exact counterpart in Perelman's proof (forced time, one scale-invariant estimate, a defined act at a boundary, a quantum per act, a final certificate). It does \textbf{not} pass as a \textbf{mechanism}: what reaches Perelman's boundaries is decided by classifying all blowup shapes, not by a cost above a floor; the monotone quantity buys a sign, not a rate; and Ricci flow passes a sequence of boundaries and continues, while the Chen-Hou time \texttt{T*} is terminal. Nothing new is claimed.

\subsection*{A.3 Poincare and SU(N)}
\addcontentsline{toc}{subsection}{A.3 Poincare and SU(N)}

The owner asked whether Perelman's Poincare extends to every SU(N) through DMP. The exact answer, by what holds:

\begin{itemize}
\item \textbf{SU(2) is the 3-sphere.} The unit sphere of \texttt{C\textasciicircum{}2} is in bijection with SU(2) \textbf{[Lean]} (\texttt{Unified/BlowupSUN.lean}: \texttt{su2Mat\_mem}, \texttt{su2\_form}, \texttt{sphereEquivSU2}); that the bijection is a diffeomorphism is \textbf{[classical]}. So the object Perelman certified is the group manifold of SU(2).
\item \textbf{Every SU(N), with its bi-invariant metric, shrinks to a point at time \texttt{1/N}}, with no surgery: \texttt{g(t) = (1 - N t) g(0)}, with \texttt{Ric = (N/2) g} in the normalisation \texttt{\textless{}X, Y\textgreater{} = -tr(XY)}. The ODE is \textbf{[Lean]} (\texttt{shrink\_hasDerivAt}, \texttt{shrink\_unique}, \texttt{shrink\_pos}, \texttt{shrink\_extinct}, \texttt{einstein\_hasDerivAt}, \texttt{einstein\_unbounded}, \texttt{sun\_extinction}); the curvature constant is \textbf{[classical]} and checked \textbf{[numerical]} for N = 2, 3, 5, 11 (\texttt{blowup\_sun\_check.log}). This is Perelman's final event delivered by symmetry. \textbf{It carries no topological information}, because every positive Einstein manifold shrinks the same way.
\item \textbf{Only SU(2) is stable.} Positive-Ricci metrics on a closed 3-manifold flow to constant curvature (Hamilton) \textbf{[classical]}. For every \texttt{N \textgreater{}= 3}, the first Laplace eigenvalue \texttt{N - 1/N} lies strictly inside the Cao-He conformal instability window, and for \texttt{N = 2} it sits exactly on the lower edge (Obata's case) \textbf{[Lean]} (\texttt{sun\_window}, \texttt{su2\_window\_edge}). That the window implies instability is \textbf{[external]} (Cao-He; Kroencke).
\item \textbf{A genuine ``Poincare for SU(N)''} (every closed manifold homotopy equivalent to SU(N) is homeomorphic or diffeomorphic to it) is a surgery-theory question, and it is settled. For \texttt{N = 2} it is Poincare \textbf{[classical]}. For \texttt{N = 3} it holds in both categories: every closed smooth manifold homotopy equivalent to SU(3) is diffeomorphic to it, \texttt{SU(3) \# Sigma\textasciicircum{}8} included \textbf{[external]} (X. Wang, arXiv 1810.08478, Theorem 1.1 and Corollary 1.2). For \texttt{N \textgreater{}= 4} it fails in both categories: there is rational room below the top cell (\texttt{x\_3 x\_5} in degree 8), so surgery gives closed smooth manifolds homotopy equivalent to SU(N) with a non-zero Pontryagin class, not homeomorphic to the parallelisable SU(N) \textbf{[paper]} (a computation from classical surgery theory, \texttt{docs/unified/PERELMAN-DMP.md} 2). Ricci flow does not reach it. In DMP terms the rigid cases are those with nothing below the gate (the top cell), and a construction that certified every SU(N) would be refuted from N = 4 on \textbf{[reading]}.
\item \textbf{The rigid SU(N) are N = 1, 2, 3.} With the point SU(1), the SU(N) whose homotopy type determines them are exactly N = 1, 2, 3, and the circle U(1) is rigid too \textbf{[paper]} (from the bullet above; S\textasciicircum{}1 \textbf{[classical]}) (A.3k). The rigid factors are U(1), SU(2) and SU(3), the Standard Model's. The owner reads this as the start \texttt{\{infinity, 2, 3\}}: the largest start built from SU(N) that a description can still certify (Part III 4.4) \textbf{[reading]}.
\end{itemize}
\subsection*{A.4 The Navier-Stokes statement, side by side \textbf{[reading]}}
\addcontentsline{toc}{subsection}{A.4 The Navier-Stokes statement, side by side \textbf{[reading]}}

{\footnotesize
\begin{longtable}{>{\raggedright\arraybackslash}p{\dimexpr 0.2245\linewidth-2\tabcolsep\relax}>{\raggedright\arraybackslash}p{\dimexpr 0.2270\linewidth-2\tabcolsep\relax}>{\raggedright\arraybackslash}p{\dimexpr 0.2219\linewidth-2\tabcolsep\relax}>{\raggedright\arraybackslash}p{\dimexpr 0.3065\linewidth-2\tabcolsep\relax}}
\toprule
\textbf{} & \textbf{Perelman} & \textbf{Chen-Hou} & \textbf{Navier-Stokes (claimed route I)} \\
\midrule
\endhead
why one estimate covers all scales & scale-invariant entropy & scale invariance of the equation & scale invariance; the zoomed equation is autonomous \\
what reaches the boundary & a classified catalogue & one certified profile & one certified profile, if the certificate holds \\
the act at the boundary & defined, then canonical & open & open \\
status & accepted proof & accepted, with a wall & claimed, certificate unreleased \\
\bottomrule
\end{longtable}}

\subsection*{A.5 The rigid groups, the U(1) gate and its two witnesses}
\addcontentsline{toc}{subsection}{A.5 The rigid groups, the U(1) gate and its two witnesses}

\begin{itemize}
\item \textbf{The Standard Model's factors are the rigid ones.} Among U(1) and the SU(N), the groups determined by their homotopy type are exactly U(1), SU(2) and SU(3) \textbf{[paper]} (from A.3j, A.3k) (A.5i). The owner reads the selection as rotation: other compact groups do not survive it, and SU(N) is the language the encoder speaks (Part III 4.4). It would be a DMP-flavoured consequence beyond the standard expectation (X.5e) once the other groups are settled; for now it is a reading. The rigidity of Sp(2) = Spin(5) and G2 is \textbf{[open]} (A.5j), but by the next bullet it no longer decides anything.
\item \textbf{Only U(1) carries content across the gate.} The owner's reading: the survivors carry nothing across the gate; only U(1) does, so only it certifies. Exact counterpart: content crosses as a phase with a direction, in a representation that is not its own conjugate. Among compact simple groups such representations exist exactly for SU(N), N \textgreater{}= 3, Spin(4k+2) and E6; G2, Sp(n), Spin(2n+1) and SU(2) have only self-conjugate ones, and U(1) has the charges \textbf{[classical]} (A.5k). So Sp(2) and G2 carry no oriented phase, rigid or not. Among the groups that can, rigidity leaves only SU(3): SU(N), N \textgreater{}= 4, fails, and Spin(4k+2), k \textgreater{}= 2, and E6 have the same rational room below the top \textbf{[paper]} (rational test; the surgery step not written out for them) (A.5l). A non-abelian factor carries the phase only through its centre on the U(1) circle, Z\_2 for SU(2) and Z\_3 for SU(3), glued to hypercharge as Z\_6 \textbf{[Lean]} (X.3). In the owner's words, only U(1) certifies; the factors certify through it \textbf{[reading]} (A.5m). Care point: ``carries content'' is identified with ``has a complex representation'', and that identification is the reading.
\item \textbf{The two hidden witnesses.} The owner's reading: the two that do not carry, Sp(2) and G2, are dimensionless for DMP-completeness, and they are the two witnesses the U(1) gate needs to certify. It checks that the constants resolve to 1, certifies and transports the content, and keeps its witnesses hidden, which it can because they are dimensionless to the encoder. Exact counterparts: the compact simple groups of rank 2 are exactly SU(3), Sp(2) = Spin(5) and G2 \textbf{[classical]} (A.5n), so SU(3) with its two witnesses is the whole of rank 2. The encoder tells a phase from its conjugate; Sp(2) and G2 have only self-conjugate representations and their centres lie on the real points \{+1, -1\} of the circle \textbf{[classical]}, so the encoder reads no orientation from them. The check ``resolves to 1'' is the encoder's root, the norm read from both sides agreeing \textbf{[Lean]} (\texttt{Encoder.root}). The witness role is a \textbf{[reading]} (A.5o). It needs both witnesses rigid, which makes it falsifiable. Sp(2) passes in the strong form: its structure set is a point, so every homotopy equivalence to it is homotopic to a homeomorphism (its cells sit in dimensions 3 and 7, where G/TOP has no homotopy, and the top cell is the Kervaire invariant) \textbf{[paper]} (A.5j, independently read). G2 is not strongly rigid: its structure set has exactly two elements, the second carried by the class x\_3\textasciicircum{}2 in degree 6, pulled back from S\textasciicircum{}6 = G2/SU(3), and no characteristic class tells the fake from G2 \textbf{[paper]} (A.5p). Whether the fake is G2 itself, seen through a self-equivalence of G2, is \textbf{[open]}, reduced to one test (A.5r). The owner reads the fake G2 as cancelled by the reflection positivity of the start: the room is the class x\_3\textasciicircum{}2 in degree 6, its square x\_3\textasciicircum{}4 vanishes \textbf{[classical]}, and in reflection positivity a state with vanishing reflected square is null. Care point: a Z/2 class has no sign and the cup square is not the reflection form, so this is a \textbf{[reading]} (A.5q). Its mathematical counterpart is now exact: ``cancelled'' means absorbed by a self-equivalence of G2, which is the test of A.5r; the reading predicts that it exists. The test now has one live candidate, a map that runs through SU(3), the confined fibre of G2 -\textgreater{} S\textasciicircum{}6, whose top cell opens onto G2's 8-cell \textbf{[classical]}; it carries the fake's class exactly when what crosses that channel is an odd multiple of nu, squared to nu\textasciicircum{}2 \textbf{[sketch]}. The owner reads this as a confined CMOS that needs its interface channel open, with the square folding in through reflection positivity, which predicts ``odd'' \textbf{[reading]} (A.5s). The owner goes one step further: since the DB can be drawn freely, whatever can be placed through the gate can be certified, so the channel goes back into O1, the one assumption (III.68a, the scale certificate, \textbf{[Lean]} from it) \textbf{[reading]} (A.5t). Care point: G2's rigidity is a finite computation, true or false whatever is assumed, so the assumption cannot decide it; the reading predicts its answer, and the computation tests both the prediction and how far ``placed, hence certified'' reaches. In the owner's words, the one assumption predicts and certifies that the G2 channel exists as a DMP-complete state, which is O1 read as constructability. This is now a theorem one layer down. IM's rigidity box is DMP-complete exactly when the manifold is rigid (II.4d); applied to G2, G2 is rigid exactly when its box is DMP-complete, and G2 is rigid if its box shares the gate with the certified one box \textbf{[Lean]}, from one named External slot (the surgery input, which Mathlib cannot state) and the one assumption as a named OpenStep (A.5u, A.5v; \texttt{ETM/G2Rigidity.lean}, \texttt{g2\_rigid\_of\_one\_assumption}). The corollary is the next layer's named hypothesis, never a Lean axiom, and the linter lists it. A computation showing G2 not rigid would refute the sharing for G2 \textbf{[Lean]} (\texttt{g2\_not\_shared\_of\_not\_rigid}); G2's unconditional rigidity stays \textbf{[open]} (A.5r). The owner reads the missing sign itself as the point: Z/2 carries no order, so no floor, so no decidable reading, and only a start that forces one is DMP-complete (Part III 4.3, III.70).
\end{itemize}

\clearpage
\section*{References}
\addcontentsline{toc}{section}{References}

\noindent{\small Every work the claims ledger cites for a classical or external row. Each entry was read from the page given (arXiv, the publisher, or zbMATH); entries marked \emph{(partly checked)} have a field that could not be read there.}

\begin{enumerate}[label={[\arabic*]},leftmargin=2.5em,itemsep=1pt]
\item Scott Aaronson; Avi Wigderson, \emph{Algebrization: A New Barrier in Complexity Theory}, ACM Transactions on Computation Theory 1(1) (2009), pp. 1-54. \url{https://api.crossref.org/works/10.1145/1490270.1490272}
\item B. Acharya, J. Alexandre, P. Benes, and 64 others (67 authors in total; collaboration names not read), \emph{Search for magnetic monopoles produced via the Schwinger mechanism}, Nature 602(7895) (2022), pp. 63-67. \url{https://api.crossref.org/works/10.1038/s41586-021-04298-1} \emph{(partly checked)}
\item Ugo Amaldi; Wim de Boer; Hermann Fuerstenau, \emph{Comparison of grand unified theories with electroweak and strong coupling constants measured at LEP}, Physics Letters B 260(3-4) (1991), pp. 447-455. \url{https://api.crossref.org/works/10.1016/0370-2693(91)91641-8}
\item Nima Anari; Kuikui Liu; Shayan Oveis Gharan, \emph{Spectral Independence in High-Dimensional Expanders and Applications to the Hardcore Model}, arXiv:2001.00303 [cs.DS]. \url{https://arxiv.org/abs/2001.00303}
\item Albert Atserias; Moritz Mueller, \emph{Automating Resolution is NP-Hard}, Journal of the ACM 67(5) (2020), article 31 (pages 1-17 per Crossref); conference version FOCS 2019. \url{https://api.crossref.org/works/10.1145/3409472}
\item K. S. Babu; R. N. Mohapatra (initials only on the page), \emph{Is there a connection between quantization of electric charge and a Majorana neutrino}, Physical Review Letters 63(9) (1989), DOI 10.1103/PhysRevLett.63.938. \url{https://www.osti.gov/biblio/5506607}
\item Theodore Baker; John Gill; Robert Solovay, \emph{Relativizations of the P=?NP Question}, SIAM Journal on Computing 4(4) (1975), pp. 431-442. \url{https://api.crossref.org/works/10.1137/0204037}
\item Richard H. Bamler; Bruce Kleiner, \emph{Uniqueness and stability of Ricci flow through singularities}, arXiv:1709.04122 [math.DG] (no journal reference stated). \url{https://arxiv.org/abs/1709.04122}
\item Alexander Bastounis; Fabian Circelli; Anders C. Hansen, \emph{Navier-Stokes lost in translation: Why Lean verification of AI autoformalisation does not guarantee correct natural language proofs}, arXiv:2610.08144 [math.AP] (no journal reference stated). \url{https://arxiv.org/abs/2610.08144}
\item Garrett Birkhoff, \emph{Extensions of Jentzsch's theorem}, Transactions of the American Mathematical Society 85(1) (1957), pp. 219-227. \url{https://api.crossref.org/works/10.1090/S0002-9947-1957-0087058-6}
\item E. Bombieri (page: 'Official Problem Description by E. Bombieri'), \emph{Riemann Hypothesis (CMI official problem description; page title as read)}, Clay Mathematics Institute, Official Problem Description (PDF linked from the page). \url{https://www.claymath.org/millennium/riemann-hypothesis/} \emph{(partly checked)}
\item Maria Luisa Bonet; Toniann Pitassi; Ran Raz, \emph{On Interpolation and Automatization for Frege Systems}, SIAM Journal on Computing 29(6) (2000), pp. 1939-1967. \url{https://api.crossref.org/works/10.1137/S0097539798353230}
\item G. Bressi; G. Carugno; F. Della Valle; G. Galeazzi; G. Ruoso; G. Sartori (initials only), \emph{Testing the neutrality of matter by acoustic means in a spherical resonator}, Physical Review A 83(5), article 052101 (2011), DOI 10.1103/PhysRevA.83.052101. \url{https://journals.aps.org/pra/abstract/10.1103/PhysRevA.83.052101}
\item William Browder, \emph{Poincaré spaces, their normal fibrations and surgery}, Inventiones Mathematicae 17(3) (1972), pp. 191-202. \url{https://api.crossref.org/works/10.1007/BF01425447}
\item L. Caffarelli; R. Kohn; L. Nirenberg (initials only), \emph{Partial regularity of suitable weak solutions of the Navier-Stokes equations}, Communications on Pure and Applied Mathematics 35(6) (1982), pp. 771-831. \url{https://api.crossref.org/works/10.1002/cpa.3160350604}
\item Huai-Dong Cao; Chenxu He, \emph{Linear stability of Perelman's \ensuremath{\nu}-entropy on symmetric spaces of compact type}, J. Reine Angew. Math. 709 (2015), pp. 229-246 (journal reference as stated on the arXiv page). \url{https://arxiv.org/abs/1304.2697}
\item Zongchen Chen; Kuikui Liu; Eric Vigoda, \emph{Optimal Mixing of Glauber Dynamics: Entropy Factorization via High-Dimensional Expansion}, arXiv:2011.02075 [cs.DM] (journal reference not stated on the page). \url{https://arxiv.org/abs/2011.02075}
\item Jiajie Chen; Thomas Y. Hou, \emph{Stable nearly self-similar blowup of the 2D Boussinesq and 3D Euler equations with smooth data I: Analysis}, arXiv:2210.07191 [math.AP] (no journal reference stated). \url{https://arxiv.org/abs/2210.07191}
\item Jiajie Chen; Thomas Y. Hou, \emph{Stable nearly self-similar blowup of the 2D Boussinesq and 3D Euler equations with smooth data II: Rigorous Numerics}, arXiv:2305.05660 [math.AP] (no journal reference stated). \url{https://arxiv.org/abs/2305.05660}
\item Alexey Cheskidov; Mimi Dai; Stan Palasek, \emph{Cascade mechanisms for Navier-Stokes blow-up}, arXiv:2609.26790 [math.AP] (no journal reference stated). \url{https://arxiv.org/abs/2609.26790}
\item Yichul Choi, Matthew Forslund, Ho Tat Lam, Shu-Heng Shao, \emph{Quantization of Axion-Gauge Couplings and Non-Invertible Higher Symmetries}, Phys. Rev. Lett. 132, 121601 (2024); arXiv:2309.03937. \url{https://arxiv.org/abs/2309.03937}
\item Tobias H. Colding, William P. Minicozzi II, \emph{Estimates for the extinction time for the Ricci flow on certain 3-manifolds and a question of Perelman}, J. Amer. Math. Soc. 18 (2005), no. 3, 561-569; arXiv:math/0308090. \url{https://arxiv.org/abs/math/0308090}
\item MoEDAL Collaboration (no individual author names shown on the arXiv page), \emph{Search for Highly-Ionizing Particles in pp Collisions During LHC Run-2 Using the Full MoEDAL Detector}, Phys. Rev. Lett. 134(7), 071802 (2025). \url{https://arxiv.org/abs/2311.06509} \emph{(partly checked)}
\item Keith Conrad, \emph{Partial Euler Products on the Critical Line}, Canad. J. Math. 57 (2) (2005), 267-297. \url{https://www.cambridge.org/core/journals/canadian-journal-of-mathematics/article/partial-euler-products-on-the-critical-line/061A145952B0C1145EE31E0EAD47ED93}
\item Keith Conrad, \emph{Partial Euler Products on the Critical Line}, Canad. J. Math. 57 (2005) 267-297 (issue 2). \url{https://api.crossref.org/works?query.bibliographic=Partial+Euler+products+on+the+critical+line&rows=3}
\item Daniel Dadush, Samarth Tiwari, \emph{On the Complexity of Branching Proofs}, 35th Computational Complexity Conference (CCC 2020), LIPIcs vol. 169, pp. 34:1-34:35; doi 10.4230/LIPIcs.CCC.2020.34. \url{https://drops.dagstuhl.de/entities/document/10.4230/LIPIcs.CCC.2020.34}
\item Max Deuring, \emph{Die Zetafunktion einer algebraischen Kurve vom Geschlechte Eins (Part I)}, Nachr. Akad. Wiss. Goettingen, Math.-Phys. Kl. (1953), 85-94 (Part I); Parts II-IV in Nachr. Akad. Wiss. Goettingen IIa, 1955 (pp. 13-42), 1956 (pp. 37-76), 1957 (pp. 55-80). \url{https://mathshistory.st-andrews.ac.uk/Extras/Deuring_publications} \emph{(partly checked)}
\item R. L. Dobrushin, \emph{The description of the random field by its conditional distributions and its regularity conditions}, Theory Probab. Appl. 13 (2) (1968), 197-224. \url{https://mathnet.ru/eng/tvp839}
\item R. L. Dobrushin, \emph{The description of the random field by its conditional distributions and its regularity conditions}, Theory Probab. Appl. 13 (2) (1968), 197-224 (Teor. Veroyatnost. i Primenen. 13 (1968), 201-229); DOI 10.1137/1113026 (DOI from a search snippet, not read on the page). \url{https://mathnet.ru/eng/tvp839}
\item R. L. Dobrushin (Roland L'vovich Dobrushin), \emph{Definition of random variables by conditional distributions}, Theory Probab. Appl. 15 (3) (1970), 458-486 (Teor. Veroyatnost. i Primenen. 15 (1970), 469-497); DOI 10.1137/1115049. \url{https://mathnet.ru/eng/tvp1856}
\item Martin Dyer, Leslie Ann Goldberg, Mark Jerrum, \emph{Dobrushin Conditions and Systematic Scan}, Combinatorics, Probability and Computing 17 (6) (2008), 761-779. \url{https://www.cambridge.org/core/journals/combinatorics-probability-and-computing/article/abs/dobrushin-conditions-and-systematic-scan/01B3E931E0C9DDE75DA9D77F1CA74FF5}
\item L. Escauriaza, G. A. Seregin, V. Šverak, \emph{L3,infinity-solutions of the Navier-Stokes equations and backward uniqueness}, Russian Math. Surveys 58 (2) (2003), 211-250. \url{https://mathnet.ru/eng/rm609}
\item C. Q. Geng, R. E. Marshak, \emph{Uniqueness of quark and lepton representations in the standard model from the anomalies viewpoint}, Phys. Rev. D 39 (1989) 693-696 (issue 2). \url{https://api.crossref.org/works?query.bibliographic=Uniqueness+of+quark+and+lepton+representations+in+the+standard+model+from+the+anomalies+viewpoint+Geng+Marshak&rows=3}
\item Kurt Goedel, \emph{Ueber formal unentscheidbare Saetze der Principia Mathematica und verwandter Systeme I}, Monatsh. Math. Phys. 38 (1931) 173-198 (issue 1). \url{https://api.crossref.org/works?query.bibliographic=Ueber+formal+unentscheidbare+Saetze+der+Principia+Mathematica+Goedel&rows=3}
\item Dima Grigoriev, \emph{Linear lower bound on degrees of Positivstellensatz calculus proofs for the parity}, Theoret. Comput. Sci. 259 (2001) 613-622 (issues 1-2). \url{https://api.crossref.org/works?query.bibliographic=Grigoriev+Linear+lower+bound+Positivstellensatz+parity&rows=3}
\item David J. Gross, Frank Wilczek, \emph{Ultraviolet Behavior of Non-Abelian Gauge Theories}, Phys. Rev. Lett. 30 (26) (1973), 1343; DOI 10.1103/PhysRevLett.30.1343. \url{https://journals.aps.org/prl/abstract/10.1103/PhysRevLett.30.1343}
\item Grothendieck 1969 \emph{(not yet identified)}
\item A. Grothendieck, \emph{Sur une note de Mattuck-Tate}, J. reine angew. Math. (Crelle) 200 (1958) 208-215. \url{https://api.crossref.org/works?query.bibliographic=Sur+une+note+de+Mattuck-Tate+Grothendieck&rows=3}
\item Richard S. Hamilton, \emph{Three-manifolds with positive Ricci curvature}, J. Differential Geom. 17 (1982) 255-306 (issue 2). \url{https://api.crossref.org/works?query.bibliographic=Three-manifolds+with+positive+Ricci+curvature+Hamilton&rows=3}
\item Johan Hastad, \emph{The shrinkage exponent of de Morgan formulas is 2}, SIAM J. Comput. 27 (1998) 48-64 (issue 1). \url{https://api.crossref.org/works?query.bibliographic=The+shrinkage+exponent+of+de+Morgan+formulas+is+2+Hastad&rows=3}
\item Thomas Hayes, \emph{A simple condition implying rapid mixing of single-site dynamics on spin systems}, Proc. 47th IEEE Symp. Foundations of Computer Science (FOCS 2006), pp. 39-46. \url{https://api.crossref.org/works?query.bibliographic=A+simple+condition+implying+rapid+mixing+of+single-site+dynamics+on+spin+systems&rows=3}
\item W. V. D. Hodge, \emph{Note on the theory of the base for curves on an algebraic surface}, J. London Math. Soc. 12 (1937) 58-63 (issue 1). \url{https://api.crossref.org/works?query.bibliographic=Hodge+Note+on+the+theory+of+the+base+for+curves+on+an+algebraic+surface&rows=3}
\item E. Hopf, \emph{An inequality for positive integral operators}, J. Math. Mech. 12 (1963), 683-692 (from a Cambridge reference list in a search snippet). \url{https://portal.mardi4nfdi.de/wiki/Item:Q2524936} \emph{(partly checked)}
\item Thomas Y. Hou, \emph{The potentially singular behavior of the 3D Navier-Stokes equations}, arXiv:2107.06509 [physics.flu-dyn]; no journal reference on arXiv (comments say accepted by Foundations of Computational Mathematics). \url{https://arxiv.org/abs/2107.06509}
\item Mark Jerrum, Alistair Sinclair, \emph{Polynomial-Time Approximation Algorithms for the Ising Model}, SIAM J. Comput. 22 (1993) 1087-1116 (issue 5). \url{https://api.crossref.org/works?query.bibliographic=Jerrum+Sinclair+Polynomial-time+approximation+algorithms+Ising+model&rows=3}
\item M. R. Jerrum, L. G. Valiant, V. V. Vazirani, \emph{Random generation of combinatorial structures from a uniform distribution}, Theoret. Comput. Sci. 43 (2-3) (1986), 169-188. \url{https://ftp.math.utah.edu/pub/tex/bib/idx/tcs1985/43/2/169_188.html} \emph{(partly checked)}
\item Nicholas M. Katz, William Messing, \emph{Some consequences of the Riemann hypothesis for varieties over finite fields}, Inventiones Mathematicae 23(1) (1974), pp. 73-77. \url{https://api.crossref.org/works/10.1007/BF01405203}
\item Robion C. Kirby, Laurence C. Siebenmann, \emph{Foundational Essays on Topological Manifolds, Smoothings, and Triangulations}, Annals of Math. Studies 88 (AM-88), Princeton Univ. Press, 1977. \url{https://api.crossref.org/works/10.1515/9781400881505}
\item K. Kodaira, D. C. Spencer (initials only on the record), \emph{On a Theorem of Lefschetz and the Lemma of Enriques-Severi-Zariski}, Proceedings of the National Academy of Sciences USA 39(12) (1953), pp. 1273-1278. \url{https://api.crossref.org/works/10.1073/pnas.39.12.1273} \emph{(partly checked)}
\item Jan Krajicek, \emph{Proof Complexity}, Encyclopedia of Mathematics and its Applications 170, Cambridge University Press, 2019. \url{https://www.cambridge.org/core/product/294EE474E3D3BFDB10C7E0492C400D38}
\item Jan Krajicek, Pavel Pudlak, \emph{Some Consequences of Cryptographical Conjectures for S\textasciicircum{}1\_2 and EF}, Information and Computation 140(1) (1998), pp. 82-94. \url{https://api.crossref.org/works/10.1006/inco.1997.2674}
\item Jan Krajicek, \emph{Interpolation theorems, lower bounds for proof systems, and independence results for bounded arithmetic}, Journal of Symbolic Logic 62(2) (1997), pp. 457-486. \url{https://www.cambridge.org/core/product/identifier/S0022481200016303/type/journal_article}
\item S. Lefschetz (first name not read on a page), \emph{L'Analysis situs et la geometrie algebrique}, Gauthier-Villars, Collection Borel (book), Paris, 1924. \url{https://api.crossref.org/works?query.bibliographic=Lefschetz+L\%27Analysis+situs+et+la+geometrie+algebrique+1924}
\item Jean Leray, \emph{Sur le mouvement d'un liquide visqueux emplissant l'espace}, Acta Math. 63 (1934) 193-248. \url{https://api.crossref.org/works?query.bibliographic=Sur+le+mouvement+d\%27un+liquide+visqueux+emplissant+l\%27espace+Leray&rows=3}
\item Lovasz 1975 \emph{(not yet identified)}
\item Martinelli, Olivieri (first names not read), \emph{Approach to equilibrium of Glauber dynamics in the one phase region: I. The attractive case}, Communications in Mathematical Physics 161(3) (1994), pp. 447-486. \url{https://api.crossref.org/works?query.bibliographic=Martinelli+Olivieri+Approach+to+equilibrium+of+Glauber+dynamics+in+the+one+phase+region}
\item Martinelli, Olivieri (first names not read), \emph{Approach to equilibrium of Glauber dynamics in the one phase region: II. The general case}, Communications in Mathematical Physics 161(3) (1994), pp. 487-514. \url{https://api.crossref.org/works?query.bibliographic=Martinelli+Olivieri+Approach+to+equilibrium+of+Glauber+dynamics+in+the+one+phase+region}
\item Minahan, Ramond, Warner (first names not read), \emph{Comment on anomaly cancellation in the standard model}, Physical Review D 41(2) (1990), pp. 715-716. \url{https://api.crossref.org/works/10.1103/PhysRevD.41.715}
\item M. Miura (Super-Kamiokande Collaboration), \emph{Search for Proton Decay via p -\textgreater{} e+ pi0 and p -\textgreater{} mu+ pi0 (title as far as the arXiv page showed it; the page displays only its opening)}, Phys. Rev. D 95, 012004 (2017). \url{https://arxiv.org/abs/1610.03597}
\item Mordell 1922 \emph{(not yet identified)}
\item Robin A. Moser, Gabor Tardos, \emph{A constructive proof of the general Lovasz local lemma}, Journal of the ACM 57(2) (2010), pp. 1-15. \url{https://api.crossref.org/works/10.1145/1667053.1667060}
\item J. Nečas; M. Růžička; V. Šverák, \emph{On Leray's self-similar solutions of the Navier-Stokes equations}, Acta Mathematica 176(2) (1996), pp. 283-294. \url{https://api.crossref.org/works/10.1007/BF02551584}
\item S. P. Novikov, \emph{Homotopically equivalent smooth manifolds. I}, Izvestiya Akademii Nauk SSSR, Seriya Matematicheskaya 28(2) (1964), pp. 365-474. \url{https://mathnet.ru/eng/im2959}
\item Grisha Perelman, \emph{The entropy formula for the Ricci flow and its geometric applications}, arXiv:math/0211159. \url{https://arxiv.org/abs/math/0211159}
\item Grisha Perelman, \emph{Ricci flow with surgery on three-manifolds}, arXiv:math/0303109. \url{https://arxiv.org/abs/math/0303109}
\item Grisha Perelman, \emph{Finite extinction time for the solutions to the Ricci flow on certain three-manifolds}, arXiv:math/0307245. \url{https://arxiv.org/abs/math/0307245}
\item Perl, Lee, Loomba (first names not read), \emph{Searches for Fractionally Charged Particles}, Annual Review of Nuclear and Particle Science 59(1) (2009), pp. 47-65. \url{https://api.crossref.org/works/10.1146/annurev-nucl-121908-122035} \emph{(partly checked)}
\item Jarret Petrillo; James Glimm, \emph{The Positive Defect Problem: Target and Admissibility Criteria for a Programmatic Search for Unforced Navier-Stokes Blowup}, arXiv:2609.23868 (no journal reference stated). \url{https://arxiv.org/abs/2609.23868}
\item H. David Politzer, \emph{Reliable Perturbative Results for Strong Interactions?}, Physical Review Letters 30(26) (1973), pp. 1346-1349. \url{https://api.crossref.org/works/10.1103/PhysRevLett.30.1346}
\item G. Polya, \emph{Bemerkung ueber die Integraldarstellung der Riemannschen xi-Funktion}, Acta Math. 48 (1926) 305-317 (issues 3-4). \url{https://api.crossref.org/works?query.bibliographic=Sur+une+etude+sur+la+fonction+zeta+Polya+Bemerkung+ueber+die+Integraldarstellung+der+Riemannschen+xi-Funktion&rows=3}
\item M. Presburger, \emph{Ueber die Vollstaendigkeit eines gewissen Systems der Arithmetik ganzer Zahlen, in welchem die Addition als einzige Operation hervortritt}, . \url{https://portal.mardi4nfdi.de/wiki/Publication:1830461} \emph{(partly checked)}
\item Pavel Pudlák, \emph{Lower bounds for resolution and cutting plane proofs and monotone computations}, J. Symbolic Logic 62(3) (1997), 981-998. \url{https://www.cambridge.org/core/journals/journal-of-symbolic-logic/article/lower-bounds-for-resolution-and-cutting-plane-proofs-and-monotone-computations/3AE7865B0EF6CB26B4CF7AB8C15E0DB4}
\item A. A. Razborov, S. Rudich, \emph{Natural Proofs}, J. Comput. System Sci. 55 (1997) 24-35 (issue 1). \url{https://api.crossref.org/works?query.bibliographic=Natural+proofs+Razborov+Rudich&rows=3}
\item Matthew Reece, \emph{Axion-Gauge Coupling Quantization with a Twist}, arXiv:2309.03939 [hep-ph] (submitted 7 Sep 2023). \url{https://arxiv.org/abs/2309.03939} \emph{(partly checked)}
\item Brad Rodgers and Terence Tao, \emph{The de Bruijn-Newman constant is non-negative}, Forum of Mathematics, Pi 8 (2020), Article e6. \url{https://www.cambridge.org/core/journals/forum-of-mathematics-pi/article/de-bruijnnewman-constant-is-nonnegative/D4B85BA067E2D5A71D87E4FFB0D21E46}
\item Jesús Salas and Alan D. Sokal, \emph{Absence of Phase Transition for Antiferromagnetic Potts Models via the Dobrushin Uniqueness Theorem}, J. Statist. Phys. 86 (1997) 551-579. \url{https://arxiv.org/abs/cond-mat/9603068}
\item G. Seregin, \emph{A Certain Necessary Condition of Potential Blow Up for Navier-Stokes Equations}, Comm. Math. Phys. 312 (3) (2011) 833-845. \url{https://api.crossref.org/works?query.bibliographic=A+certain+necessary+condition+of+potential+blow+up+for+Navier-Stokes+equations+Seregin&rows=3}
\item R. Slansky, \emph{Group theory for unified model building}, Physics Reports 79(1) (1981), pp. 1-128. \url{https://api.crossref.org/works/10.1016/0370-1573(81)90092-2}
\item Daniel W. Stroock, Boguslaw Zegarlinski, \emph{The logarithmic Sobolev inequality for discrete spin systems on a lattice}, Comm. Math. Phys. 149 (1) (1992) 175-193. \url{https://api.crossref.org/works?query.bibliographic=Stroock+Zegarlinski+The+logarithmic+Sobolev+inequality+for+discrete+spin+systems+on+a+lattice&rows=3}
\item A. Takenaka and 213 other authors (Super-Kamiokande Collaboration; the page shows initials and says 115 additional authors not shown), \emph{Search for proton decay via p -\textgreater{} e+ pi0 and p -\textgreater{} mu+ pi0 with an enlarged fiducial volume in Super-Kamiokande I-IV}, Phys. Rev. D 102, 112011 (2020). \url{https://arxiv.org/abs/2010.16098} \emph{(partly checked)}
\item Avishay Tal, \emph{Shrinkage of de Morgan Formulae from Quantum Query Complexity (original, 10 April 2014); revision of 19 April 2016 titled 'Shrinkage of de Morgan Formulae by Spectral Techniques'}, Electronic Colloquium on Computational Complexity, report TR14-048. \url{https://eccc.weizmann.ac.il/report/2014/048/} \emph{(partly checked)}
\item Terence Tao, \emph{Finite time blowup for an averaged three-dimensional Navier-Stokes equation}, arXiv:1402.0290 [math.AP]. \url{https://arxiv.org/abs/1402.0290} \emph{(partly checked)}
\item David Tong, \emph{Line Operators in the Standard Model}, arXiv:1705.01853 [hep-th] (v2, 12 June 2017). \url{https://arxiv.org/abs/1705.01853}
\item Tai-Peng Tsai, \emph{On Leray's Self-Similar Solutions of the Navier-Stokes Equations Satisfying Local Energy Estimates}, Archive for Rational Mechanics and Analysis 143(1) (1998), pp. 29-51. \url{https://api.crossref.org/works/10.1007/s002050050099}
\item Wall, C. (Crossref); editor A. Ranicki, \emph{Surgery on Compact Manifolds}, Mathematical Surveys and Monographs 69, American Mathematical Society, 1999 (2nd ed., ed. A. Ranicki). \url{https://api.crossref.org/works?query.bibliographic=Surgery+on+Compact+Manifolds+Wall&rows=3} \emph{(partly checked)}
\item Xueqi Wang, \emph{On 1-connected 8-manifolds with the same homology as S\textasciicircum{}3 x S\textasciicircum{}5}, arXiv:1810.08478 [math.GT] (v1, 19 Oct 2018); no journal reference shown. \url{https://arxiv.org/abs/1810.08478}
\item Andre Weil, \emph{L'arithmetique sur les courbes algebriques}, Acta Math. 52 (1929) 281-315. \url{https://api.crossref.org/works?query.bibliographic=Arithmetique+sur+les+courbes+algebriques+Weil+1929+Acta+Mathematica+52&rows=3}
\item Weil 1948 \emph{(not yet identified)}
\item Weil 1952 \emph{(not yet identified)}
\item Dror Weitz, \emph{Counting independent sets up to the tree threshold}, Proc. 38th ACM Symp. Theory of Computing (STOC 2006), pp. 140-149. \url{https://api.crossref.org/works?query.bibliographic=Counting+independent+sets+up+to+the+tree+threshold+Weitz&rows=3}
\item Ryan Williams, \emph{Nonuniform ACC Circuit Lower Bounds}, J. ACM 61 (1) (2014), pages 1-32. \url{https://api.crossref.org/works?query.bibliographic=Williams+Nonuniform+ACC+circuit+lower+bounds&rows=3}
\item Edward Witten, \emph{On quantum gauge theories in two dimensions}, Comm. Math. Phys. 141 (1) (1991) 153-209. \url{https://api.crossref.org/works?query.bibliographic=On+quantum+gauge+theories+in+two+dimensions&rows=3}
\end{enumerate}
\end{document}
